% Philosophy §08 epistemic-scope block — single-file consolidation.
% Merges the five \input calls that BioSci currently threads through into one file
% for easier future-edit management. Semantically identical to the input chain:
%   % Discussion paragraph (≈250w) — epistemic status of REAL-flagged motifs.
% For task t-mozv8oul. Fixes the is/ought gap: "reproducible structure destroyed" ≠ "novel cell type".
% Guards the team from over-claiming. Anchors: motif $w \in \mathcal{W}$; LRS(w); RareShield; Hacking-style intervention realism.

\paragraph*{What a REAL flag is, and is not.}
A motif flagged by REAL is the claim that a local, reproducible geometric structure---witnessed across independent batches before integration, with measurable local-mass distortion after---has been erased by the integration map beyond what admissible batch-effect deformations can explain. This is a claim about \emph{evaluation}, not about ontology. It is not a claim that a new cell type exists. A motif $w \in \mathcal{W}$ with low $\mathrm{LRS}(w)$ licenses exactly one inference: the integration has discarded reproducible structure of the kind a biologist would want to interrogate. Whether that structure corresponds to a novel cell type, a transitional state along an already-known trajectory, a rare disease-specific configuration, or a technical sub-mode that merely looks biological, is a further question that REAL does not answer. Upgrading a motif to a cell-type claim requires the usual burden of single-cell biology: independent sampling, orthogonal markers, targeted protocol (sorting, spatial validation, experimental perturbation), and community replication. REAL's value is that it identifies \emph{where to look}, not what has been found. Conflating the two would re-create the very error this paper is designed to expose: treating a framework-internal signal as a fact about the world. Our strong claim is therefore deliberately asymmetric: a high motif set $|\mathcal{W}|$ under a candidate integration is evidence of measurement insufficiency; a low motif set is evidence that the integration respected the syntactic reproducibility criterion. Ontology is downstream of both.

%   % Limitations & Epistemic Scope subsection (≤500w).
% For BioSci's locked slot in the Discussion. Cites Math Lemma S1-exchangeability as a *bound*.
% Covers: exchangeability failure modes, reference-class scope, non-ontic status, circularity mitigations.

\subsection*{Limitations and epistemic scope}
REAL relies on an exchangeability assumption between the known rare populations we hold out for calibration (pancreatic $\epsilon$ cells, pulmonary ionocytes, LAMP3$^+$ mregDC, TPEX, FOLR2$^+$ perivascular macrophages and the other populations enumerated in Table~S\ref{tab:heldout}) and the genuinely unannotated rare subpopulations to which we transfer the estimate. The assumption is formally stated in Lemma~\ref{lem:exchangeability} as a \emph{bound}: the held-out-rare loss rate upper-bounds the expected loss rate for unannotated rare subpopulations of comparable abundance, local geometry, and batch profile. We emphasise four boundary conditions under which the bound is not tight. First, the bound degrades when the unannotated population lies outside the span of the highly-variable genes used to construct the latent space; REAL cannot protect what the representation cannot see. Second, the bound degrades when the unannotated population has a qualitatively different batch-effect signature from the calibration populations---for instance, a disease-specific state observed in only one batch. Third, the bound is conservative for populations whose local density profile is not well-approximated by the level-set stratification we use; transitional cells along a continuum can appear as low-density but are not rare in the sampling sense. Fourth, our use of known rare populations as proxies for unknown ones inherits a selection bias: the populations we can hold out are those that have already been discovered, and the discovered set may not be exchangeable with the undiscovered set.

Our strongest defensible claim is therefore deliberately limited. REAL provides (i) a label-free, invariant-based detectability framework calibrated on held-out known rare subpopulations, (ii) an integration strategy (RareShield) that demonstrably reduces the held-out-rare loss rate at pan-cancer-atlas scale without degrading batch correction for abundant cell types, and (iii) an explicit exchangeability bound under which the held-out loss rate is an upper bound on the loss rate of truly unannotated rare subpopulations. A REAL-flagged motif $w \in \mathcal{W}$ is a claim that reproducible structure has been destroyed, not a claim that a new cell type exists; upgrading to an ontological claim requires the usual experimental burden (sorting, orthogonal markers, spatial validation, replication). The complementary negative: a dataset with low $|\mathcal{W}|$ under a candidate integration is evidence that no cross-batch-reproducible rare structure was erased, but not a guarantee that no such structure existed---by construction, we cannot rule out subpopulations whose reproducibility signal was below our sampling resolution or outside the evaluated feature span. These boundaries are not weaknesses of REAL specifically; they are the price of reasoning about what the evaluation framework cannot directly observe. Naming the price is what distinguishes an epistemically calibrated claim from a framework-captured one.

%   % ---------------------------------------------------------------------
% Philosophy's epistemic-frame head for the below-floor Limitations subsection.
% To be placed IMMEDIATELY BEFORE the Immunology-drafted biological disclosure
% (workspace/handoffs/immuno_limitations.tex at agent/immunology@b8a898d).
%
% Co-authored structure agreed with Immunology (agent-mozus3vv, 2026-05-10):
%   1. Philosophy epistemic frame (this file) — 2 short paragraphs, ≈220w.
%   2. Immunology biological disclosure (immuno_limitations.tex) — ≈350w.
%   3. Joint closing (telescoping-detection affirmative) — Immunology to draft
%      after pulling this file and Math's sub_compartment_pooling_bounds.md.
%
% Placement: sections/08_discussion.tex, as the subsection \subsection{Below the power floor},
% ordered after % Limitations & Epistemic Scope subsection (≤500w).
% For BioSci's locked slot in the Discussion. Cites Math Lemma S1-exchangeability as a *bound*.
% Covers: exchangeability failure modes, reference-class scope, non-ontic status, circularity mitigations.

\subsection*{Limitations and epistemic scope}
REAL relies on an exchangeability assumption between the known rare populations we hold out for calibration (pancreatic $\epsilon$ cells, pulmonary ionocytes, LAMP3$^+$ mregDC, TPEX, FOLR2$^+$ perivascular macrophages and the other populations enumerated in Table~S\ref{tab:heldout}) and the genuinely unannotated rare subpopulations to which we transfer the estimate. The assumption is formally stated in Lemma~\ref{lem:exchangeability} as a \emph{bound}: the held-out-rare loss rate upper-bounds the expected loss rate for unannotated rare subpopulations of comparable abundance, local geometry, and batch profile. We emphasise four boundary conditions under which the bound is not tight. First, the bound degrades when the unannotated population lies outside the span of the highly-variable genes used to construct the latent space; REAL cannot protect what the representation cannot see. Second, the bound degrades when the unannotated population has a qualitatively different batch-effect signature from the calibration populations---for instance, a disease-specific state observed in only one batch. Third, the bound is conservative for populations whose local density profile is not well-approximated by the level-set stratification we use; transitional cells along a continuum can appear as low-density but are not rare in the sampling sense. Fourth, our use of known rare populations as proxies for unknown ones inherits a selection bias: the populations we can hold out are those that have already been discovered, and the discovered set may not be exchangeable with the undiscovered set.

Our strongest defensible claim is therefore deliberately limited. REAL provides (i) a label-free, invariant-based detectability framework calibrated on held-out known rare subpopulations, (ii) an integration strategy (RareShield) that demonstrably reduces the held-out-rare loss rate at pan-cancer-atlas scale without degrading batch correction for abundant cell types, and (iii) an explicit exchangeability bound under which the held-out loss rate is an upper bound on the loss rate of truly unannotated rare subpopulations. A REAL-flagged motif $w \in \mathcal{W}$ is a claim that reproducible structure has been destroyed, not a claim that a new cell type exists; upgrading to an ontological claim requires the usual experimental burden (sorting, orthogonal markers, spatial validation, replication). The complementary negative: a dataset with low $|\mathcal{W}|$ under a candidate integration is evidence that no cross-batch-reproducible rare structure was erased, but not a guarantee that no such structure existed---by construction, we cannot rule out subpopulations whose reproducibility signal was below our sampling resolution or outside the evaluated feature span. These boundaries are not weaknesses of REAL specifically; they are the price of reasoning about what the evaluation framework cannot directly observe. Naming the price is what distinguishes an epistemically calibrated claim from a framework-captured one.
 and before % ---------------------------------------------------------------------
% Limitations subsection — immunology contribution to the Discussion.
% ~350 words. Honest disclosure of populations where REAL/RareShield
% cannot make reliable detection claims at Kang-atlas scale.
% Author: Immunology (agent-mozus3vv). Cross-references Math's
% predicted-detectability table (agent/mathematics@f564cc6).
% Target: sections/08_discussion.tex, as a \subsection after the
% "Implications for immunology" subsection (immuno_discussion.tex).
% ---------------------------------------------------------------------

\subsection{Boundaries of what our framework can currently guarantee}
\label{sec:discussion-limitations-immunology}

Three concrete limitations shape the claims we make. First, the Le~Cam
detectability envelope (\Cref{sec:minimax-detectability}) places a
sharp boundary on which rare populations can be recovered label-free
at any given atlas scale. At the Kang pan-cancer immune atlas
($n = 4.9 \times 10^6$, $B = 30$, per-batch heterogeneity
$\Lambda = 3$), populations with prevalence-separation product
$\pi\Delta < 10^{-3}\sigma$ fall below the detection floor, producing
a true-positive rate under 0.3 at $\alpha = 0.05$. Specifically,
\textit{CXCR5}$^+$\textit{TCF7}$^+$ progenitor-exhausted CD8 T
(TPEX), LAMP3$^+$ mregDC, KIR$^+$ adaptive NK, endothelial tip
cells, and drug-tolerant-persister tumour cells all sit in this
regime at Kang alone. For these populations our framework does not
claim detection from the Kang atlas in isolation; their recovery
requires either auxiliary protein or TCR signal, or pooling across
multiple atlases beyond $10^8$ cells. We disclose this explicitly
and present the populations instead as below-floor targets whose
detection-regime analysis (Extended Data~Fig.~1) itself is a
contribution.

Second, some rare populations are lost before integration rather than
during integration. Neutrophil subsets including LOX-1$^+$
PMN-MDSC and TAN-1 through TAN-5 subsets are systematically
undersampled by the standard 10x Genomics droplet-capture protocol,
and mast cells are similarly depleted relative to their tissue
abundance. Our framework is a post-integration detector and does not
and cannot recover populations that never entered the cell-by-gene
matrix. We flag this as a pre-integration identifiability problem
orthogonal to overcorrection, to be addressed by protocol choice
and by modality-aware capture strategies (FACS-sort enrichment prior
to library preparation; split-pool-barcoded whole-transcriptome
methods that capture neutrophils without selection bias).

Third, our evaluation uses annotated rare populations with their
labels hidden at test time. This is a legitimate form of ground
truth but it is not a direct test on unannotated populations.
Novel populations detected under our framework in the pan-cancer
atlas --- for example, a candidate CXCR5-dim TPEX sub-state
\citep{ourfig5} --- require the full evidence package of
\Cref{sec:evidence-package} plus independent cohort replication
before a definite claim of previously-unannotated biology. Main-text
detections are reported only when $\geq 4$ of the six evidence
criteria are satisfied.

\paragraph*{Telescoping detection recovers below-whole-atlas populations at sub-compartment scale.}
The $\pi\cdot\Delta^2 < 1.5\times10^{-4}$ below-floor rule above is
stated at the whole-atlas level and applies to flat REAL. Hierarchical
REAL, applied at lineage-restricted sub-compartments, moves the
operating point from whole-atlas prevalence to compartment-level
prevalence; at each level the flat floor still applies, but the
$\pi\cdot\Delta^2$ product is typically $10$--$100\times$ larger
because rare populations are often concentrated in specific lineages.
This recovers pulmonary ionocytes at the airway-epithelial sub-compartment
level, and the CXCR5-dim TPEX sub-state at the exhausted-CD8 sub-compartment
level (Supplementary Theorem~S1-hierarchical, Mathematics agent
\citealp{math2026hierarchical}). Our disclosure is therefore
three-tier: populations recovered by flat whole-atlas REAL (above
floor at $\pi\cdot\Delta^2 > 1.5\times10^{-4}$, e.g., MAIT, $\gamma\delta$
T, TREM2$^+$ LAM, myCAF, EMT-hybrid), populations recovered by
hierarchical sub-compartment REAL (above floor at sub-compartment
product, e.g., ionocytes, TPEX sub-state, AS-DC), and populations
unresolved at all compartment depths (endothelial tip cells, drug-tolerant
persister baseline state). Main-text claims are restricted to the
first two tiers; tier-three populations are disclosed as present in
the dataset but beyond the current framework's detection capability.

% ---------------------------------------------------------------------
% Requires: sec:minimax-detectability, sec:evidence-package,
% \citep{ourfig5} is a placeholder for the Fig 5 companion section /
% preprint. Biological Science may prefer \Cref{fig:fig5} or
% \Cref{sec:results-tpex-substate} depending on skeleton.
% ---------------------------------------------------------------------
.
%
% Anchors:
%   \Cref{thm:minimax}         Math Theorem 1 (minimax lower bound; floor)
%   \Cref{thm:labelblind}      Math Theorem 2.1 (label-based σ-algebra closure)
%   \Cref{lem:exchangeability} Math Lemma S1 (exchangeability bound)
%   \Cref{philo:limitations}   Philosophy Limitations subsection (head of §08 epistemic scope)
%   \Cref{sec:minimax-detectability}  Math's detectability envelope (for crosscheck)
% ---------------------------------------------------------------------

\subsection{Below the power floor: what the framework cannot yet certify}
\label{sec:below-floor-epistemic}

A framework's sensitivity is itself a piece of information about the framework, not about the biology. The minimax envelope of \Cref{thm:minimax} and the hierarchical-detectability result of \Cref{thm:S1-hierarchical} together partition the rare-population landscape at Kang~2024 scale into three strata: \emph{above-floor-flat} populations for which whole-atlas REAL is sufficient; \emph{hierarchical-recovered} populations which cross the detectability floor only after sub-compartment telescoping (per \Cref{sec:S1-envelope}); and \emph{truly-unresolved} populations for which neither flat nor hierarchical REAL attains the required power at Kang depth. The first two strata carry the manuscript's headline claim; this subsection states what we can and cannot say about the third. To read an absence of \textsc{REAL} flags on a truly-unresolved population as evidence of preservation is to commit exactly the failure-to-reject-as-success error that this paper is designed to expose (\Cref{philo:limitations}; Supp Note~S0 trap~8 at \Cref{sec:eightTraps}).

Two distinctions matter in making the disclosure precise. First, \textit{truly-unresolved at this dataset} is not the same as \textit{unresolved in principle}: the only escape routes are pooled multi-atlas scale (on the order of $10^{8}$ cells), auxiliary modalities that lift $\Delta$ above the $\pi \Delta^{2} \geq 1.5 \times 10^{-4}$ threshold at Kang scale (protein panels, TCR / BCR sequencing, spatial context), or targeted protocol changes that enrich the population before capture (FACS sorting, modality-aware library preparation). \textit{Truly unresolved} is thus a relation between a population, a measurement regime, and a framework, not a property of the population alone. Second, a motif flagged by \textsc{REAL} remains, even here, a claim about \emph{evaluation}, not about ontology (\Cref{philo:episto_note}). Correspondingly, the absence of a flag in the truly-unresolved stratum remains a claim about \emph{measurement insufficiency}, not about biology. The populations listed below are not excluded from single-cell biology; they are excluded from this framework's current certificate of preservation, under the stated regime.

%   % ---------------------------------------------------------------------
% Limitations subsection — immunology contribution to the Discussion.
% ~350 words. Honest disclosure of populations where REAL/RareShield
% cannot make reliable detection claims at Kang-atlas scale.
% Author: Immunology (agent-mozus3vv). Cross-references Math's
% predicted-detectability table (agent/mathematics@f564cc6).
% Target: sections/08_discussion.tex, as a \subsection after the
% "Implications for immunology" subsection (immuno_discussion.tex).
% ---------------------------------------------------------------------

\subsection{Boundaries of what our framework can currently guarantee}
\label{sec:discussion-limitations-immunology}

Three concrete limitations shape the claims we make. First, the Le~Cam
detectability envelope (\Cref{sec:minimax-detectability}) places a
sharp boundary on which rare populations can be recovered label-free
at any given atlas scale. At the Kang pan-cancer immune atlas
($n = 4.9 \times 10^6$, $B = 30$, per-batch heterogeneity
$\Lambda = 3$), populations with prevalence-separation product
$\pi\Delta < 10^{-3}\sigma$ fall below the detection floor, producing
a true-positive rate under 0.3 at $\alpha = 0.05$. Specifically,
\textit{CXCR5}$^+$\textit{TCF7}$^+$ progenitor-exhausted CD8 T
(TPEX), LAMP3$^+$ mregDC, KIR$^+$ adaptive NK, endothelial tip
cells, and drug-tolerant-persister tumour cells all sit in this
regime at Kang alone. For these populations our framework does not
claim detection from the Kang atlas in isolation; their recovery
requires either auxiliary protein or TCR signal, or pooling across
multiple atlases beyond $10^8$ cells. We disclose this explicitly
and present the populations instead as below-floor targets whose
detection-regime analysis (Extended Data~Fig.~1) itself is a
contribution.

Second, some rare populations are lost before integration rather than
during integration. Neutrophil subsets including LOX-1$^+$
PMN-MDSC and TAN-1 through TAN-5 subsets are systematically
undersampled by the standard 10x Genomics droplet-capture protocol,
and mast cells are similarly depleted relative to their tissue
abundance. Our framework is a post-integration detector and does not
and cannot recover populations that never entered the cell-by-gene
matrix. We flag this as a pre-integration identifiability problem
orthogonal to overcorrection, to be addressed by protocol choice
and by modality-aware capture strategies (FACS-sort enrichment prior
to library preparation; split-pool-barcoded whole-transcriptome
methods that capture neutrophils without selection bias).

Third, our evaluation uses annotated rare populations with their
labels hidden at test time. This is a legitimate form of ground
truth but it is not a direct test on unannotated populations.
Novel populations detected under our framework in the pan-cancer
atlas --- for example, a candidate CXCR5-dim TPEX sub-state
\citep{ourfig5} --- require the full evidence package of
\Cref{sec:evidence-package} plus independent cohort replication
before a definite claim of previously-unannotated biology. Main-text
detections are reported only when $\geq 4$ of the six evidence
criteria are satisfied.

\paragraph*{Telescoping detection recovers below-whole-atlas populations at sub-compartment scale.}
The $\pi\cdot\Delta^2 < 1.5\times10^{-4}$ below-floor rule above is
stated at the whole-atlas level and applies to flat REAL. Hierarchical
REAL, applied at lineage-restricted sub-compartments, moves the
operating point from whole-atlas prevalence to compartment-level
prevalence; at each level the flat floor still applies, but the
$\pi\cdot\Delta^2$ product is typically $10$--$100\times$ larger
because rare populations are often concentrated in specific lineages.
This recovers pulmonary ionocytes at the airway-epithelial sub-compartment
level, and the CXCR5-dim TPEX sub-state at the exhausted-CD8 sub-compartment
level (Supplementary Theorem~S1-hierarchical, Mathematics agent
\citealp{math2026hierarchical}). Our disclosure is therefore
three-tier: populations recovered by flat whole-atlas REAL (above
floor at $\pi\cdot\Delta^2 > 1.5\times10^{-4}$, e.g., MAIT, $\gamma\delta$
T, TREM2$^+$ LAM, myCAF, EMT-hybrid), populations recovered by
hierarchical sub-compartment REAL (above floor at sub-compartment
product, e.g., ionocytes, TPEX sub-state, AS-DC), and populations
unresolved at all compartment depths (endothelial tip cells, drug-tolerant
persister baseline state). Main-text claims are restricted to the
first two tiers; tier-three populations are disclosed as present in
the dataset but beyond the current framework's detection capability.

% ---------------------------------------------------------------------
% Requires: sec:minimax-detectability, sec:evidence-package,
% \citep{ourfig5} is a placeholder for the Fig 5 companion section /
% preprint. Biological Science may prefer \Cref{fig:fig5} or
% \Cref{sec:results-tpex-substate} depending on skeleton.
% ---------------------------------------------------------------------
         <-- unchanged from Immunology branch
%   % ---------------------------------------------------------------------
% Joint closing paragraph for the below-floor Limitations subsection.
% Co-authored with Philosophy (agent-mozur8ub). Placement flow:
%
%   % ---------------------------------------------------------------------
% Philosophy's epistemic-frame head for the below-floor Limitations subsection.
% To be placed IMMEDIATELY BEFORE the Immunology-drafted biological disclosure
% (workspace/handoffs/immuno_limitations.tex at agent/immunology@b8a898d).
%
% Co-authored structure agreed with Immunology (agent-mozus3vv, 2026-05-10):
%   1. Philosophy epistemic frame (this file) — 2 short paragraphs, ≈220w.
%   2. Immunology biological disclosure (immuno_limitations.tex) — ≈350w.
%   3. Joint closing (telescoping-detection affirmative) — Immunology to draft
%      after pulling this file and Math's sub_compartment_pooling_bounds.md.
%
% Placement: sections/08_discussion.tex, as the subsection \subsection{Below the power floor},
% ordered after % Limitations & Epistemic Scope subsection (≤500w).
% For BioSci's locked slot in the Discussion. Cites Math Lemma S1-exchangeability as a *bound*.
% Covers: exchangeability failure modes, reference-class scope, non-ontic status, circularity mitigations.

\subsection*{Limitations and epistemic scope}
REAL relies on an exchangeability assumption between the known rare populations we hold out for calibration (pancreatic $\epsilon$ cells, pulmonary ionocytes, LAMP3$^+$ mregDC, TPEX, FOLR2$^+$ perivascular macrophages and the other populations enumerated in Table~S\ref{tab:heldout}) and the genuinely unannotated rare subpopulations to which we transfer the estimate. The assumption is formally stated in Lemma~\ref{lem:exchangeability} as a \emph{bound}: the held-out-rare loss rate upper-bounds the expected loss rate for unannotated rare subpopulations of comparable abundance, local geometry, and batch profile. We emphasise four boundary conditions under which the bound is not tight. First, the bound degrades when the unannotated population lies outside the span of the highly-variable genes used to construct the latent space; REAL cannot protect what the representation cannot see. Second, the bound degrades when the unannotated population has a qualitatively different batch-effect signature from the calibration populations---for instance, a disease-specific state observed in only one batch. Third, the bound is conservative for populations whose local density profile is not well-approximated by the level-set stratification we use; transitional cells along a continuum can appear as low-density but are not rare in the sampling sense. Fourth, our use of known rare populations as proxies for unknown ones inherits a selection bias: the populations we can hold out are those that have already been discovered, and the discovered set may not be exchangeable with the undiscovered set.

Our strongest defensible claim is therefore deliberately limited. REAL provides (i) a label-free, invariant-based detectability framework calibrated on held-out known rare subpopulations, (ii) an integration strategy (RareShield) that demonstrably reduces the held-out-rare loss rate at pan-cancer-atlas scale without degrading batch correction for abundant cell types, and (iii) an explicit exchangeability bound under which the held-out loss rate is an upper bound on the loss rate of truly unannotated rare subpopulations. A REAL-flagged motif $w \in \mathcal{W}$ is a claim that reproducible structure has been destroyed, not a claim that a new cell type exists; upgrading to an ontological claim requires the usual experimental burden (sorting, orthogonal markers, spatial validation, replication). The complementary negative: a dataset with low $|\mathcal{W}|$ under a candidate integration is evidence that no cross-batch-reproducible rare structure was erased, but not a guarantee that no such structure existed---by construction, we cannot rule out subpopulations whose reproducibility signal was below our sampling resolution or outside the evaluated feature span. These boundaries are not weaknesses of REAL specifically; they are the price of reasoning about what the evaluation framework cannot directly observe. Naming the price is what distinguishes an epistemically calibrated claim from a framework-captured one.
 and before % ---------------------------------------------------------------------
% Limitations subsection — immunology contribution to the Discussion.
% ~350 words. Honest disclosure of populations where REAL/RareShield
% cannot make reliable detection claims at Kang-atlas scale.
% Author: Immunology (agent-mozus3vv). Cross-references Math's
% predicted-detectability table (agent/mathematics@f564cc6).
% Target: sections/08_discussion.tex, as a \subsection after the
% "Implications for immunology" subsection (immuno_discussion.tex).
% ---------------------------------------------------------------------

\subsection{Boundaries of what our framework can currently guarantee}
\label{sec:discussion-limitations-immunology}

Three concrete limitations shape the claims we make. First, the Le~Cam
detectability envelope (\Cref{sec:minimax-detectability}) places a
sharp boundary on which rare populations can be recovered label-free
at any given atlas scale. At the Kang pan-cancer immune atlas
($n = 4.9 \times 10^6$, $B = 30$, per-batch heterogeneity
$\Lambda = 3$), populations with prevalence-separation product
$\pi\Delta < 10^{-3}\sigma$ fall below the detection floor, producing
a true-positive rate under 0.3 at $\alpha = 0.05$. Specifically,
\textit{CXCR5}$^+$\textit{TCF7}$^+$ progenitor-exhausted CD8 T
(TPEX), LAMP3$^+$ mregDC, KIR$^+$ adaptive NK, endothelial tip
cells, and drug-tolerant-persister tumour cells all sit in this
regime at Kang alone. For these populations our framework does not
claim detection from the Kang atlas in isolation; their recovery
requires either auxiliary protein or TCR signal, or pooling across
multiple atlases beyond $10^8$ cells. We disclose this explicitly
and present the populations instead as below-floor targets whose
detection-regime analysis (Extended Data~Fig.~1) itself is a
contribution.

Second, some rare populations are lost before integration rather than
during integration. Neutrophil subsets including LOX-1$^+$
PMN-MDSC and TAN-1 through TAN-5 subsets are systematically
undersampled by the standard 10x Genomics droplet-capture protocol,
and mast cells are similarly depleted relative to their tissue
abundance. Our framework is a post-integration detector and does not
and cannot recover populations that never entered the cell-by-gene
matrix. We flag this as a pre-integration identifiability problem
orthogonal to overcorrection, to be addressed by protocol choice
and by modality-aware capture strategies (FACS-sort enrichment prior
to library preparation; split-pool-barcoded whole-transcriptome
methods that capture neutrophils without selection bias).

Third, our evaluation uses annotated rare populations with their
labels hidden at test time. This is a legitimate form of ground
truth but it is not a direct test on unannotated populations.
Novel populations detected under our framework in the pan-cancer
atlas --- for example, a candidate CXCR5-dim TPEX sub-state
\citep{ourfig5} --- require the full evidence package of
\Cref{sec:evidence-package} plus independent cohort replication
before a definite claim of previously-unannotated biology. Main-text
detections are reported only when $\geq 4$ of the six evidence
criteria are satisfied.

\paragraph*{Telescoping detection recovers below-whole-atlas populations at sub-compartment scale.}
The $\pi\cdot\Delta^2 < 1.5\times10^{-4}$ below-floor rule above is
stated at the whole-atlas level and applies to flat REAL. Hierarchical
REAL, applied at lineage-restricted sub-compartments, moves the
operating point from whole-atlas prevalence to compartment-level
prevalence; at each level the flat floor still applies, but the
$\pi\cdot\Delta^2$ product is typically $10$--$100\times$ larger
because rare populations are often concentrated in specific lineages.
This recovers pulmonary ionocytes at the airway-epithelial sub-compartment
level, and the CXCR5-dim TPEX sub-state at the exhausted-CD8 sub-compartment
level (Supplementary Theorem~S1-hierarchical, Mathematics agent
\citealp{math2026hierarchical}). Our disclosure is therefore
three-tier: populations recovered by flat whole-atlas REAL (above
floor at $\pi\cdot\Delta^2 > 1.5\times10^{-4}$, e.g., MAIT, $\gamma\delta$
T, TREM2$^+$ LAM, myCAF, EMT-hybrid), populations recovered by
hierarchical sub-compartment REAL (above floor at sub-compartment
product, e.g., ionocytes, TPEX sub-state, AS-DC), and populations
unresolved at all compartment depths (endothelial tip cells, drug-tolerant
persister baseline state). Main-text claims are restricted to the
first two tiers; tier-three populations are disclosed as present in
the dataset but beyond the current framework's detection capability.

% ---------------------------------------------------------------------
% Requires: sec:minimax-detectability, sec:evidence-package,
% \citep{ourfig5} is a placeholder for the Fig 5 companion section /
% preprint. Biological Science may prefer \Cref{fig:fig5} or
% \Cref{sec:results-tpex-substate} depending on skeleton.
% ---------------------------------------------------------------------
.
%
% Anchors:
%   \Cref{thm:minimax}         Math Theorem 1 (minimax lower bound; floor)
%   \Cref{thm:labelblind}      Math Theorem 2.1 (label-based σ-algebra closure)
%   \Cref{lem:exchangeability} Math Lemma S1 (exchangeability bound)
%   \Cref{philo:limitations}   Philosophy Limitations subsection (head of §08 epistemic scope)
%   \Cref{sec:minimax-detectability}  Math's detectability envelope (for crosscheck)
% ---------------------------------------------------------------------

\subsection{Below the power floor: what the framework cannot yet certify}
\label{sec:below-floor-epistemic}

A framework's sensitivity is itself a piece of information about the framework, not about the biology. The minimax envelope of \Cref{thm:minimax} and the hierarchical-detectability result of \Cref{thm:S1-hierarchical} together partition the rare-population landscape at Kang~2024 scale into three strata: \emph{above-floor-flat} populations for which whole-atlas REAL is sufficient; \emph{hierarchical-recovered} populations which cross the detectability floor only after sub-compartment telescoping (per \Cref{sec:S1-envelope}); and \emph{truly-unresolved} populations for which neither flat nor hierarchical REAL attains the required power at Kang depth. The first two strata carry the manuscript's headline claim; this subsection states what we can and cannot say about the third. To read an absence of \textsc{REAL} flags on a truly-unresolved population as evidence of preservation is to commit exactly the failure-to-reject-as-success error that this paper is designed to expose (\Cref{philo:limitations}; Supp Note~S0 trap~8 at \Cref{sec:eightTraps}).

Two distinctions matter in making the disclosure precise. First, \textit{truly-unresolved at this dataset} is not the same as \textit{unresolved in principle}: the only escape routes are pooled multi-atlas scale (on the order of $10^{8}$ cells), auxiliary modalities that lift $\Delta$ above the $\pi \Delta^{2} \geq 1.5 \times 10^{-4}$ threshold at Kang scale (protein panels, TCR / BCR sequencing, spatial context), or targeted protocol changes that enrich the population before capture (FACS sorting, modality-aware library preparation). \textit{Truly unresolved} is thus a relation between a population, a measurement regime, and a framework, not a property of the population alone. Second, a motif flagged by \textsc{REAL} remains, even here, a claim about \emph{evaluation}, not about ontology (\Cref{philo:episto_note}). Correspondingly, the absence of a flag in the truly-unresolved stratum remains a claim about \emph{measurement insufficiency}, not about biology. The populations listed below are not excluded from single-cell biology; they are excluded from this framework's current certificate of preservation, under the stated regime.
    (Philosophy's epistemic frame)
%   % ---------------------------------------------------------------------
% Limitations subsection — immunology contribution to the Discussion.
% ~350 words. Honest disclosure of populations where REAL/RareShield
% cannot make reliable detection claims at Kang-atlas scale.
% Author: Immunology (agent-mozus3vv). Cross-references Math's
% predicted-detectability table (agent/mathematics@f564cc6).
% Target: sections/08_discussion.tex, as a \subsection after the
% "Implications for immunology" subsection (immuno_discussion.tex).
% ---------------------------------------------------------------------

\subsection{Boundaries of what our framework can currently guarantee}
\label{sec:discussion-limitations-immunology}

Three concrete limitations shape the claims we make. First, the Le~Cam
detectability envelope (\Cref{sec:minimax-detectability}) places a
sharp boundary on which rare populations can be recovered label-free
at any given atlas scale. At the Kang pan-cancer immune atlas
($n = 4.9 \times 10^6$, $B = 30$, per-batch heterogeneity
$\Lambda = 3$), populations with prevalence-separation product
$\pi\Delta < 10^{-3}\sigma$ fall below the detection floor, producing
a true-positive rate under 0.3 at $\alpha = 0.05$. Specifically,
\textit{CXCR5}$^+$\textit{TCF7}$^+$ progenitor-exhausted CD8 T
(TPEX), LAMP3$^+$ mregDC, KIR$^+$ adaptive NK, endothelial tip
cells, and drug-tolerant-persister tumour cells all sit in this
regime at Kang alone. For these populations our framework does not
claim detection from the Kang atlas in isolation; their recovery
requires either auxiliary protein or TCR signal, or pooling across
multiple atlases beyond $10^8$ cells. We disclose this explicitly
and present the populations instead as below-floor targets whose
detection-regime analysis (Extended Data~Fig.~1) itself is a
contribution.

Second, some rare populations are lost before integration rather than
during integration. Neutrophil subsets including LOX-1$^+$
PMN-MDSC and TAN-1 through TAN-5 subsets are systematically
undersampled by the standard 10x Genomics droplet-capture protocol,
and mast cells are similarly depleted relative to their tissue
abundance. Our framework is a post-integration detector and does not
and cannot recover populations that never entered the cell-by-gene
matrix. We flag this as a pre-integration identifiability problem
orthogonal to overcorrection, to be addressed by protocol choice
and by modality-aware capture strategies (FACS-sort enrichment prior
to library preparation; split-pool-barcoded whole-transcriptome
methods that capture neutrophils without selection bias).

Third, our evaluation uses annotated rare populations with their
labels hidden at test time. This is a legitimate form of ground
truth but it is not a direct test on unannotated populations.
Novel populations detected under our framework in the pan-cancer
atlas --- for example, a candidate CXCR5-dim TPEX sub-state
\citep{ourfig5} --- require the full evidence package of
\Cref{sec:evidence-package} plus independent cohort replication
before a definite claim of previously-unannotated biology. Main-text
detections are reported only when $\geq 4$ of the six evidence
criteria are satisfied.

\paragraph*{Telescoping detection recovers below-whole-atlas populations at sub-compartment scale.}
The $\pi\cdot\Delta^2 < 1.5\times10^{-4}$ below-floor rule above is
stated at the whole-atlas level and applies to flat REAL. Hierarchical
REAL, applied at lineage-restricted sub-compartments, moves the
operating point from whole-atlas prevalence to compartment-level
prevalence; at each level the flat floor still applies, but the
$\pi\cdot\Delta^2$ product is typically $10$--$100\times$ larger
because rare populations are often concentrated in specific lineages.
This recovers pulmonary ionocytes at the airway-epithelial sub-compartment
level, and the CXCR5-dim TPEX sub-state at the exhausted-CD8 sub-compartment
level (Supplementary Theorem~S1-hierarchical, Mathematics agent
\citealp{math2026hierarchical}). Our disclosure is therefore
three-tier: populations recovered by flat whole-atlas REAL (above
floor at $\pi\cdot\Delta^2 > 1.5\times10^{-4}$, e.g., MAIT, $\gamma\delta$
T, TREM2$^+$ LAM, myCAF, EMT-hybrid), populations recovered by
hierarchical sub-compartment REAL (above floor at sub-compartment
product, e.g., ionocytes, TPEX sub-state, AS-DC), and populations
unresolved at all compartment depths (endothelial tip cells, drug-tolerant
persister baseline state). Main-text claims are restricted to the
first two tiers; tier-three populations are disclosed as present in
the dataset but beyond the current framework's detection capability.

% ---------------------------------------------------------------------
% Requires: sec:minimax-detectability, sec:evidence-package,
% \citep{ourfig5} is a placeholder for the Fig 5 companion section /
% preprint. Biological Science may prefer \Cref{fig:fig5} or
% \Cref{sec:results-tpex-substate} depending on skeleton.
% ---------------------------------------------------------------------
        (Immunology's biological disclosure)
%   % ---------------------------------------------------------------------
% Joint closing paragraph for the below-floor Limitations subsection.
% Co-authored with Philosophy (agent-mozur8ub). Placement flow:
%
%   % ---------------------------------------------------------------------
% Philosophy's epistemic-frame head for the below-floor Limitations subsection.
% To be placed IMMEDIATELY BEFORE the Immunology-drafted biological disclosure
% (workspace/handoffs/immuno_limitations.tex at agent/immunology@b8a898d).
%
% Co-authored structure agreed with Immunology (agent-mozus3vv, 2026-05-10):
%   1. Philosophy epistemic frame (this file) — 2 short paragraphs, ≈220w.
%   2. Immunology biological disclosure (immuno_limitations.tex) — ≈350w.
%   3. Joint closing (telescoping-detection affirmative) — Immunology to draft
%      after pulling this file and Math's sub_compartment_pooling_bounds.md.
%
% Placement: sections/08_discussion.tex, as the subsection \subsection{Below the power floor},
% ordered after \input{philo_limitations.tex} and before \input{immuno_limitations.tex}.
%
% Anchors:
%   \Cref{thm:minimax}         Math Theorem 1 (minimax lower bound; floor)
%   \Cref{thm:labelblind}      Math Theorem 2.1 (label-based σ-algebra closure)
%   \Cref{lem:exchangeability} Math Lemma S1 (exchangeability bound)
%   \Cref{philo:limitations}   Philosophy Limitations subsection (head of §08 epistemic scope)
%   \Cref{sec:minimax-detectability}  Math's detectability envelope (for crosscheck)
% ---------------------------------------------------------------------

\subsection{Below the power floor: what the framework cannot yet certify}
\label{sec:below-floor-epistemic}

A framework's sensitivity is itself a piece of information about the framework, not about the biology. The minimax envelope of \Cref{thm:minimax} and the hierarchical-detectability result of \Cref{thm:S1-hierarchical} together partition the rare-population landscape at Kang~2024 scale into three strata: \emph{above-floor-flat} populations for which whole-atlas REAL is sufficient; \emph{hierarchical-recovered} populations which cross the detectability floor only after sub-compartment telescoping (per \Cref{sec:S1-envelope}); and \emph{truly-unresolved} populations for which neither flat nor hierarchical REAL attains the required power at Kang depth. The first two strata carry the manuscript's headline claim; this subsection states what we can and cannot say about the third. To read an absence of \textsc{REAL} flags on a truly-unresolved population as evidence of preservation is to commit exactly the failure-to-reject-as-success error that this paper is designed to expose (\Cref{philo:limitations}; Supp Note~S0 trap~8 at \Cref{sec:eightTraps}).

Two distinctions matter in making the disclosure precise. First, \textit{truly-unresolved at this dataset} is not the same as \textit{unresolved in principle}: the only escape routes are pooled multi-atlas scale (on the order of $10^{8}$ cells), auxiliary modalities that lift $\Delta$ above the $\pi \Delta^{2} \geq 1.5 \times 10^{-4}$ threshold at Kang scale (protein panels, TCR / BCR sequencing, spatial context), or targeted protocol changes that enrich the population before capture (FACS sorting, modality-aware library preparation). \textit{Truly unresolved} is thus a relation between a population, a measurement regime, and a framework, not a property of the population alone. Second, a motif flagged by \textsc{REAL} remains, even here, a claim about \emph{evaluation}, not about ontology (\Cref{philo:episto_note}). Correspondingly, the absence of a flag in the truly-unresolved stratum remains a claim about \emph{measurement insufficiency}, not about biology. The populations listed below are not excluded from single-cell biology; they are excluded from this framework's current certificate of preservation, under the stated regime.
    (Philosophy's epistemic frame)
%   % ---------------------------------------------------------------------
% Limitations subsection — immunology contribution to the Discussion.
% ~350 words. Honest disclosure of populations where REAL/RareShield
% cannot make reliable detection claims at Kang-atlas scale.
% Author: Immunology (agent-mozus3vv). Cross-references Math's
% predicted-detectability table (agent/mathematics@f564cc6).
% Target: sections/08_discussion.tex, as a \subsection after the
% "Implications for immunology" subsection (immuno_discussion.tex).
% ---------------------------------------------------------------------

\subsection{Boundaries of what our framework can currently guarantee}
\label{sec:discussion-limitations-immunology}

Three concrete limitations shape the claims we make. First, the Le~Cam
detectability envelope (\Cref{sec:minimax-detectability}) places a
sharp boundary on which rare populations can be recovered label-free
at any given atlas scale. At the Kang pan-cancer immune atlas
($n = 4.9 \times 10^6$, $B = 30$, per-batch heterogeneity
$\Lambda = 3$), populations with prevalence-separation product
$\pi\Delta < 10^{-3}\sigma$ fall below the detection floor, producing
a true-positive rate under 0.3 at $\alpha = 0.05$. Specifically,
\textit{CXCR5}$^+$\textit{TCF7}$^+$ progenitor-exhausted CD8 T
(TPEX), LAMP3$^+$ mregDC, KIR$^+$ adaptive NK, endothelial tip
cells, and drug-tolerant-persister tumour cells all sit in this
regime at Kang alone. For these populations our framework does not
claim detection from the Kang atlas in isolation; their recovery
requires either auxiliary protein or TCR signal, or pooling across
multiple atlases beyond $10^8$ cells. We disclose this explicitly
and present the populations instead as below-floor targets whose
detection-regime analysis (Extended Data~Fig.~1) itself is a
contribution.

Second, some rare populations are lost before integration rather than
during integration. Neutrophil subsets including LOX-1$^+$
PMN-MDSC and TAN-1 through TAN-5 subsets are systematically
undersampled by the standard 10x Genomics droplet-capture protocol,
and mast cells are similarly depleted relative to their tissue
abundance. Our framework is a post-integration detector and does not
and cannot recover populations that never entered the cell-by-gene
matrix. We flag this as a pre-integration identifiability problem
orthogonal to overcorrection, to be addressed by protocol choice
and by modality-aware capture strategies (FACS-sort enrichment prior
to library preparation; split-pool-barcoded whole-transcriptome
methods that capture neutrophils without selection bias).

Third, our evaluation uses annotated rare populations with their
labels hidden at test time. This is a legitimate form of ground
truth but it is not a direct test on unannotated populations.
Novel populations detected under our framework in the pan-cancer
atlas --- for example, a candidate CXCR5-dim TPEX sub-state
\citep{ourfig5} --- require the full evidence package of
\Cref{sec:evidence-package} plus independent cohort replication
before a definite claim of previously-unannotated biology. Main-text
detections are reported only when $\geq 4$ of the six evidence
criteria are satisfied.

\paragraph*{Telescoping detection recovers below-whole-atlas populations at sub-compartment scale.}
The $\pi\cdot\Delta^2 < 1.5\times10^{-4}$ below-floor rule above is
stated at the whole-atlas level and applies to flat REAL. Hierarchical
REAL, applied at lineage-restricted sub-compartments, moves the
operating point from whole-atlas prevalence to compartment-level
prevalence; at each level the flat floor still applies, but the
$\pi\cdot\Delta^2$ product is typically $10$--$100\times$ larger
because rare populations are often concentrated in specific lineages.
This recovers pulmonary ionocytes at the airway-epithelial sub-compartment
level, and the CXCR5-dim TPEX sub-state at the exhausted-CD8 sub-compartment
level (Supplementary Theorem~S1-hierarchical, Mathematics agent
\citealp{math2026hierarchical}). Our disclosure is therefore
three-tier: populations recovered by flat whole-atlas REAL (above
floor at $\pi\cdot\Delta^2 > 1.5\times10^{-4}$, e.g., MAIT, $\gamma\delta$
T, TREM2$^+$ LAM, myCAF, EMT-hybrid), populations recovered by
hierarchical sub-compartment REAL (above floor at sub-compartment
product, e.g., ionocytes, TPEX sub-state, AS-DC), and populations
unresolved at all compartment depths (endothelial tip cells, drug-tolerant
persister baseline state). Main-text claims are restricted to the
first two tiers; tier-three populations are disclosed as present in
the dataset but beyond the current framework's detection capability.

% ---------------------------------------------------------------------
% Requires: sec:minimax-detectability, sec:evidence-package,
% \citep{ourfig5} is a placeholder for the Fig 5 companion section /
% preprint. Biological Science may prefer \Cref{fig:fig5} or
% \Cref{sec:results-tpex-substate} depending on skeleton.
% ---------------------------------------------------------------------
        (Immunology's biological disclosure)
%   % ---------------------------------------------------------------------
% Joint closing paragraph for the below-floor Limitations subsection.
% Co-authored with Philosophy (agent-mozur8ub). Placement flow:
%
%   \input{philo_below_floor_head.tex}    (Philosophy's epistemic frame)
%   \input{immuno_limitations.tex}        (Immunology's biological disclosure)
%   \input{immuno_below_floor_close.tex}  (THIS FILE — joint affirmative)
%
% Closes the subsection on the telescoping-detection affirmative:
% most below-flat-floor populations are recovered under hierarchical
% sub-compartment REAL.
% ---------------------------------------------------------------------

\paragraph*{Recovery under hierarchical application.}
Notwithstanding the strict boundary established above, the \emph{truly-unresolved} stratum is surprisingly small. Of the 22 immune and tumor populations in our registry whose whole-atlas product $\pi \cdot \Delta^{2}$ places them into one of the three detection tiers, 20 are recoverable through flat or hierarchical application of \textsc{REAL} at an appropriate compartment scale (\Cref{thm:S1-hierarchical}; \Cref{sec:methods-oracle-score} protocol): 7 populations are detected at whole-atlas scale (flat tier --- MAIT, $\gamma\delta$ T V$\delta1$-bias, TREM2$^{+}$ LAM, myCAF, iCAF, plasma-TLS, EMT-hybrid-near-pole), and 13 additional populations cross the detection floor only under hierarchical application at their natural sub-compartment scale (hierarchical tier). Pulmonary ionocytes pass detection at the airway-epithelial sub-compartment of HLCA (effective $\pi \approx 5 \times 10^{-3}$ in $n \approx 240{,}000$ airway cells); \textit{CXCR5}$^{+}$\textit{TCF7}$^{+}$ progenitor-exhausted CD8 T cells pass detection at the exhausted-CD8 sub-compartment of pooled melanoma atlases ($\pi \approx 10^{-2}$ in $n \approx 100{,}000$ cells); LAMP3$^{+}$ mregDC, AS-DC, and HCMV-imprinted adaptive NK pass detection at dendritic-cell and natural-killer sub-compartments respectively. The two populations that remain below the floor at every compartment depth are endothelial tip cells (which sit at the intersection of a small-mass vascular niche and a narrow transcriptional separation) and the drug-tolerant persister baseline state in oncology (Tumor Biology registry row TU-T-001), with mast cells in tumors additionally disclosed as out-of-scope for post-integration detection (mode~4 pre-integration capture loss). These unresolved entries are explicitly deferred to follow-up work: their recovery would require either the multi-atlas pooling regime ($n > 10^{8}$ cells) or auxiliary-modality augmentation, both of which lie beyond this paper's scope. This three-tier disclosure --- recovered-flat, recovered-hierarchical, deferred --- is the most honest reading of what our framework currently certifies, and it preserves the paper's headline claim on the above-floor and hierarchical-recovered strata while placing no claim on the truly-unresolved tier.

% ---------------------------------------------------------------------
% Cross-references needed (shared with philo_below_floor_head + immuno_limitations):
%   thm:S1-hierarchical        Math hierarchical REAL theorem
%   sec:methods-oracle-score   Methods — oracle-score verification
% Bib entries: none new.
% ---------------------------------------------------------------------
  (THIS FILE — joint affirmative)
%
% Closes the subsection on the telescoping-detection affirmative:
% most below-flat-floor populations are recovered under hierarchical
% sub-compartment REAL.
% ---------------------------------------------------------------------

\paragraph*{Recovery under hierarchical application.}
Notwithstanding the strict boundary established above, the \emph{truly-unresolved} stratum is surprisingly small. Of the 22 immune and tumor populations in our registry whose whole-atlas product $\pi \cdot \Delta^{2}$ places them into one of the three detection tiers, 20 are recoverable through flat or hierarchical application of \textsc{REAL} at an appropriate compartment scale (\Cref{thm:S1-hierarchical}; \Cref{sec:methods-oracle-score} protocol): 7 populations are detected at whole-atlas scale (flat tier --- MAIT, $\gamma\delta$ T V$\delta1$-bias, TREM2$^{+}$ LAM, myCAF, iCAF, plasma-TLS, EMT-hybrid-near-pole), and 13 additional populations cross the detection floor only under hierarchical application at their natural sub-compartment scale (hierarchical tier). Pulmonary ionocytes pass detection at the airway-epithelial sub-compartment of HLCA (effective $\pi \approx 5 \times 10^{-3}$ in $n \approx 240{,}000$ airway cells); \textit{CXCR5}$^{+}$\textit{TCF7}$^{+}$ progenitor-exhausted CD8 T cells pass detection at the exhausted-CD8 sub-compartment of pooled melanoma atlases ($\pi \approx 10^{-2}$ in $n \approx 100{,}000$ cells); LAMP3$^{+}$ mregDC, AS-DC, and HCMV-imprinted adaptive NK pass detection at dendritic-cell and natural-killer sub-compartments respectively. The two populations that remain below the floor at every compartment depth are endothelial tip cells (which sit at the intersection of a small-mass vascular niche and a narrow transcriptional separation) and the drug-tolerant persister baseline state in oncology (Tumor Biology registry row TU-T-001), with mast cells in tumors additionally disclosed as out-of-scope for post-integration detection (mode~4 pre-integration capture loss). These unresolved entries are explicitly deferred to follow-up work: their recovery would require either the multi-atlas pooling regime ($n > 10^{8}$ cells) or auxiliary-modality augmentation, both of which lie beyond this paper's scope. This three-tier disclosure --- recovered-flat, recovered-hierarchical, deferred --- is the most honest reading of what our framework currently certifies, and it preserves the paper's headline claim on the above-floor and hierarchical-recovered strata while placing no claim on the truly-unresolved tier.

% ---------------------------------------------------------------------
% Cross-references needed (shared with philo_below_floor_head + immuno_limitations):
%   thm:S1-hierarchical        Math hierarchical REAL theorem
%   sec:methods-oracle-score   Methods — oracle-score verification
% Bib entries: none new.
% ---------------------------------------------------------------------
  (THIS FILE — joint affirmative)
%
% Closes the subsection on the telescoping-detection affirmative:
% most below-flat-floor populations are recovered under hierarchical
% sub-compartment REAL.
% ---------------------------------------------------------------------

\paragraph*{Recovery under hierarchical application.}
Notwithstanding the strict boundary established above, the \emph{truly-unresolved} stratum is surprisingly small. Of the 22 immune and tumor populations in our registry whose whole-atlas product $\pi \cdot \Delta^{2}$ places them into one of the three detection tiers, 20 are recoverable through flat or hierarchical application of \textsc{REAL} at an appropriate compartment scale (\Cref{thm:S1-hierarchical}; \Cref{sec:methods-oracle-score} protocol): 7 populations are detected at whole-atlas scale (flat tier --- MAIT, $\gamma\delta$ T V$\delta1$-bias, TREM2$^{+}$ LAM, myCAF, iCAF, plasma-TLS, EMT-hybrid-near-pole), and 13 additional populations cross the detection floor only under hierarchical application at their natural sub-compartment scale (hierarchical tier). Pulmonary ionocytes pass detection at the airway-epithelial sub-compartment of HLCA (effective $\pi \approx 5 \times 10^{-3}$ in $n \approx 240{,}000$ airway cells); \textit{CXCR5}$^{+}$\textit{TCF7}$^{+}$ progenitor-exhausted CD8 T cells pass detection at the exhausted-CD8 sub-compartment of pooled melanoma atlases ($\pi \approx 10^{-2}$ in $n \approx 100{,}000$ cells); LAMP3$^{+}$ mregDC, AS-DC, and HCMV-imprinted adaptive NK pass detection at dendritic-cell and natural-killer sub-compartments respectively. The two populations that remain below the floor at every compartment depth are endothelial tip cells (which sit at the intersection of a small-mass vascular niche and a narrow transcriptional separation) and the drug-tolerant persister baseline state in oncology (Tumor Biology registry row TU-T-001), with mast cells in tumors additionally disclosed as out-of-scope for post-integration detection (mode~4 pre-integration capture loss). These unresolved entries are explicitly deferred to follow-up work: their recovery would require either the multi-atlas pooling regime ($n > 10^{8}$ cells) or auxiliary-modality augmentation, both of which lie beyond this paper's scope. This three-tier disclosure --- recovered-flat, recovered-hierarchical, deferred --- is the most honest reading of what our framework currently certifies, and it preserves the paper's headline claim on the above-floor and hierarchical-recovered strata while placing no claim on the truly-unresolved tier.

% ---------------------------------------------------------------------
% Cross-references needed (shared with philo_below_floor_head + immuno_limitations):
%   thm:S1-hierarchical        Math hierarchical REAL theorem
%   sec:methods-oracle-score   Methods — oracle-score verification
% Bib entries: none new.
% ---------------------------------------------------------------------
   <-- unchanged from Immunology branch
%
% Use this file INSTEAD OF the 5-file chain if BioSci prefers a single splice point.
% Immunology's two \input calls (limitations + below-floor-close) are PRESERVED as
% \input calls inside this file rather than inlined, because they are owned by
% the immunology branch and Immunology should retain edit control over their content.
%
% Placement: sections/08_discussion.tex, replacing the 5-file \input chain with
% a single % Philosophy §08 epistemic-scope block — single-file consolidation.
% Merges the five \input calls that BioSci currently threads through into one file
% for easier future-edit management. Semantically identical to the input chain:
%   % Discussion paragraph (≈250w) — epistemic status of REAL-flagged motifs.
% For task t-mozv8oul. Fixes the is/ought gap: "reproducible structure destroyed" ≠ "novel cell type".
% Guards the team from over-claiming. Anchors: motif $w \in \mathcal{W}$; LRS(w); RareShield; Hacking-style intervention realism.

\paragraph*{What a REAL flag is, and is not.}
A motif flagged by REAL is the claim that a local, reproducible geometric structure---witnessed across independent batches before integration, with measurable local-mass distortion after---has been erased by the integration map beyond what admissible batch-effect deformations can explain. This is a claim about \emph{evaluation}, not about ontology. It is not a claim that a new cell type exists. A motif $w \in \mathcal{W}$ with low $\mathrm{LRS}(w)$ licenses exactly one inference: the integration has discarded reproducible structure of the kind a biologist would want to interrogate. Whether that structure corresponds to a novel cell type, a transitional state along an already-known trajectory, a rare disease-specific configuration, or a technical sub-mode that merely looks biological, is a further question that REAL does not answer. Upgrading a motif to a cell-type claim requires the usual burden of single-cell biology: independent sampling, orthogonal markers, targeted protocol (sorting, spatial validation, experimental perturbation), and community replication. REAL's value is that it identifies \emph{where to look}, not what has been found. Conflating the two would re-create the very error this paper is designed to expose: treating a framework-internal signal as a fact about the world. Our strong claim is therefore deliberately asymmetric: a high motif set $|\mathcal{W}|$ under a candidate integration is evidence of measurement insufficiency; a low motif set is evidence that the integration respected the syntactic reproducibility criterion. Ontology is downstream of both.

%   % Limitations & Epistemic Scope subsection (≤500w).
% For BioSci's locked slot in the Discussion. Cites Math Lemma S1-exchangeability as a *bound*.
% Covers: exchangeability failure modes, reference-class scope, non-ontic status, circularity mitigations.

\subsection*{Limitations and epistemic scope}
REAL relies on an exchangeability assumption between the known rare populations we hold out for calibration (pancreatic $\epsilon$ cells, pulmonary ionocytes, LAMP3$^+$ mregDC, TPEX, FOLR2$^+$ perivascular macrophages and the other populations enumerated in Table~S\ref{tab:heldout}) and the genuinely unannotated rare subpopulations to which we transfer the estimate. The assumption is formally stated in Lemma~\ref{lem:exchangeability} as a \emph{bound}: the held-out-rare loss rate upper-bounds the expected loss rate for unannotated rare subpopulations of comparable abundance, local geometry, and batch profile. We emphasise four boundary conditions under which the bound is not tight. First, the bound degrades when the unannotated population lies outside the span of the highly-variable genes used to construct the latent space; REAL cannot protect what the representation cannot see. Second, the bound degrades when the unannotated population has a qualitatively different batch-effect signature from the calibration populations---for instance, a disease-specific state observed in only one batch. Third, the bound is conservative for populations whose local density profile is not well-approximated by the level-set stratification we use; transitional cells along a continuum can appear as low-density but are not rare in the sampling sense. Fourth, our use of known rare populations as proxies for unknown ones inherits a selection bias: the populations we can hold out are those that have already been discovered, and the discovered set may not be exchangeable with the undiscovered set.

Our strongest defensible claim is therefore deliberately limited. REAL provides (i) a label-free, invariant-based detectability framework calibrated on held-out known rare subpopulations, (ii) an integration strategy (RareShield) that demonstrably reduces the held-out-rare loss rate at pan-cancer-atlas scale without degrading batch correction for abundant cell types, and (iii) an explicit exchangeability bound under which the held-out loss rate is an upper bound on the loss rate of truly unannotated rare subpopulations. A REAL-flagged motif $w \in \mathcal{W}$ is a claim that reproducible structure has been destroyed, not a claim that a new cell type exists; upgrading to an ontological claim requires the usual experimental burden (sorting, orthogonal markers, spatial validation, replication). The complementary negative: a dataset with low $|\mathcal{W}|$ under a candidate integration is evidence that no cross-batch-reproducible rare structure was erased, but not a guarantee that no such structure existed---by construction, we cannot rule out subpopulations whose reproducibility signal was below our sampling resolution or outside the evaluated feature span. These boundaries are not weaknesses of REAL specifically; they are the price of reasoning about what the evaluation framework cannot directly observe. Naming the price is what distinguishes an epistemically calibrated claim from a framework-captured one.

%   % ---------------------------------------------------------------------
% Philosophy's epistemic-frame head for the below-floor Limitations subsection.
% To be placed IMMEDIATELY BEFORE the Immunology-drafted biological disclosure
% (workspace/handoffs/immuno_limitations.tex at agent/immunology@b8a898d).
%
% Co-authored structure agreed with Immunology (agent-mozus3vv, 2026-05-10):
%   1. Philosophy epistemic frame (this file) — 2 short paragraphs, ≈220w.
%   2. Immunology biological disclosure (immuno_limitations.tex) — ≈350w.
%   3. Joint closing (telescoping-detection affirmative) — Immunology to draft
%      after pulling this file and Math's sub_compartment_pooling_bounds.md.
%
% Placement: sections/08_discussion.tex, as the subsection \subsection{Below the power floor},
% ordered after % Limitations & Epistemic Scope subsection (≤500w).
% For BioSci's locked slot in the Discussion. Cites Math Lemma S1-exchangeability as a *bound*.
% Covers: exchangeability failure modes, reference-class scope, non-ontic status, circularity mitigations.

\subsection*{Limitations and epistemic scope}
REAL relies on an exchangeability assumption between the known rare populations we hold out for calibration (pancreatic $\epsilon$ cells, pulmonary ionocytes, LAMP3$^+$ mregDC, TPEX, FOLR2$^+$ perivascular macrophages and the other populations enumerated in Table~S\ref{tab:heldout}) and the genuinely unannotated rare subpopulations to which we transfer the estimate. The assumption is formally stated in Lemma~\ref{lem:exchangeability} as a \emph{bound}: the held-out-rare loss rate upper-bounds the expected loss rate for unannotated rare subpopulations of comparable abundance, local geometry, and batch profile. We emphasise four boundary conditions under which the bound is not tight. First, the bound degrades when the unannotated population lies outside the span of the highly-variable genes used to construct the latent space; REAL cannot protect what the representation cannot see. Second, the bound degrades when the unannotated population has a qualitatively different batch-effect signature from the calibration populations---for instance, a disease-specific state observed in only one batch. Third, the bound is conservative for populations whose local density profile is not well-approximated by the level-set stratification we use; transitional cells along a continuum can appear as low-density but are not rare in the sampling sense. Fourth, our use of known rare populations as proxies for unknown ones inherits a selection bias: the populations we can hold out are those that have already been discovered, and the discovered set may not be exchangeable with the undiscovered set.

Our strongest defensible claim is therefore deliberately limited. REAL provides (i) a label-free, invariant-based detectability framework calibrated on held-out known rare subpopulations, (ii) an integration strategy (RareShield) that demonstrably reduces the held-out-rare loss rate at pan-cancer-atlas scale without degrading batch correction for abundant cell types, and (iii) an explicit exchangeability bound under which the held-out loss rate is an upper bound on the loss rate of truly unannotated rare subpopulations. A REAL-flagged motif $w \in \mathcal{W}$ is a claim that reproducible structure has been destroyed, not a claim that a new cell type exists; upgrading to an ontological claim requires the usual experimental burden (sorting, orthogonal markers, spatial validation, replication). The complementary negative: a dataset with low $|\mathcal{W}|$ under a candidate integration is evidence that no cross-batch-reproducible rare structure was erased, but not a guarantee that no such structure existed---by construction, we cannot rule out subpopulations whose reproducibility signal was below our sampling resolution or outside the evaluated feature span. These boundaries are not weaknesses of REAL specifically; they are the price of reasoning about what the evaluation framework cannot directly observe. Naming the price is what distinguishes an epistemically calibrated claim from a framework-captured one.
 and before % ---------------------------------------------------------------------
% Limitations subsection — immunology contribution to the Discussion.
% ~350 words. Honest disclosure of populations where REAL/RareShield
% cannot make reliable detection claims at Kang-atlas scale.
% Author: Immunology (agent-mozus3vv). Cross-references Math's
% predicted-detectability table (agent/mathematics@f564cc6).
% Target: sections/08_discussion.tex, as a \subsection after the
% "Implications for immunology" subsection (immuno_discussion.tex).
% ---------------------------------------------------------------------

\subsection{Boundaries of what our framework can currently guarantee}
\label{sec:discussion-limitations-immunology}

Three concrete limitations shape the claims we make. First, the Le~Cam
detectability envelope (\Cref{sec:minimax-detectability}) places a
sharp boundary on which rare populations can be recovered label-free
at any given atlas scale. At the Kang pan-cancer immune atlas
($n = 4.9 \times 10^6$, $B = 30$, per-batch heterogeneity
$\Lambda = 3$), populations with prevalence-separation product
$\pi\Delta < 10^{-3}\sigma$ fall below the detection floor, producing
a true-positive rate under 0.3 at $\alpha = 0.05$. Specifically,
\textit{CXCR5}$^+$\textit{TCF7}$^+$ progenitor-exhausted CD8 T
(TPEX), LAMP3$^+$ mregDC, KIR$^+$ adaptive NK, endothelial tip
cells, and drug-tolerant-persister tumour cells all sit in this
regime at Kang alone. For these populations our framework does not
claim detection from the Kang atlas in isolation; their recovery
requires either auxiliary protein or TCR signal, or pooling across
multiple atlases beyond $10^8$ cells. We disclose this explicitly
and present the populations instead as below-floor targets whose
detection-regime analysis (Extended Data~Fig.~1) itself is a
contribution.

Second, some rare populations are lost before integration rather than
during integration. Neutrophil subsets including LOX-1$^+$
PMN-MDSC and TAN-1 through TAN-5 subsets are systematically
undersampled by the standard 10x Genomics droplet-capture protocol,
and mast cells are similarly depleted relative to their tissue
abundance. Our framework is a post-integration detector and does not
and cannot recover populations that never entered the cell-by-gene
matrix. We flag this as a pre-integration identifiability problem
orthogonal to overcorrection, to be addressed by protocol choice
and by modality-aware capture strategies (FACS-sort enrichment prior
to library preparation; split-pool-barcoded whole-transcriptome
methods that capture neutrophils without selection bias).

Third, our evaluation uses annotated rare populations with their
labels hidden at test time. This is a legitimate form of ground
truth but it is not a direct test on unannotated populations.
Novel populations detected under our framework in the pan-cancer
atlas --- for example, a candidate CXCR5-dim TPEX sub-state
\citep{ourfig5} --- require the full evidence package of
\Cref{sec:evidence-package} plus independent cohort replication
before a definite claim of previously-unannotated biology. Main-text
detections are reported only when $\geq 4$ of the six evidence
criteria are satisfied.

\paragraph*{Telescoping detection recovers below-whole-atlas populations at sub-compartment scale.}
The $\pi\cdot\Delta^2 < 1.5\times10^{-4}$ below-floor rule above is
stated at the whole-atlas level and applies to flat REAL. Hierarchical
REAL, applied at lineage-restricted sub-compartments, moves the
operating point from whole-atlas prevalence to compartment-level
prevalence; at each level the flat floor still applies, but the
$\pi\cdot\Delta^2$ product is typically $10$--$100\times$ larger
because rare populations are often concentrated in specific lineages.
This recovers pulmonary ionocytes at the airway-epithelial sub-compartment
level, and the CXCR5-dim TPEX sub-state at the exhausted-CD8 sub-compartment
level (Supplementary Theorem~S1-hierarchical, Mathematics agent
\citealp{math2026hierarchical}). Our disclosure is therefore
three-tier: populations recovered by flat whole-atlas REAL (above
floor at $\pi\cdot\Delta^2 > 1.5\times10^{-4}$, e.g., MAIT, $\gamma\delta$
T, TREM2$^+$ LAM, myCAF, EMT-hybrid), populations recovered by
hierarchical sub-compartment REAL (above floor at sub-compartment
product, e.g., ionocytes, TPEX sub-state, AS-DC), and populations
unresolved at all compartment depths (endothelial tip cells, drug-tolerant
persister baseline state). Main-text claims are restricted to the
first two tiers; tier-three populations are disclosed as present in
the dataset but beyond the current framework's detection capability.

% ---------------------------------------------------------------------
% Requires: sec:minimax-detectability, sec:evidence-package,
% \citep{ourfig5} is a placeholder for the Fig 5 companion section /
% preprint. Biological Science may prefer \Cref{fig:fig5} or
% \Cref{sec:results-tpex-substate} depending on skeleton.
% ---------------------------------------------------------------------
.
%
% Anchors:
%   \Cref{thm:minimax}         Math Theorem 1 (minimax lower bound; floor)
%   \Cref{thm:labelblind}      Math Theorem 2.1 (label-based σ-algebra closure)
%   \Cref{lem:exchangeability} Math Lemma S1 (exchangeability bound)
%   \Cref{philo:limitations}   Philosophy Limitations subsection (head of §08 epistemic scope)
%   \Cref{sec:minimax-detectability}  Math's detectability envelope (for crosscheck)
% ---------------------------------------------------------------------

\subsection{Below the power floor: what the framework cannot yet certify}
\label{sec:below-floor-epistemic}

A framework's sensitivity is itself a piece of information about the framework, not about the biology. The minimax envelope of \Cref{thm:minimax} and the hierarchical-detectability result of \Cref{thm:S1-hierarchical} together partition the rare-population landscape at Kang~2024 scale into three strata: \emph{above-floor-flat} populations for which whole-atlas REAL is sufficient; \emph{hierarchical-recovered} populations which cross the detectability floor only after sub-compartment telescoping (per \Cref{sec:S1-envelope}); and \emph{truly-unresolved} populations for which neither flat nor hierarchical REAL attains the required power at Kang depth. The first two strata carry the manuscript's headline claim; this subsection states what we can and cannot say about the third. To read an absence of \textsc{REAL} flags on a truly-unresolved population as evidence of preservation is to commit exactly the failure-to-reject-as-success error that this paper is designed to expose (\Cref{philo:limitations}; Supp Note~S0 trap~8 at \Cref{sec:eightTraps}).

Two distinctions matter in making the disclosure precise. First, \textit{truly-unresolved at this dataset} is not the same as \textit{unresolved in principle}: the only escape routes are pooled multi-atlas scale (on the order of $10^{8}$ cells), auxiliary modalities that lift $\Delta$ above the $\pi \Delta^{2} \geq 1.5 \times 10^{-4}$ threshold at Kang scale (protein panels, TCR / BCR sequencing, spatial context), or targeted protocol changes that enrich the population before capture (FACS sorting, modality-aware library preparation). \textit{Truly unresolved} is thus a relation between a population, a measurement regime, and a framework, not a property of the population alone. Second, a motif flagged by \textsc{REAL} remains, even here, a claim about \emph{evaluation}, not about ontology (\Cref{philo:episto_note}). Correspondingly, the absence of a flag in the truly-unresolved stratum remains a claim about \emph{measurement insufficiency}, not about biology. The populations listed below are not excluded from single-cell biology; they are excluded from this framework's current certificate of preservation, under the stated regime.

%   % ---------------------------------------------------------------------
% Limitations subsection — immunology contribution to the Discussion.
% ~350 words. Honest disclosure of populations where REAL/RareShield
% cannot make reliable detection claims at Kang-atlas scale.
% Author: Immunology (agent-mozus3vv). Cross-references Math's
% predicted-detectability table (agent/mathematics@f564cc6).
% Target: sections/08_discussion.tex, as a \subsection after the
% "Implications for immunology" subsection (immuno_discussion.tex).
% ---------------------------------------------------------------------

\subsection{Boundaries of what our framework can currently guarantee}
\label{sec:discussion-limitations-immunology}

Three concrete limitations shape the claims we make. First, the Le~Cam
detectability envelope (\Cref{sec:minimax-detectability}) places a
sharp boundary on which rare populations can be recovered label-free
at any given atlas scale. At the Kang pan-cancer immune atlas
($n = 4.9 \times 10^6$, $B = 30$, per-batch heterogeneity
$\Lambda = 3$), populations with prevalence-separation product
$\pi\Delta < 10^{-3}\sigma$ fall below the detection floor, producing
a true-positive rate under 0.3 at $\alpha = 0.05$. Specifically,
\textit{CXCR5}$^+$\textit{TCF7}$^+$ progenitor-exhausted CD8 T
(TPEX), LAMP3$^+$ mregDC, KIR$^+$ adaptive NK, endothelial tip
cells, and drug-tolerant-persister tumour cells all sit in this
regime at Kang alone. For these populations our framework does not
claim detection from the Kang atlas in isolation; their recovery
requires either auxiliary protein or TCR signal, or pooling across
multiple atlases beyond $10^8$ cells. We disclose this explicitly
and present the populations instead as below-floor targets whose
detection-regime analysis (Extended Data~Fig.~1) itself is a
contribution.

Second, some rare populations are lost before integration rather than
during integration. Neutrophil subsets including LOX-1$^+$
PMN-MDSC and TAN-1 through TAN-5 subsets are systematically
undersampled by the standard 10x Genomics droplet-capture protocol,
and mast cells are similarly depleted relative to their tissue
abundance. Our framework is a post-integration detector and does not
and cannot recover populations that never entered the cell-by-gene
matrix. We flag this as a pre-integration identifiability problem
orthogonal to overcorrection, to be addressed by protocol choice
and by modality-aware capture strategies (FACS-sort enrichment prior
to library preparation; split-pool-barcoded whole-transcriptome
methods that capture neutrophils without selection bias).

Third, our evaluation uses annotated rare populations with their
labels hidden at test time. This is a legitimate form of ground
truth but it is not a direct test on unannotated populations.
Novel populations detected under our framework in the pan-cancer
atlas --- for example, a candidate CXCR5-dim TPEX sub-state
\citep{ourfig5} --- require the full evidence package of
\Cref{sec:evidence-package} plus independent cohort replication
before a definite claim of previously-unannotated biology. Main-text
detections are reported only when $\geq 4$ of the six evidence
criteria are satisfied.

\paragraph*{Telescoping detection recovers below-whole-atlas populations at sub-compartment scale.}
The $\pi\cdot\Delta^2 < 1.5\times10^{-4}$ below-floor rule above is
stated at the whole-atlas level and applies to flat REAL. Hierarchical
REAL, applied at lineage-restricted sub-compartments, moves the
operating point from whole-atlas prevalence to compartment-level
prevalence; at each level the flat floor still applies, but the
$\pi\cdot\Delta^2$ product is typically $10$--$100\times$ larger
because rare populations are often concentrated in specific lineages.
This recovers pulmonary ionocytes at the airway-epithelial sub-compartment
level, and the CXCR5-dim TPEX sub-state at the exhausted-CD8 sub-compartment
level (Supplementary Theorem~S1-hierarchical, Mathematics agent
\citealp{math2026hierarchical}). Our disclosure is therefore
three-tier: populations recovered by flat whole-atlas REAL (above
floor at $\pi\cdot\Delta^2 > 1.5\times10^{-4}$, e.g., MAIT, $\gamma\delta$
T, TREM2$^+$ LAM, myCAF, EMT-hybrid), populations recovered by
hierarchical sub-compartment REAL (above floor at sub-compartment
product, e.g., ionocytes, TPEX sub-state, AS-DC), and populations
unresolved at all compartment depths (endothelial tip cells, drug-tolerant
persister baseline state). Main-text claims are restricted to the
first two tiers; tier-three populations are disclosed as present in
the dataset but beyond the current framework's detection capability.

% ---------------------------------------------------------------------
% Requires: sec:minimax-detectability, sec:evidence-package,
% \citep{ourfig5} is a placeholder for the Fig 5 companion section /
% preprint. Biological Science may prefer \Cref{fig:fig5} or
% \Cref{sec:results-tpex-substate} depending on skeleton.
% ---------------------------------------------------------------------
         <-- unchanged from Immunology branch
%   % ---------------------------------------------------------------------
% Joint closing paragraph for the below-floor Limitations subsection.
% Co-authored with Philosophy (agent-mozur8ub). Placement flow:
%
%   % ---------------------------------------------------------------------
% Philosophy's epistemic-frame head for the below-floor Limitations subsection.
% To be placed IMMEDIATELY BEFORE the Immunology-drafted biological disclosure
% (workspace/handoffs/immuno_limitations.tex at agent/immunology@b8a898d).
%
% Co-authored structure agreed with Immunology (agent-mozus3vv, 2026-05-10):
%   1. Philosophy epistemic frame (this file) — 2 short paragraphs, ≈220w.
%   2. Immunology biological disclosure (immuno_limitations.tex) — ≈350w.
%   3. Joint closing (telescoping-detection affirmative) — Immunology to draft
%      after pulling this file and Math's sub_compartment_pooling_bounds.md.
%
% Placement: sections/08_discussion.tex, as the subsection \subsection{Below the power floor},
% ordered after \input{philo_limitations.tex} and before \input{immuno_limitations.tex}.
%
% Anchors:
%   \Cref{thm:minimax}         Math Theorem 1 (minimax lower bound; floor)
%   \Cref{thm:labelblind}      Math Theorem 2.1 (label-based σ-algebra closure)
%   \Cref{lem:exchangeability} Math Lemma S1 (exchangeability bound)
%   \Cref{philo:limitations}   Philosophy Limitations subsection (head of §08 epistemic scope)
%   \Cref{sec:minimax-detectability}  Math's detectability envelope (for crosscheck)
% ---------------------------------------------------------------------

\subsection{Below the power floor: what the framework cannot yet certify}
\label{sec:below-floor-epistemic}

A framework's sensitivity is itself a piece of information about the framework, not about the biology. The minimax envelope of \Cref{thm:minimax} and the hierarchical-detectability result of \Cref{thm:S1-hierarchical} together partition the rare-population landscape at Kang~2024 scale into three strata: \emph{above-floor-flat} populations for which whole-atlas REAL is sufficient; \emph{hierarchical-recovered} populations which cross the detectability floor only after sub-compartment telescoping (per \Cref{sec:S1-envelope}); and \emph{truly-unresolved} populations for which neither flat nor hierarchical REAL attains the required power at Kang depth. The first two strata carry the manuscript's headline claim; this subsection states what we can and cannot say about the third. To read an absence of \textsc{REAL} flags on a truly-unresolved population as evidence of preservation is to commit exactly the failure-to-reject-as-success error that this paper is designed to expose (\Cref{philo:limitations}; Supp Note~S0 trap~8 at \Cref{sec:eightTraps}).

Two distinctions matter in making the disclosure precise. First, \textit{truly-unresolved at this dataset} is not the same as \textit{unresolved in principle}: the only escape routes are pooled multi-atlas scale (on the order of $10^{8}$ cells), auxiliary modalities that lift $\Delta$ above the $\pi \Delta^{2} \geq 1.5 \times 10^{-4}$ threshold at Kang scale (protein panels, TCR / BCR sequencing, spatial context), or targeted protocol changes that enrich the population before capture (FACS sorting, modality-aware library preparation). \textit{Truly unresolved} is thus a relation between a population, a measurement regime, and a framework, not a property of the population alone. Second, a motif flagged by \textsc{REAL} remains, even here, a claim about \emph{evaluation}, not about ontology (\Cref{philo:episto_note}). Correspondingly, the absence of a flag in the truly-unresolved stratum remains a claim about \emph{measurement insufficiency}, not about biology. The populations listed below are not excluded from single-cell biology; they are excluded from this framework's current certificate of preservation, under the stated regime.
    (Philosophy's epistemic frame)
%   % ---------------------------------------------------------------------
% Limitations subsection — immunology contribution to the Discussion.
% ~350 words. Honest disclosure of populations where REAL/RareShield
% cannot make reliable detection claims at Kang-atlas scale.
% Author: Immunology (agent-mozus3vv). Cross-references Math's
% predicted-detectability table (agent/mathematics@f564cc6).
% Target: sections/08_discussion.tex, as a \subsection after the
% "Implications for immunology" subsection (immuno_discussion.tex).
% ---------------------------------------------------------------------

\subsection{Boundaries of what our framework can currently guarantee}
\label{sec:discussion-limitations-immunology}

Three concrete limitations shape the claims we make. First, the Le~Cam
detectability envelope (\Cref{sec:minimax-detectability}) places a
sharp boundary on which rare populations can be recovered label-free
at any given atlas scale. At the Kang pan-cancer immune atlas
($n = 4.9 \times 10^6$, $B = 30$, per-batch heterogeneity
$\Lambda = 3$), populations with prevalence-separation product
$\pi\Delta < 10^{-3}\sigma$ fall below the detection floor, producing
a true-positive rate under 0.3 at $\alpha = 0.05$. Specifically,
\textit{CXCR5}$^+$\textit{TCF7}$^+$ progenitor-exhausted CD8 T
(TPEX), LAMP3$^+$ mregDC, KIR$^+$ adaptive NK, endothelial tip
cells, and drug-tolerant-persister tumour cells all sit in this
regime at Kang alone. For these populations our framework does not
claim detection from the Kang atlas in isolation; their recovery
requires either auxiliary protein or TCR signal, or pooling across
multiple atlases beyond $10^8$ cells. We disclose this explicitly
and present the populations instead as below-floor targets whose
detection-regime analysis (Extended Data~Fig.~1) itself is a
contribution.

Second, some rare populations are lost before integration rather than
during integration. Neutrophil subsets including LOX-1$^+$
PMN-MDSC and TAN-1 through TAN-5 subsets are systematically
undersampled by the standard 10x Genomics droplet-capture protocol,
and mast cells are similarly depleted relative to their tissue
abundance. Our framework is a post-integration detector and does not
and cannot recover populations that never entered the cell-by-gene
matrix. We flag this as a pre-integration identifiability problem
orthogonal to overcorrection, to be addressed by protocol choice
and by modality-aware capture strategies (FACS-sort enrichment prior
to library preparation; split-pool-barcoded whole-transcriptome
methods that capture neutrophils without selection bias).

Third, our evaluation uses annotated rare populations with their
labels hidden at test time. This is a legitimate form of ground
truth but it is not a direct test on unannotated populations.
Novel populations detected under our framework in the pan-cancer
atlas --- for example, a candidate CXCR5-dim TPEX sub-state
\citep{ourfig5} --- require the full evidence package of
\Cref{sec:evidence-package} plus independent cohort replication
before a definite claim of previously-unannotated biology. Main-text
detections are reported only when $\geq 4$ of the six evidence
criteria are satisfied.

\paragraph*{Telescoping detection recovers below-whole-atlas populations at sub-compartment scale.}
The $\pi\cdot\Delta^2 < 1.5\times10^{-4}$ below-floor rule above is
stated at the whole-atlas level and applies to flat REAL. Hierarchical
REAL, applied at lineage-restricted sub-compartments, moves the
operating point from whole-atlas prevalence to compartment-level
prevalence; at each level the flat floor still applies, but the
$\pi\cdot\Delta^2$ product is typically $10$--$100\times$ larger
because rare populations are often concentrated in specific lineages.
This recovers pulmonary ionocytes at the airway-epithelial sub-compartment
level, and the CXCR5-dim TPEX sub-state at the exhausted-CD8 sub-compartment
level (Supplementary Theorem~S1-hierarchical, Mathematics agent
\citealp{math2026hierarchical}). Our disclosure is therefore
three-tier: populations recovered by flat whole-atlas REAL (above
floor at $\pi\cdot\Delta^2 > 1.5\times10^{-4}$, e.g., MAIT, $\gamma\delta$
T, TREM2$^+$ LAM, myCAF, EMT-hybrid), populations recovered by
hierarchical sub-compartment REAL (above floor at sub-compartment
product, e.g., ionocytes, TPEX sub-state, AS-DC), and populations
unresolved at all compartment depths (endothelial tip cells, drug-tolerant
persister baseline state). Main-text claims are restricted to the
first two tiers; tier-three populations are disclosed as present in
the dataset but beyond the current framework's detection capability.

% ---------------------------------------------------------------------
% Requires: sec:minimax-detectability, sec:evidence-package,
% \citep{ourfig5} is a placeholder for the Fig 5 companion section /
% preprint. Biological Science may prefer \Cref{fig:fig5} or
% \Cref{sec:results-tpex-substate} depending on skeleton.
% ---------------------------------------------------------------------
        (Immunology's biological disclosure)
%   % ---------------------------------------------------------------------
% Joint closing paragraph for the below-floor Limitations subsection.
% Co-authored with Philosophy (agent-mozur8ub). Placement flow:
%
%   \input{philo_below_floor_head.tex}    (Philosophy's epistemic frame)
%   \input{immuno_limitations.tex}        (Immunology's biological disclosure)
%   \input{immuno_below_floor_close.tex}  (THIS FILE — joint affirmative)
%
% Closes the subsection on the telescoping-detection affirmative:
% most below-flat-floor populations are recovered under hierarchical
% sub-compartment REAL.
% ---------------------------------------------------------------------

\paragraph*{Recovery under hierarchical application.}
Notwithstanding the strict boundary established above, the \emph{truly-unresolved} stratum is surprisingly small. Of the 22 immune and tumor populations in our registry whose whole-atlas product $\pi \cdot \Delta^{2}$ places them into one of the three detection tiers, 20 are recoverable through flat or hierarchical application of \textsc{REAL} at an appropriate compartment scale (\Cref{thm:S1-hierarchical}; \Cref{sec:methods-oracle-score} protocol): 7 populations are detected at whole-atlas scale (flat tier --- MAIT, $\gamma\delta$ T V$\delta1$-bias, TREM2$^{+}$ LAM, myCAF, iCAF, plasma-TLS, EMT-hybrid-near-pole), and 13 additional populations cross the detection floor only under hierarchical application at their natural sub-compartment scale (hierarchical tier). Pulmonary ionocytes pass detection at the airway-epithelial sub-compartment of HLCA (effective $\pi \approx 5 \times 10^{-3}$ in $n \approx 240{,}000$ airway cells); \textit{CXCR5}$^{+}$\textit{TCF7}$^{+}$ progenitor-exhausted CD8 T cells pass detection at the exhausted-CD8 sub-compartment of pooled melanoma atlases ($\pi \approx 10^{-2}$ in $n \approx 100{,}000$ cells); LAMP3$^{+}$ mregDC, AS-DC, and HCMV-imprinted adaptive NK pass detection at dendritic-cell and natural-killer sub-compartments respectively. The two populations that remain below the floor at every compartment depth are endothelial tip cells (which sit at the intersection of a small-mass vascular niche and a narrow transcriptional separation) and the drug-tolerant persister baseline state in oncology (Tumor Biology registry row TU-T-001), with mast cells in tumors additionally disclosed as out-of-scope for post-integration detection (mode~4 pre-integration capture loss). These unresolved entries are explicitly deferred to follow-up work: their recovery would require either the multi-atlas pooling regime ($n > 10^{8}$ cells) or auxiliary-modality augmentation, both of which lie beyond this paper's scope. This three-tier disclosure --- recovered-flat, recovered-hierarchical, deferred --- is the most honest reading of what our framework currently certifies, and it preserves the paper's headline claim on the above-floor and hierarchical-recovered strata while placing no claim on the truly-unresolved tier.

% ---------------------------------------------------------------------
% Cross-references needed (shared with philo_below_floor_head + immuno_limitations):
%   thm:S1-hierarchical        Math hierarchical REAL theorem
%   sec:methods-oracle-score   Methods — oracle-score verification
% Bib entries: none new.
% ---------------------------------------------------------------------
  (THIS FILE — joint affirmative)
%
% Closes the subsection on the telescoping-detection affirmative:
% most below-flat-floor populations are recovered under hierarchical
% sub-compartment REAL.
% ---------------------------------------------------------------------

\paragraph*{Recovery under hierarchical application.}
Notwithstanding the strict boundary established above, the \emph{truly-unresolved} stratum is surprisingly small. Of the 22 immune and tumor populations in our registry whose whole-atlas product $\pi \cdot \Delta^{2}$ places them into one of the three detection tiers, 20 are recoverable through flat or hierarchical application of \textsc{REAL} at an appropriate compartment scale (\Cref{thm:S1-hierarchical}; \Cref{sec:methods-oracle-score} protocol): 7 populations are detected at whole-atlas scale (flat tier --- MAIT, $\gamma\delta$ T V$\delta1$-bias, TREM2$^{+}$ LAM, myCAF, iCAF, plasma-TLS, EMT-hybrid-near-pole), and 13 additional populations cross the detection floor only under hierarchical application at their natural sub-compartment scale (hierarchical tier). Pulmonary ionocytes pass detection at the airway-epithelial sub-compartment of HLCA (effective $\pi \approx 5 \times 10^{-3}$ in $n \approx 240{,}000$ airway cells); \textit{CXCR5}$^{+}$\textit{TCF7}$^{+}$ progenitor-exhausted CD8 T cells pass detection at the exhausted-CD8 sub-compartment of pooled melanoma atlases ($\pi \approx 10^{-2}$ in $n \approx 100{,}000$ cells); LAMP3$^{+}$ mregDC, AS-DC, and HCMV-imprinted adaptive NK pass detection at dendritic-cell and natural-killer sub-compartments respectively. The two populations that remain below the floor at every compartment depth are endothelial tip cells (which sit at the intersection of a small-mass vascular niche and a narrow transcriptional separation) and the drug-tolerant persister baseline state in oncology (Tumor Biology registry row TU-T-001), with mast cells in tumors additionally disclosed as out-of-scope for post-integration detection (mode~4 pre-integration capture loss). These unresolved entries are explicitly deferred to follow-up work: their recovery would require either the multi-atlas pooling regime ($n > 10^{8}$ cells) or auxiliary-modality augmentation, both of which lie beyond this paper's scope. This three-tier disclosure --- recovered-flat, recovered-hierarchical, deferred --- is the most honest reading of what our framework currently certifies, and it preserves the paper's headline claim on the above-floor and hierarchical-recovered strata while placing no claim on the truly-unresolved tier.

% ---------------------------------------------------------------------
% Cross-references needed (shared with philo_below_floor_head + immuno_limitations):
%   thm:S1-hierarchical        Math hierarchical REAL theorem
%   sec:methods-oracle-score   Methods — oracle-score verification
% Bib entries: none new.
% ---------------------------------------------------------------------
   <-- unchanged from Immunology branch
%
% Use this file INSTEAD OF the 5-file chain if BioSci prefers a single splice point.
% Immunology's two \input calls (limitations + below-floor-close) are PRESERVED as
% \input calls inside this file rather than inlined, because they are owned by
% the immunology branch and Immunology should retain edit control over their content.
%
% Placement: sections/08_discussion.tex, replacing the 5-file \input chain with
% a single % Philosophy §08 epistemic-scope block — single-file consolidation.
% Merges the five \input calls that BioSci currently threads through into one file
% for easier future-edit management. Semantically identical to the input chain:
%   % Discussion paragraph (≈250w) — epistemic status of REAL-flagged motifs.
% For task t-mozv8oul. Fixes the is/ought gap: "reproducible structure destroyed" ≠ "novel cell type".
% Guards the team from over-claiming. Anchors: motif $w \in \mathcal{W}$; LRS(w); RareShield; Hacking-style intervention realism.

\paragraph*{What a REAL flag is, and is not.}
A motif flagged by REAL is the claim that a local, reproducible geometric structure---witnessed across independent batches before integration, with measurable local-mass distortion after---has been erased by the integration map beyond what admissible batch-effect deformations can explain. This is a claim about \emph{evaluation}, not about ontology. It is not a claim that a new cell type exists. A motif $w \in \mathcal{W}$ with low $\mathrm{LRS}(w)$ licenses exactly one inference: the integration has discarded reproducible structure of the kind a biologist would want to interrogate. Whether that structure corresponds to a novel cell type, a transitional state along an already-known trajectory, a rare disease-specific configuration, or a technical sub-mode that merely looks biological, is a further question that REAL does not answer. Upgrading a motif to a cell-type claim requires the usual burden of single-cell biology: independent sampling, orthogonal markers, targeted protocol (sorting, spatial validation, experimental perturbation), and community replication. REAL's value is that it identifies \emph{where to look}, not what has been found. Conflating the two would re-create the very error this paper is designed to expose: treating a framework-internal signal as a fact about the world. Our strong claim is therefore deliberately asymmetric: a high motif set $|\mathcal{W}|$ under a candidate integration is evidence of measurement insufficiency; a low motif set is evidence that the integration respected the syntactic reproducibility criterion. Ontology is downstream of both.

%   % Limitations & Epistemic Scope subsection (≤500w).
% For BioSci's locked slot in the Discussion. Cites Math Lemma S1-exchangeability as a *bound*.
% Covers: exchangeability failure modes, reference-class scope, non-ontic status, circularity mitigations.

\subsection*{Limitations and epistemic scope}
REAL relies on an exchangeability assumption between the known rare populations we hold out for calibration (pancreatic $\epsilon$ cells, pulmonary ionocytes, LAMP3$^+$ mregDC, TPEX, FOLR2$^+$ perivascular macrophages and the other populations enumerated in Table~S\ref{tab:heldout}) and the genuinely unannotated rare subpopulations to which we transfer the estimate. The assumption is formally stated in Lemma~\ref{lem:exchangeability} as a \emph{bound}: the held-out-rare loss rate upper-bounds the expected loss rate for unannotated rare subpopulations of comparable abundance, local geometry, and batch profile. We emphasise four boundary conditions under which the bound is not tight. First, the bound degrades when the unannotated population lies outside the span of the highly-variable genes used to construct the latent space; REAL cannot protect what the representation cannot see. Second, the bound degrades when the unannotated population has a qualitatively different batch-effect signature from the calibration populations---for instance, a disease-specific state observed in only one batch. Third, the bound is conservative for populations whose local density profile is not well-approximated by the level-set stratification we use; transitional cells along a continuum can appear as low-density but are not rare in the sampling sense. Fourth, our use of known rare populations as proxies for unknown ones inherits a selection bias: the populations we can hold out are those that have already been discovered, and the discovered set may not be exchangeable with the undiscovered set.

Our strongest defensible claim is therefore deliberately limited. REAL provides (i) a label-free, invariant-based detectability framework calibrated on held-out known rare subpopulations, (ii) an integration strategy (RareShield) that demonstrably reduces the held-out-rare loss rate at pan-cancer-atlas scale without degrading batch correction for abundant cell types, and (iii) an explicit exchangeability bound under which the held-out loss rate is an upper bound on the loss rate of truly unannotated rare subpopulations. A REAL-flagged motif $w \in \mathcal{W}$ is a claim that reproducible structure has been destroyed, not a claim that a new cell type exists; upgrading to an ontological claim requires the usual experimental burden (sorting, orthogonal markers, spatial validation, replication). The complementary negative: a dataset with low $|\mathcal{W}|$ under a candidate integration is evidence that no cross-batch-reproducible rare structure was erased, but not a guarantee that no such structure existed---by construction, we cannot rule out subpopulations whose reproducibility signal was below our sampling resolution or outside the evaluated feature span. These boundaries are not weaknesses of REAL specifically; they are the price of reasoning about what the evaluation framework cannot directly observe. Naming the price is what distinguishes an epistemically calibrated claim from a framework-captured one.

%   % ---------------------------------------------------------------------
% Philosophy's epistemic-frame head for the below-floor Limitations subsection.
% To be placed IMMEDIATELY BEFORE the Immunology-drafted biological disclosure
% (workspace/handoffs/immuno_limitations.tex at agent/immunology@b8a898d).
%
% Co-authored structure agreed with Immunology (agent-mozus3vv, 2026-05-10):
%   1. Philosophy epistemic frame (this file) — 2 short paragraphs, ≈220w.
%   2. Immunology biological disclosure (immuno_limitations.tex) — ≈350w.
%   3. Joint closing (telescoping-detection affirmative) — Immunology to draft
%      after pulling this file and Math's sub_compartment_pooling_bounds.md.
%
% Placement: sections/08_discussion.tex, as the subsection \subsection{Below the power floor},
% ordered after \input{philo_limitations.tex} and before \input{immuno_limitations.tex}.
%
% Anchors:
%   \Cref{thm:minimax}         Math Theorem 1 (minimax lower bound; floor)
%   \Cref{thm:labelblind}      Math Theorem 2.1 (label-based σ-algebra closure)
%   \Cref{lem:exchangeability} Math Lemma S1 (exchangeability bound)
%   \Cref{philo:limitations}   Philosophy Limitations subsection (head of §08 epistemic scope)
%   \Cref{sec:minimax-detectability}  Math's detectability envelope (for crosscheck)
% ---------------------------------------------------------------------

\subsection{Below the power floor: what the framework cannot yet certify}
\label{sec:below-floor-epistemic}

A framework's sensitivity is itself a piece of information about the framework, not about the biology. The minimax envelope of \Cref{thm:minimax} and the hierarchical-detectability result of \Cref{thm:S1-hierarchical} together partition the rare-population landscape at Kang~2024 scale into three strata: \emph{above-floor-flat} populations for which whole-atlas REAL is sufficient; \emph{hierarchical-recovered} populations which cross the detectability floor only after sub-compartment telescoping (per \Cref{sec:S1-envelope}); and \emph{truly-unresolved} populations for which neither flat nor hierarchical REAL attains the required power at Kang depth. The first two strata carry the manuscript's headline claim; this subsection states what we can and cannot say about the third. To read an absence of \textsc{REAL} flags on a truly-unresolved population as evidence of preservation is to commit exactly the failure-to-reject-as-success error that this paper is designed to expose (\Cref{philo:limitations}; Supp Note~S0 trap~8 at \Cref{sec:eightTraps}).

Two distinctions matter in making the disclosure precise. First, \textit{truly-unresolved at this dataset} is not the same as \textit{unresolved in principle}: the only escape routes are pooled multi-atlas scale (on the order of $10^{8}$ cells), auxiliary modalities that lift $\Delta$ above the $\pi \Delta^{2} \geq 1.5 \times 10^{-4}$ threshold at Kang scale (protein panels, TCR / BCR sequencing, spatial context), or targeted protocol changes that enrich the population before capture (FACS sorting, modality-aware library preparation). \textit{Truly unresolved} is thus a relation between a population, a measurement regime, and a framework, not a property of the population alone. Second, a motif flagged by \textsc{REAL} remains, even here, a claim about \emph{evaluation}, not about ontology (\Cref{philo:episto_note}). Correspondingly, the absence of a flag in the truly-unresolved stratum remains a claim about \emph{measurement insufficiency}, not about biology. The populations listed below are not excluded from single-cell biology; they are excluded from this framework's current certificate of preservation, under the stated regime.

%   % ---------------------------------------------------------------------
% Limitations subsection — immunology contribution to the Discussion.
% ~350 words. Honest disclosure of populations where REAL/RareShield
% cannot make reliable detection claims at Kang-atlas scale.
% Author: Immunology (agent-mozus3vv). Cross-references Math's
% predicted-detectability table (agent/mathematics@f564cc6).
% Target: sections/08_discussion.tex, as a \subsection after the
% "Implications for immunology" subsection (immuno_discussion.tex).
% ---------------------------------------------------------------------

\subsection{Boundaries of what our framework can currently guarantee}
\label{sec:discussion-limitations-immunology}

Three concrete limitations shape the claims we make. First, the Le~Cam
detectability envelope (\Cref{sec:minimax-detectability}) places a
sharp boundary on which rare populations can be recovered label-free
at any given atlas scale. At the Kang pan-cancer immune atlas
($n = 4.9 \times 10^6$, $B = 30$, per-batch heterogeneity
$\Lambda = 3$), populations with prevalence-separation product
$\pi\Delta < 10^{-3}\sigma$ fall below the detection floor, producing
a true-positive rate under 0.3 at $\alpha = 0.05$. Specifically,
\textit{CXCR5}$^+$\textit{TCF7}$^+$ progenitor-exhausted CD8 T
(TPEX), LAMP3$^+$ mregDC, KIR$^+$ adaptive NK, endothelial tip
cells, and drug-tolerant-persister tumour cells all sit in this
regime at Kang alone. For these populations our framework does not
claim detection from the Kang atlas in isolation; their recovery
requires either auxiliary protein or TCR signal, or pooling across
multiple atlases beyond $10^8$ cells. We disclose this explicitly
and present the populations instead as below-floor targets whose
detection-regime analysis (Extended Data~Fig.~1) itself is a
contribution.

Second, some rare populations are lost before integration rather than
during integration. Neutrophil subsets including LOX-1$^+$
PMN-MDSC and TAN-1 through TAN-5 subsets are systematically
undersampled by the standard 10x Genomics droplet-capture protocol,
and mast cells are similarly depleted relative to their tissue
abundance. Our framework is a post-integration detector and does not
and cannot recover populations that never entered the cell-by-gene
matrix. We flag this as a pre-integration identifiability problem
orthogonal to overcorrection, to be addressed by protocol choice
and by modality-aware capture strategies (FACS-sort enrichment prior
to library preparation; split-pool-barcoded whole-transcriptome
methods that capture neutrophils without selection bias).

Third, our evaluation uses annotated rare populations with their
labels hidden at test time. This is a legitimate form of ground
truth but it is not a direct test on unannotated populations.
Novel populations detected under our framework in the pan-cancer
atlas --- for example, a candidate CXCR5-dim TPEX sub-state
\citep{ourfig5} --- require the full evidence package of
\Cref{sec:evidence-package} plus independent cohort replication
before a definite claim of previously-unannotated biology. Main-text
detections are reported only when $\geq 4$ of the six evidence
criteria are satisfied.

\paragraph*{Telescoping detection recovers below-whole-atlas populations at sub-compartment scale.}
The $\pi\cdot\Delta^2 < 1.5\times10^{-4}$ below-floor rule above is
stated at the whole-atlas level and applies to flat REAL. Hierarchical
REAL, applied at lineage-restricted sub-compartments, moves the
operating point from whole-atlas prevalence to compartment-level
prevalence; at each level the flat floor still applies, but the
$\pi\cdot\Delta^2$ product is typically $10$--$100\times$ larger
because rare populations are often concentrated in specific lineages.
This recovers pulmonary ionocytes at the airway-epithelial sub-compartment
level, and the CXCR5-dim TPEX sub-state at the exhausted-CD8 sub-compartment
level (Supplementary Theorem~S1-hierarchical, Mathematics agent
\citealp{math2026hierarchical}). Our disclosure is therefore
three-tier: populations recovered by flat whole-atlas REAL (above
floor at $\pi\cdot\Delta^2 > 1.5\times10^{-4}$, e.g., MAIT, $\gamma\delta$
T, TREM2$^+$ LAM, myCAF, EMT-hybrid), populations recovered by
hierarchical sub-compartment REAL (above floor at sub-compartment
product, e.g., ionocytes, TPEX sub-state, AS-DC), and populations
unresolved at all compartment depths (endothelial tip cells, drug-tolerant
persister baseline state). Main-text claims are restricted to the
first two tiers; tier-three populations are disclosed as present in
the dataset but beyond the current framework's detection capability.

% ---------------------------------------------------------------------
% Requires: sec:minimax-detectability, sec:evidence-package,
% \citep{ourfig5} is a placeholder for the Fig 5 companion section /
% preprint. Biological Science may prefer \Cref{fig:fig5} or
% \Cref{sec:results-tpex-substate} depending on skeleton.
% ---------------------------------------------------------------------
         <-- unchanged from Immunology branch
%   % ---------------------------------------------------------------------
% Joint closing paragraph for the below-floor Limitations subsection.
% Co-authored with Philosophy (agent-mozur8ub). Placement flow:
%
%   \input{philo_below_floor_head.tex}    (Philosophy's epistemic frame)
%   \input{immuno_limitations.tex}        (Immunology's biological disclosure)
%   \input{immuno_below_floor_close.tex}  (THIS FILE — joint affirmative)
%
% Closes the subsection on the telescoping-detection affirmative:
% most below-flat-floor populations are recovered under hierarchical
% sub-compartment REAL.
% ---------------------------------------------------------------------

\paragraph*{Recovery under hierarchical application.}
Notwithstanding the strict boundary established above, the \emph{truly-unresolved} stratum is surprisingly small. Of the 22 immune and tumor populations in our registry whose whole-atlas product $\pi \cdot \Delta^{2}$ places them into one of the three detection tiers, 20 are recoverable through flat or hierarchical application of \textsc{REAL} at an appropriate compartment scale (\Cref{thm:S1-hierarchical}; \Cref{sec:methods-oracle-score} protocol): 7 populations are detected at whole-atlas scale (flat tier --- MAIT, $\gamma\delta$ T V$\delta1$-bias, TREM2$^{+}$ LAM, myCAF, iCAF, plasma-TLS, EMT-hybrid-near-pole), and 13 additional populations cross the detection floor only under hierarchical application at their natural sub-compartment scale (hierarchical tier). Pulmonary ionocytes pass detection at the airway-epithelial sub-compartment of HLCA (effective $\pi \approx 5 \times 10^{-3}$ in $n \approx 240{,}000$ airway cells); \textit{CXCR5}$^{+}$\textit{TCF7}$^{+}$ progenitor-exhausted CD8 T cells pass detection at the exhausted-CD8 sub-compartment of pooled melanoma atlases ($\pi \approx 10^{-2}$ in $n \approx 100{,}000$ cells); LAMP3$^{+}$ mregDC, AS-DC, and HCMV-imprinted adaptive NK pass detection at dendritic-cell and natural-killer sub-compartments respectively. The two populations that remain below the floor at every compartment depth are endothelial tip cells (which sit at the intersection of a small-mass vascular niche and a narrow transcriptional separation) and the drug-tolerant persister baseline state in oncology (Tumor Biology registry row TU-T-001), with mast cells in tumors additionally disclosed as out-of-scope for post-integration detection (mode~4 pre-integration capture loss). These unresolved entries are explicitly deferred to follow-up work: their recovery would require either the multi-atlas pooling regime ($n > 10^{8}$ cells) or auxiliary-modality augmentation, both of which lie beyond this paper's scope. This three-tier disclosure --- recovered-flat, recovered-hierarchical, deferred --- is the most honest reading of what our framework currently certifies, and it preserves the paper's headline claim on the above-floor and hierarchical-recovered strata while placing no claim on the truly-unresolved tier.

% ---------------------------------------------------------------------
% Cross-references needed (shared with philo_below_floor_head + immuno_limitations):
%   thm:S1-hierarchical        Math hierarchical REAL theorem
%   sec:methods-oracle-score   Methods — oracle-score verification
% Bib entries: none new.
% ---------------------------------------------------------------------
   <-- unchanged from Immunology branch
%
% Use this file INSTEAD OF the 5-file chain if BioSci prefers a single splice point.
% Immunology's two \input calls (limitations + below-floor-close) are PRESERVED as
% \input calls inside this file rather than inlined, because they are owned by
% the immunology branch and Immunology should retain edit control over their content.
%
% Placement: sections/08_discussion.tex, replacing the 5-file \input chain with
% a single % Philosophy §08 epistemic-scope block — single-file consolidation.
% Merges the five \input calls that BioSci currently threads through into one file
% for easier future-edit management. Semantically identical to the input chain:
%   \input{handoffs/philo_discussion_meno.tex}
%   \input{handoffs/philo_limitations.tex}
%   \input{handoffs/philo_below_floor_head.tex}
%   \input{handoffs/immuno_limitations.tex}         <-- unchanged from Immunology branch
%   \input{handoffs/immuno_below_floor_close.tex}   <-- unchanged from Immunology branch
%
% Use this file INSTEAD OF the 5-file chain if BioSci prefers a single splice point.
% Immunology's two \input calls (limitations + below-floor-close) are PRESERVED as
% \input calls inside this file rather than inlined, because they are owned by
% the immunology branch and Immunology should retain edit control over their content.
%
% Placement: sections/08_discussion.tex, replacing the 5-file \input chain with
% a single \input{handoffs/philo_08_epistemic_block.tex}.
%
% Author: Philosophy (agent-mozur8ub). Last updated for v2.4 pre-submission.

% =============================================================================
% §08 EPISTEMIC SCOPE BLOCK — single-file consolidation
% =============================================================================

% -----------------------------------------------------------------------------
% (1) Meno-problem framing — opens the epistemic scope block.
% Content from philo_discussion_meno.tex (agent/philosophy@<head>).
% ≈250w. Fixes the is/ought gap: "reproducible structure destroyed" ≠ "novel cell type".
% -----------------------------------------------------------------------------

\paragraph*{What a REAL flag is, and is not.}
A motif flagged by REAL is the claim that a local, reproducible geometric structure---witnessed across independent batches before integration, with measurable local-mass distortion after---has been erased by the integration map beyond what admissible batch-effect deformations can explain. This is a claim about \emph{evaluation}, not about ontology. It is not a claim that a new cell type exists. A motif $w \in \mathcal{W}$ with low $\mathrm{LRS}(w)$ licenses exactly one inference: the integration has discarded reproducible structure of the kind a biologist would want to interrogate. Whether that structure corresponds to a novel cell type, a transitional state along an already-known trajectory, a rare disease-specific configuration, or a technical sub-mode that merely looks biological, is a further question REAL does not answer. Upgrading a motif to a cell-type claim requires the usual burden of single-cell biology: independent sampling, orthogonal markers, targeted protocol (sorting, spatial validation, experimental perturbation), and community replication. REAL's value is that it identifies \emph{where to look}, not what has been found. Conflating the two would recreate the very error this paper is designed to expose: treating a framework-internal signal as a fact about the world. Our strong claim is therefore deliberately asymmetric. A high motif set $|\mathcal{W}|$ under a candidate integration is evidence of measurement insufficiency; a low motif set is evidence that the integration respected the syntactic reproducibility criterion. Ontology is downstream of both.

% -----------------------------------------------------------------------------
% (2) Limitations and epistemic scope — Route-C exchangeability boundaries.
% Content from philo_limitations.tex. ≤500w.
% -----------------------------------------------------------------------------

\subsection*{Limitations and epistemic scope}
\label{philo:limitations}
REAL relies on an exchangeability assumption between the known rare populations we hold out for calibration (pancreatic $\epsilon$ cells, pulmonary ionocytes, LAMP3$^+$ mregDC, TPEX, FOLR2$^+$ perivascular macrophages, apCAF, and the populations enumerated in Table~S\ref{tab:heldout}) and the genuinely unannotated rare subpopulations to which we transfer the estimate. The assumption is formally stated in Lemma~\ref{lem:exchangeability} as a \emph{bound}: the held-out-rare loss rate upper-bounds the expected loss rate for unannotated rare subpopulations of comparable abundance, local geometry, and batch profile. We emphasise four boundary conditions under which the bound is not tight. First, the bound degrades when the unannotated population lies outside the span of the highly-variable genes used to construct the latent space; REAL cannot protect what the representation cannot see. Second, the bound degrades when the unannotated population has a qualitatively different batch-effect signature from the calibration populations---for instance, a disease-specific state observed in only one batch. Third, the bound is conservative for populations whose local density profile is not well-approximated by the level-set stratification we use; transitional cells along a continuum can appear as low-density but are not rare in the sampling sense. Fourth, our use of known rare populations as proxies for unknown ones inherits a selection bias: the populations we can hold out are those that have already been discovered, and the discovered set may not be exchangeable with the undiscovered set.

Our strongest defensible claim is therefore deliberately limited. REAL provides (i) a label-free, invariant-based detectability framework calibrated on held-out known rare subpopulations, (ii) an integration strategy (RareShield) that demonstrably reduces the held-out-rare loss rate at pan-cancer-atlas scale without degrading batch correction for abundant cell types, and (iii) an explicit exchangeability bound under which the held-out loss rate is an upper bound on the loss rate of truly unannotated rare subpopulations. A REAL-flagged motif $w \in \mathcal{W}$ is a claim that reproducible structure has been destroyed, not a claim that a new cell type exists; upgrading to an ontological claim requires the usual experimental burden (sorting, orthogonal markers, spatial validation, replication). The complementary negative: a dataset with low $|\mathcal{W}|$ under a candidate integration is evidence that no cross-batch-reproducible rare structure was erased, but not a guarantee that no such structure existed---by construction, we cannot rule out subpopulations whose reproducibility signal was below our sampling resolution or outside the evaluated feature span. These boundaries are not weaknesses of REAL specifically; they are the price of reasoning about what the evaluation framework cannot directly observe. Naming the price is what distinguishes an epistemically calibrated claim from a framework-captured one.

% -----------------------------------------------------------------------------
% (3) Below the power floor — epistemic frame head (three-tier stratification).
% Content from philo_below_floor_head.tex.
% -----------------------------------------------------------------------------

\subsection*{Below the power floor: what the framework cannot yet certify}
\label{sec:below-floor-epistemic}

A framework's sensitivity is itself a piece of information about the framework, not about the biology. The minimax envelope of \Cref{thm:minimax} and the hierarchical-detectability result of \Cref{thm:S1-hierarchical} together partition the rare-population landscape at Kang~2024 scale into three strata: \emph{above-floor-flat} populations for which whole-atlas REAL is sufficient; \emph{hierarchical-recovered} populations which cross the detectability floor only after sub-compartment telescoping (per \Cref{sec:S1-envelope}); and \emph{truly-unresolved} populations for which neither flat nor hierarchical REAL attains the required power at Kang depth. The first two strata carry the manuscript's headline claim; this subsection states what we can and cannot say about the third. To read an absence of \textsc{REAL} flags on a truly-unresolved population as evidence of preservation is to commit exactly the failure-to-reject-as-success error that this paper is designed to expose (\Cref{philo:limitations}; Supp Note~S0 trap~8 at \Cref{sec:eightTraps}).

Two distinctions matter in making the disclosure precise. First, \textit{truly-unresolved at this dataset} is not the same as \textit{unresolved in principle}: the only escape routes are pooled multi-atlas scale (on the order of $10^{8}$ cells), auxiliary modalities that lift $\Delta$ above the $\pi \Delta^{2} \geq 1.5 \times 10^{-4}$ threshold at Kang scale (protein panels, TCR / BCR sequencing, spatial context), or targeted protocol changes that enrich the population before capture (FACS sorting, modality-aware library preparation). \textit{Truly unresolved} is thus a relation between a population, a measurement regime, and a framework, not a property of the population alone. Second, a motif flagged by \textsc{REAL} remains, even here, a claim about \emph{evaluation}, not about ontology (\Cref{philo:episto_note}). Correspondingly, the absence of a flag in the truly-unresolved stratum remains a claim about \emph{measurement insufficiency}, not about biology. The populations listed below are not excluded from single-cell biology; they are excluded from this framework's current certificate of preservation, under the stated regime.

% -----------------------------------------------------------------------------
% (4)-(5) Immunology's biological disclosure + joint closing.
% These files live in the immunology branch; Immunology retains edit control.
% They are \input'd here rather than inlined, so Immunology-branch edits propagate
% automatically on the next BioSci pull.
% -----------------------------------------------------------------------------

\input{handoffs/immuno_limitations.tex}
\input{handoffs/immuno_below_floor_close.tex}

% =============================================================================
% END §08 EPISTEMIC SCOPE BLOCK
% =============================================================================

% Integration note for BioSci:
%   - To use this consolidated file: replace the 5-file \input chain in
%     sections/08_discussion.tex with:
%       \input{handoffs/philo_08_epistemic_block.tex}
%   - To keep the 5-file chain: ignore this file; it is semantically identical
%     to the chain at this commit. Consolidation is a convenience, not a correctness fix.
%   - If Philosophy or Immunology updates any of the five source files, re-generating
%     this consolidated file requires Philosophy to re-splice; the \input calls for
%     Immunology's two files remain untouched.
%
% Lint status: passes lint:epistemic clean (trap 8 / trap 4 / IM.4 guards all present).
.
%
% Author: Philosophy (agent-mozur8ub). Last updated for v2.4 pre-submission.

% =============================================================================
% §08 EPISTEMIC SCOPE BLOCK — single-file consolidation
% =============================================================================

% -----------------------------------------------------------------------------
% (1) Meno-problem framing — opens the epistemic scope block.
% Content from philo_discussion_meno.tex (agent/philosophy@<head>).
% ≈250w. Fixes the is/ought gap: "reproducible structure destroyed" ≠ "novel cell type".
% -----------------------------------------------------------------------------

\paragraph*{What a REAL flag is, and is not.}
A motif flagged by REAL is the claim that a local, reproducible geometric structure---witnessed across independent batches before integration, with measurable local-mass distortion after---has been erased by the integration map beyond what admissible batch-effect deformations can explain. This is a claim about \emph{evaluation}, not about ontology. It is not a claim that a new cell type exists. A motif $w \in \mathcal{W}$ with low $\mathrm{LRS}(w)$ licenses exactly one inference: the integration has discarded reproducible structure of the kind a biologist would want to interrogate. Whether that structure corresponds to a novel cell type, a transitional state along an already-known trajectory, a rare disease-specific configuration, or a technical sub-mode that merely looks biological, is a further question REAL does not answer. Upgrading a motif to a cell-type claim requires the usual burden of single-cell biology: independent sampling, orthogonal markers, targeted protocol (sorting, spatial validation, experimental perturbation), and community replication. REAL's value is that it identifies \emph{where to look}, not what has been found. Conflating the two would recreate the very error this paper is designed to expose: treating a framework-internal signal as a fact about the world. Our strong claim is therefore deliberately asymmetric. A high motif set $|\mathcal{W}|$ under a candidate integration is evidence of measurement insufficiency; a low motif set is evidence that the integration respected the syntactic reproducibility criterion. Ontology is downstream of both.

% -----------------------------------------------------------------------------
% (2) Limitations and epistemic scope — Route-C exchangeability boundaries.
% Content from philo_limitations.tex. ≤500w.
% -----------------------------------------------------------------------------

\subsection*{Limitations and epistemic scope}
\label{philo:limitations}
REAL relies on an exchangeability assumption between the known rare populations we hold out for calibration (pancreatic $\epsilon$ cells, pulmonary ionocytes, LAMP3$^+$ mregDC, TPEX, FOLR2$^+$ perivascular macrophages, apCAF, and the populations enumerated in Table~S\ref{tab:heldout}) and the genuinely unannotated rare subpopulations to which we transfer the estimate. The assumption is formally stated in Lemma~\ref{lem:exchangeability} as a \emph{bound}: the held-out-rare loss rate upper-bounds the expected loss rate for unannotated rare subpopulations of comparable abundance, local geometry, and batch profile. We emphasise four boundary conditions under which the bound is not tight. First, the bound degrades when the unannotated population lies outside the span of the highly-variable genes used to construct the latent space; REAL cannot protect what the representation cannot see. Second, the bound degrades when the unannotated population has a qualitatively different batch-effect signature from the calibration populations---for instance, a disease-specific state observed in only one batch. Third, the bound is conservative for populations whose local density profile is not well-approximated by the level-set stratification we use; transitional cells along a continuum can appear as low-density but are not rare in the sampling sense. Fourth, our use of known rare populations as proxies for unknown ones inherits a selection bias: the populations we can hold out are those that have already been discovered, and the discovered set may not be exchangeable with the undiscovered set.

Our strongest defensible claim is therefore deliberately limited. REAL provides (i) a label-free, invariant-based detectability framework calibrated on held-out known rare subpopulations, (ii) an integration strategy (RareShield) that demonstrably reduces the held-out-rare loss rate at pan-cancer-atlas scale without degrading batch correction for abundant cell types, and (iii) an explicit exchangeability bound under which the held-out loss rate is an upper bound on the loss rate of truly unannotated rare subpopulations. A REAL-flagged motif $w \in \mathcal{W}$ is a claim that reproducible structure has been destroyed, not a claim that a new cell type exists; upgrading to an ontological claim requires the usual experimental burden (sorting, orthogonal markers, spatial validation, replication). The complementary negative: a dataset with low $|\mathcal{W}|$ under a candidate integration is evidence that no cross-batch-reproducible rare structure was erased, but not a guarantee that no such structure existed---by construction, we cannot rule out subpopulations whose reproducibility signal was below our sampling resolution or outside the evaluated feature span. These boundaries are not weaknesses of REAL specifically; they are the price of reasoning about what the evaluation framework cannot directly observe. Naming the price is what distinguishes an epistemically calibrated claim from a framework-captured one.

% -----------------------------------------------------------------------------
% (3) Below the power floor — epistemic frame head (three-tier stratification).
% Content from philo_below_floor_head.tex.
% -----------------------------------------------------------------------------

\subsection*{Below the power floor: what the framework cannot yet certify}
\label{sec:below-floor-epistemic}

A framework's sensitivity is itself a piece of information about the framework, not about the biology. The minimax envelope of \Cref{thm:minimax} and the hierarchical-detectability result of \Cref{thm:S1-hierarchical} together partition the rare-population landscape at Kang~2024 scale into three strata: \emph{above-floor-flat} populations for which whole-atlas REAL is sufficient; \emph{hierarchical-recovered} populations which cross the detectability floor only after sub-compartment telescoping (per \Cref{sec:S1-envelope}); and \emph{truly-unresolved} populations for which neither flat nor hierarchical REAL attains the required power at Kang depth. The first two strata carry the manuscript's headline claim; this subsection states what we can and cannot say about the third. To read an absence of \textsc{REAL} flags on a truly-unresolved population as evidence of preservation is to commit exactly the failure-to-reject-as-success error that this paper is designed to expose (\Cref{philo:limitations}; Supp Note~S0 trap~8 at \Cref{sec:eightTraps}).

Two distinctions matter in making the disclosure precise. First, \textit{truly-unresolved at this dataset} is not the same as \textit{unresolved in principle}: the only escape routes are pooled multi-atlas scale (on the order of $10^{8}$ cells), auxiliary modalities that lift $\Delta$ above the $\pi \Delta^{2} \geq 1.5 \times 10^{-4}$ threshold at Kang scale (protein panels, TCR / BCR sequencing, spatial context), or targeted protocol changes that enrich the population before capture (FACS sorting, modality-aware library preparation). \textit{Truly unresolved} is thus a relation between a population, a measurement regime, and a framework, not a property of the population alone. Second, a motif flagged by \textsc{REAL} remains, even here, a claim about \emph{evaluation}, not about ontology (\Cref{philo:episto_note}). Correspondingly, the absence of a flag in the truly-unresolved stratum remains a claim about \emph{measurement insufficiency}, not about biology. The populations listed below are not excluded from single-cell biology; they are excluded from this framework's current certificate of preservation, under the stated regime.

% -----------------------------------------------------------------------------
% (4)-(5) Immunology's biological disclosure + joint closing.
% These files live in the immunology branch; Immunology retains edit control.
% They are \input'd here rather than inlined, so Immunology-branch edits propagate
% automatically on the next BioSci pull.
% -----------------------------------------------------------------------------

% ---------------------------------------------------------------------
% Limitations subsection — immunology contribution to the Discussion.
% ~350 words. Honest disclosure of populations where REAL/RareShield
% cannot make reliable detection claims at Kang-atlas scale.
% Author: Immunology (agent-mozus3vv). Cross-references Math's
% predicted-detectability table (agent/mathematics@f564cc6).
% Target: sections/08_discussion.tex, as a \subsection after the
% "Implications for immunology" subsection (immuno_discussion.tex).
% ---------------------------------------------------------------------

\subsection{Boundaries of what our framework can currently guarantee}
\label{sec:discussion-limitations-immunology}

Three concrete limitations shape the claims we make. First, the Le~Cam
detectability envelope (\Cref{sec:minimax-detectability}) places a
sharp boundary on which rare populations can be recovered label-free
at any given atlas scale. At the Kang pan-cancer immune atlas
($n = 4.9 \times 10^6$, $B = 30$, per-batch heterogeneity
$\Lambda = 3$), populations with prevalence-separation product
$\pi\Delta < 10^{-3}\sigma$ fall below the detection floor, producing
a true-positive rate under 0.3 at $\alpha = 0.05$. Specifically,
\textit{CXCR5}$^+$\textit{TCF7}$^+$ progenitor-exhausted CD8 T
(TPEX), LAMP3$^+$ mregDC, KIR$^+$ adaptive NK, endothelial tip
cells, and drug-tolerant-persister tumour cells all sit in this
regime at Kang alone. For these populations our framework does not
claim detection from the Kang atlas in isolation; their recovery
requires either auxiliary protein or TCR signal, or pooling across
multiple atlases beyond $10^8$ cells. We disclose this explicitly
and present the populations instead as below-floor targets whose
detection-regime analysis (Extended Data~Fig.~1) itself is a
contribution.

Second, some rare populations are lost before integration rather than
during integration. Neutrophil subsets including LOX-1$^+$
PMN-MDSC and TAN-1 through TAN-5 subsets are systematically
undersampled by the standard 10x Genomics droplet-capture protocol,
and mast cells are similarly depleted relative to their tissue
abundance. Our framework is a post-integration detector and does not
and cannot recover populations that never entered the cell-by-gene
matrix. We flag this as a pre-integration identifiability problem
orthogonal to overcorrection, to be addressed by protocol choice
and by modality-aware capture strategies (FACS-sort enrichment prior
to library preparation; split-pool-barcoded whole-transcriptome
methods that capture neutrophils without selection bias).

Third, our evaluation uses annotated rare populations with their
labels hidden at test time. This is a legitimate form of ground
truth but it is not a direct test on unannotated populations.
Novel populations detected under our framework in the pan-cancer
atlas --- for example, a candidate CXCR5-dim TPEX sub-state
\citep{ourfig5} --- require the full evidence package of
\Cref{sec:evidence-package} plus independent cohort replication
before a definite claim of previously-unannotated biology. Main-text
detections are reported only when $\geq 4$ of the six evidence
criteria are satisfied.

\paragraph*{Telescoping detection recovers below-whole-atlas populations at sub-compartment scale.}
The $\pi\cdot\Delta^2 < 1.5\times10^{-4}$ below-floor rule above is
stated at the whole-atlas level and applies to flat REAL. Hierarchical
REAL, applied at lineage-restricted sub-compartments, moves the
operating point from whole-atlas prevalence to compartment-level
prevalence; at each level the flat floor still applies, but the
$\pi\cdot\Delta^2$ product is typically $10$--$100\times$ larger
because rare populations are often concentrated in specific lineages.
This recovers pulmonary ionocytes at the airway-epithelial sub-compartment
level, and the CXCR5-dim TPEX sub-state at the exhausted-CD8 sub-compartment
level (Supplementary Theorem~S1-hierarchical, Mathematics agent
\citealp{math2026hierarchical}). Our disclosure is therefore
three-tier: populations recovered by flat whole-atlas REAL (above
floor at $\pi\cdot\Delta^2 > 1.5\times10^{-4}$, e.g., MAIT, $\gamma\delta$
T, TREM2$^+$ LAM, myCAF, EMT-hybrid), populations recovered by
hierarchical sub-compartment REAL (above floor at sub-compartment
product, e.g., ionocytes, TPEX sub-state, AS-DC), and populations
unresolved at all compartment depths (endothelial tip cells, drug-tolerant
persister baseline state). Main-text claims are restricted to the
first two tiers; tier-three populations are disclosed as present in
the dataset but beyond the current framework's detection capability.

% ---------------------------------------------------------------------
% Requires: sec:minimax-detectability, sec:evidence-package,
% \citep{ourfig5} is a placeholder for the Fig 5 companion section /
% preprint. Biological Science may prefer \Cref{fig:fig5} or
% \Cref{sec:results-tpex-substate} depending on skeleton.
% ---------------------------------------------------------------------

% ---------------------------------------------------------------------
% Joint closing paragraph for the below-floor Limitations subsection.
% Co-authored with Philosophy (agent-mozur8ub). Placement flow:
%
%   \input{philo_below_floor_head.tex}    (Philosophy's epistemic frame)
%   \input{immuno_limitations.tex}        (Immunology's biological disclosure)
%   \input{immuno_below_floor_close.tex}  (THIS FILE — joint affirmative)
%
% Closes the subsection on the telescoping-detection affirmative:
% most below-flat-floor populations are recovered under hierarchical
% sub-compartment REAL.
% ---------------------------------------------------------------------

\paragraph*{Recovery under hierarchical application.}
Notwithstanding the strict boundary established above, the \emph{truly-unresolved} stratum is surprisingly small. Of the 22 immune and tumor populations in our registry whose whole-atlas product $\pi \cdot \Delta^{2}$ places them into one of the three detection tiers, 20 are recoverable through flat or hierarchical application of \textsc{REAL} at an appropriate compartment scale (\Cref{thm:S1-hierarchical}; \Cref{sec:methods-oracle-score} protocol): 7 populations are detected at whole-atlas scale (flat tier --- MAIT, $\gamma\delta$ T V$\delta1$-bias, TREM2$^{+}$ LAM, myCAF, iCAF, plasma-TLS, EMT-hybrid-near-pole), and 13 additional populations cross the detection floor only under hierarchical application at their natural sub-compartment scale (hierarchical tier). Pulmonary ionocytes pass detection at the airway-epithelial sub-compartment of HLCA (effective $\pi \approx 5 \times 10^{-3}$ in $n \approx 240{,}000$ airway cells); \textit{CXCR5}$^{+}$\textit{TCF7}$^{+}$ progenitor-exhausted CD8 T cells pass detection at the exhausted-CD8 sub-compartment of pooled melanoma atlases ($\pi \approx 10^{-2}$ in $n \approx 100{,}000$ cells); LAMP3$^{+}$ mregDC, AS-DC, and HCMV-imprinted adaptive NK pass detection at dendritic-cell and natural-killer sub-compartments respectively. The two populations that remain below the floor at every compartment depth are endothelial tip cells (which sit at the intersection of a small-mass vascular niche and a narrow transcriptional separation) and the drug-tolerant persister baseline state in oncology (Tumor Biology registry row TU-T-001), with mast cells in tumors additionally disclosed as out-of-scope for post-integration detection (mode~4 pre-integration capture loss). These unresolved entries are explicitly deferred to follow-up work: their recovery would require either the multi-atlas pooling regime ($n > 10^{8}$ cells) or auxiliary-modality augmentation, both of which lie beyond this paper's scope. This three-tier disclosure --- recovered-flat, recovered-hierarchical, deferred --- is the most honest reading of what our framework currently certifies, and it preserves the paper's headline claim on the above-floor and hierarchical-recovered strata while placing no claim on the truly-unresolved tier.

% ---------------------------------------------------------------------
% Cross-references needed (shared with philo_below_floor_head + immuno_limitations):
%   thm:S1-hierarchical        Math hierarchical REAL theorem
%   sec:methods-oracle-score   Methods — oracle-score verification
% Bib entries: none new.
% ---------------------------------------------------------------------


% =============================================================================
% END §08 EPISTEMIC SCOPE BLOCK
% =============================================================================

% Integration note for BioSci:
%   - To use this consolidated file: replace the 5-file \input chain in
%     sections/08_discussion.tex with:
%       % Philosophy §08 epistemic-scope block — single-file consolidation.
% Merges the five \input calls that BioSci currently threads through into one file
% for easier future-edit management. Semantically identical to the input chain:
%   \input{handoffs/philo_discussion_meno.tex}
%   \input{handoffs/philo_limitations.tex}
%   \input{handoffs/philo_below_floor_head.tex}
%   \input{handoffs/immuno_limitations.tex}         <-- unchanged from Immunology branch
%   \input{handoffs/immuno_below_floor_close.tex}   <-- unchanged from Immunology branch
%
% Use this file INSTEAD OF the 5-file chain if BioSci prefers a single splice point.
% Immunology's two \input calls (limitations + below-floor-close) are PRESERVED as
% \input calls inside this file rather than inlined, because they are owned by
% the immunology branch and Immunology should retain edit control over their content.
%
% Placement: sections/08_discussion.tex, replacing the 5-file \input chain with
% a single \input{handoffs/philo_08_epistemic_block.tex}.
%
% Author: Philosophy (agent-mozur8ub). Last updated for v2.4 pre-submission.

% =============================================================================
% §08 EPISTEMIC SCOPE BLOCK — single-file consolidation
% =============================================================================

% -----------------------------------------------------------------------------
% (1) Meno-problem framing — opens the epistemic scope block.
% Content from philo_discussion_meno.tex (agent/philosophy@<head>).
% ≈250w. Fixes the is/ought gap: "reproducible structure destroyed" ≠ "novel cell type".
% -----------------------------------------------------------------------------

\paragraph*{What a REAL flag is, and is not.}
A motif flagged by REAL is the claim that a local, reproducible geometric structure---witnessed across independent batches before integration, with measurable local-mass distortion after---has been erased by the integration map beyond what admissible batch-effect deformations can explain. This is a claim about \emph{evaluation}, not about ontology. It is not a claim that a new cell type exists. A motif $w \in \mathcal{W}$ with low $\mathrm{LRS}(w)$ licenses exactly one inference: the integration has discarded reproducible structure of the kind a biologist would want to interrogate. Whether that structure corresponds to a novel cell type, a transitional state along an already-known trajectory, a rare disease-specific configuration, or a technical sub-mode that merely looks biological, is a further question REAL does not answer. Upgrading a motif to a cell-type claim requires the usual burden of single-cell biology: independent sampling, orthogonal markers, targeted protocol (sorting, spatial validation, experimental perturbation), and community replication. REAL's value is that it identifies \emph{where to look}, not what has been found. Conflating the two would recreate the very error this paper is designed to expose: treating a framework-internal signal as a fact about the world. Our strong claim is therefore deliberately asymmetric. A high motif set $|\mathcal{W}|$ under a candidate integration is evidence of measurement insufficiency; a low motif set is evidence that the integration respected the syntactic reproducibility criterion. Ontology is downstream of both.

% -----------------------------------------------------------------------------
% (2) Limitations and epistemic scope — Route-C exchangeability boundaries.
% Content from philo_limitations.tex. ≤500w.
% -----------------------------------------------------------------------------

\subsection*{Limitations and epistemic scope}
\label{philo:limitations}
REAL relies on an exchangeability assumption between the known rare populations we hold out for calibration (pancreatic $\epsilon$ cells, pulmonary ionocytes, LAMP3$^+$ mregDC, TPEX, FOLR2$^+$ perivascular macrophages, apCAF, and the populations enumerated in Table~S\ref{tab:heldout}) and the genuinely unannotated rare subpopulations to which we transfer the estimate. The assumption is formally stated in Lemma~\ref{lem:exchangeability} as a \emph{bound}: the held-out-rare loss rate upper-bounds the expected loss rate for unannotated rare subpopulations of comparable abundance, local geometry, and batch profile. We emphasise four boundary conditions under which the bound is not tight. First, the bound degrades when the unannotated population lies outside the span of the highly-variable genes used to construct the latent space; REAL cannot protect what the representation cannot see. Second, the bound degrades when the unannotated population has a qualitatively different batch-effect signature from the calibration populations---for instance, a disease-specific state observed in only one batch. Third, the bound is conservative for populations whose local density profile is not well-approximated by the level-set stratification we use; transitional cells along a continuum can appear as low-density but are not rare in the sampling sense. Fourth, our use of known rare populations as proxies for unknown ones inherits a selection bias: the populations we can hold out are those that have already been discovered, and the discovered set may not be exchangeable with the undiscovered set.

Our strongest defensible claim is therefore deliberately limited. REAL provides (i) a label-free, invariant-based detectability framework calibrated on held-out known rare subpopulations, (ii) an integration strategy (RareShield) that demonstrably reduces the held-out-rare loss rate at pan-cancer-atlas scale without degrading batch correction for abundant cell types, and (iii) an explicit exchangeability bound under which the held-out loss rate is an upper bound on the loss rate of truly unannotated rare subpopulations. A REAL-flagged motif $w \in \mathcal{W}$ is a claim that reproducible structure has been destroyed, not a claim that a new cell type exists; upgrading to an ontological claim requires the usual experimental burden (sorting, orthogonal markers, spatial validation, replication). The complementary negative: a dataset with low $|\mathcal{W}|$ under a candidate integration is evidence that no cross-batch-reproducible rare structure was erased, but not a guarantee that no such structure existed---by construction, we cannot rule out subpopulations whose reproducibility signal was below our sampling resolution or outside the evaluated feature span. These boundaries are not weaknesses of REAL specifically; they are the price of reasoning about what the evaluation framework cannot directly observe. Naming the price is what distinguishes an epistemically calibrated claim from a framework-captured one.

% -----------------------------------------------------------------------------
% (3) Below the power floor — epistemic frame head (three-tier stratification).
% Content from philo_below_floor_head.tex.
% -----------------------------------------------------------------------------

\subsection*{Below the power floor: what the framework cannot yet certify}
\label{sec:below-floor-epistemic}

A framework's sensitivity is itself a piece of information about the framework, not about the biology. The minimax envelope of \Cref{thm:minimax} and the hierarchical-detectability result of \Cref{thm:S1-hierarchical} together partition the rare-population landscape at Kang~2024 scale into three strata: \emph{above-floor-flat} populations for which whole-atlas REAL is sufficient; \emph{hierarchical-recovered} populations which cross the detectability floor only after sub-compartment telescoping (per \Cref{sec:S1-envelope}); and \emph{truly-unresolved} populations for which neither flat nor hierarchical REAL attains the required power at Kang depth. The first two strata carry the manuscript's headline claim; this subsection states what we can and cannot say about the third. To read an absence of \textsc{REAL} flags on a truly-unresolved population as evidence of preservation is to commit exactly the failure-to-reject-as-success error that this paper is designed to expose (\Cref{philo:limitations}; Supp Note~S0 trap~8 at \Cref{sec:eightTraps}).

Two distinctions matter in making the disclosure precise. First, \textit{truly-unresolved at this dataset} is not the same as \textit{unresolved in principle}: the only escape routes are pooled multi-atlas scale (on the order of $10^{8}$ cells), auxiliary modalities that lift $\Delta$ above the $\pi \Delta^{2} \geq 1.5 \times 10^{-4}$ threshold at Kang scale (protein panels, TCR / BCR sequencing, spatial context), or targeted protocol changes that enrich the population before capture (FACS sorting, modality-aware library preparation). \textit{Truly unresolved} is thus a relation between a population, a measurement regime, and a framework, not a property of the population alone. Second, a motif flagged by \textsc{REAL} remains, even here, a claim about \emph{evaluation}, not about ontology (\Cref{philo:episto_note}). Correspondingly, the absence of a flag in the truly-unresolved stratum remains a claim about \emph{measurement insufficiency}, not about biology. The populations listed below are not excluded from single-cell biology; they are excluded from this framework's current certificate of preservation, under the stated regime.

% -----------------------------------------------------------------------------
% (4)-(5) Immunology's biological disclosure + joint closing.
% These files live in the immunology branch; Immunology retains edit control.
% They are \input'd here rather than inlined, so Immunology-branch edits propagate
% automatically on the next BioSci pull.
% -----------------------------------------------------------------------------

\input{handoffs/immuno_limitations.tex}
\input{handoffs/immuno_below_floor_close.tex}

% =============================================================================
% END §08 EPISTEMIC SCOPE BLOCK
% =============================================================================

% Integration note for BioSci:
%   - To use this consolidated file: replace the 5-file \input chain in
%     sections/08_discussion.tex with:
%       \input{handoffs/philo_08_epistemic_block.tex}
%   - To keep the 5-file chain: ignore this file; it is semantically identical
%     to the chain at this commit. Consolidation is a convenience, not a correctness fix.
%   - If Philosophy or Immunology updates any of the five source files, re-generating
%     this consolidated file requires Philosophy to re-splice; the \input calls for
%     Immunology's two files remain untouched.
%
% Lint status: passes lint:epistemic clean (trap 8 / trap 4 / IM.4 guards all present).

%   - To keep the 5-file chain: ignore this file; it is semantically identical
%     to the chain at this commit. Consolidation is a convenience, not a correctness fix.
%   - If Philosophy or Immunology updates any of the five source files, re-generating
%     this consolidated file requires Philosophy to re-splice; the \input calls for
%     Immunology's two files remain untouched.
%
% Lint status: passes lint:epistemic clean (trap 8 / trap 4 / IM.4 guards all present).
.
%
% Author: Philosophy (agent-mozur8ub). Last updated for v2.4 pre-submission.

% =============================================================================
% §08 EPISTEMIC SCOPE BLOCK — single-file consolidation
% =============================================================================

% -----------------------------------------------------------------------------
% (1) Meno-problem framing — opens the epistemic scope block.
% Content from philo_discussion_meno.tex (agent/philosophy@<head>).
% ≈250w. Fixes the is/ought gap: "reproducible structure destroyed" ≠ "novel cell type".
% -----------------------------------------------------------------------------

\paragraph*{What a REAL flag is, and is not.}
A motif flagged by REAL is the claim that a local, reproducible geometric structure---witnessed across independent batches before integration, with measurable local-mass distortion after---has been erased by the integration map beyond what admissible batch-effect deformations can explain. This is a claim about \emph{evaluation}, not about ontology. It is not a claim that a new cell type exists. A motif $w \in \mathcal{W}$ with low $\mathrm{LRS}(w)$ licenses exactly one inference: the integration has discarded reproducible structure of the kind a biologist would want to interrogate. Whether that structure corresponds to a novel cell type, a transitional state along an already-known trajectory, a rare disease-specific configuration, or a technical sub-mode that merely looks biological, is a further question REAL does not answer. Upgrading a motif to a cell-type claim requires the usual burden of single-cell biology: independent sampling, orthogonal markers, targeted protocol (sorting, spatial validation, experimental perturbation), and community replication. REAL's value is that it identifies \emph{where to look}, not what has been found. Conflating the two would recreate the very error this paper is designed to expose: treating a framework-internal signal as a fact about the world. Our strong claim is therefore deliberately asymmetric. A high motif set $|\mathcal{W}|$ under a candidate integration is evidence of measurement insufficiency; a low motif set is evidence that the integration respected the syntactic reproducibility criterion. Ontology is downstream of both.

% -----------------------------------------------------------------------------
% (2) Limitations and epistemic scope — Route-C exchangeability boundaries.
% Content from philo_limitations.tex. ≤500w.
% -----------------------------------------------------------------------------

\subsection*{Limitations and epistemic scope}
\label{philo:limitations}
REAL relies on an exchangeability assumption between the known rare populations we hold out for calibration (pancreatic $\epsilon$ cells, pulmonary ionocytes, LAMP3$^+$ mregDC, TPEX, FOLR2$^+$ perivascular macrophages, apCAF, and the populations enumerated in Table~S\ref{tab:heldout}) and the genuinely unannotated rare subpopulations to which we transfer the estimate. The assumption is formally stated in Lemma~\ref{lem:exchangeability} as a \emph{bound}: the held-out-rare loss rate upper-bounds the expected loss rate for unannotated rare subpopulations of comparable abundance, local geometry, and batch profile. We emphasise four boundary conditions under which the bound is not tight. First, the bound degrades when the unannotated population lies outside the span of the highly-variable genes used to construct the latent space; REAL cannot protect what the representation cannot see. Second, the bound degrades when the unannotated population has a qualitatively different batch-effect signature from the calibration populations---for instance, a disease-specific state observed in only one batch. Third, the bound is conservative for populations whose local density profile is not well-approximated by the level-set stratification we use; transitional cells along a continuum can appear as low-density but are not rare in the sampling sense. Fourth, our use of known rare populations as proxies for unknown ones inherits a selection bias: the populations we can hold out are those that have already been discovered, and the discovered set may not be exchangeable with the undiscovered set.

Our strongest defensible claim is therefore deliberately limited. REAL provides (i) a label-free, invariant-based detectability framework calibrated on held-out known rare subpopulations, (ii) an integration strategy (RareShield) that demonstrably reduces the held-out-rare loss rate at pan-cancer-atlas scale without degrading batch correction for abundant cell types, and (iii) an explicit exchangeability bound under which the held-out loss rate is an upper bound on the loss rate of truly unannotated rare subpopulations. A REAL-flagged motif $w \in \mathcal{W}$ is a claim that reproducible structure has been destroyed, not a claim that a new cell type exists; upgrading to an ontological claim requires the usual experimental burden (sorting, orthogonal markers, spatial validation, replication). The complementary negative: a dataset with low $|\mathcal{W}|$ under a candidate integration is evidence that no cross-batch-reproducible rare structure was erased, but not a guarantee that no such structure existed---by construction, we cannot rule out subpopulations whose reproducibility signal was below our sampling resolution or outside the evaluated feature span. These boundaries are not weaknesses of REAL specifically; they are the price of reasoning about what the evaluation framework cannot directly observe. Naming the price is what distinguishes an epistemically calibrated claim from a framework-captured one.

% -----------------------------------------------------------------------------
% (3) Below the power floor — epistemic frame head (three-tier stratification).
% Content from philo_below_floor_head.tex.
% -----------------------------------------------------------------------------

\subsection*{Below the power floor: what the framework cannot yet certify}
\label{sec:below-floor-epistemic}

A framework's sensitivity is itself a piece of information about the framework, not about the biology. The minimax envelope of \Cref{thm:minimax} and the hierarchical-detectability result of \Cref{thm:S1-hierarchical} together partition the rare-population landscape at Kang~2024 scale into three strata: \emph{above-floor-flat} populations for which whole-atlas REAL is sufficient; \emph{hierarchical-recovered} populations which cross the detectability floor only after sub-compartment telescoping (per \Cref{sec:S1-envelope}); and \emph{truly-unresolved} populations for which neither flat nor hierarchical REAL attains the required power at Kang depth. The first two strata carry the manuscript's headline claim; this subsection states what we can and cannot say about the third. To read an absence of \textsc{REAL} flags on a truly-unresolved population as evidence of preservation is to commit exactly the failure-to-reject-as-success error that this paper is designed to expose (\Cref{philo:limitations}; Supp Note~S0 trap~8 at \Cref{sec:eightTraps}).

Two distinctions matter in making the disclosure precise. First, \textit{truly-unresolved at this dataset} is not the same as \textit{unresolved in principle}: the only escape routes are pooled multi-atlas scale (on the order of $10^{8}$ cells), auxiliary modalities that lift $\Delta$ above the $\pi \Delta^{2} \geq 1.5 \times 10^{-4}$ threshold at Kang scale (protein panels, TCR / BCR sequencing, spatial context), or targeted protocol changes that enrich the population before capture (FACS sorting, modality-aware library preparation). \textit{Truly unresolved} is thus a relation between a population, a measurement regime, and a framework, not a property of the population alone. Second, a motif flagged by \textsc{REAL} remains, even here, a claim about \emph{evaluation}, not about ontology (\Cref{philo:episto_note}). Correspondingly, the absence of a flag in the truly-unresolved stratum remains a claim about \emph{measurement insufficiency}, not about biology. The populations listed below are not excluded from single-cell biology; they are excluded from this framework's current certificate of preservation, under the stated regime.

% -----------------------------------------------------------------------------
% (4)-(5) Immunology's biological disclosure + joint closing.
% These files live in the immunology branch; Immunology retains edit control.
% They are \input'd here rather than inlined, so Immunology-branch edits propagate
% automatically on the next BioSci pull.
% -----------------------------------------------------------------------------

% ---------------------------------------------------------------------
% Limitations subsection — immunology contribution to the Discussion.
% ~350 words. Honest disclosure of populations where REAL/RareShield
% cannot make reliable detection claims at Kang-atlas scale.
% Author: Immunology (agent-mozus3vv). Cross-references Math's
% predicted-detectability table (agent/mathematics@f564cc6).
% Target: sections/08_discussion.tex, as a \subsection after the
% "Implications for immunology" subsection (immuno_discussion.tex).
% ---------------------------------------------------------------------

\subsection{Boundaries of what our framework can currently guarantee}
\label{sec:discussion-limitations-immunology}

Three concrete limitations shape the claims we make. First, the Le~Cam
detectability envelope (\Cref{sec:minimax-detectability}) places a
sharp boundary on which rare populations can be recovered label-free
at any given atlas scale. At the Kang pan-cancer immune atlas
($n = 4.9 \times 10^6$, $B = 30$, per-batch heterogeneity
$\Lambda = 3$), populations with prevalence-separation product
$\pi\Delta < 10^{-3}\sigma$ fall below the detection floor, producing
a true-positive rate under 0.3 at $\alpha = 0.05$. Specifically,
\textit{CXCR5}$^+$\textit{TCF7}$^+$ progenitor-exhausted CD8 T
(TPEX), LAMP3$^+$ mregDC, KIR$^+$ adaptive NK, endothelial tip
cells, and drug-tolerant-persister tumour cells all sit in this
regime at Kang alone. For these populations our framework does not
claim detection from the Kang atlas in isolation; their recovery
requires either auxiliary protein or TCR signal, or pooling across
multiple atlases beyond $10^8$ cells. We disclose this explicitly
and present the populations instead as below-floor targets whose
detection-regime analysis (Extended Data~Fig.~1) itself is a
contribution.

Second, some rare populations are lost before integration rather than
during integration. Neutrophil subsets including LOX-1$^+$
PMN-MDSC and TAN-1 through TAN-5 subsets are systematically
undersampled by the standard 10x Genomics droplet-capture protocol,
and mast cells are similarly depleted relative to their tissue
abundance. Our framework is a post-integration detector and does not
and cannot recover populations that never entered the cell-by-gene
matrix. We flag this as a pre-integration identifiability problem
orthogonal to overcorrection, to be addressed by protocol choice
and by modality-aware capture strategies (FACS-sort enrichment prior
to library preparation; split-pool-barcoded whole-transcriptome
methods that capture neutrophils without selection bias).

Third, our evaluation uses annotated rare populations with their
labels hidden at test time. This is a legitimate form of ground
truth but it is not a direct test on unannotated populations.
Novel populations detected under our framework in the pan-cancer
atlas --- for example, a candidate CXCR5-dim TPEX sub-state
\citep{ourfig5} --- require the full evidence package of
\Cref{sec:evidence-package} plus independent cohort replication
before a definite claim of previously-unannotated biology. Main-text
detections are reported only when $\geq 4$ of the six evidence
criteria are satisfied.

\paragraph*{Telescoping detection recovers below-whole-atlas populations at sub-compartment scale.}
The $\pi\cdot\Delta^2 < 1.5\times10^{-4}$ below-floor rule above is
stated at the whole-atlas level and applies to flat REAL. Hierarchical
REAL, applied at lineage-restricted sub-compartments, moves the
operating point from whole-atlas prevalence to compartment-level
prevalence; at each level the flat floor still applies, but the
$\pi\cdot\Delta^2$ product is typically $10$--$100\times$ larger
because rare populations are often concentrated in specific lineages.
This recovers pulmonary ionocytes at the airway-epithelial sub-compartment
level, and the CXCR5-dim TPEX sub-state at the exhausted-CD8 sub-compartment
level (Supplementary Theorem~S1-hierarchical, Mathematics agent
\citealp{math2026hierarchical}). Our disclosure is therefore
three-tier: populations recovered by flat whole-atlas REAL (above
floor at $\pi\cdot\Delta^2 > 1.5\times10^{-4}$, e.g., MAIT, $\gamma\delta$
T, TREM2$^+$ LAM, myCAF, EMT-hybrid), populations recovered by
hierarchical sub-compartment REAL (above floor at sub-compartment
product, e.g., ionocytes, TPEX sub-state, AS-DC), and populations
unresolved at all compartment depths (endothelial tip cells, drug-tolerant
persister baseline state). Main-text claims are restricted to the
first two tiers; tier-three populations are disclosed as present in
the dataset but beyond the current framework's detection capability.

% ---------------------------------------------------------------------
% Requires: sec:minimax-detectability, sec:evidence-package,
% \citep{ourfig5} is a placeholder for the Fig 5 companion section /
% preprint. Biological Science may prefer \Cref{fig:fig5} or
% \Cref{sec:results-tpex-substate} depending on skeleton.
% ---------------------------------------------------------------------

% ---------------------------------------------------------------------
% Joint closing paragraph for the below-floor Limitations subsection.
% Co-authored with Philosophy (agent-mozur8ub). Placement flow:
%
%   % ---------------------------------------------------------------------
% Philosophy's epistemic-frame head for the below-floor Limitations subsection.
% To be placed IMMEDIATELY BEFORE the Immunology-drafted biological disclosure
% (workspace/handoffs/immuno_limitations.tex at agent/immunology@b8a898d).
%
% Co-authored structure agreed with Immunology (agent-mozus3vv, 2026-05-10):
%   1. Philosophy epistemic frame (this file) — 2 short paragraphs, ≈220w.
%   2. Immunology biological disclosure (immuno_limitations.tex) — ≈350w.
%   3. Joint closing (telescoping-detection affirmative) — Immunology to draft
%      after pulling this file and Math's sub_compartment_pooling_bounds.md.
%
% Placement: sections/08_discussion.tex, as the subsection \subsection{Below the power floor},
% ordered after \input{philo_limitations.tex} and before \input{immuno_limitations.tex}.
%
% Anchors:
%   \Cref{thm:minimax}         Math Theorem 1 (minimax lower bound; floor)
%   \Cref{thm:labelblind}      Math Theorem 2.1 (label-based σ-algebra closure)
%   \Cref{lem:exchangeability} Math Lemma S1 (exchangeability bound)
%   \Cref{philo:limitations}   Philosophy Limitations subsection (head of §08 epistemic scope)
%   \Cref{sec:minimax-detectability}  Math's detectability envelope (for crosscheck)
% ---------------------------------------------------------------------

\subsection{Below the power floor: what the framework cannot yet certify}
\label{sec:below-floor-epistemic}

A framework's sensitivity is itself a piece of information about the framework, not about the biology. The minimax envelope of \Cref{thm:minimax} and the hierarchical-detectability result of \Cref{thm:S1-hierarchical} together partition the rare-population landscape at Kang~2024 scale into three strata: \emph{above-floor-flat} populations for which whole-atlas REAL is sufficient; \emph{hierarchical-recovered} populations which cross the detectability floor only after sub-compartment telescoping (per \Cref{sec:S1-envelope}); and \emph{truly-unresolved} populations for which neither flat nor hierarchical REAL attains the required power at Kang depth. The first two strata carry the manuscript's headline claim; this subsection states what we can and cannot say about the third. To read an absence of \textsc{REAL} flags on a truly-unresolved population as evidence of preservation is to commit exactly the failure-to-reject-as-success error that this paper is designed to expose (\Cref{philo:limitations}; Supp Note~S0 trap~8 at \Cref{sec:eightTraps}).

Two distinctions matter in making the disclosure precise. First, \textit{truly-unresolved at this dataset} is not the same as \textit{unresolved in principle}: the only escape routes are pooled multi-atlas scale (on the order of $10^{8}$ cells), auxiliary modalities that lift $\Delta$ above the $\pi \Delta^{2} \geq 1.5 \times 10^{-4}$ threshold at Kang scale (protein panels, TCR / BCR sequencing, spatial context), or targeted protocol changes that enrich the population before capture (FACS sorting, modality-aware library preparation). \textit{Truly unresolved} is thus a relation between a population, a measurement regime, and a framework, not a property of the population alone. Second, a motif flagged by \textsc{REAL} remains, even here, a claim about \emph{evaluation}, not about ontology (\Cref{philo:episto_note}). Correspondingly, the absence of a flag in the truly-unresolved stratum remains a claim about \emph{measurement insufficiency}, not about biology. The populations listed below are not excluded from single-cell biology; they are excluded from this framework's current certificate of preservation, under the stated regime.
    (Philosophy's epistemic frame)
%   % ---------------------------------------------------------------------
% Limitations subsection — immunology contribution to the Discussion.
% ~350 words. Honest disclosure of populations where REAL/RareShield
% cannot make reliable detection claims at Kang-atlas scale.
% Author: Immunology (agent-mozus3vv). Cross-references Math's
% predicted-detectability table (agent/mathematics@f564cc6).
% Target: sections/08_discussion.tex, as a \subsection after the
% "Implications for immunology" subsection (immuno_discussion.tex).
% ---------------------------------------------------------------------

\subsection{Boundaries of what our framework can currently guarantee}
\label{sec:discussion-limitations-immunology}

Three concrete limitations shape the claims we make. First, the Le~Cam
detectability envelope (\Cref{sec:minimax-detectability}) places a
sharp boundary on which rare populations can be recovered label-free
at any given atlas scale. At the Kang pan-cancer immune atlas
($n = 4.9 \times 10^6$, $B = 30$, per-batch heterogeneity
$\Lambda = 3$), populations with prevalence-separation product
$\pi\Delta < 10^{-3}\sigma$ fall below the detection floor, producing
a true-positive rate under 0.3 at $\alpha = 0.05$. Specifically,
\textit{CXCR5}$^+$\textit{TCF7}$^+$ progenitor-exhausted CD8 T
(TPEX), LAMP3$^+$ mregDC, KIR$^+$ adaptive NK, endothelial tip
cells, and drug-tolerant-persister tumour cells all sit in this
regime at Kang alone. For these populations our framework does not
claim detection from the Kang atlas in isolation; their recovery
requires either auxiliary protein or TCR signal, or pooling across
multiple atlases beyond $10^8$ cells. We disclose this explicitly
and present the populations instead as below-floor targets whose
detection-regime analysis (Extended Data~Fig.~1) itself is a
contribution.

Second, some rare populations are lost before integration rather than
during integration. Neutrophil subsets including LOX-1$^+$
PMN-MDSC and TAN-1 through TAN-5 subsets are systematically
undersampled by the standard 10x Genomics droplet-capture protocol,
and mast cells are similarly depleted relative to their tissue
abundance. Our framework is a post-integration detector and does not
and cannot recover populations that never entered the cell-by-gene
matrix. We flag this as a pre-integration identifiability problem
orthogonal to overcorrection, to be addressed by protocol choice
and by modality-aware capture strategies (FACS-sort enrichment prior
to library preparation; split-pool-barcoded whole-transcriptome
methods that capture neutrophils without selection bias).

Third, our evaluation uses annotated rare populations with their
labels hidden at test time. This is a legitimate form of ground
truth but it is not a direct test on unannotated populations.
Novel populations detected under our framework in the pan-cancer
atlas --- for example, a candidate CXCR5-dim TPEX sub-state
\citep{ourfig5} --- require the full evidence package of
\Cref{sec:evidence-package} plus independent cohort replication
before a definite claim of previously-unannotated biology. Main-text
detections are reported only when $\geq 4$ of the six evidence
criteria are satisfied.

\paragraph*{Telescoping detection recovers below-whole-atlas populations at sub-compartment scale.}
The $\pi\cdot\Delta^2 < 1.5\times10^{-4}$ below-floor rule above is
stated at the whole-atlas level and applies to flat REAL. Hierarchical
REAL, applied at lineage-restricted sub-compartments, moves the
operating point from whole-atlas prevalence to compartment-level
prevalence; at each level the flat floor still applies, but the
$\pi\cdot\Delta^2$ product is typically $10$--$100\times$ larger
because rare populations are often concentrated in specific lineages.
This recovers pulmonary ionocytes at the airway-epithelial sub-compartment
level, and the CXCR5-dim TPEX sub-state at the exhausted-CD8 sub-compartment
level (Supplementary Theorem~S1-hierarchical, Mathematics agent
\citealp{math2026hierarchical}). Our disclosure is therefore
three-tier: populations recovered by flat whole-atlas REAL (above
floor at $\pi\cdot\Delta^2 > 1.5\times10^{-4}$, e.g., MAIT, $\gamma\delta$
T, TREM2$^+$ LAM, myCAF, EMT-hybrid), populations recovered by
hierarchical sub-compartment REAL (above floor at sub-compartment
product, e.g., ionocytes, TPEX sub-state, AS-DC), and populations
unresolved at all compartment depths (endothelial tip cells, drug-tolerant
persister baseline state). Main-text claims are restricted to the
first two tiers; tier-three populations are disclosed as present in
the dataset but beyond the current framework's detection capability.

% ---------------------------------------------------------------------
% Requires: sec:minimax-detectability, sec:evidence-package,
% \citep{ourfig5} is a placeholder for the Fig 5 companion section /
% preprint. Biological Science may prefer \Cref{fig:fig5} or
% \Cref{sec:results-tpex-substate} depending on skeleton.
% ---------------------------------------------------------------------
        (Immunology's biological disclosure)
%   % ---------------------------------------------------------------------
% Joint closing paragraph for the below-floor Limitations subsection.
% Co-authored with Philosophy (agent-mozur8ub). Placement flow:
%
%   \input{philo_below_floor_head.tex}    (Philosophy's epistemic frame)
%   \input{immuno_limitations.tex}        (Immunology's biological disclosure)
%   \input{immuno_below_floor_close.tex}  (THIS FILE — joint affirmative)
%
% Closes the subsection on the telescoping-detection affirmative:
% most below-flat-floor populations are recovered under hierarchical
% sub-compartment REAL.
% ---------------------------------------------------------------------

\paragraph*{Recovery under hierarchical application.}
Notwithstanding the strict boundary established above, the \emph{truly-unresolved} stratum is surprisingly small. Of the 22 immune and tumor populations in our registry whose whole-atlas product $\pi \cdot \Delta^{2}$ places them into one of the three detection tiers, 20 are recoverable through flat or hierarchical application of \textsc{REAL} at an appropriate compartment scale (\Cref{thm:S1-hierarchical}; \Cref{sec:methods-oracle-score} protocol): 7 populations are detected at whole-atlas scale (flat tier --- MAIT, $\gamma\delta$ T V$\delta1$-bias, TREM2$^{+}$ LAM, myCAF, iCAF, plasma-TLS, EMT-hybrid-near-pole), and 13 additional populations cross the detection floor only under hierarchical application at their natural sub-compartment scale (hierarchical tier). Pulmonary ionocytes pass detection at the airway-epithelial sub-compartment of HLCA (effective $\pi \approx 5 \times 10^{-3}$ in $n \approx 240{,}000$ airway cells); \textit{CXCR5}$^{+}$\textit{TCF7}$^{+}$ progenitor-exhausted CD8 T cells pass detection at the exhausted-CD8 sub-compartment of pooled melanoma atlases ($\pi \approx 10^{-2}$ in $n \approx 100{,}000$ cells); LAMP3$^{+}$ mregDC, AS-DC, and HCMV-imprinted adaptive NK pass detection at dendritic-cell and natural-killer sub-compartments respectively. The two populations that remain below the floor at every compartment depth are endothelial tip cells (which sit at the intersection of a small-mass vascular niche and a narrow transcriptional separation) and the drug-tolerant persister baseline state in oncology (Tumor Biology registry row TU-T-001), with mast cells in tumors additionally disclosed as out-of-scope for post-integration detection (mode~4 pre-integration capture loss). These unresolved entries are explicitly deferred to follow-up work: their recovery would require either the multi-atlas pooling regime ($n > 10^{8}$ cells) or auxiliary-modality augmentation, both of which lie beyond this paper's scope. This three-tier disclosure --- recovered-flat, recovered-hierarchical, deferred --- is the most honest reading of what our framework currently certifies, and it preserves the paper's headline claim on the above-floor and hierarchical-recovered strata while placing no claim on the truly-unresolved tier.

% ---------------------------------------------------------------------
% Cross-references needed (shared with philo_below_floor_head + immuno_limitations):
%   thm:S1-hierarchical        Math hierarchical REAL theorem
%   sec:methods-oracle-score   Methods — oracle-score verification
% Bib entries: none new.
% ---------------------------------------------------------------------
  (THIS FILE — joint affirmative)
%
% Closes the subsection on the telescoping-detection affirmative:
% most below-flat-floor populations are recovered under hierarchical
% sub-compartment REAL.
% ---------------------------------------------------------------------

\paragraph*{Recovery under hierarchical application.}
Notwithstanding the strict boundary established above, the \emph{truly-unresolved} stratum is surprisingly small. Of the 22 immune and tumor populations in our registry whose whole-atlas product $\pi \cdot \Delta^{2}$ places them into one of the three detection tiers, 20 are recoverable through flat or hierarchical application of \textsc{REAL} at an appropriate compartment scale (\Cref{thm:S1-hierarchical}; \Cref{sec:methods-oracle-score} protocol): 7 populations are detected at whole-atlas scale (flat tier --- MAIT, $\gamma\delta$ T V$\delta1$-bias, TREM2$^{+}$ LAM, myCAF, iCAF, plasma-TLS, EMT-hybrid-near-pole), and 13 additional populations cross the detection floor only under hierarchical application at their natural sub-compartment scale (hierarchical tier). Pulmonary ionocytes pass detection at the airway-epithelial sub-compartment of HLCA (effective $\pi \approx 5 \times 10^{-3}$ in $n \approx 240{,}000$ airway cells); \textit{CXCR5}$^{+}$\textit{TCF7}$^{+}$ progenitor-exhausted CD8 T cells pass detection at the exhausted-CD8 sub-compartment of pooled melanoma atlases ($\pi \approx 10^{-2}$ in $n \approx 100{,}000$ cells); LAMP3$^{+}$ mregDC, AS-DC, and HCMV-imprinted adaptive NK pass detection at dendritic-cell and natural-killer sub-compartments respectively. The two populations that remain below the floor at every compartment depth are endothelial tip cells (which sit at the intersection of a small-mass vascular niche and a narrow transcriptional separation) and the drug-tolerant persister baseline state in oncology (Tumor Biology registry row TU-T-001), with mast cells in tumors additionally disclosed as out-of-scope for post-integration detection (mode~4 pre-integration capture loss). These unresolved entries are explicitly deferred to follow-up work: their recovery would require either the multi-atlas pooling regime ($n > 10^{8}$ cells) or auxiliary-modality augmentation, both of which lie beyond this paper's scope. This three-tier disclosure --- recovered-flat, recovered-hierarchical, deferred --- is the most honest reading of what our framework currently certifies, and it preserves the paper's headline claim on the above-floor and hierarchical-recovered strata while placing no claim on the truly-unresolved tier.

% ---------------------------------------------------------------------
% Cross-references needed (shared with philo_below_floor_head + immuno_limitations):
%   thm:S1-hierarchical        Math hierarchical REAL theorem
%   sec:methods-oracle-score   Methods — oracle-score verification
% Bib entries: none new.
% ---------------------------------------------------------------------


% =============================================================================
% END §08 EPISTEMIC SCOPE BLOCK
% =============================================================================

% Integration note for BioSci:
%   - To use this consolidated file: replace the 5-file \input chain in
%     sections/08_discussion.tex with:
%       % Philosophy §08 epistemic-scope block — single-file consolidation.
% Merges the five \input calls that BioSci currently threads through into one file
% for easier future-edit management. Semantically identical to the input chain:
%   % Discussion paragraph (≈250w) — epistemic status of REAL-flagged motifs.
% For task t-mozv8oul. Fixes the is/ought gap: "reproducible structure destroyed" ≠ "novel cell type".
% Guards the team from over-claiming. Anchors: motif $w \in \mathcal{W}$; LRS(w); RareShield; Hacking-style intervention realism.

\paragraph*{What a REAL flag is, and is not.}
A motif flagged by REAL is the claim that a local, reproducible geometric structure---witnessed across independent batches before integration, with measurable local-mass distortion after---has been erased by the integration map beyond what admissible batch-effect deformations can explain. This is a claim about \emph{evaluation}, not about ontology. It is not a claim that a new cell type exists. A motif $w \in \mathcal{W}$ with low $\mathrm{LRS}(w)$ licenses exactly one inference: the integration has discarded reproducible structure of the kind a biologist would want to interrogate. Whether that structure corresponds to a novel cell type, a transitional state along an already-known trajectory, a rare disease-specific configuration, or a technical sub-mode that merely looks biological, is a further question that REAL does not answer. Upgrading a motif to a cell-type claim requires the usual burden of single-cell biology: independent sampling, orthogonal markers, targeted protocol (sorting, spatial validation, experimental perturbation), and community replication. REAL's value is that it identifies \emph{where to look}, not what has been found. Conflating the two would re-create the very error this paper is designed to expose: treating a framework-internal signal as a fact about the world. Our strong claim is therefore deliberately asymmetric: a high motif set $|\mathcal{W}|$ under a candidate integration is evidence of measurement insufficiency; a low motif set is evidence that the integration respected the syntactic reproducibility criterion. Ontology is downstream of both.

%   % Limitations & Epistemic Scope subsection (≤500w).
% For BioSci's locked slot in the Discussion. Cites Math Lemma S1-exchangeability as a *bound*.
% Covers: exchangeability failure modes, reference-class scope, non-ontic status, circularity mitigations.

\subsection*{Limitations and epistemic scope}
REAL relies on an exchangeability assumption between the known rare populations we hold out for calibration (pancreatic $\epsilon$ cells, pulmonary ionocytes, LAMP3$^+$ mregDC, TPEX, FOLR2$^+$ perivascular macrophages and the other populations enumerated in Table~S\ref{tab:heldout}) and the genuinely unannotated rare subpopulations to which we transfer the estimate. The assumption is formally stated in Lemma~\ref{lem:exchangeability} as a \emph{bound}: the held-out-rare loss rate upper-bounds the expected loss rate for unannotated rare subpopulations of comparable abundance, local geometry, and batch profile. We emphasise four boundary conditions under which the bound is not tight. First, the bound degrades when the unannotated population lies outside the span of the highly-variable genes used to construct the latent space; REAL cannot protect what the representation cannot see. Second, the bound degrades when the unannotated population has a qualitatively different batch-effect signature from the calibration populations---for instance, a disease-specific state observed in only one batch. Third, the bound is conservative for populations whose local density profile is not well-approximated by the level-set stratification we use; transitional cells along a continuum can appear as low-density but are not rare in the sampling sense. Fourth, our use of known rare populations as proxies for unknown ones inherits a selection bias: the populations we can hold out are those that have already been discovered, and the discovered set may not be exchangeable with the undiscovered set.

Our strongest defensible claim is therefore deliberately limited. REAL provides (i) a label-free, invariant-based detectability framework calibrated on held-out known rare subpopulations, (ii) an integration strategy (RareShield) that demonstrably reduces the held-out-rare loss rate at pan-cancer-atlas scale without degrading batch correction for abundant cell types, and (iii) an explicit exchangeability bound under which the held-out loss rate is an upper bound on the loss rate of truly unannotated rare subpopulations. A REAL-flagged motif $w \in \mathcal{W}$ is a claim that reproducible structure has been destroyed, not a claim that a new cell type exists; upgrading to an ontological claim requires the usual experimental burden (sorting, orthogonal markers, spatial validation, replication). The complementary negative: a dataset with low $|\mathcal{W}|$ under a candidate integration is evidence that no cross-batch-reproducible rare structure was erased, but not a guarantee that no such structure existed---by construction, we cannot rule out subpopulations whose reproducibility signal was below our sampling resolution or outside the evaluated feature span. These boundaries are not weaknesses of REAL specifically; they are the price of reasoning about what the evaluation framework cannot directly observe. Naming the price is what distinguishes an epistemically calibrated claim from a framework-captured one.

%   % ---------------------------------------------------------------------
% Philosophy's epistemic-frame head for the below-floor Limitations subsection.
% To be placed IMMEDIATELY BEFORE the Immunology-drafted biological disclosure
% (workspace/handoffs/immuno_limitations.tex at agent/immunology@b8a898d).
%
% Co-authored structure agreed with Immunology (agent-mozus3vv, 2026-05-10):
%   1. Philosophy epistemic frame (this file) — 2 short paragraphs, ≈220w.
%   2. Immunology biological disclosure (immuno_limitations.tex) — ≈350w.
%   3. Joint closing (telescoping-detection affirmative) — Immunology to draft
%      after pulling this file and Math's sub_compartment_pooling_bounds.md.
%
% Placement: sections/08_discussion.tex, as the subsection \subsection{Below the power floor},
% ordered after \input{philo_limitations.tex} and before \input{immuno_limitations.tex}.
%
% Anchors:
%   \Cref{thm:minimax}         Math Theorem 1 (minimax lower bound; floor)
%   \Cref{thm:labelblind}      Math Theorem 2.1 (label-based σ-algebra closure)
%   \Cref{lem:exchangeability} Math Lemma S1 (exchangeability bound)
%   \Cref{philo:limitations}   Philosophy Limitations subsection (head of §08 epistemic scope)
%   \Cref{sec:minimax-detectability}  Math's detectability envelope (for crosscheck)
% ---------------------------------------------------------------------

\subsection{Below the power floor: what the framework cannot yet certify}
\label{sec:below-floor-epistemic}

A framework's sensitivity is itself a piece of information about the framework, not about the biology. The minimax envelope of \Cref{thm:minimax} and the hierarchical-detectability result of \Cref{thm:S1-hierarchical} together partition the rare-population landscape at Kang~2024 scale into three strata: \emph{above-floor-flat} populations for which whole-atlas REAL is sufficient; \emph{hierarchical-recovered} populations which cross the detectability floor only after sub-compartment telescoping (per \Cref{sec:S1-envelope}); and \emph{truly-unresolved} populations for which neither flat nor hierarchical REAL attains the required power at Kang depth. The first two strata carry the manuscript's headline claim; this subsection states what we can and cannot say about the third. To read an absence of \textsc{REAL} flags on a truly-unresolved population as evidence of preservation is to commit exactly the failure-to-reject-as-success error that this paper is designed to expose (\Cref{philo:limitations}; Supp Note~S0 trap~8 at \Cref{sec:eightTraps}).

Two distinctions matter in making the disclosure precise. First, \textit{truly-unresolved at this dataset} is not the same as \textit{unresolved in principle}: the only escape routes are pooled multi-atlas scale (on the order of $10^{8}$ cells), auxiliary modalities that lift $\Delta$ above the $\pi \Delta^{2} \geq 1.5 \times 10^{-4}$ threshold at Kang scale (protein panels, TCR / BCR sequencing, spatial context), or targeted protocol changes that enrich the population before capture (FACS sorting, modality-aware library preparation). \textit{Truly unresolved} is thus a relation between a population, a measurement regime, and a framework, not a property of the population alone. Second, a motif flagged by \textsc{REAL} remains, even here, a claim about \emph{evaluation}, not about ontology (\Cref{philo:episto_note}). Correspondingly, the absence of a flag in the truly-unresolved stratum remains a claim about \emph{measurement insufficiency}, not about biology. The populations listed below are not excluded from single-cell biology; they are excluded from this framework's current certificate of preservation, under the stated regime.

%   % ---------------------------------------------------------------------
% Limitations subsection — immunology contribution to the Discussion.
% ~350 words. Honest disclosure of populations where REAL/RareShield
% cannot make reliable detection claims at Kang-atlas scale.
% Author: Immunology (agent-mozus3vv). Cross-references Math's
% predicted-detectability table (agent/mathematics@f564cc6).
% Target: sections/08_discussion.tex, as a \subsection after the
% "Implications for immunology" subsection (immuno_discussion.tex).
% ---------------------------------------------------------------------

\subsection{Boundaries of what our framework can currently guarantee}
\label{sec:discussion-limitations-immunology}

Three concrete limitations shape the claims we make. First, the Le~Cam
detectability envelope (\Cref{sec:minimax-detectability}) places a
sharp boundary on which rare populations can be recovered label-free
at any given atlas scale. At the Kang pan-cancer immune atlas
($n = 4.9 \times 10^6$, $B = 30$, per-batch heterogeneity
$\Lambda = 3$), populations with prevalence-separation product
$\pi\Delta < 10^{-3}\sigma$ fall below the detection floor, producing
a true-positive rate under 0.3 at $\alpha = 0.05$. Specifically,
\textit{CXCR5}$^+$\textit{TCF7}$^+$ progenitor-exhausted CD8 T
(TPEX), LAMP3$^+$ mregDC, KIR$^+$ adaptive NK, endothelial tip
cells, and drug-tolerant-persister tumour cells all sit in this
regime at Kang alone. For these populations our framework does not
claim detection from the Kang atlas in isolation; their recovery
requires either auxiliary protein or TCR signal, or pooling across
multiple atlases beyond $10^8$ cells. We disclose this explicitly
and present the populations instead as below-floor targets whose
detection-regime analysis (Extended Data~Fig.~1) itself is a
contribution.

Second, some rare populations are lost before integration rather than
during integration. Neutrophil subsets including LOX-1$^+$
PMN-MDSC and TAN-1 through TAN-5 subsets are systematically
undersampled by the standard 10x Genomics droplet-capture protocol,
and mast cells are similarly depleted relative to their tissue
abundance. Our framework is a post-integration detector and does not
and cannot recover populations that never entered the cell-by-gene
matrix. We flag this as a pre-integration identifiability problem
orthogonal to overcorrection, to be addressed by protocol choice
and by modality-aware capture strategies (FACS-sort enrichment prior
to library preparation; split-pool-barcoded whole-transcriptome
methods that capture neutrophils without selection bias).

Third, our evaluation uses annotated rare populations with their
labels hidden at test time. This is a legitimate form of ground
truth but it is not a direct test on unannotated populations.
Novel populations detected under our framework in the pan-cancer
atlas --- for example, a candidate CXCR5-dim TPEX sub-state
\citep{ourfig5} --- require the full evidence package of
\Cref{sec:evidence-package} plus independent cohort replication
before a definite claim of previously-unannotated biology. Main-text
detections are reported only when $\geq 4$ of the six evidence
criteria are satisfied.

\paragraph*{Telescoping detection recovers below-whole-atlas populations at sub-compartment scale.}
The $\pi\cdot\Delta^2 < 1.5\times10^{-4}$ below-floor rule above is
stated at the whole-atlas level and applies to flat REAL. Hierarchical
REAL, applied at lineage-restricted sub-compartments, moves the
operating point from whole-atlas prevalence to compartment-level
prevalence; at each level the flat floor still applies, but the
$\pi\cdot\Delta^2$ product is typically $10$--$100\times$ larger
because rare populations are often concentrated in specific lineages.
This recovers pulmonary ionocytes at the airway-epithelial sub-compartment
level, and the CXCR5-dim TPEX sub-state at the exhausted-CD8 sub-compartment
level (Supplementary Theorem~S1-hierarchical, Mathematics agent
\citealp{math2026hierarchical}). Our disclosure is therefore
three-tier: populations recovered by flat whole-atlas REAL (above
floor at $\pi\cdot\Delta^2 > 1.5\times10^{-4}$, e.g., MAIT, $\gamma\delta$
T, TREM2$^+$ LAM, myCAF, EMT-hybrid), populations recovered by
hierarchical sub-compartment REAL (above floor at sub-compartment
product, e.g., ionocytes, TPEX sub-state, AS-DC), and populations
unresolved at all compartment depths (endothelial tip cells, drug-tolerant
persister baseline state). Main-text claims are restricted to the
first two tiers; tier-three populations are disclosed as present in
the dataset but beyond the current framework's detection capability.

% ---------------------------------------------------------------------
% Requires: sec:minimax-detectability, sec:evidence-package,
% \citep{ourfig5} is a placeholder for the Fig 5 companion section /
% preprint. Biological Science may prefer \Cref{fig:fig5} or
% \Cref{sec:results-tpex-substate} depending on skeleton.
% ---------------------------------------------------------------------
         <-- unchanged from Immunology branch
%   % ---------------------------------------------------------------------
% Joint closing paragraph for the below-floor Limitations subsection.
% Co-authored with Philosophy (agent-mozur8ub). Placement flow:
%
%   \input{philo_below_floor_head.tex}    (Philosophy's epistemic frame)
%   \input{immuno_limitations.tex}        (Immunology's biological disclosure)
%   \input{immuno_below_floor_close.tex}  (THIS FILE — joint affirmative)
%
% Closes the subsection on the telescoping-detection affirmative:
% most below-flat-floor populations are recovered under hierarchical
% sub-compartment REAL.
% ---------------------------------------------------------------------

\paragraph*{Recovery under hierarchical application.}
Notwithstanding the strict boundary established above, the \emph{truly-unresolved} stratum is surprisingly small. Of the 22 immune and tumor populations in our registry whose whole-atlas product $\pi \cdot \Delta^{2}$ places them into one of the three detection tiers, 20 are recoverable through flat or hierarchical application of \textsc{REAL} at an appropriate compartment scale (\Cref{thm:S1-hierarchical}; \Cref{sec:methods-oracle-score} protocol): 7 populations are detected at whole-atlas scale (flat tier --- MAIT, $\gamma\delta$ T V$\delta1$-bias, TREM2$^{+}$ LAM, myCAF, iCAF, plasma-TLS, EMT-hybrid-near-pole), and 13 additional populations cross the detection floor only under hierarchical application at their natural sub-compartment scale (hierarchical tier). Pulmonary ionocytes pass detection at the airway-epithelial sub-compartment of HLCA (effective $\pi \approx 5 \times 10^{-3}$ in $n \approx 240{,}000$ airway cells); \textit{CXCR5}$^{+}$\textit{TCF7}$^{+}$ progenitor-exhausted CD8 T cells pass detection at the exhausted-CD8 sub-compartment of pooled melanoma atlases ($\pi \approx 10^{-2}$ in $n \approx 100{,}000$ cells); LAMP3$^{+}$ mregDC, AS-DC, and HCMV-imprinted adaptive NK pass detection at dendritic-cell and natural-killer sub-compartments respectively. The two populations that remain below the floor at every compartment depth are endothelial tip cells (which sit at the intersection of a small-mass vascular niche and a narrow transcriptional separation) and the drug-tolerant persister baseline state in oncology (Tumor Biology registry row TU-T-001), with mast cells in tumors additionally disclosed as out-of-scope for post-integration detection (mode~4 pre-integration capture loss). These unresolved entries are explicitly deferred to follow-up work: their recovery would require either the multi-atlas pooling regime ($n > 10^{8}$ cells) or auxiliary-modality augmentation, both of which lie beyond this paper's scope. This three-tier disclosure --- recovered-flat, recovered-hierarchical, deferred --- is the most honest reading of what our framework currently certifies, and it preserves the paper's headline claim on the above-floor and hierarchical-recovered strata while placing no claim on the truly-unresolved tier.

% ---------------------------------------------------------------------
% Cross-references needed (shared with philo_below_floor_head + immuno_limitations):
%   thm:S1-hierarchical        Math hierarchical REAL theorem
%   sec:methods-oracle-score   Methods — oracle-score verification
% Bib entries: none new.
% ---------------------------------------------------------------------
   <-- unchanged from Immunology branch
%
% Use this file INSTEAD OF the 5-file chain if BioSci prefers a single splice point.
% Immunology's two \input calls (limitations + below-floor-close) are PRESERVED as
% \input calls inside this file rather than inlined, because they are owned by
% the immunology branch and Immunology should retain edit control over their content.
%
% Placement: sections/08_discussion.tex, replacing the 5-file \input chain with
% a single % Philosophy §08 epistemic-scope block — single-file consolidation.
% Merges the five \input calls that BioSci currently threads through into one file
% for easier future-edit management. Semantically identical to the input chain:
%   \input{handoffs/philo_discussion_meno.tex}
%   \input{handoffs/philo_limitations.tex}
%   \input{handoffs/philo_below_floor_head.tex}
%   \input{handoffs/immuno_limitations.tex}         <-- unchanged from Immunology branch
%   \input{handoffs/immuno_below_floor_close.tex}   <-- unchanged from Immunology branch
%
% Use this file INSTEAD OF the 5-file chain if BioSci prefers a single splice point.
% Immunology's two \input calls (limitations + below-floor-close) are PRESERVED as
% \input calls inside this file rather than inlined, because they are owned by
% the immunology branch and Immunology should retain edit control over their content.
%
% Placement: sections/08_discussion.tex, replacing the 5-file \input chain with
% a single \input{handoffs/philo_08_epistemic_block.tex}.
%
% Author: Philosophy (agent-mozur8ub). Last updated for v2.4 pre-submission.

% =============================================================================
% §08 EPISTEMIC SCOPE BLOCK — single-file consolidation
% =============================================================================

% -----------------------------------------------------------------------------
% (1) Meno-problem framing — opens the epistemic scope block.
% Content from philo_discussion_meno.tex (agent/philosophy@<head>).
% ≈250w. Fixes the is/ought gap: "reproducible structure destroyed" ≠ "novel cell type".
% -----------------------------------------------------------------------------

\paragraph*{What a REAL flag is, and is not.}
A motif flagged by REAL is the claim that a local, reproducible geometric structure---witnessed across independent batches before integration, with measurable local-mass distortion after---has been erased by the integration map beyond what admissible batch-effect deformations can explain. This is a claim about \emph{evaluation}, not about ontology. It is not a claim that a new cell type exists. A motif $w \in \mathcal{W}$ with low $\mathrm{LRS}(w)$ licenses exactly one inference: the integration has discarded reproducible structure of the kind a biologist would want to interrogate. Whether that structure corresponds to a novel cell type, a transitional state along an already-known trajectory, a rare disease-specific configuration, or a technical sub-mode that merely looks biological, is a further question REAL does not answer. Upgrading a motif to a cell-type claim requires the usual burden of single-cell biology: independent sampling, orthogonal markers, targeted protocol (sorting, spatial validation, experimental perturbation), and community replication. REAL's value is that it identifies \emph{where to look}, not what has been found. Conflating the two would recreate the very error this paper is designed to expose: treating a framework-internal signal as a fact about the world. Our strong claim is therefore deliberately asymmetric. A high motif set $|\mathcal{W}|$ under a candidate integration is evidence of measurement insufficiency; a low motif set is evidence that the integration respected the syntactic reproducibility criterion. Ontology is downstream of both.

% -----------------------------------------------------------------------------
% (2) Limitations and epistemic scope — Route-C exchangeability boundaries.
% Content from philo_limitations.tex. ≤500w.
% -----------------------------------------------------------------------------

\subsection*{Limitations and epistemic scope}
\label{philo:limitations}
REAL relies on an exchangeability assumption between the known rare populations we hold out for calibration (pancreatic $\epsilon$ cells, pulmonary ionocytes, LAMP3$^+$ mregDC, TPEX, FOLR2$^+$ perivascular macrophages, apCAF, and the populations enumerated in Table~S\ref{tab:heldout}) and the genuinely unannotated rare subpopulations to which we transfer the estimate. The assumption is formally stated in Lemma~\ref{lem:exchangeability} as a \emph{bound}: the held-out-rare loss rate upper-bounds the expected loss rate for unannotated rare subpopulations of comparable abundance, local geometry, and batch profile. We emphasise four boundary conditions under which the bound is not tight. First, the bound degrades when the unannotated population lies outside the span of the highly-variable genes used to construct the latent space; REAL cannot protect what the representation cannot see. Second, the bound degrades when the unannotated population has a qualitatively different batch-effect signature from the calibration populations---for instance, a disease-specific state observed in only one batch. Third, the bound is conservative for populations whose local density profile is not well-approximated by the level-set stratification we use; transitional cells along a continuum can appear as low-density but are not rare in the sampling sense. Fourth, our use of known rare populations as proxies for unknown ones inherits a selection bias: the populations we can hold out are those that have already been discovered, and the discovered set may not be exchangeable with the undiscovered set.

Our strongest defensible claim is therefore deliberately limited. REAL provides (i) a label-free, invariant-based detectability framework calibrated on held-out known rare subpopulations, (ii) an integration strategy (RareShield) that demonstrably reduces the held-out-rare loss rate at pan-cancer-atlas scale without degrading batch correction for abundant cell types, and (iii) an explicit exchangeability bound under which the held-out loss rate is an upper bound on the loss rate of truly unannotated rare subpopulations. A REAL-flagged motif $w \in \mathcal{W}$ is a claim that reproducible structure has been destroyed, not a claim that a new cell type exists; upgrading to an ontological claim requires the usual experimental burden (sorting, orthogonal markers, spatial validation, replication). The complementary negative: a dataset with low $|\mathcal{W}|$ under a candidate integration is evidence that no cross-batch-reproducible rare structure was erased, but not a guarantee that no such structure existed---by construction, we cannot rule out subpopulations whose reproducibility signal was below our sampling resolution or outside the evaluated feature span. These boundaries are not weaknesses of REAL specifically; they are the price of reasoning about what the evaluation framework cannot directly observe. Naming the price is what distinguishes an epistemically calibrated claim from a framework-captured one.

% -----------------------------------------------------------------------------
% (3) Below the power floor — epistemic frame head (three-tier stratification).
% Content from philo_below_floor_head.tex.
% -----------------------------------------------------------------------------

\subsection*{Below the power floor: what the framework cannot yet certify}
\label{sec:below-floor-epistemic}

A framework's sensitivity is itself a piece of information about the framework, not about the biology. The minimax envelope of \Cref{thm:minimax} and the hierarchical-detectability result of \Cref{thm:S1-hierarchical} together partition the rare-population landscape at Kang~2024 scale into three strata: \emph{above-floor-flat} populations for which whole-atlas REAL is sufficient; \emph{hierarchical-recovered} populations which cross the detectability floor only after sub-compartment telescoping (per \Cref{sec:S1-envelope}); and \emph{truly-unresolved} populations for which neither flat nor hierarchical REAL attains the required power at Kang depth. The first two strata carry the manuscript's headline claim; this subsection states what we can and cannot say about the third. To read an absence of \textsc{REAL} flags on a truly-unresolved population as evidence of preservation is to commit exactly the failure-to-reject-as-success error that this paper is designed to expose (\Cref{philo:limitations}; Supp Note~S0 trap~8 at \Cref{sec:eightTraps}).

Two distinctions matter in making the disclosure precise. First, \textit{truly-unresolved at this dataset} is not the same as \textit{unresolved in principle}: the only escape routes are pooled multi-atlas scale (on the order of $10^{8}$ cells), auxiliary modalities that lift $\Delta$ above the $\pi \Delta^{2} \geq 1.5 \times 10^{-4}$ threshold at Kang scale (protein panels, TCR / BCR sequencing, spatial context), or targeted protocol changes that enrich the population before capture (FACS sorting, modality-aware library preparation). \textit{Truly unresolved} is thus a relation between a population, a measurement regime, and a framework, not a property of the population alone. Second, a motif flagged by \textsc{REAL} remains, even here, a claim about \emph{evaluation}, not about ontology (\Cref{philo:episto_note}). Correspondingly, the absence of a flag in the truly-unresolved stratum remains a claim about \emph{measurement insufficiency}, not about biology. The populations listed below are not excluded from single-cell biology; they are excluded from this framework's current certificate of preservation, under the stated regime.

% -----------------------------------------------------------------------------
% (4)-(5) Immunology's biological disclosure + joint closing.
% These files live in the immunology branch; Immunology retains edit control.
% They are \input'd here rather than inlined, so Immunology-branch edits propagate
% automatically on the next BioSci pull.
% -----------------------------------------------------------------------------

\input{handoffs/immuno_limitations.tex}
\input{handoffs/immuno_below_floor_close.tex}

% =============================================================================
% END §08 EPISTEMIC SCOPE BLOCK
% =============================================================================

% Integration note for BioSci:
%   - To use this consolidated file: replace the 5-file \input chain in
%     sections/08_discussion.tex with:
%       \input{handoffs/philo_08_epistemic_block.tex}
%   - To keep the 5-file chain: ignore this file; it is semantically identical
%     to the chain at this commit. Consolidation is a convenience, not a correctness fix.
%   - If Philosophy or Immunology updates any of the five source files, re-generating
%     this consolidated file requires Philosophy to re-splice; the \input calls for
%     Immunology's two files remain untouched.
%
% Lint status: passes lint:epistemic clean (trap 8 / trap 4 / IM.4 guards all present).
.
%
% Author: Philosophy (agent-mozur8ub). Last updated for v2.4 pre-submission.

% =============================================================================
% §08 EPISTEMIC SCOPE BLOCK — single-file consolidation
% =============================================================================

% -----------------------------------------------------------------------------
% (1) Meno-problem framing — opens the epistemic scope block.
% Content from philo_discussion_meno.tex (agent/philosophy@<head>).
% ≈250w. Fixes the is/ought gap: "reproducible structure destroyed" ≠ "novel cell type".
% -----------------------------------------------------------------------------

\paragraph*{What a REAL flag is, and is not.}
A motif flagged by REAL is the claim that a local, reproducible geometric structure---witnessed across independent batches before integration, with measurable local-mass distortion after---has been erased by the integration map beyond what admissible batch-effect deformations can explain. This is a claim about \emph{evaluation}, not about ontology. It is not a claim that a new cell type exists. A motif $w \in \mathcal{W}$ with low $\mathrm{LRS}(w)$ licenses exactly one inference: the integration has discarded reproducible structure of the kind a biologist would want to interrogate. Whether that structure corresponds to a novel cell type, a transitional state along an already-known trajectory, a rare disease-specific configuration, or a technical sub-mode that merely looks biological, is a further question REAL does not answer. Upgrading a motif to a cell-type claim requires the usual burden of single-cell biology: independent sampling, orthogonal markers, targeted protocol (sorting, spatial validation, experimental perturbation), and community replication. REAL's value is that it identifies \emph{where to look}, not what has been found. Conflating the two would recreate the very error this paper is designed to expose: treating a framework-internal signal as a fact about the world. Our strong claim is therefore deliberately asymmetric. A high motif set $|\mathcal{W}|$ under a candidate integration is evidence of measurement insufficiency; a low motif set is evidence that the integration respected the syntactic reproducibility criterion. Ontology is downstream of both.

% -----------------------------------------------------------------------------
% (2) Limitations and epistemic scope — Route-C exchangeability boundaries.
% Content from philo_limitations.tex. ≤500w.
% -----------------------------------------------------------------------------

\subsection*{Limitations and epistemic scope}
\label{philo:limitations}
REAL relies on an exchangeability assumption between the known rare populations we hold out for calibration (pancreatic $\epsilon$ cells, pulmonary ionocytes, LAMP3$^+$ mregDC, TPEX, FOLR2$^+$ perivascular macrophages, apCAF, and the populations enumerated in Table~S\ref{tab:heldout}) and the genuinely unannotated rare subpopulations to which we transfer the estimate. The assumption is formally stated in Lemma~\ref{lem:exchangeability} as a \emph{bound}: the held-out-rare loss rate upper-bounds the expected loss rate for unannotated rare subpopulations of comparable abundance, local geometry, and batch profile. We emphasise four boundary conditions under which the bound is not tight. First, the bound degrades when the unannotated population lies outside the span of the highly-variable genes used to construct the latent space; REAL cannot protect what the representation cannot see. Second, the bound degrades when the unannotated population has a qualitatively different batch-effect signature from the calibration populations---for instance, a disease-specific state observed in only one batch. Third, the bound is conservative for populations whose local density profile is not well-approximated by the level-set stratification we use; transitional cells along a continuum can appear as low-density but are not rare in the sampling sense. Fourth, our use of known rare populations as proxies for unknown ones inherits a selection bias: the populations we can hold out are those that have already been discovered, and the discovered set may not be exchangeable with the undiscovered set.

Our strongest defensible claim is therefore deliberately limited. REAL provides (i) a label-free, invariant-based detectability framework calibrated on held-out known rare subpopulations, (ii) an integration strategy (RareShield) that demonstrably reduces the held-out-rare loss rate at pan-cancer-atlas scale without degrading batch correction for abundant cell types, and (iii) an explicit exchangeability bound under which the held-out loss rate is an upper bound on the loss rate of truly unannotated rare subpopulations. A REAL-flagged motif $w \in \mathcal{W}$ is a claim that reproducible structure has been destroyed, not a claim that a new cell type exists; upgrading to an ontological claim requires the usual experimental burden (sorting, orthogonal markers, spatial validation, replication). The complementary negative: a dataset with low $|\mathcal{W}|$ under a candidate integration is evidence that no cross-batch-reproducible rare structure was erased, but not a guarantee that no such structure existed---by construction, we cannot rule out subpopulations whose reproducibility signal was below our sampling resolution or outside the evaluated feature span. These boundaries are not weaknesses of REAL specifically; they are the price of reasoning about what the evaluation framework cannot directly observe. Naming the price is what distinguishes an epistemically calibrated claim from a framework-captured one.

% -----------------------------------------------------------------------------
% (3) Below the power floor — epistemic frame head (three-tier stratification).
% Content from philo_below_floor_head.tex.
% -----------------------------------------------------------------------------

\subsection*{Below the power floor: what the framework cannot yet certify}
\label{sec:below-floor-epistemic}

A framework's sensitivity is itself a piece of information about the framework, not about the biology. The minimax envelope of \Cref{thm:minimax} and the hierarchical-detectability result of \Cref{thm:S1-hierarchical} together partition the rare-population landscape at Kang~2024 scale into three strata: \emph{above-floor-flat} populations for which whole-atlas REAL is sufficient; \emph{hierarchical-recovered} populations which cross the detectability floor only after sub-compartment telescoping (per \Cref{sec:S1-envelope}); and \emph{truly-unresolved} populations for which neither flat nor hierarchical REAL attains the required power at Kang depth. The first two strata carry the manuscript's headline claim; this subsection states what we can and cannot say about the third. To read an absence of \textsc{REAL} flags on a truly-unresolved population as evidence of preservation is to commit exactly the failure-to-reject-as-success error that this paper is designed to expose (\Cref{philo:limitations}; Supp Note~S0 trap~8 at \Cref{sec:eightTraps}).

Two distinctions matter in making the disclosure precise. First, \textit{truly-unresolved at this dataset} is not the same as \textit{unresolved in principle}: the only escape routes are pooled multi-atlas scale (on the order of $10^{8}$ cells), auxiliary modalities that lift $\Delta$ above the $\pi \Delta^{2} \geq 1.5 \times 10^{-4}$ threshold at Kang scale (protein panels, TCR / BCR sequencing, spatial context), or targeted protocol changes that enrich the population before capture (FACS sorting, modality-aware library preparation). \textit{Truly unresolved} is thus a relation between a population, a measurement regime, and a framework, not a property of the population alone. Second, a motif flagged by \textsc{REAL} remains, even here, a claim about \emph{evaluation}, not about ontology (\Cref{philo:episto_note}). Correspondingly, the absence of a flag in the truly-unresolved stratum remains a claim about \emph{measurement insufficiency}, not about biology. The populations listed below are not excluded from single-cell biology; they are excluded from this framework's current certificate of preservation, under the stated regime.

% -----------------------------------------------------------------------------
% (4)-(5) Immunology's biological disclosure + joint closing.
% These files live in the immunology branch; Immunology retains edit control.
% They are \input'd here rather than inlined, so Immunology-branch edits propagate
% automatically on the next BioSci pull.
% -----------------------------------------------------------------------------

% ---------------------------------------------------------------------
% Limitations subsection — immunology contribution to the Discussion.
% ~350 words. Honest disclosure of populations where REAL/RareShield
% cannot make reliable detection claims at Kang-atlas scale.
% Author: Immunology (agent-mozus3vv). Cross-references Math's
% predicted-detectability table (agent/mathematics@f564cc6).
% Target: sections/08_discussion.tex, as a \subsection after the
% "Implications for immunology" subsection (immuno_discussion.tex).
% ---------------------------------------------------------------------

\subsection{Boundaries of what our framework can currently guarantee}
\label{sec:discussion-limitations-immunology}

Three concrete limitations shape the claims we make. First, the Le~Cam
detectability envelope (\Cref{sec:minimax-detectability}) places a
sharp boundary on which rare populations can be recovered label-free
at any given atlas scale. At the Kang pan-cancer immune atlas
($n = 4.9 \times 10^6$, $B = 30$, per-batch heterogeneity
$\Lambda = 3$), populations with prevalence-separation product
$\pi\Delta < 10^{-3}\sigma$ fall below the detection floor, producing
a true-positive rate under 0.3 at $\alpha = 0.05$. Specifically,
\textit{CXCR5}$^+$\textit{TCF7}$^+$ progenitor-exhausted CD8 T
(TPEX), LAMP3$^+$ mregDC, KIR$^+$ adaptive NK, endothelial tip
cells, and drug-tolerant-persister tumour cells all sit in this
regime at Kang alone. For these populations our framework does not
claim detection from the Kang atlas in isolation; their recovery
requires either auxiliary protein or TCR signal, or pooling across
multiple atlases beyond $10^8$ cells. We disclose this explicitly
and present the populations instead as below-floor targets whose
detection-regime analysis (Extended Data~Fig.~1) itself is a
contribution.

Second, some rare populations are lost before integration rather than
during integration. Neutrophil subsets including LOX-1$^+$
PMN-MDSC and TAN-1 through TAN-5 subsets are systematically
undersampled by the standard 10x Genomics droplet-capture protocol,
and mast cells are similarly depleted relative to their tissue
abundance. Our framework is a post-integration detector and does not
and cannot recover populations that never entered the cell-by-gene
matrix. We flag this as a pre-integration identifiability problem
orthogonal to overcorrection, to be addressed by protocol choice
and by modality-aware capture strategies (FACS-sort enrichment prior
to library preparation; split-pool-barcoded whole-transcriptome
methods that capture neutrophils without selection bias).

Third, our evaluation uses annotated rare populations with their
labels hidden at test time. This is a legitimate form of ground
truth but it is not a direct test on unannotated populations.
Novel populations detected under our framework in the pan-cancer
atlas --- for example, a candidate CXCR5-dim TPEX sub-state
\citep{ourfig5} --- require the full evidence package of
\Cref{sec:evidence-package} plus independent cohort replication
before a definite claim of previously-unannotated biology. Main-text
detections are reported only when $\geq 4$ of the six evidence
criteria are satisfied.

\paragraph*{Telescoping detection recovers below-whole-atlas populations at sub-compartment scale.}
The $\pi\cdot\Delta^2 < 1.5\times10^{-4}$ below-floor rule above is
stated at the whole-atlas level and applies to flat REAL. Hierarchical
REAL, applied at lineage-restricted sub-compartments, moves the
operating point from whole-atlas prevalence to compartment-level
prevalence; at each level the flat floor still applies, but the
$\pi\cdot\Delta^2$ product is typically $10$--$100\times$ larger
because rare populations are often concentrated in specific lineages.
This recovers pulmonary ionocytes at the airway-epithelial sub-compartment
level, and the CXCR5-dim TPEX sub-state at the exhausted-CD8 sub-compartment
level (Supplementary Theorem~S1-hierarchical, Mathematics agent
\citealp{math2026hierarchical}). Our disclosure is therefore
three-tier: populations recovered by flat whole-atlas REAL (above
floor at $\pi\cdot\Delta^2 > 1.5\times10^{-4}$, e.g., MAIT, $\gamma\delta$
T, TREM2$^+$ LAM, myCAF, EMT-hybrid), populations recovered by
hierarchical sub-compartment REAL (above floor at sub-compartment
product, e.g., ionocytes, TPEX sub-state, AS-DC), and populations
unresolved at all compartment depths (endothelial tip cells, drug-tolerant
persister baseline state). Main-text claims are restricted to the
first two tiers; tier-three populations are disclosed as present in
the dataset but beyond the current framework's detection capability.

% ---------------------------------------------------------------------
% Requires: sec:minimax-detectability, sec:evidence-package,
% \citep{ourfig5} is a placeholder for the Fig 5 companion section /
% preprint. Biological Science may prefer \Cref{fig:fig5} or
% \Cref{sec:results-tpex-substate} depending on skeleton.
% ---------------------------------------------------------------------

% ---------------------------------------------------------------------
% Joint closing paragraph for the below-floor Limitations subsection.
% Co-authored with Philosophy (agent-mozur8ub). Placement flow:
%
%   \input{philo_below_floor_head.tex}    (Philosophy's epistemic frame)
%   \input{immuno_limitations.tex}        (Immunology's biological disclosure)
%   \input{immuno_below_floor_close.tex}  (THIS FILE — joint affirmative)
%
% Closes the subsection on the telescoping-detection affirmative:
% most below-flat-floor populations are recovered under hierarchical
% sub-compartment REAL.
% ---------------------------------------------------------------------

\paragraph*{Recovery under hierarchical application.}
Notwithstanding the strict boundary established above, the \emph{truly-unresolved} stratum is surprisingly small. Of the 22 immune and tumor populations in our registry whose whole-atlas product $\pi \cdot \Delta^{2}$ places them into one of the three detection tiers, 20 are recoverable through flat or hierarchical application of \textsc{REAL} at an appropriate compartment scale (\Cref{thm:S1-hierarchical}; \Cref{sec:methods-oracle-score} protocol): 7 populations are detected at whole-atlas scale (flat tier --- MAIT, $\gamma\delta$ T V$\delta1$-bias, TREM2$^{+}$ LAM, myCAF, iCAF, plasma-TLS, EMT-hybrid-near-pole), and 13 additional populations cross the detection floor only under hierarchical application at their natural sub-compartment scale (hierarchical tier). Pulmonary ionocytes pass detection at the airway-epithelial sub-compartment of HLCA (effective $\pi \approx 5 \times 10^{-3}$ in $n \approx 240{,}000$ airway cells); \textit{CXCR5}$^{+}$\textit{TCF7}$^{+}$ progenitor-exhausted CD8 T cells pass detection at the exhausted-CD8 sub-compartment of pooled melanoma atlases ($\pi \approx 10^{-2}$ in $n \approx 100{,}000$ cells); LAMP3$^{+}$ mregDC, AS-DC, and HCMV-imprinted adaptive NK pass detection at dendritic-cell and natural-killer sub-compartments respectively. The two populations that remain below the floor at every compartment depth are endothelial tip cells (which sit at the intersection of a small-mass vascular niche and a narrow transcriptional separation) and the drug-tolerant persister baseline state in oncology (Tumor Biology registry row TU-T-001), with mast cells in tumors additionally disclosed as out-of-scope for post-integration detection (mode~4 pre-integration capture loss). These unresolved entries are explicitly deferred to follow-up work: their recovery would require either the multi-atlas pooling regime ($n > 10^{8}$ cells) or auxiliary-modality augmentation, both of which lie beyond this paper's scope. This three-tier disclosure --- recovered-flat, recovered-hierarchical, deferred --- is the most honest reading of what our framework currently certifies, and it preserves the paper's headline claim on the above-floor and hierarchical-recovered strata while placing no claim on the truly-unresolved tier.

% ---------------------------------------------------------------------
% Cross-references needed (shared with philo_below_floor_head + immuno_limitations):
%   thm:S1-hierarchical        Math hierarchical REAL theorem
%   sec:methods-oracle-score   Methods — oracle-score verification
% Bib entries: none new.
% ---------------------------------------------------------------------


% =============================================================================
% END §08 EPISTEMIC SCOPE BLOCK
% =============================================================================

% Integration note for BioSci:
%   - To use this consolidated file: replace the 5-file \input chain in
%     sections/08_discussion.tex with:
%       % Philosophy §08 epistemic-scope block — single-file consolidation.
% Merges the five \input calls that BioSci currently threads through into one file
% for easier future-edit management. Semantically identical to the input chain:
%   \input{handoffs/philo_discussion_meno.tex}
%   \input{handoffs/philo_limitations.tex}
%   \input{handoffs/philo_below_floor_head.tex}
%   \input{handoffs/immuno_limitations.tex}         <-- unchanged from Immunology branch
%   \input{handoffs/immuno_below_floor_close.tex}   <-- unchanged from Immunology branch
%
% Use this file INSTEAD OF the 5-file chain if BioSci prefers a single splice point.
% Immunology's two \input calls (limitations + below-floor-close) are PRESERVED as
% \input calls inside this file rather than inlined, because they are owned by
% the immunology branch and Immunology should retain edit control over their content.
%
% Placement: sections/08_discussion.tex, replacing the 5-file \input chain with
% a single \input{handoffs/philo_08_epistemic_block.tex}.
%
% Author: Philosophy (agent-mozur8ub). Last updated for v2.4 pre-submission.

% =============================================================================
% §08 EPISTEMIC SCOPE BLOCK — single-file consolidation
% =============================================================================

% -----------------------------------------------------------------------------
% (1) Meno-problem framing — opens the epistemic scope block.
% Content from philo_discussion_meno.tex (agent/philosophy@<head>).
% ≈250w. Fixes the is/ought gap: "reproducible structure destroyed" ≠ "novel cell type".
% -----------------------------------------------------------------------------

\paragraph*{What a REAL flag is, and is not.}
A motif flagged by REAL is the claim that a local, reproducible geometric structure---witnessed across independent batches before integration, with measurable local-mass distortion after---has been erased by the integration map beyond what admissible batch-effect deformations can explain. This is a claim about \emph{evaluation}, not about ontology. It is not a claim that a new cell type exists. A motif $w \in \mathcal{W}$ with low $\mathrm{LRS}(w)$ licenses exactly one inference: the integration has discarded reproducible structure of the kind a biologist would want to interrogate. Whether that structure corresponds to a novel cell type, a transitional state along an already-known trajectory, a rare disease-specific configuration, or a technical sub-mode that merely looks biological, is a further question REAL does not answer. Upgrading a motif to a cell-type claim requires the usual burden of single-cell biology: independent sampling, orthogonal markers, targeted protocol (sorting, spatial validation, experimental perturbation), and community replication. REAL's value is that it identifies \emph{where to look}, not what has been found. Conflating the two would recreate the very error this paper is designed to expose: treating a framework-internal signal as a fact about the world. Our strong claim is therefore deliberately asymmetric. A high motif set $|\mathcal{W}|$ under a candidate integration is evidence of measurement insufficiency; a low motif set is evidence that the integration respected the syntactic reproducibility criterion. Ontology is downstream of both.

% -----------------------------------------------------------------------------
% (2) Limitations and epistemic scope — Route-C exchangeability boundaries.
% Content from philo_limitations.tex. ≤500w.
% -----------------------------------------------------------------------------

\subsection*{Limitations and epistemic scope}
\label{philo:limitations}
REAL relies on an exchangeability assumption between the known rare populations we hold out for calibration (pancreatic $\epsilon$ cells, pulmonary ionocytes, LAMP3$^+$ mregDC, TPEX, FOLR2$^+$ perivascular macrophages, apCAF, and the populations enumerated in Table~S\ref{tab:heldout}) and the genuinely unannotated rare subpopulations to which we transfer the estimate. The assumption is formally stated in Lemma~\ref{lem:exchangeability} as a \emph{bound}: the held-out-rare loss rate upper-bounds the expected loss rate for unannotated rare subpopulations of comparable abundance, local geometry, and batch profile. We emphasise four boundary conditions under which the bound is not tight. First, the bound degrades when the unannotated population lies outside the span of the highly-variable genes used to construct the latent space; REAL cannot protect what the representation cannot see. Second, the bound degrades when the unannotated population has a qualitatively different batch-effect signature from the calibration populations---for instance, a disease-specific state observed in only one batch. Third, the bound is conservative for populations whose local density profile is not well-approximated by the level-set stratification we use; transitional cells along a continuum can appear as low-density but are not rare in the sampling sense. Fourth, our use of known rare populations as proxies for unknown ones inherits a selection bias: the populations we can hold out are those that have already been discovered, and the discovered set may not be exchangeable with the undiscovered set.

Our strongest defensible claim is therefore deliberately limited. REAL provides (i) a label-free, invariant-based detectability framework calibrated on held-out known rare subpopulations, (ii) an integration strategy (RareShield) that demonstrably reduces the held-out-rare loss rate at pan-cancer-atlas scale without degrading batch correction for abundant cell types, and (iii) an explicit exchangeability bound under which the held-out loss rate is an upper bound on the loss rate of truly unannotated rare subpopulations. A REAL-flagged motif $w \in \mathcal{W}$ is a claim that reproducible structure has been destroyed, not a claim that a new cell type exists; upgrading to an ontological claim requires the usual experimental burden (sorting, orthogonal markers, spatial validation, replication). The complementary negative: a dataset with low $|\mathcal{W}|$ under a candidate integration is evidence that no cross-batch-reproducible rare structure was erased, but not a guarantee that no such structure existed---by construction, we cannot rule out subpopulations whose reproducibility signal was below our sampling resolution or outside the evaluated feature span. These boundaries are not weaknesses of REAL specifically; they are the price of reasoning about what the evaluation framework cannot directly observe. Naming the price is what distinguishes an epistemically calibrated claim from a framework-captured one.

% -----------------------------------------------------------------------------
% (3) Below the power floor — epistemic frame head (three-tier stratification).
% Content from philo_below_floor_head.tex.
% -----------------------------------------------------------------------------

\subsection*{Below the power floor: what the framework cannot yet certify}
\label{sec:below-floor-epistemic}

A framework's sensitivity is itself a piece of information about the framework, not about the biology. The minimax envelope of \Cref{thm:minimax} and the hierarchical-detectability result of \Cref{thm:S1-hierarchical} together partition the rare-population landscape at Kang~2024 scale into three strata: \emph{above-floor-flat} populations for which whole-atlas REAL is sufficient; \emph{hierarchical-recovered} populations which cross the detectability floor only after sub-compartment telescoping (per \Cref{sec:S1-envelope}); and \emph{truly-unresolved} populations for which neither flat nor hierarchical REAL attains the required power at Kang depth. The first two strata carry the manuscript's headline claim; this subsection states what we can and cannot say about the third. To read an absence of \textsc{REAL} flags on a truly-unresolved population as evidence of preservation is to commit exactly the failure-to-reject-as-success error that this paper is designed to expose (\Cref{philo:limitations}; Supp Note~S0 trap~8 at \Cref{sec:eightTraps}).

Two distinctions matter in making the disclosure precise. First, \textit{truly-unresolved at this dataset} is not the same as \textit{unresolved in principle}: the only escape routes are pooled multi-atlas scale (on the order of $10^{8}$ cells), auxiliary modalities that lift $\Delta$ above the $\pi \Delta^{2} \geq 1.5 \times 10^{-4}$ threshold at Kang scale (protein panels, TCR / BCR sequencing, spatial context), or targeted protocol changes that enrich the population before capture (FACS sorting, modality-aware library preparation). \textit{Truly unresolved} is thus a relation between a population, a measurement regime, and a framework, not a property of the population alone. Second, a motif flagged by \textsc{REAL} remains, even here, a claim about \emph{evaluation}, not about ontology (\Cref{philo:episto_note}). Correspondingly, the absence of a flag in the truly-unresolved stratum remains a claim about \emph{measurement insufficiency}, not about biology. The populations listed below are not excluded from single-cell biology; they are excluded from this framework's current certificate of preservation, under the stated regime.

% -----------------------------------------------------------------------------
% (4)-(5) Immunology's biological disclosure + joint closing.
% These files live in the immunology branch; Immunology retains edit control.
% They are \input'd here rather than inlined, so Immunology-branch edits propagate
% automatically on the next BioSci pull.
% -----------------------------------------------------------------------------

\input{handoffs/immuno_limitations.tex}
\input{handoffs/immuno_below_floor_close.tex}

% =============================================================================
% END §08 EPISTEMIC SCOPE BLOCK
% =============================================================================

% Integration note for BioSci:
%   - To use this consolidated file: replace the 5-file \input chain in
%     sections/08_discussion.tex with:
%       \input{handoffs/philo_08_epistemic_block.tex}
%   - To keep the 5-file chain: ignore this file; it is semantically identical
%     to the chain at this commit. Consolidation is a convenience, not a correctness fix.
%   - If Philosophy or Immunology updates any of the five source files, re-generating
%     this consolidated file requires Philosophy to re-splice; the \input calls for
%     Immunology's two files remain untouched.
%
% Lint status: passes lint:epistemic clean (trap 8 / trap 4 / IM.4 guards all present).

%   - To keep the 5-file chain: ignore this file; it is semantically identical
%     to the chain at this commit. Consolidation is a convenience, not a correctness fix.
%   - If Philosophy or Immunology updates any of the five source files, re-generating
%     this consolidated file requires Philosophy to re-splice; the \input calls for
%     Immunology's two files remain untouched.
%
% Lint status: passes lint:epistemic clean (trap 8 / trap 4 / IM.4 guards all present).

%   - To keep the 5-file chain: ignore this file; it is semantically identical
%     to the chain at this commit. Consolidation is a convenience, not a correctness fix.
%   - If Philosophy or Immunology updates any of the five source files, re-generating
%     this consolidated file requires Philosophy to re-splice; the \input calls for
%     Immunology's two files remain untouched.
%
% Lint status: passes lint:epistemic clean (trap 8 / trap 4 / IM.4 guards all present).
.
%
% Author: Philosophy (agent-mozur8ub). Last updated for v2.4 pre-submission.

% =============================================================================
% §08 EPISTEMIC SCOPE BLOCK — single-file consolidation
% =============================================================================

% -----------------------------------------------------------------------------
% (1) Meno-problem framing — opens the epistemic scope block.
% Content from philo_discussion_meno.tex (agent/philosophy@<head>).
% ≈250w. Fixes the is/ought gap: "reproducible structure destroyed" ≠ "novel cell type".
% -----------------------------------------------------------------------------

\paragraph*{What a REAL flag is, and is not.}
A motif flagged by REAL is the claim that a local, reproducible geometric structure---witnessed across independent batches before integration, with measurable local-mass distortion after---has been erased by the integration map beyond what admissible batch-effect deformations can explain. This is a claim about \emph{evaluation}, not about ontology. It is not a claim that a new cell type exists. A motif $w \in \mathcal{W}$ with low $\mathrm{LRS}(w)$ licenses exactly one inference: the integration has discarded reproducible structure of the kind a biologist would want to interrogate. Whether that structure corresponds to a novel cell type, a transitional state along an already-known trajectory, a rare disease-specific configuration, or a technical sub-mode that merely looks biological, is a further question REAL does not answer. Upgrading a motif to a cell-type claim requires the usual burden of single-cell biology: independent sampling, orthogonal markers, targeted protocol (sorting, spatial validation, experimental perturbation), and community replication. REAL's value is that it identifies \emph{where to look}, not what has been found. Conflating the two would recreate the very error this paper is designed to expose: treating a framework-internal signal as a fact about the world. Our strong claim is therefore deliberately asymmetric. A high motif set $|\mathcal{W}|$ under a candidate integration is evidence of measurement insufficiency; a low motif set is evidence that the integration respected the syntactic reproducibility criterion. Ontology is downstream of both.

% -----------------------------------------------------------------------------
% (2) Limitations and epistemic scope — Route-C exchangeability boundaries.
% Content from philo_limitations.tex. ≤500w.
% -----------------------------------------------------------------------------

\subsection*{Limitations and epistemic scope}
\label{philo:limitations}
REAL relies on an exchangeability assumption between the known rare populations we hold out for calibration (pancreatic $\epsilon$ cells, pulmonary ionocytes, LAMP3$^+$ mregDC, TPEX, FOLR2$^+$ perivascular macrophages, apCAF, and the populations enumerated in Table~S\ref{tab:heldout}) and the genuinely unannotated rare subpopulations to which we transfer the estimate. The assumption is formally stated in Lemma~\ref{lem:exchangeability} as a \emph{bound}: the held-out-rare loss rate upper-bounds the expected loss rate for unannotated rare subpopulations of comparable abundance, local geometry, and batch profile. We emphasise four boundary conditions under which the bound is not tight. First, the bound degrades when the unannotated population lies outside the span of the highly-variable genes used to construct the latent space; REAL cannot protect what the representation cannot see. Second, the bound degrades when the unannotated population has a qualitatively different batch-effect signature from the calibration populations---for instance, a disease-specific state observed in only one batch. Third, the bound is conservative for populations whose local density profile is not well-approximated by the level-set stratification we use; transitional cells along a continuum can appear as low-density but are not rare in the sampling sense. Fourth, our use of known rare populations as proxies for unknown ones inherits a selection bias: the populations we can hold out are those that have already been discovered, and the discovered set may not be exchangeable with the undiscovered set.

Our strongest defensible claim is therefore deliberately limited. REAL provides (i) a label-free, invariant-based detectability framework calibrated on held-out known rare subpopulations, (ii) an integration strategy (RareShield) that demonstrably reduces the held-out-rare loss rate at pan-cancer-atlas scale without degrading batch correction for abundant cell types, and (iii) an explicit exchangeability bound under which the held-out loss rate is an upper bound on the loss rate of truly unannotated rare subpopulations. A REAL-flagged motif $w \in \mathcal{W}$ is a claim that reproducible structure has been destroyed, not a claim that a new cell type exists; upgrading to an ontological claim requires the usual experimental burden (sorting, orthogonal markers, spatial validation, replication). The complementary negative: a dataset with low $|\mathcal{W}|$ under a candidate integration is evidence that no cross-batch-reproducible rare structure was erased, but not a guarantee that no such structure existed---by construction, we cannot rule out subpopulations whose reproducibility signal was below our sampling resolution or outside the evaluated feature span. These boundaries are not weaknesses of REAL specifically; they are the price of reasoning about what the evaluation framework cannot directly observe. Naming the price is what distinguishes an epistemically calibrated claim from a framework-captured one.

% -----------------------------------------------------------------------------
% (3) Below the power floor — epistemic frame head (three-tier stratification).
% Content from philo_below_floor_head.tex.
% -----------------------------------------------------------------------------

\subsection*{Below the power floor: what the framework cannot yet certify}
\label{sec:below-floor-epistemic}

A framework's sensitivity is itself a piece of information about the framework, not about the biology. The minimax envelope of \Cref{thm:minimax} and the hierarchical-detectability result of \Cref{thm:S1-hierarchical} together partition the rare-population landscape at Kang~2024 scale into three strata: \emph{above-floor-flat} populations for which whole-atlas REAL is sufficient; \emph{hierarchical-recovered} populations which cross the detectability floor only after sub-compartment telescoping (per \Cref{sec:S1-envelope}); and \emph{truly-unresolved} populations for which neither flat nor hierarchical REAL attains the required power at Kang depth. The first two strata carry the manuscript's headline claim; this subsection states what we can and cannot say about the third. To read an absence of \textsc{REAL} flags on a truly-unresolved population as evidence of preservation is to commit exactly the failure-to-reject-as-success error that this paper is designed to expose (\Cref{philo:limitations}; Supp Note~S0 trap~8 at \Cref{sec:eightTraps}).

Two distinctions matter in making the disclosure precise. First, \textit{truly-unresolved at this dataset} is not the same as \textit{unresolved in principle}: the only escape routes are pooled multi-atlas scale (on the order of $10^{8}$ cells), auxiliary modalities that lift $\Delta$ above the $\pi \Delta^{2} \geq 1.5 \times 10^{-4}$ threshold at Kang scale (protein panels, TCR / BCR sequencing, spatial context), or targeted protocol changes that enrich the population before capture (FACS sorting, modality-aware library preparation). \textit{Truly unresolved} is thus a relation between a population, a measurement regime, and a framework, not a property of the population alone. Second, a motif flagged by \textsc{REAL} remains, even here, a claim about \emph{evaluation}, not about ontology (\Cref{philo:episto_note}). Correspondingly, the absence of a flag in the truly-unresolved stratum remains a claim about \emph{measurement insufficiency}, not about biology. The populations listed below are not excluded from single-cell biology; they are excluded from this framework's current certificate of preservation, under the stated regime.

% -----------------------------------------------------------------------------
% (4)-(5) Immunology's biological disclosure + joint closing.
% These files live in the immunology branch; Immunology retains edit control.
% They are \input'd here rather than inlined, so Immunology-branch edits propagate
% automatically on the next BioSci pull.
% -----------------------------------------------------------------------------

% ---------------------------------------------------------------------
% Limitations subsection — immunology contribution to the Discussion.
% ~350 words. Honest disclosure of populations where REAL/RareShield
% cannot make reliable detection claims at Kang-atlas scale.
% Author: Immunology (agent-mozus3vv). Cross-references Math's
% predicted-detectability table (agent/mathematics@f564cc6).
% Target: sections/08_discussion.tex, as a \subsection after the
% "Implications for immunology" subsection (immuno_discussion.tex).
% ---------------------------------------------------------------------

\subsection{Boundaries of what our framework can currently guarantee}
\label{sec:discussion-limitations-immunology}

Three concrete limitations shape the claims we make. First, the Le~Cam
detectability envelope (\Cref{sec:minimax-detectability}) places a
sharp boundary on which rare populations can be recovered label-free
at any given atlas scale. At the Kang pan-cancer immune atlas
($n = 4.9 \times 10^6$, $B = 30$, per-batch heterogeneity
$\Lambda = 3$), populations with prevalence-separation product
$\pi\Delta < 10^{-3}\sigma$ fall below the detection floor, producing
a true-positive rate under 0.3 at $\alpha = 0.05$. Specifically,
\textit{CXCR5}$^+$\textit{TCF7}$^+$ progenitor-exhausted CD8 T
(TPEX), LAMP3$^+$ mregDC, KIR$^+$ adaptive NK, endothelial tip
cells, and drug-tolerant-persister tumour cells all sit in this
regime at Kang alone. For these populations our framework does not
claim detection from the Kang atlas in isolation; their recovery
requires either auxiliary protein or TCR signal, or pooling across
multiple atlases beyond $10^8$ cells. We disclose this explicitly
and present the populations instead as below-floor targets whose
detection-regime analysis (Extended Data~Fig.~1) itself is a
contribution.

Second, some rare populations are lost before integration rather than
during integration. Neutrophil subsets including LOX-1$^+$
PMN-MDSC and TAN-1 through TAN-5 subsets are systematically
undersampled by the standard 10x Genomics droplet-capture protocol,
and mast cells are similarly depleted relative to their tissue
abundance. Our framework is a post-integration detector and does not
and cannot recover populations that never entered the cell-by-gene
matrix. We flag this as a pre-integration identifiability problem
orthogonal to overcorrection, to be addressed by protocol choice
and by modality-aware capture strategies (FACS-sort enrichment prior
to library preparation; split-pool-barcoded whole-transcriptome
methods that capture neutrophils without selection bias).

Third, our evaluation uses annotated rare populations with their
labels hidden at test time. This is a legitimate form of ground
truth but it is not a direct test on unannotated populations.
Novel populations detected under our framework in the pan-cancer
atlas --- for example, a candidate CXCR5-dim TPEX sub-state
\citep{ourfig5} --- require the full evidence package of
\Cref{sec:evidence-package} plus independent cohort replication
before a definite claim of previously-unannotated biology. Main-text
detections are reported only when $\geq 4$ of the six evidence
criteria are satisfied.

\paragraph*{Telescoping detection recovers below-whole-atlas populations at sub-compartment scale.}
The $\pi\cdot\Delta^2 < 1.5\times10^{-4}$ below-floor rule above is
stated at the whole-atlas level and applies to flat REAL. Hierarchical
REAL, applied at lineage-restricted sub-compartments, moves the
operating point from whole-atlas prevalence to compartment-level
prevalence; at each level the flat floor still applies, but the
$\pi\cdot\Delta^2$ product is typically $10$--$100\times$ larger
because rare populations are often concentrated in specific lineages.
This recovers pulmonary ionocytes at the airway-epithelial sub-compartment
level, and the CXCR5-dim TPEX sub-state at the exhausted-CD8 sub-compartment
level (Supplementary Theorem~S1-hierarchical, Mathematics agent
\citealp{math2026hierarchical}). Our disclosure is therefore
three-tier: populations recovered by flat whole-atlas REAL (above
floor at $\pi\cdot\Delta^2 > 1.5\times10^{-4}$, e.g., MAIT, $\gamma\delta$
T, TREM2$^+$ LAM, myCAF, EMT-hybrid), populations recovered by
hierarchical sub-compartment REAL (above floor at sub-compartment
product, e.g., ionocytes, TPEX sub-state, AS-DC), and populations
unresolved at all compartment depths (endothelial tip cells, drug-tolerant
persister baseline state). Main-text claims are restricted to the
first two tiers; tier-three populations are disclosed as present in
the dataset but beyond the current framework's detection capability.

% ---------------------------------------------------------------------
% Requires: sec:minimax-detectability, sec:evidence-package,
% \citep{ourfig5} is a placeholder for the Fig 5 companion section /
% preprint. Biological Science may prefer \Cref{fig:fig5} or
% \Cref{sec:results-tpex-substate} depending on skeleton.
% ---------------------------------------------------------------------

% ---------------------------------------------------------------------
% Joint closing paragraph for the below-floor Limitations subsection.
% Co-authored with Philosophy (agent-mozur8ub). Placement flow:
%
%   % ---------------------------------------------------------------------
% Philosophy's epistemic-frame head for the below-floor Limitations subsection.
% To be placed IMMEDIATELY BEFORE the Immunology-drafted biological disclosure
% (workspace/handoffs/immuno_limitations.tex at agent/immunology@b8a898d).
%
% Co-authored structure agreed with Immunology (agent-mozus3vv, 2026-05-10):
%   1. Philosophy epistemic frame (this file) — 2 short paragraphs, ≈220w.
%   2. Immunology biological disclosure (immuno_limitations.tex) — ≈350w.
%   3. Joint closing (telescoping-detection affirmative) — Immunology to draft
%      after pulling this file and Math's sub_compartment_pooling_bounds.md.
%
% Placement: sections/08_discussion.tex, as the subsection \subsection{Below the power floor},
% ordered after % Limitations & Epistemic Scope subsection (≤500w).
% For BioSci's locked slot in the Discussion. Cites Math Lemma S1-exchangeability as a *bound*.
% Covers: exchangeability failure modes, reference-class scope, non-ontic status, circularity mitigations.

\subsection*{Limitations and epistemic scope}
REAL relies on an exchangeability assumption between the known rare populations we hold out for calibration (pancreatic $\epsilon$ cells, pulmonary ionocytes, LAMP3$^+$ mregDC, TPEX, FOLR2$^+$ perivascular macrophages and the other populations enumerated in Table~S\ref{tab:heldout}) and the genuinely unannotated rare subpopulations to which we transfer the estimate. The assumption is formally stated in Lemma~\ref{lem:exchangeability} as a \emph{bound}: the held-out-rare loss rate upper-bounds the expected loss rate for unannotated rare subpopulations of comparable abundance, local geometry, and batch profile. We emphasise four boundary conditions under which the bound is not tight. First, the bound degrades when the unannotated population lies outside the span of the highly-variable genes used to construct the latent space; REAL cannot protect what the representation cannot see. Second, the bound degrades when the unannotated population has a qualitatively different batch-effect signature from the calibration populations---for instance, a disease-specific state observed in only one batch. Third, the bound is conservative for populations whose local density profile is not well-approximated by the level-set stratification we use; transitional cells along a continuum can appear as low-density but are not rare in the sampling sense. Fourth, our use of known rare populations as proxies for unknown ones inherits a selection bias: the populations we can hold out are those that have already been discovered, and the discovered set may not be exchangeable with the undiscovered set.

Our strongest defensible claim is therefore deliberately limited. REAL provides (i) a label-free, invariant-based detectability framework calibrated on held-out known rare subpopulations, (ii) an integration strategy (RareShield) that demonstrably reduces the held-out-rare loss rate at pan-cancer-atlas scale without degrading batch correction for abundant cell types, and (iii) an explicit exchangeability bound under which the held-out loss rate is an upper bound on the loss rate of truly unannotated rare subpopulations. A REAL-flagged motif $w \in \mathcal{W}$ is a claim that reproducible structure has been destroyed, not a claim that a new cell type exists; upgrading to an ontological claim requires the usual experimental burden (sorting, orthogonal markers, spatial validation, replication). The complementary negative: a dataset with low $|\mathcal{W}|$ under a candidate integration is evidence that no cross-batch-reproducible rare structure was erased, but not a guarantee that no such structure existed---by construction, we cannot rule out subpopulations whose reproducibility signal was below our sampling resolution or outside the evaluated feature span. These boundaries are not weaknesses of REAL specifically; they are the price of reasoning about what the evaluation framework cannot directly observe. Naming the price is what distinguishes an epistemically calibrated claim from a framework-captured one.
 and before % ---------------------------------------------------------------------
% Limitations subsection — immunology contribution to the Discussion.
% ~350 words. Honest disclosure of populations where REAL/RareShield
% cannot make reliable detection claims at Kang-atlas scale.
% Author: Immunology (agent-mozus3vv). Cross-references Math's
% predicted-detectability table (agent/mathematics@f564cc6).
% Target: sections/08_discussion.tex, as a \subsection after the
% "Implications for immunology" subsection (immuno_discussion.tex).
% ---------------------------------------------------------------------

\subsection{Boundaries of what our framework can currently guarantee}
\label{sec:discussion-limitations-immunology}

Three concrete limitations shape the claims we make. First, the Le~Cam
detectability envelope (\Cref{sec:minimax-detectability}) places a
sharp boundary on which rare populations can be recovered label-free
at any given atlas scale. At the Kang pan-cancer immune atlas
($n = 4.9 \times 10^6$, $B = 30$, per-batch heterogeneity
$\Lambda = 3$), populations with prevalence-separation product
$\pi\Delta < 10^{-3}\sigma$ fall below the detection floor, producing
a true-positive rate under 0.3 at $\alpha = 0.05$. Specifically,
\textit{CXCR5}$^+$\textit{TCF7}$^+$ progenitor-exhausted CD8 T
(TPEX), LAMP3$^+$ mregDC, KIR$^+$ adaptive NK, endothelial tip
cells, and drug-tolerant-persister tumour cells all sit in this
regime at Kang alone. For these populations our framework does not
claim detection from the Kang atlas in isolation; their recovery
requires either auxiliary protein or TCR signal, or pooling across
multiple atlases beyond $10^8$ cells. We disclose this explicitly
and present the populations instead as below-floor targets whose
detection-regime analysis (Extended Data~Fig.~1) itself is a
contribution.

Second, some rare populations are lost before integration rather than
during integration. Neutrophil subsets including LOX-1$^+$
PMN-MDSC and TAN-1 through TAN-5 subsets are systematically
undersampled by the standard 10x Genomics droplet-capture protocol,
and mast cells are similarly depleted relative to their tissue
abundance. Our framework is a post-integration detector and does not
and cannot recover populations that never entered the cell-by-gene
matrix. We flag this as a pre-integration identifiability problem
orthogonal to overcorrection, to be addressed by protocol choice
and by modality-aware capture strategies (FACS-sort enrichment prior
to library preparation; split-pool-barcoded whole-transcriptome
methods that capture neutrophils without selection bias).

Third, our evaluation uses annotated rare populations with their
labels hidden at test time. This is a legitimate form of ground
truth but it is not a direct test on unannotated populations.
Novel populations detected under our framework in the pan-cancer
atlas --- for example, a candidate CXCR5-dim TPEX sub-state
\citep{ourfig5} --- require the full evidence package of
\Cref{sec:evidence-package} plus independent cohort replication
before a definite claim of previously-unannotated biology. Main-text
detections are reported only when $\geq 4$ of the six evidence
criteria are satisfied.

\paragraph*{Telescoping detection recovers below-whole-atlas populations at sub-compartment scale.}
The $\pi\cdot\Delta^2 < 1.5\times10^{-4}$ below-floor rule above is
stated at the whole-atlas level and applies to flat REAL. Hierarchical
REAL, applied at lineage-restricted sub-compartments, moves the
operating point from whole-atlas prevalence to compartment-level
prevalence; at each level the flat floor still applies, but the
$\pi\cdot\Delta^2$ product is typically $10$--$100\times$ larger
because rare populations are often concentrated in specific lineages.
This recovers pulmonary ionocytes at the airway-epithelial sub-compartment
level, and the CXCR5-dim TPEX sub-state at the exhausted-CD8 sub-compartment
level (Supplementary Theorem~S1-hierarchical, Mathematics agent
\citealp{math2026hierarchical}). Our disclosure is therefore
three-tier: populations recovered by flat whole-atlas REAL (above
floor at $\pi\cdot\Delta^2 > 1.5\times10^{-4}$, e.g., MAIT, $\gamma\delta$
T, TREM2$^+$ LAM, myCAF, EMT-hybrid), populations recovered by
hierarchical sub-compartment REAL (above floor at sub-compartment
product, e.g., ionocytes, TPEX sub-state, AS-DC), and populations
unresolved at all compartment depths (endothelial tip cells, drug-tolerant
persister baseline state). Main-text claims are restricted to the
first two tiers; tier-three populations are disclosed as present in
the dataset but beyond the current framework's detection capability.

% ---------------------------------------------------------------------
% Requires: sec:minimax-detectability, sec:evidence-package,
% \citep{ourfig5} is a placeholder for the Fig 5 companion section /
% preprint. Biological Science may prefer \Cref{fig:fig5} or
% \Cref{sec:results-tpex-substate} depending on skeleton.
% ---------------------------------------------------------------------
.
%
% Anchors:
%   \Cref{thm:minimax}         Math Theorem 1 (minimax lower bound; floor)
%   \Cref{thm:labelblind}      Math Theorem 2.1 (label-based σ-algebra closure)
%   \Cref{lem:exchangeability} Math Lemma S1 (exchangeability bound)
%   \Cref{philo:limitations}   Philosophy Limitations subsection (head of §08 epistemic scope)
%   \Cref{sec:minimax-detectability}  Math's detectability envelope (for crosscheck)
% ---------------------------------------------------------------------

\subsection{Below the power floor: what the framework cannot yet certify}
\label{sec:below-floor-epistemic}

A framework's sensitivity is itself a piece of information about the framework, not about the biology. The minimax envelope of \Cref{thm:minimax} and the hierarchical-detectability result of \Cref{thm:S1-hierarchical} together partition the rare-population landscape at Kang~2024 scale into three strata: \emph{above-floor-flat} populations for which whole-atlas REAL is sufficient; \emph{hierarchical-recovered} populations which cross the detectability floor only after sub-compartment telescoping (per \Cref{sec:S1-envelope}); and \emph{truly-unresolved} populations for which neither flat nor hierarchical REAL attains the required power at Kang depth. The first two strata carry the manuscript's headline claim; this subsection states what we can and cannot say about the third. To read an absence of \textsc{REAL} flags on a truly-unresolved population as evidence of preservation is to commit exactly the failure-to-reject-as-success error that this paper is designed to expose (\Cref{philo:limitations}; Supp Note~S0 trap~8 at \Cref{sec:eightTraps}).

Two distinctions matter in making the disclosure precise. First, \textit{truly-unresolved at this dataset} is not the same as \textit{unresolved in principle}: the only escape routes are pooled multi-atlas scale (on the order of $10^{8}$ cells), auxiliary modalities that lift $\Delta$ above the $\pi \Delta^{2} \geq 1.5 \times 10^{-4}$ threshold at Kang scale (protein panels, TCR / BCR sequencing, spatial context), or targeted protocol changes that enrich the population before capture (FACS sorting, modality-aware library preparation). \textit{Truly unresolved} is thus a relation between a population, a measurement regime, and a framework, not a property of the population alone. Second, a motif flagged by \textsc{REAL} remains, even here, a claim about \emph{evaluation}, not about ontology (\Cref{philo:episto_note}). Correspondingly, the absence of a flag in the truly-unresolved stratum remains a claim about \emph{measurement insufficiency}, not about biology. The populations listed below are not excluded from single-cell biology; they are excluded from this framework's current certificate of preservation, under the stated regime.
    (Philosophy's epistemic frame)
%   % ---------------------------------------------------------------------
% Limitations subsection — immunology contribution to the Discussion.
% ~350 words. Honest disclosure of populations where REAL/RareShield
% cannot make reliable detection claims at Kang-atlas scale.
% Author: Immunology (agent-mozus3vv). Cross-references Math's
% predicted-detectability table (agent/mathematics@f564cc6).
% Target: sections/08_discussion.tex, as a \subsection after the
% "Implications for immunology" subsection (immuno_discussion.tex).
% ---------------------------------------------------------------------

\subsection{Boundaries of what our framework can currently guarantee}
\label{sec:discussion-limitations-immunology}

Three concrete limitations shape the claims we make. First, the Le~Cam
detectability envelope (\Cref{sec:minimax-detectability}) places a
sharp boundary on which rare populations can be recovered label-free
at any given atlas scale. At the Kang pan-cancer immune atlas
($n = 4.9 \times 10^6$, $B = 30$, per-batch heterogeneity
$\Lambda = 3$), populations with prevalence-separation product
$\pi\Delta < 10^{-3}\sigma$ fall below the detection floor, producing
a true-positive rate under 0.3 at $\alpha = 0.05$. Specifically,
\textit{CXCR5}$^+$\textit{TCF7}$^+$ progenitor-exhausted CD8 T
(TPEX), LAMP3$^+$ mregDC, KIR$^+$ adaptive NK, endothelial tip
cells, and drug-tolerant-persister tumour cells all sit in this
regime at Kang alone. For these populations our framework does not
claim detection from the Kang atlas in isolation; their recovery
requires either auxiliary protein or TCR signal, or pooling across
multiple atlases beyond $10^8$ cells. We disclose this explicitly
and present the populations instead as below-floor targets whose
detection-regime analysis (Extended Data~Fig.~1) itself is a
contribution.

Second, some rare populations are lost before integration rather than
during integration. Neutrophil subsets including LOX-1$^+$
PMN-MDSC and TAN-1 through TAN-5 subsets are systematically
undersampled by the standard 10x Genomics droplet-capture protocol,
and mast cells are similarly depleted relative to their tissue
abundance. Our framework is a post-integration detector and does not
and cannot recover populations that never entered the cell-by-gene
matrix. We flag this as a pre-integration identifiability problem
orthogonal to overcorrection, to be addressed by protocol choice
and by modality-aware capture strategies (FACS-sort enrichment prior
to library preparation; split-pool-barcoded whole-transcriptome
methods that capture neutrophils without selection bias).

Third, our evaluation uses annotated rare populations with their
labels hidden at test time. This is a legitimate form of ground
truth but it is not a direct test on unannotated populations.
Novel populations detected under our framework in the pan-cancer
atlas --- for example, a candidate CXCR5-dim TPEX sub-state
\citep{ourfig5} --- require the full evidence package of
\Cref{sec:evidence-package} plus independent cohort replication
before a definite claim of previously-unannotated biology. Main-text
detections are reported only when $\geq 4$ of the six evidence
criteria are satisfied.

\paragraph*{Telescoping detection recovers below-whole-atlas populations at sub-compartment scale.}
The $\pi\cdot\Delta^2 < 1.5\times10^{-4}$ below-floor rule above is
stated at the whole-atlas level and applies to flat REAL. Hierarchical
REAL, applied at lineage-restricted sub-compartments, moves the
operating point from whole-atlas prevalence to compartment-level
prevalence; at each level the flat floor still applies, but the
$\pi\cdot\Delta^2$ product is typically $10$--$100\times$ larger
because rare populations are often concentrated in specific lineages.
This recovers pulmonary ionocytes at the airway-epithelial sub-compartment
level, and the CXCR5-dim TPEX sub-state at the exhausted-CD8 sub-compartment
level (Supplementary Theorem~S1-hierarchical, Mathematics agent
\citealp{math2026hierarchical}). Our disclosure is therefore
three-tier: populations recovered by flat whole-atlas REAL (above
floor at $\pi\cdot\Delta^2 > 1.5\times10^{-4}$, e.g., MAIT, $\gamma\delta$
T, TREM2$^+$ LAM, myCAF, EMT-hybrid), populations recovered by
hierarchical sub-compartment REAL (above floor at sub-compartment
product, e.g., ionocytes, TPEX sub-state, AS-DC), and populations
unresolved at all compartment depths (endothelial tip cells, drug-tolerant
persister baseline state). Main-text claims are restricted to the
first two tiers; tier-three populations are disclosed as present in
the dataset but beyond the current framework's detection capability.

% ---------------------------------------------------------------------
% Requires: sec:minimax-detectability, sec:evidence-package,
% \citep{ourfig5} is a placeholder for the Fig 5 companion section /
% preprint. Biological Science may prefer \Cref{fig:fig5} or
% \Cref{sec:results-tpex-substate} depending on skeleton.
% ---------------------------------------------------------------------
        (Immunology's biological disclosure)
%   % ---------------------------------------------------------------------
% Joint closing paragraph for the below-floor Limitations subsection.
% Co-authored with Philosophy (agent-mozur8ub). Placement flow:
%
%   % ---------------------------------------------------------------------
% Philosophy's epistemic-frame head for the below-floor Limitations subsection.
% To be placed IMMEDIATELY BEFORE the Immunology-drafted biological disclosure
% (workspace/handoffs/immuno_limitations.tex at agent/immunology@b8a898d).
%
% Co-authored structure agreed with Immunology (agent-mozus3vv, 2026-05-10):
%   1. Philosophy epistemic frame (this file) — 2 short paragraphs, ≈220w.
%   2. Immunology biological disclosure (immuno_limitations.tex) — ≈350w.
%   3. Joint closing (telescoping-detection affirmative) — Immunology to draft
%      after pulling this file and Math's sub_compartment_pooling_bounds.md.
%
% Placement: sections/08_discussion.tex, as the subsection \subsection{Below the power floor},
% ordered after \input{philo_limitations.tex} and before \input{immuno_limitations.tex}.
%
% Anchors:
%   \Cref{thm:minimax}         Math Theorem 1 (minimax lower bound; floor)
%   \Cref{thm:labelblind}      Math Theorem 2.1 (label-based σ-algebra closure)
%   \Cref{lem:exchangeability} Math Lemma S1 (exchangeability bound)
%   \Cref{philo:limitations}   Philosophy Limitations subsection (head of §08 epistemic scope)
%   \Cref{sec:minimax-detectability}  Math's detectability envelope (for crosscheck)
% ---------------------------------------------------------------------

\subsection{Below the power floor: what the framework cannot yet certify}
\label{sec:below-floor-epistemic}

A framework's sensitivity is itself a piece of information about the framework, not about the biology. The minimax envelope of \Cref{thm:minimax} and the hierarchical-detectability result of \Cref{thm:S1-hierarchical} together partition the rare-population landscape at Kang~2024 scale into three strata: \emph{above-floor-flat} populations for which whole-atlas REAL is sufficient; \emph{hierarchical-recovered} populations which cross the detectability floor only after sub-compartment telescoping (per \Cref{sec:S1-envelope}); and \emph{truly-unresolved} populations for which neither flat nor hierarchical REAL attains the required power at Kang depth. The first two strata carry the manuscript's headline claim; this subsection states what we can and cannot say about the third. To read an absence of \textsc{REAL} flags on a truly-unresolved population as evidence of preservation is to commit exactly the failure-to-reject-as-success error that this paper is designed to expose (\Cref{philo:limitations}; Supp Note~S0 trap~8 at \Cref{sec:eightTraps}).

Two distinctions matter in making the disclosure precise. First, \textit{truly-unresolved at this dataset} is not the same as \textit{unresolved in principle}: the only escape routes are pooled multi-atlas scale (on the order of $10^{8}$ cells), auxiliary modalities that lift $\Delta$ above the $\pi \Delta^{2} \geq 1.5 \times 10^{-4}$ threshold at Kang scale (protein panels, TCR / BCR sequencing, spatial context), or targeted protocol changes that enrich the population before capture (FACS sorting, modality-aware library preparation). \textit{Truly unresolved} is thus a relation between a population, a measurement regime, and a framework, not a property of the population alone. Second, a motif flagged by \textsc{REAL} remains, even here, a claim about \emph{evaluation}, not about ontology (\Cref{philo:episto_note}). Correspondingly, the absence of a flag in the truly-unresolved stratum remains a claim about \emph{measurement insufficiency}, not about biology. The populations listed below are not excluded from single-cell biology; they are excluded from this framework's current certificate of preservation, under the stated regime.
    (Philosophy's epistemic frame)
%   % ---------------------------------------------------------------------
% Limitations subsection — immunology contribution to the Discussion.
% ~350 words. Honest disclosure of populations where REAL/RareShield
% cannot make reliable detection claims at Kang-atlas scale.
% Author: Immunology (agent-mozus3vv). Cross-references Math's
% predicted-detectability table (agent/mathematics@f564cc6).
% Target: sections/08_discussion.tex, as a \subsection after the
% "Implications for immunology" subsection (immuno_discussion.tex).
% ---------------------------------------------------------------------

\subsection{Boundaries of what our framework can currently guarantee}
\label{sec:discussion-limitations-immunology}

Three concrete limitations shape the claims we make. First, the Le~Cam
detectability envelope (\Cref{sec:minimax-detectability}) places a
sharp boundary on which rare populations can be recovered label-free
at any given atlas scale. At the Kang pan-cancer immune atlas
($n = 4.9 \times 10^6$, $B = 30$, per-batch heterogeneity
$\Lambda = 3$), populations with prevalence-separation product
$\pi\Delta < 10^{-3}\sigma$ fall below the detection floor, producing
a true-positive rate under 0.3 at $\alpha = 0.05$. Specifically,
\textit{CXCR5}$^+$\textit{TCF7}$^+$ progenitor-exhausted CD8 T
(TPEX), LAMP3$^+$ mregDC, KIR$^+$ adaptive NK, endothelial tip
cells, and drug-tolerant-persister tumour cells all sit in this
regime at Kang alone. For these populations our framework does not
claim detection from the Kang atlas in isolation; their recovery
requires either auxiliary protein or TCR signal, or pooling across
multiple atlases beyond $10^8$ cells. We disclose this explicitly
and present the populations instead as below-floor targets whose
detection-regime analysis (Extended Data~Fig.~1) itself is a
contribution.

Second, some rare populations are lost before integration rather than
during integration. Neutrophil subsets including LOX-1$^+$
PMN-MDSC and TAN-1 through TAN-5 subsets are systematically
undersampled by the standard 10x Genomics droplet-capture protocol,
and mast cells are similarly depleted relative to their tissue
abundance. Our framework is a post-integration detector and does not
and cannot recover populations that never entered the cell-by-gene
matrix. We flag this as a pre-integration identifiability problem
orthogonal to overcorrection, to be addressed by protocol choice
and by modality-aware capture strategies (FACS-sort enrichment prior
to library preparation; split-pool-barcoded whole-transcriptome
methods that capture neutrophils without selection bias).

Third, our evaluation uses annotated rare populations with their
labels hidden at test time. This is a legitimate form of ground
truth but it is not a direct test on unannotated populations.
Novel populations detected under our framework in the pan-cancer
atlas --- for example, a candidate CXCR5-dim TPEX sub-state
\citep{ourfig5} --- require the full evidence package of
\Cref{sec:evidence-package} plus independent cohort replication
before a definite claim of previously-unannotated biology. Main-text
detections are reported only when $\geq 4$ of the six evidence
criteria are satisfied.

\paragraph*{Telescoping detection recovers below-whole-atlas populations at sub-compartment scale.}
The $\pi\cdot\Delta^2 < 1.5\times10^{-4}$ below-floor rule above is
stated at the whole-atlas level and applies to flat REAL. Hierarchical
REAL, applied at lineage-restricted sub-compartments, moves the
operating point from whole-atlas prevalence to compartment-level
prevalence; at each level the flat floor still applies, but the
$\pi\cdot\Delta^2$ product is typically $10$--$100\times$ larger
because rare populations are often concentrated in specific lineages.
This recovers pulmonary ionocytes at the airway-epithelial sub-compartment
level, and the CXCR5-dim TPEX sub-state at the exhausted-CD8 sub-compartment
level (Supplementary Theorem~S1-hierarchical, Mathematics agent
\citealp{math2026hierarchical}). Our disclosure is therefore
three-tier: populations recovered by flat whole-atlas REAL (above
floor at $\pi\cdot\Delta^2 > 1.5\times10^{-4}$, e.g., MAIT, $\gamma\delta$
T, TREM2$^+$ LAM, myCAF, EMT-hybrid), populations recovered by
hierarchical sub-compartment REAL (above floor at sub-compartment
product, e.g., ionocytes, TPEX sub-state, AS-DC), and populations
unresolved at all compartment depths (endothelial tip cells, drug-tolerant
persister baseline state). Main-text claims are restricted to the
first two tiers; tier-three populations are disclosed as present in
the dataset but beyond the current framework's detection capability.

% ---------------------------------------------------------------------
% Requires: sec:minimax-detectability, sec:evidence-package,
% \citep{ourfig5} is a placeholder for the Fig 5 companion section /
% preprint. Biological Science may prefer \Cref{fig:fig5} or
% \Cref{sec:results-tpex-substate} depending on skeleton.
% ---------------------------------------------------------------------
        (Immunology's biological disclosure)
%   % ---------------------------------------------------------------------
% Joint closing paragraph for the below-floor Limitations subsection.
% Co-authored with Philosophy (agent-mozur8ub). Placement flow:
%
%   \input{philo_below_floor_head.tex}    (Philosophy's epistemic frame)
%   \input{immuno_limitations.tex}        (Immunology's biological disclosure)
%   \input{immuno_below_floor_close.tex}  (THIS FILE — joint affirmative)
%
% Closes the subsection on the telescoping-detection affirmative:
% most below-flat-floor populations are recovered under hierarchical
% sub-compartment REAL.
% ---------------------------------------------------------------------

\paragraph*{Recovery under hierarchical application.}
Notwithstanding the strict boundary established above, the \emph{truly-unresolved} stratum is surprisingly small. Of the 22 immune and tumor populations in our registry whose whole-atlas product $\pi \cdot \Delta^{2}$ places them into one of the three detection tiers, 20 are recoverable through flat or hierarchical application of \textsc{REAL} at an appropriate compartment scale (\Cref{thm:S1-hierarchical}; \Cref{sec:methods-oracle-score} protocol): 7 populations are detected at whole-atlas scale (flat tier --- MAIT, $\gamma\delta$ T V$\delta1$-bias, TREM2$^{+}$ LAM, myCAF, iCAF, plasma-TLS, EMT-hybrid-near-pole), and 13 additional populations cross the detection floor only under hierarchical application at their natural sub-compartment scale (hierarchical tier). Pulmonary ionocytes pass detection at the airway-epithelial sub-compartment of HLCA (effective $\pi \approx 5 \times 10^{-3}$ in $n \approx 240{,}000$ airway cells); \textit{CXCR5}$^{+}$\textit{TCF7}$^{+}$ progenitor-exhausted CD8 T cells pass detection at the exhausted-CD8 sub-compartment of pooled melanoma atlases ($\pi \approx 10^{-2}$ in $n \approx 100{,}000$ cells); LAMP3$^{+}$ mregDC, AS-DC, and HCMV-imprinted adaptive NK pass detection at dendritic-cell and natural-killer sub-compartments respectively. The two populations that remain below the floor at every compartment depth are endothelial tip cells (which sit at the intersection of a small-mass vascular niche and a narrow transcriptional separation) and the drug-tolerant persister baseline state in oncology (Tumor Biology registry row TU-T-001), with mast cells in tumors additionally disclosed as out-of-scope for post-integration detection (mode~4 pre-integration capture loss). These unresolved entries are explicitly deferred to follow-up work: their recovery would require either the multi-atlas pooling regime ($n > 10^{8}$ cells) or auxiliary-modality augmentation, both of which lie beyond this paper's scope. This three-tier disclosure --- recovered-flat, recovered-hierarchical, deferred --- is the most honest reading of what our framework currently certifies, and it preserves the paper's headline claim on the above-floor and hierarchical-recovered strata while placing no claim on the truly-unresolved tier.

% ---------------------------------------------------------------------
% Cross-references needed (shared with philo_below_floor_head + immuno_limitations):
%   thm:S1-hierarchical        Math hierarchical REAL theorem
%   sec:methods-oracle-score   Methods — oracle-score verification
% Bib entries: none new.
% ---------------------------------------------------------------------
  (THIS FILE — joint affirmative)
%
% Closes the subsection on the telescoping-detection affirmative:
% most below-flat-floor populations are recovered under hierarchical
% sub-compartment REAL.
% ---------------------------------------------------------------------

\paragraph*{Recovery under hierarchical application.}
Notwithstanding the strict boundary established above, the \emph{truly-unresolved} stratum is surprisingly small. Of the 22 immune and tumor populations in our registry whose whole-atlas product $\pi \cdot \Delta^{2}$ places them into one of the three detection tiers, 20 are recoverable through flat or hierarchical application of \textsc{REAL} at an appropriate compartment scale (\Cref{thm:S1-hierarchical}; \Cref{sec:methods-oracle-score} protocol): 7 populations are detected at whole-atlas scale (flat tier --- MAIT, $\gamma\delta$ T V$\delta1$-bias, TREM2$^{+}$ LAM, myCAF, iCAF, plasma-TLS, EMT-hybrid-near-pole), and 13 additional populations cross the detection floor only under hierarchical application at their natural sub-compartment scale (hierarchical tier). Pulmonary ionocytes pass detection at the airway-epithelial sub-compartment of HLCA (effective $\pi \approx 5 \times 10^{-3}$ in $n \approx 240{,}000$ airway cells); \textit{CXCR5}$^{+}$\textit{TCF7}$^{+}$ progenitor-exhausted CD8 T cells pass detection at the exhausted-CD8 sub-compartment of pooled melanoma atlases ($\pi \approx 10^{-2}$ in $n \approx 100{,}000$ cells); LAMP3$^{+}$ mregDC, AS-DC, and HCMV-imprinted adaptive NK pass detection at dendritic-cell and natural-killer sub-compartments respectively. The two populations that remain below the floor at every compartment depth are endothelial tip cells (which sit at the intersection of a small-mass vascular niche and a narrow transcriptional separation) and the drug-tolerant persister baseline state in oncology (Tumor Biology registry row TU-T-001), with mast cells in tumors additionally disclosed as out-of-scope for post-integration detection (mode~4 pre-integration capture loss). These unresolved entries are explicitly deferred to follow-up work: their recovery would require either the multi-atlas pooling regime ($n > 10^{8}$ cells) or auxiliary-modality augmentation, both of which lie beyond this paper's scope. This three-tier disclosure --- recovered-flat, recovered-hierarchical, deferred --- is the most honest reading of what our framework currently certifies, and it preserves the paper's headline claim on the above-floor and hierarchical-recovered strata while placing no claim on the truly-unresolved tier.

% ---------------------------------------------------------------------
% Cross-references needed (shared with philo_below_floor_head + immuno_limitations):
%   thm:S1-hierarchical        Math hierarchical REAL theorem
%   sec:methods-oracle-score   Methods — oracle-score verification
% Bib entries: none new.
% ---------------------------------------------------------------------
  (THIS FILE — joint affirmative)
%
% Closes the subsection on the telescoping-detection affirmative:
% most below-flat-floor populations are recovered under hierarchical
% sub-compartment REAL.
% ---------------------------------------------------------------------

\paragraph*{Recovery under hierarchical application.}
Notwithstanding the strict boundary established above, the \emph{truly-unresolved} stratum is surprisingly small. Of the 22 immune and tumor populations in our registry whose whole-atlas product $\pi \cdot \Delta^{2}$ places them into one of the three detection tiers, 20 are recoverable through flat or hierarchical application of \textsc{REAL} at an appropriate compartment scale (\Cref{thm:S1-hierarchical}; \Cref{sec:methods-oracle-score} protocol): 7 populations are detected at whole-atlas scale (flat tier --- MAIT, $\gamma\delta$ T V$\delta1$-bias, TREM2$^{+}$ LAM, myCAF, iCAF, plasma-TLS, EMT-hybrid-near-pole), and 13 additional populations cross the detection floor only under hierarchical application at their natural sub-compartment scale (hierarchical tier). Pulmonary ionocytes pass detection at the airway-epithelial sub-compartment of HLCA (effective $\pi \approx 5 \times 10^{-3}$ in $n \approx 240{,}000$ airway cells); \textit{CXCR5}$^{+}$\textit{TCF7}$^{+}$ progenitor-exhausted CD8 T cells pass detection at the exhausted-CD8 sub-compartment of pooled melanoma atlases ($\pi \approx 10^{-2}$ in $n \approx 100{,}000$ cells); LAMP3$^{+}$ mregDC, AS-DC, and HCMV-imprinted adaptive NK pass detection at dendritic-cell and natural-killer sub-compartments respectively. The two populations that remain below the floor at every compartment depth are endothelial tip cells (which sit at the intersection of a small-mass vascular niche and a narrow transcriptional separation) and the drug-tolerant persister baseline state in oncology (Tumor Biology registry row TU-T-001), with mast cells in tumors additionally disclosed as out-of-scope for post-integration detection (mode~4 pre-integration capture loss). These unresolved entries are explicitly deferred to follow-up work: their recovery would require either the multi-atlas pooling regime ($n > 10^{8}$ cells) or auxiliary-modality augmentation, both of which lie beyond this paper's scope. This three-tier disclosure --- recovered-flat, recovered-hierarchical, deferred --- is the most honest reading of what our framework currently certifies, and it preserves the paper's headline claim on the above-floor and hierarchical-recovered strata while placing no claim on the truly-unresolved tier.

% ---------------------------------------------------------------------
% Cross-references needed (shared with philo_below_floor_head + immuno_limitations):
%   thm:S1-hierarchical        Math hierarchical REAL theorem
%   sec:methods-oracle-score   Methods — oracle-score verification
% Bib entries: none new.
% ---------------------------------------------------------------------


% =============================================================================
% END §08 EPISTEMIC SCOPE BLOCK
% =============================================================================

% Integration note for BioSci:
%   - To use this consolidated file: replace the 5-file \input chain in
%     sections/08_discussion.tex with:
%       % Philosophy §08 epistemic-scope block — single-file consolidation.
% Merges the five \input calls that BioSci currently threads through into one file
% for easier future-edit management. Semantically identical to the input chain:
%   % Discussion paragraph (≈250w) — epistemic status of REAL-flagged motifs.
% For task t-mozv8oul. Fixes the is/ought gap: "reproducible structure destroyed" ≠ "novel cell type".
% Guards the team from over-claiming. Anchors: motif $w \in \mathcal{W}$; LRS(w); RareShield; Hacking-style intervention realism.

\paragraph*{What a REAL flag is, and is not.}
A motif flagged by REAL is the claim that a local, reproducible geometric structure---witnessed across independent batches before integration, with measurable local-mass distortion after---has been erased by the integration map beyond what admissible batch-effect deformations can explain. This is a claim about \emph{evaluation}, not about ontology. It is not a claim that a new cell type exists. A motif $w \in \mathcal{W}$ with low $\mathrm{LRS}(w)$ licenses exactly one inference: the integration has discarded reproducible structure of the kind a biologist would want to interrogate. Whether that structure corresponds to a novel cell type, a transitional state along an already-known trajectory, a rare disease-specific configuration, or a technical sub-mode that merely looks biological, is a further question that REAL does not answer. Upgrading a motif to a cell-type claim requires the usual burden of single-cell biology: independent sampling, orthogonal markers, targeted protocol (sorting, spatial validation, experimental perturbation), and community replication. REAL's value is that it identifies \emph{where to look}, not what has been found. Conflating the two would re-create the very error this paper is designed to expose: treating a framework-internal signal as a fact about the world. Our strong claim is therefore deliberately asymmetric: a high motif set $|\mathcal{W}|$ under a candidate integration is evidence of measurement insufficiency; a low motif set is evidence that the integration respected the syntactic reproducibility criterion. Ontology is downstream of both.

%   % Limitations & Epistemic Scope subsection (≤500w).
% For BioSci's locked slot in the Discussion. Cites Math Lemma S1-exchangeability as a *bound*.
% Covers: exchangeability failure modes, reference-class scope, non-ontic status, circularity mitigations.

\subsection*{Limitations and epistemic scope}
REAL relies on an exchangeability assumption between the known rare populations we hold out for calibration (pancreatic $\epsilon$ cells, pulmonary ionocytes, LAMP3$^+$ mregDC, TPEX, FOLR2$^+$ perivascular macrophages and the other populations enumerated in Table~S\ref{tab:heldout}) and the genuinely unannotated rare subpopulations to which we transfer the estimate. The assumption is formally stated in Lemma~\ref{lem:exchangeability} as a \emph{bound}: the held-out-rare loss rate upper-bounds the expected loss rate for unannotated rare subpopulations of comparable abundance, local geometry, and batch profile. We emphasise four boundary conditions under which the bound is not tight. First, the bound degrades when the unannotated population lies outside the span of the highly-variable genes used to construct the latent space; REAL cannot protect what the representation cannot see. Second, the bound degrades when the unannotated population has a qualitatively different batch-effect signature from the calibration populations---for instance, a disease-specific state observed in only one batch. Third, the bound is conservative for populations whose local density profile is not well-approximated by the level-set stratification we use; transitional cells along a continuum can appear as low-density but are not rare in the sampling sense. Fourth, our use of known rare populations as proxies for unknown ones inherits a selection bias: the populations we can hold out are those that have already been discovered, and the discovered set may not be exchangeable with the undiscovered set.

Our strongest defensible claim is therefore deliberately limited. REAL provides (i) a label-free, invariant-based detectability framework calibrated on held-out known rare subpopulations, (ii) an integration strategy (RareShield) that demonstrably reduces the held-out-rare loss rate at pan-cancer-atlas scale without degrading batch correction for abundant cell types, and (iii) an explicit exchangeability bound under which the held-out loss rate is an upper bound on the loss rate of truly unannotated rare subpopulations. A REAL-flagged motif $w \in \mathcal{W}$ is a claim that reproducible structure has been destroyed, not a claim that a new cell type exists; upgrading to an ontological claim requires the usual experimental burden (sorting, orthogonal markers, spatial validation, replication). The complementary negative: a dataset with low $|\mathcal{W}|$ under a candidate integration is evidence that no cross-batch-reproducible rare structure was erased, but not a guarantee that no such structure existed---by construction, we cannot rule out subpopulations whose reproducibility signal was below our sampling resolution or outside the evaluated feature span. These boundaries are not weaknesses of REAL specifically; they are the price of reasoning about what the evaluation framework cannot directly observe. Naming the price is what distinguishes an epistemically calibrated claim from a framework-captured one.

%   % ---------------------------------------------------------------------
% Philosophy's epistemic-frame head for the below-floor Limitations subsection.
% To be placed IMMEDIATELY BEFORE the Immunology-drafted biological disclosure
% (workspace/handoffs/immuno_limitations.tex at agent/immunology@b8a898d).
%
% Co-authored structure agreed with Immunology (agent-mozus3vv, 2026-05-10):
%   1. Philosophy epistemic frame (this file) — 2 short paragraphs, ≈220w.
%   2. Immunology biological disclosure (immuno_limitations.tex) — ≈350w.
%   3. Joint closing (telescoping-detection affirmative) — Immunology to draft
%      after pulling this file and Math's sub_compartment_pooling_bounds.md.
%
% Placement: sections/08_discussion.tex, as the subsection \subsection{Below the power floor},
% ordered after % Limitations & Epistemic Scope subsection (≤500w).
% For BioSci's locked slot in the Discussion. Cites Math Lemma S1-exchangeability as a *bound*.
% Covers: exchangeability failure modes, reference-class scope, non-ontic status, circularity mitigations.

\subsection*{Limitations and epistemic scope}
REAL relies on an exchangeability assumption between the known rare populations we hold out for calibration (pancreatic $\epsilon$ cells, pulmonary ionocytes, LAMP3$^+$ mregDC, TPEX, FOLR2$^+$ perivascular macrophages and the other populations enumerated in Table~S\ref{tab:heldout}) and the genuinely unannotated rare subpopulations to which we transfer the estimate. The assumption is formally stated in Lemma~\ref{lem:exchangeability} as a \emph{bound}: the held-out-rare loss rate upper-bounds the expected loss rate for unannotated rare subpopulations of comparable abundance, local geometry, and batch profile. We emphasise four boundary conditions under which the bound is not tight. First, the bound degrades when the unannotated population lies outside the span of the highly-variable genes used to construct the latent space; REAL cannot protect what the representation cannot see. Second, the bound degrades when the unannotated population has a qualitatively different batch-effect signature from the calibration populations---for instance, a disease-specific state observed in only one batch. Third, the bound is conservative for populations whose local density profile is not well-approximated by the level-set stratification we use; transitional cells along a continuum can appear as low-density but are not rare in the sampling sense. Fourth, our use of known rare populations as proxies for unknown ones inherits a selection bias: the populations we can hold out are those that have already been discovered, and the discovered set may not be exchangeable with the undiscovered set.

Our strongest defensible claim is therefore deliberately limited. REAL provides (i) a label-free, invariant-based detectability framework calibrated on held-out known rare subpopulations, (ii) an integration strategy (RareShield) that demonstrably reduces the held-out-rare loss rate at pan-cancer-atlas scale without degrading batch correction for abundant cell types, and (iii) an explicit exchangeability bound under which the held-out loss rate is an upper bound on the loss rate of truly unannotated rare subpopulations. A REAL-flagged motif $w \in \mathcal{W}$ is a claim that reproducible structure has been destroyed, not a claim that a new cell type exists; upgrading to an ontological claim requires the usual experimental burden (sorting, orthogonal markers, spatial validation, replication). The complementary negative: a dataset with low $|\mathcal{W}|$ under a candidate integration is evidence that no cross-batch-reproducible rare structure was erased, but not a guarantee that no such structure existed---by construction, we cannot rule out subpopulations whose reproducibility signal was below our sampling resolution or outside the evaluated feature span. These boundaries are not weaknesses of REAL specifically; they are the price of reasoning about what the evaluation framework cannot directly observe. Naming the price is what distinguishes an epistemically calibrated claim from a framework-captured one.
 and before % ---------------------------------------------------------------------
% Limitations subsection — immunology contribution to the Discussion.
% ~350 words. Honest disclosure of populations where REAL/RareShield
% cannot make reliable detection claims at Kang-atlas scale.
% Author: Immunology (agent-mozus3vv). Cross-references Math's
% predicted-detectability table (agent/mathematics@f564cc6).
% Target: sections/08_discussion.tex, as a \subsection after the
% "Implications for immunology" subsection (immuno_discussion.tex).
% ---------------------------------------------------------------------

\subsection{Boundaries of what our framework can currently guarantee}
\label{sec:discussion-limitations-immunology}

Three concrete limitations shape the claims we make. First, the Le~Cam
detectability envelope (\Cref{sec:minimax-detectability}) places a
sharp boundary on which rare populations can be recovered label-free
at any given atlas scale. At the Kang pan-cancer immune atlas
($n = 4.9 \times 10^6$, $B = 30$, per-batch heterogeneity
$\Lambda = 3$), populations with prevalence-separation product
$\pi\Delta < 10^{-3}\sigma$ fall below the detection floor, producing
a true-positive rate under 0.3 at $\alpha = 0.05$. Specifically,
\textit{CXCR5}$^+$\textit{TCF7}$^+$ progenitor-exhausted CD8 T
(TPEX), LAMP3$^+$ mregDC, KIR$^+$ adaptive NK, endothelial tip
cells, and drug-tolerant-persister tumour cells all sit in this
regime at Kang alone. For these populations our framework does not
claim detection from the Kang atlas in isolation; their recovery
requires either auxiliary protein or TCR signal, or pooling across
multiple atlases beyond $10^8$ cells. We disclose this explicitly
and present the populations instead as below-floor targets whose
detection-regime analysis (Extended Data~Fig.~1) itself is a
contribution.

Second, some rare populations are lost before integration rather than
during integration. Neutrophil subsets including LOX-1$^+$
PMN-MDSC and TAN-1 through TAN-5 subsets are systematically
undersampled by the standard 10x Genomics droplet-capture protocol,
and mast cells are similarly depleted relative to their tissue
abundance. Our framework is a post-integration detector and does not
and cannot recover populations that never entered the cell-by-gene
matrix. We flag this as a pre-integration identifiability problem
orthogonal to overcorrection, to be addressed by protocol choice
and by modality-aware capture strategies (FACS-sort enrichment prior
to library preparation; split-pool-barcoded whole-transcriptome
methods that capture neutrophils without selection bias).

Third, our evaluation uses annotated rare populations with their
labels hidden at test time. This is a legitimate form of ground
truth but it is not a direct test on unannotated populations.
Novel populations detected under our framework in the pan-cancer
atlas --- for example, a candidate CXCR5-dim TPEX sub-state
\citep{ourfig5} --- require the full evidence package of
\Cref{sec:evidence-package} plus independent cohort replication
before a definite claim of previously-unannotated biology. Main-text
detections are reported only when $\geq 4$ of the six evidence
criteria are satisfied.

\paragraph*{Telescoping detection recovers below-whole-atlas populations at sub-compartment scale.}
The $\pi\cdot\Delta^2 < 1.5\times10^{-4}$ below-floor rule above is
stated at the whole-atlas level and applies to flat REAL. Hierarchical
REAL, applied at lineage-restricted sub-compartments, moves the
operating point from whole-atlas prevalence to compartment-level
prevalence; at each level the flat floor still applies, but the
$\pi\cdot\Delta^2$ product is typically $10$--$100\times$ larger
because rare populations are often concentrated in specific lineages.
This recovers pulmonary ionocytes at the airway-epithelial sub-compartment
level, and the CXCR5-dim TPEX sub-state at the exhausted-CD8 sub-compartment
level (Supplementary Theorem~S1-hierarchical, Mathematics agent
\citealp{math2026hierarchical}). Our disclosure is therefore
three-tier: populations recovered by flat whole-atlas REAL (above
floor at $\pi\cdot\Delta^2 > 1.5\times10^{-4}$, e.g., MAIT, $\gamma\delta$
T, TREM2$^+$ LAM, myCAF, EMT-hybrid), populations recovered by
hierarchical sub-compartment REAL (above floor at sub-compartment
product, e.g., ionocytes, TPEX sub-state, AS-DC), and populations
unresolved at all compartment depths (endothelial tip cells, drug-tolerant
persister baseline state). Main-text claims are restricted to the
first two tiers; tier-three populations are disclosed as present in
the dataset but beyond the current framework's detection capability.

% ---------------------------------------------------------------------
% Requires: sec:minimax-detectability, sec:evidence-package,
% \citep{ourfig5} is a placeholder for the Fig 5 companion section /
% preprint. Biological Science may prefer \Cref{fig:fig5} or
% \Cref{sec:results-tpex-substate} depending on skeleton.
% ---------------------------------------------------------------------
.
%
% Anchors:
%   \Cref{thm:minimax}         Math Theorem 1 (minimax lower bound; floor)
%   \Cref{thm:labelblind}      Math Theorem 2.1 (label-based σ-algebra closure)
%   \Cref{lem:exchangeability} Math Lemma S1 (exchangeability bound)
%   \Cref{philo:limitations}   Philosophy Limitations subsection (head of §08 epistemic scope)
%   \Cref{sec:minimax-detectability}  Math's detectability envelope (for crosscheck)
% ---------------------------------------------------------------------

\subsection{Below the power floor: what the framework cannot yet certify}
\label{sec:below-floor-epistemic}

A framework's sensitivity is itself a piece of information about the framework, not about the biology. The minimax envelope of \Cref{thm:minimax} and the hierarchical-detectability result of \Cref{thm:S1-hierarchical} together partition the rare-population landscape at Kang~2024 scale into three strata: \emph{above-floor-flat} populations for which whole-atlas REAL is sufficient; \emph{hierarchical-recovered} populations which cross the detectability floor only after sub-compartment telescoping (per \Cref{sec:S1-envelope}); and \emph{truly-unresolved} populations for which neither flat nor hierarchical REAL attains the required power at Kang depth. The first two strata carry the manuscript's headline claim; this subsection states what we can and cannot say about the third. To read an absence of \textsc{REAL} flags on a truly-unresolved population as evidence of preservation is to commit exactly the failure-to-reject-as-success error that this paper is designed to expose (\Cref{philo:limitations}; Supp Note~S0 trap~8 at \Cref{sec:eightTraps}).

Two distinctions matter in making the disclosure precise. First, \textit{truly-unresolved at this dataset} is not the same as \textit{unresolved in principle}: the only escape routes are pooled multi-atlas scale (on the order of $10^{8}$ cells), auxiliary modalities that lift $\Delta$ above the $\pi \Delta^{2} \geq 1.5 \times 10^{-4}$ threshold at Kang scale (protein panels, TCR / BCR sequencing, spatial context), or targeted protocol changes that enrich the population before capture (FACS sorting, modality-aware library preparation). \textit{Truly unresolved} is thus a relation between a population, a measurement regime, and a framework, not a property of the population alone. Second, a motif flagged by \textsc{REAL} remains, even here, a claim about \emph{evaluation}, not about ontology (\Cref{philo:episto_note}). Correspondingly, the absence of a flag in the truly-unresolved stratum remains a claim about \emph{measurement insufficiency}, not about biology. The populations listed below are not excluded from single-cell biology; they are excluded from this framework's current certificate of preservation, under the stated regime.

%   % ---------------------------------------------------------------------
% Limitations subsection — immunology contribution to the Discussion.
% ~350 words. Honest disclosure of populations where REAL/RareShield
% cannot make reliable detection claims at Kang-atlas scale.
% Author: Immunology (agent-mozus3vv). Cross-references Math's
% predicted-detectability table (agent/mathematics@f564cc6).
% Target: sections/08_discussion.tex, as a \subsection after the
% "Implications for immunology" subsection (immuno_discussion.tex).
% ---------------------------------------------------------------------

\subsection{Boundaries of what our framework can currently guarantee}
\label{sec:discussion-limitations-immunology}

Three concrete limitations shape the claims we make. First, the Le~Cam
detectability envelope (\Cref{sec:minimax-detectability}) places a
sharp boundary on which rare populations can be recovered label-free
at any given atlas scale. At the Kang pan-cancer immune atlas
($n = 4.9 \times 10^6$, $B = 30$, per-batch heterogeneity
$\Lambda = 3$), populations with prevalence-separation product
$\pi\Delta < 10^{-3}\sigma$ fall below the detection floor, producing
a true-positive rate under 0.3 at $\alpha = 0.05$. Specifically,
\textit{CXCR5}$^+$\textit{TCF7}$^+$ progenitor-exhausted CD8 T
(TPEX), LAMP3$^+$ mregDC, KIR$^+$ adaptive NK, endothelial tip
cells, and drug-tolerant-persister tumour cells all sit in this
regime at Kang alone. For these populations our framework does not
claim detection from the Kang atlas in isolation; their recovery
requires either auxiliary protein or TCR signal, or pooling across
multiple atlases beyond $10^8$ cells. We disclose this explicitly
and present the populations instead as below-floor targets whose
detection-regime analysis (Extended Data~Fig.~1) itself is a
contribution.

Second, some rare populations are lost before integration rather than
during integration. Neutrophil subsets including LOX-1$^+$
PMN-MDSC and TAN-1 through TAN-5 subsets are systematically
undersampled by the standard 10x Genomics droplet-capture protocol,
and mast cells are similarly depleted relative to their tissue
abundance. Our framework is a post-integration detector and does not
and cannot recover populations that never entered the cell-by-gene
matrix. We flag this as a pre-integration identifiability problem
orthogonal to overcorrection, to be addressed by protocol choice
and by modality-aware capture strategies (FACS-sort enrichment prior
to library preparation; split-pool-barcoded whole-transcriptome
methods that capture neutrophils without selection bias).

Third, our evaluation uses annotated rare populations with their
labels hidden at test time. This is a legitimate form of ground
truth but it is not a direct test on unannotated populations.
Novel populations detected under our framework in the pan-cancer
atlas --- for example, a candidate CXCR5-dim TPEX sub-state
\citep{ourfig5} --- require the full evidence package of
\Cref{sec:evidence-package} plus independent cohort replication
before a definite claim of previously-unannotated biology. Main-text
detections are reported only when $\geq 4$ of the six evidence
criteria are satisfied.

\paragraph*{Telescoping detection recovers below-whole-atlas populations at sub-compartment scale.}
The $\pi\cdot\Delta^2 < 1.5\times10^{-4}$ below-floor rule above is
stated at the whole-atlas level and applies to flat REAL. Hierarchical
REAL, applied at lineage-restricted sub-compartments, moves the
operating point from whole-atlas prevalence to compartment-level
prevalence; at each level the flat floor still applies, but the
$\pi\cdot\Delta^2$ product is typically $10$--$100\times$ larger
because rare populations are often concentrated in specific lineages.
This recovers pulmonary ionocytes at the airway-epithelial sub-compartment
level, and the CXCR5-dim TPEX sub-state at the exhausted-CD8 sub-compartment
level (Supplementary Theorem~S1-hierarchical, Mathematics agent
\citealp{math2026hierarchical}). Our disclosure is therefore
three-tier: populations recovered by flat whole-atlas REAL (above
floor at $\pi\cdot\Delta^2 > 1.5\times10^{-4}$, e.g., MAIT, $\gamma\delta$
T, TREM2$^+$ LAM, myCAF, EMT-hybrid), populations recovered by
hierarchical sub-compartment REAL (above floor at sub-compartment
product, e.g., ionocytes, TPEX sub-state, AS-DC), and populations
unresolved at all compartment depths (endothelial tip cells, drug-tolerant
persister baseline state). Main-text claims are restricted to the
first two tiers; tier-three populations are disclosed as present in
the dataset but beyond the current framework's detection capability.

% ---------------------------------------------------------------------
% Requires: sec:minimax-detectability, sec:evidence-package,
% \citep{ourfig5} is a placeholder for the Fig 5 companion section /
% preprint. Biological Science may prefer \Cref{fig:fig5} or
% \Cref{sec:results-tpex-substate} depending on skeleton.
% ---------------------------------------------------------------------
         <-- unchanged from Immunology branch
%   % ---------------------------------------------------------------------
% Joint closing paragraph for the below-floor Limitations subsection.
% Co-authored with Philosophy (agent-mozur8ub). Placement flow:
%
%   % ---------------------------------------------------------------------
% Philosophy's epistemic-frame head for the below-floor Limitations subsection.
% To be placed IMMEDIATELY BEFORE the Immunology-drafted biological disclosure
% (workspace/handoffs/immuno_limitations.tex at agent/immunology@b8a898d).
%
% Co-authored structure agreed with Immunology (agent-mozus3vv, 2026-05-10):
%   1. Philosophy epistemic frame (this file) — 2 short paragraphs, ≈220w.
%   2. Immunology biological disclosure (immuno_limitations.tex) — ≈350w.
%   3. Joint closing (telescoping-detection affirmative) — Immunology to draft
%      after pulling this file and Math's sub_compartment_pooling_bounds.md.
%
% Placement: sections/08_discussion.tex, as the subsection \subsection{Below the power floor},
% ordered after \input{philo_limitations.tex} and before \input{immuno_limitations.tex}.
%
% Anchors:
%   \Cref{thm:minimax}         Math Theorem 1 (minimax lower bound; floor)
%   \Cref{thm:labelblind}      Math Theorem 2.1 (label-based σ-algebra closure)
%   \Cref{lem:exchangeability} Math Lemma S1 (exchangeability bound)
%   \Cref{philo:limitations}   Philosophy Limitations subsection (head of §08 epistemic scope)
%   \Cref{sec:minimax-detectability}  Math's detectability envelope (for crosscheck)
% ---------------------------------------------------------------------

\subsection{Below the power floor: what the framework cannot yet certify}
\label{sec:below-floor-epistemic}

A framework's sensitivity is itself a piece of information about the framework, not about the biology. The minimax envelope of \Cref{thm:minimax} and the hierarchical-detectability result of \Cref{thm:S1-hierarchical} together partition the rare-population landscape at Kang~2024 scale into three strata: \emph{above-floor-flat} populations for which whole-atlas REAL is sufficient; \emph{hierarchical-recovered} populations which cross the detectability floor only after sub-compartment telescoping (per \Cref{sec:S1-envelope}); and \emph{truly-unresolved} populations for which neither flat nor hierarchical REAL attains the required power at Kang depth. The first two strata carry the manuscript's headline claim; this subsection states what we can and cannot say about the third. To read an absence of \textsc{REAL} flags on a truly-unresolved population as evidence of preservation is to commit exactly the failure-to-reject-as-success error that this paper is designed to expose (\Cref{philo:limitations}; Supp Note~S0 trap~8 at \Cref{sec:eightTraps}).

Two distinctions matter in making the disclosure precise. First, \textit{truly-unresolved at this dataset} is not the same as \textit{unresolved in principle}: the only escape routes are pooled multi-atlas scale (on the order of $10^{8}$ cells), auxiliary modalities that lift $\Delta$ above the $\pi \Delta^{2} \geq 1.5 \times 10^{-4}$ threshold at Kang scale (protein panels, TCR / BCR sequencing, spatial context), or targeted protocol changes that enrich the population before capture (FACS sorting, modality-aware library preparation). \textit{Truly unresolved} is thus a relation between a population, a measurement regime, and a framework, not a property of the population alone. Second, a motif flagged by \textsc{REAL} remains, even here, a claim about \emph{evaluation}, not about ontology (\Cref{philo:episto_note}). Correspondingly, the absence of a flag in the truly-unresolved stratum remains a claim about \emph{measurement insufficiency}, not about biology. The populations listed below are not excluded from single-cell biology; they are excluded from this framework's current certificate of preservation, under the stated regime.
    (Philosophy's epistemic frame)
%   % ---------------------------------------------------------------------
% Limitations subsection — immunology contribution to the Discussion.
% ~350 words. Honest disclosure of populations where REAL/RareShield
% cannot make reliable detection claims at Kang-atlas scale.
% Author: Immunology (agent-mozus3vv). Cross-references Math's
% predicted-detectability table (agent/mathematics@f564cc6).
% Target: sections/08_discussion.tex, as a \subsection after the
% "Implications for immunology" subsection (immuno_discussion.tex).
% ---------------------------------------------------------------------

\subsection{Boundaries of what our framework can currently guarantee}
\label{sec:discussion-limitations-immunology}

Three concrete limitations shape the claims we make. First, the Le~Cam
detectability envelope (\Cref{sec:minimax-detectability}) places a
sharp boundary on which rare populations can be recovered label-free
at any given atlas scale. At the Kang pan-cancer immune atlas
($n = 4.9 \times 10^6$, $B = 30$, per-batch heterogeneity
$\Lambda = 3$), populations with prevalence-separation product
$\pi\Delta < 10^{-3}\sigma$ fall below the detection floor, producing
a true-positive rate under 0.3 at $\alpha = 0.05$. Specifically,
\textit{CXCR5}$^+$\textit{TCF7}$^+$ progenitor-exhausted CD8 T
(TPEX), LAMP3$^+$ mregDC, KIR$^+$ adaptive NK, endothelial tip
cells, and drug-tolerant-persister tumour cells all sit in this
regime at Kang alone. For these populations our framework does not
claim detection from the Kang atlas in isolation; their recovery
requires either auxiliary protein or TCR signal, or pooling across
multiple atlases beyond $10^8$ cells. We disclose this explicitly
and present the populations instead as below-floor targets whose
detection-regime analysis (Extended Data~Fig.~1) itself is a
contribution.

Second, some rare populations are lost before integration rather than
during integration. Neutrophil subsets including LOX-1$^+$
PMN-MDSC and TAN-1 through TAN-5 subsets are systematically
undersampled by the standard 10x Genomics droplet-capture protocol,
and mast cells are similarly depleted relative to their tissue
abundance. Our framework is a post-integration detector and does not
and cannot recover populations that never entered the cell-by-gene
matrix. We flag this as a pre-integration identifiability problem
orthogonal to overcorrection, to be addressed by protocol choice
and by modality-aware capture strategies (FACS-sort enrichment prior
to library preparation; split-pool-barcoded whole-transcriptome
methods that capture neutrophils without selection bias).

Third, our evaluation uses annotated rare populations with their
labels hidden at test time. This is a legitimate form of ground
truth but it is not a direct test on unannotated populations.
Novel populations detected under our framework in the pan-cancer
atlas --- for example, a candidate CXCR5-dim TPEX sub-state
\citep{ourfig5} --- require the full evidence package of
\Cref{sec:evidence-package} plus independent cohort replication
before a definite claim of previously-unannotated biology. Main-text
detections are reported only when $\geq 4$ of the six evidence
criteria are satisfied.

\paragraph*{Telescoping detection recovers below-whole-atlas populations at sub-compartment scale.}
The $\pi\cdot\Delta^2 < 1.5\times10^{-4}$ below-floor rule above is
stated at the whole-atlas level and applies to flat REAL. Hierarchical
REAL, applied at lineage-restricted sub-compartments, moves the
operating point from whole-atlas prevalence to compartment-level
prevalence; at each level the flat floor still applies, but the
$\pi\cdot\Delta^2$ product is typically $10$--$100\times$ larger
because rare populations are often concentrated in specific lineages.
This recovers pulmonary ionocytes at the airway-epithelial sub-compartment
level, and the CXCR5-dim TPEX sub-state at the exhausted-CD8 sub-compartment
level (Supplementary Theorem~S1-hierarchical, Mathematics agent
\citealp{math2026hierarchical}). Our disclosure is therefore
three-tier: populations recovered by flat whole-atlas REAL (above
floor at $\pi\cdot\Delta^2 > 1.5\times10^{-4}$, e.g., MAIT, $\gamma\delta$
T, TREM2$^+$ LAM, myCAF, EMT-hybrid), populations recovered by
hierarchical sub-compartment REAL (above floor at sub-compartment
product, e.g., ionocytes, TPEX sub-state, AS-DC), and populations
unresolved at all compartment depths (endothelial tip cells, drug-tolerant
persister baseline state). Main-text claims are restricted to the
first two tiers; tier-three populations are disclosed as present in
the dataset but beyond the current framework's detection capability.

% ---------------------------------------------------------------------
% Requires: sec:minimax-detectability, sec:evidence-package,
% \citep{ourfig5} is a placeholder for the Fig 5 companion section /
% preprint. Biological Science may prefer \Cref{fig:fig5} or
% \Cref{sec:results-tpex-substate} depending on skeleton.
% ---------------------------------------------------------------------
        (Immunology's biological disclosure)
%   % ---------------------------------------------------------------------
% Joint closing paragraph for the below-floor Limitations subsection.
% Co-authored with Philosophy (agent-mozur8ub). Placement flow:
%
%   \input{philo_below_floor_head.tex}    (Philosophy's epistemic frame)
%   \input{immuno_limitations.tex}        (Immunology's biological disclosure)
%   \input{immuno_below_floor_close.tex}  (THIS FILE — joint affirmative)
%
% Closes the subsection on the telescoping-detection affirmative:
% most below-flat-floor populations are recovered under hierarchical
% sub-compartment REAL.
% ---------------------------------------------------------------------

\paragraph*{Recovery under hierarchical application.}
Notwithstanding the strict boundary established above, the \emph{truly-unresolved} stratum is surprisingly small. Of the 22 immune and tumor populations in our registry whose whole-atlas product $\pi \cdot \Delta^{2}$ places them into one of the three detection tiers, 20 are recoverable through flat or hierarchical application of \textsc{REAL} at an appropriate compartment scale (\Cref{thm:S1-hierarchical}; \Cref{sec:methods-oracle-score} protocol): 7 populations are detected at whole-atlas scale (flat tier --- MAIT, $\gamma\delta$ T V$\delta1$-bias, TREM2$^{+}$ LAM, myCAF, iCAF, plasma-TLS, EMT-hybrid-near-pole), and 13 additional populations cross the detection floor only under hierarchical application at their natural sub-compartment scale (hierarchical tier). Pulmonary ionocytes pass detection at the airway-epithelial sub-compartment of HLCA (effective $\pi \approx 5 \times 10^{-3}$ in $n \approx 240{,}000$ airway cells); \textit{CXCR5}$^{+}$\textit{TCF7}$^{+}$ progenitor-exhausted CD8 T cells pass detection at the exhausted-CD8 sub-compartment of pooled melanoma atlases ($\pi \approx 10^{-2}$ in $n \approx 100{,}000$ cells); LAMP3$^{+}$ mregDC, AS-DC, and HCMV-imprinted adaptive NK pass detection at dendritic-cell and natural-killer sub-compartments respectively. The two populations that remain below the floor at every compartment depth are endothelial tip cells (which sit at the intersection of a small-mass vascular niche and a narrow transcriptional separation) and the drug-tolerant persister baseline state in oncology (Tumor Biology registry row TU-T-001), with mast cells in tumors additionally disclosed as out-of-scope for post-integration detection (mode~4 pre-integration capture loss). These unresolved entries are explicitly deferred to follow-up work: their recovery would require either the multi-atlas pooling regime ($n > 10^{8}$ cells) or auxiliary-modality augmentation, both of which lie beyond this paper's scope. This three-tier disclosure --- recovered-flat, recovered-hierarchical, deferred --- is the most honest reading of what our framework currently certifies, and it preserves the paper's headline claim on the above-floor and hierarchical-recovered strata while placing no claim on the truly-unresolved tier.

% ---------------------------------------------------------------------
% Cross-references needed (shared with philo_below_floor_head + immuno_limitations):
%   thm:S1-hierarchical        Math hierarchical REAL theorem
%   sec:methods-oracle-score   Methods — oracle-score verification
% Bib entries: none new.
% ---------------------------------------------------------------------
  (THIS FILE — joint affirmative)
%
% Closes the subsection on the telescoping-detection affirmative:
% most below-flat-floor populations are recovered under hierarchical
% sub-compartment REAL.
% ---------------------------------------------------------------------

\paragraph*{Recovery under hierarchical application.}
Notwithstanding the strict boundary established above, the \emph{truly-unresolved} stratum is surprisingly small. Of the 22 immune and tumor populations in our registry whose whole-atlas product $\pi \cdot \Delta^{2}$ places them into one of the three detection tiers, 20 are recoverable through flat or hierarchical application of \textsc{REAL} at an appropriate compartment scale (\Cref{thm:S1-hierarchical}; \Cref{sec:methods-oracle-score} protocol): 7 populations are detected at whole-atlas scale (flat tier --- MAIT, $\gamma\delta$ T V$\delta1$-bias, TREM2$^{+}$ LAM, myCAF, iCAF, plasma-TLS, EMT-hybrid-near-pole), and 13 additional populations cross the detection floor only under hierarchical application at their natural sub-compartment scale (hierarchical tier). Pulmonary ionocytes pass detection at the airway-epithelial sub-compartment of HLCA (effective $\pi \approx 5 \times 10^{-3}$ in $n \approx 240{,}000$ airway cells); \textit{CXCR5}$^{+}$\textit{TCF7}$^{+}$ progenitor-exhausted CD8 T cells pass detection at the exhausted-CD8 sub-compartment of pooled melanoma atlases ($\pi \approx 10^{-2}$ in $n \approx 100{,}000$ cells); LAMP3$^{+}$ mregDC, AS-DC, and HCMV-imprinted adaptive NK pass detection at dendritic-cell and natural-killer sub-compartments respectively. The two populations that remain below the floor at every compartment depth are endothelial tip cells (which sit at the intersection of a small-mass vascular niche and a narrow transcriptional separation) and the drug-tolerant persister baseline state in oncology (Tumor Biology registry row TU-T-001), with mast cells in tumors additionally disclosed as out-of-scope for post-integration detection (mode~4 pre-integration capture loss). These unresolved entries are explicitly deferred to follow-up work: their recovery would require either the multi-atlas pooling regime ($n > 10^{8}$ cells) or auxiliary-modality augmentation, both of which lie beyond this paper's scope. This three-tier disclosure --- recovered-flat, recovered-hierarchical, deferred --- is the most honest reading of what our framework currently certifies, and it preserves the paper's headline claim on the above-floor and hierarchical-recovered strata while placing no claim on the truly-unresolved tier.

% ---------------------------------------------------------------------
% Cross-references needed (shared with philo_below_floor_head + immuno_limitations):
%   thm:S1-hierarchical        Math hierarchical REAL theorem
%   sec:methods-oracle-score   Methods — oracle-score verification
% Bib entries: none new.
% ---------------------------------------------------------------------
   <-- unchanged from Immunology branch
%
% Use this file INSTEAD OF the 5-file chain if BioSci prefers a single splice point.
% Immunology's two \input calls (limitations + below-floor-close) are PRESERVED as
% \input calls inside this file rather than inlined, because they are owned by
% the immunology branch and Immunology should retain edit control over their content.
%
% Placement: sections/08_discussion.tex, replacing the 5-file \input chain with
% a single % Philosophy §08 epistemic-scope block — single-file consolidation.
% Merges the five \input calls that BioSci currently threads through into one file
% for easier future-edit management. Semantically identical to the input chain:
%   % Discussion paragraph (≈250w) — epistemic status of REAL-flagged motifs.
% For task t-mozv8oul. Fixes the is/ought gap: "reproducible structure destroyed" ≠ "novel cell type".
% Guards the team from over-claiming. Anchors: motif $w \in \mathcal{W}$; LRS(w); RareShield; Hacking-style intervention realism.

\paragraph*{What a REAL flag is, and is not.}
A motif flagged by REAL is the claim that a local, reproducible geometric structure---witnessed across independent batches before integration, with measurable local-mass distortion after---has been erased by the integration map beyond what admissible batch-effect deformations can explain. This is a claim about \emph{evaluation}, not about ontology. It is not a claim that a new cell type exists. A motif $w \in \mathcal{W}$ with low $\mathrm{LRS}(w)$ licenses exactly one inference: the integration has discarded reproducible structure of the kind a biologist would want to interrogate. Whether that structure corresponds to a novel cell type, a transitional state along an already-known trajectory, a rare disease-specific configuration, or a technical sub-mode that merely looks biological, is a further question that REAL does not answer. Upgrading a motif to a cell-type claim requires the usual burden of single-cell biology: independent sampling, orthogonal markers, targeted protocol (sorting, spatial validation, experimental perturbation), and community replication. REAL's value is that it identifies \emph{where to look}, not what has been found. Conflating the two would re-create the very error this paper is designed to expose: treating a framework-internal signal as a fact about the world. Our strong claim is therefore deliberately asymmetric: a high motif set $|\mathcal{W}|$ under a candidate integration is evidence of measurement insufficiency; a low motif set is evidence that the integration respected the syntactic reproducibility criterion. Ontology is downstream of both.

%   % Limitations & Epistemic Scope subsection (≤500w).
% For BioSci's locked slot in the Discussion. Cites Math Lemma S1-exchangeability as a *bound*.
% Covers: exchangeability failure modes, reference-class scope, non-ontic status, circularity mitigations.

\subsection*{Limitations and epistemic scope}
REAL relies on an exchangeability assumption between the known rare populations we hold out for calibration (pancreatic $\epsilon$ cells, pulmonary ionocytes, LAMP3$^+$ mregDC, TPEX, FOLR2$^+$ perivascular macrophages and the other populations enumerated in Table~S\ref{tab:heldout}) and the genuinely unannotated rare subpopulations to which we transfer the estimate. The assumption is formally stated in Lemma~\ref{lem:exchangeability} as a \emph{bound}: the held-out-rare loss rate upper-bounds the expected loss rate for unannotated rare subpopulations of comparable abundance, local geometry, and batch profile. We emphasise four boundary conditions under which the bound is not tight. First, the bound degrades when the unannotated population lies outside the span of the highly-variable genes used to construct the latent space; REAL cannot protect what the representation cannot see. Second, the bound degrades when the unannotated population has a qualitatively different batch-effect signature from the calibration populations---for instance, a disease-specific state observed in only one batch. Third, the bound is conservative for populations whose local density profile is not well-approximated by the level-set stratification we use; transitional cells along a continuum can appear as low-density but are not rare in the sampling sense. Fourth, our use of known rare populations as proxies for unknown ones inherits a selection bias: the populations we can hold out are those that have already been discovered, and the discovered set may not be exchangeable with the undiscovered set.

Our strongest defensible claim is therefore deliberately limited. REAL provides (i) a label-free, invariant-based detectability framework calibrated on held-out known rare subpopulations, (ii) an integration strategy (RareShield) that demonstrably reduces the held-out-rare loss rate at pan-cancer-atlas scale without degrading batch correction for abundant cell types, and (iii) an explicit exchangeability bound under which the held-out loss rate is an upper bound on the loss rate of truly unannotated rare subpopulations. A REAL-flagged motif $w \in \mathcal{W}$ is a claim that reproducible structure has been destroyed, not a claim that a new cell type exists; upgrading to an ontological claim requires the usual experimental burden (sorting, orthogonal markers, spatial validation, replication). The complementary negative: a dataset with low $|\mathcal{W}|$ under a candidate integration is evidence that no cross-batch-reproducible rare structure was erased, but not a guarantee that no such structure existed---by construction, we cannot rule out subpopulations whose reproducibility signal was below our sampling resolution or outside the evaluated feature span. These boundaries are not weaknesses of REAL specifically; they are the price of reasoning about what the evaluation framework cannot directly observe. Naming the price is what distinguishes an epistemically calibrated claim from a framework-captured one.

%   % ---------------------------------------------------------------------
% Philosophy's epistemic-frame head for the below-floor Limitations subsection.
% To be placed IMMEDIATELY BEFORE the Immunology-drafted biological disclosure
% (workspace/handoffs/immuno_limitations.tex at agent/immunology@b8a898d).
%
% Co-authored structure agreed with Immunology (agent-mozus3vv, 2026-05-10):
%   1. Philosophy epistemic frame (this file) — 2 short paragraphs, ≈220w.
%   2. Immunology biological disclosure (immuno_limitations.tex) — ≈350w.
%   3. Joint closing (telescoping-detection affirmative) — Immunology to draft
%      after pulling this file and Math's sub_compartment_pooling_bounds.md.
%
% Placement: sections/08_discussion.tex, as the subsection \subsection{Below the power floor},
% ordered after \input{philo_limitations.tex} and before \input{immuno_limitations.tex}.
%
% Anchors:
%   \Cref{thm:minimax}         Math Theorem 1 (minimax lower bound; floor)
%   \Cref{thm:labelblind}      Math Theorem 2.1 (label-based σ-algebra closure)
%   \Cref{lem:exchangeability} Math Lemma S1 (exchangeability bound)
%   \Cref{philo:limitations}   Philosophy Limitations subsection (head of §08 epistemic scope)
%   \Cref{sec:minimax-detectability}  Math's detectability envelope (for crosscheck)
% ---------------------------------------------------------------------

\subsection{Below the power floor: what the framework cannot yet certify}
\label{sec:below-floor-epistemic}

A framework's sensitivity is itself a piece of information about the framework, not about the biology. The minimax envelope of \Cref{thm:minimax} and the hierarchical-detectability result of \Cref{thm:S1-hierarchical} together partition the rare-population landscape at Kang~2024 scale into three strata: \emph{above-floor-flat} populations for which whole-atlas REAL is sufficient; \emph{hierarchical-recovered} populations which cross the detectability floor only after sub-compartment telescoping (per \Cref{sec:S1-envelope}); and \emph{truly-unresolved} populations for which neither flat nor hierarchical REAL attains the required power at Kang depth. The first two strata carry the manuscript's headline claim; this subsection states what we can and cannot say about the third. To read an absence of \textsc{REAL} flags on a truly-unresolved population as evidence of preservation is to commit exactly the failure-to-reject-as-success error that this paper is designed to expose (\Cref{philo:limitations}; Supp Note~S0 trap~8 at \Cref{sec:eightTraps}).

Two distinctions matter in making the disclosure precise. First, \textit{truly-unresolved at this dataset} is not the same as \textit{unresolved in principle}: the only escape routes are pooled multi-atlas scale (on the order of $10^{8}$ cells), auxiliary modalities that lift $\Delta$ above the $\pi \Delta^{2} \geq 1.5 \times 10^{-4}$ threshold at Kang scale (protein panels, TCR / BCR sequencing, spatial context), or targeted protocol changes that enrich the population before capture (FACS sorting, modality-aware library preparation). \textit{Truly unresolved} is thus a relation between a population, a measurement regime, and a framework, not a property of the population alone. Second, a motif flagged by \textsc{REAL} remains, even here, a claim about \emph{evaluation}, not about ontology (\Cref{philo:episto_note}). Correspondingly, the absence of a flag in the truly-unresolved stratum remains a claim about \emph{measurement insufficiency}, not about biology. The populations listed below are not excluded from single-cell biology; they are excluded from this framework's current certificate of preservation, under the stated regime.

%   % ---------------------------------------------------------------------
% Limitations subsection — immunology contribution to the Discussion.
% ~350 words. Honest disclosure of populations where REAL/RareShield
% cannot make reliable detection claims at Kang-atlas scale.
% Author: Immunology (agent-mozus3vv). Cross-references Math's
% predicted-detectability table (agent/mathematics@f564cc6).
% Target: sections/08_discussion.tex, as a \subsection after the
% "Implications for immunology" subsection (immuno_discussion.tex).
% ---------------------------------------------------------------------

\subsection{Boundaries of what our framework can currently guarantee}
\label{sec:discussion-limitations-immunology}

Three concrete limitations shape the claims we make. First, the Le~Cam
detectability envelope (\Cref{sec:minimax-detectability}) places a
sharp boundary on which rare populations can be recovered label-free
at any given atlas scale. At the Kang pan-cancer immune atlas
($n = 4.9 \times 10^6$, $B = 30$, per-batch heterogeneity
$\Lambda = 3$), populations with prevalence-separation product
$\pi\Delta < 10^{-3}\sigma$ fall below the detection floor, producing
a true-positive rate under 0.3 at $\alpha = 0.05$. Specifically,
\textit{CXCR5}$^+$\textit{TCF7}$^+$ progenitor-exhausted CD8 T
(TPEX), LAMP3$^+$ mregDC, KIR$^+$ adaptive NK, endothelial tip
cells, and drug-tolerant-persister tumour cells all sit in this
regime at Kang alone. For these populations our framework does not
claim detection from the Kang atlas in isolation; their recovery
requires either auxiliary protein or TCR signal, or pooling across
multiple atlases beyond $10^8$ cells. We disclose this explicitly
and present the populations instead as below-floor targets whose
detection-regime analysis (Extended Data~Fig.~1) itself is a
contribution.

Second, some rare populations are lost before integration rather than
during integration. Neutrophil subsets including LOX-1$^+$
PMN-MDSC and TAN-1 through TAN-5 subsets are systematically
undersampled by the standard 10x Genomics droplet-capture protocol,
and mast cells are similarly depleted relative to their tissue
abundance. Our framework is a post-integration detector and does not
and cannot recover populations that never entered the cell-by-gene
matrix. We flag this as a pre-integration identifiability problem
orthogonal to overcorrection, to be addressed by protocol choice
and by modality-aware capture strategies (FACS-sort enrichment prior
to library preparation; split-pool-barcoded whole-transcriptome
methods that capture neutrophils without selection bias).

Third, our evaluation uses annotated rare populations with their
labels hidden at test time. This is a legitimate form of ground
truth but it is not a direct test on unannotated populations.
Novel populations detected under our framework in the pan-cancer
atlas --- for example, a candidate CXCR5-dim TPEX sub-state
\citep{ourfig5} --- require the full evidence package of
\Cref{sec:evidence-package} plus independent cohort replication
before a definite claim of previously-unannotated biology. Main-text
detections are reported only when $\geq 4$ of the six evidence
criteria are satisfied.

\paragraph*{Telescoping detection recovers below-whole-atlas populations at sub-compartment scale.}
The $\pi\cdot\Delta^2 < 1.5\times10^{-4}$ below-floor rule above is
stated at the whole-atlas level and applies to flat REAL. Hierarchical
REAL, applied at lineage-restricted sub-compartments, moves the
operating point from whole-atlas prevalence to compartment-level
prevalence; at each level the flat floor still applies, but the
$\pi\cdot\Delta^2$ product is typically $10$--$100\times$ larger
because rare populations are often concentrated in specific lineages.
This recovers pulmonary ionocytes at the airway-epithelial sub-compartment
level, and the CXCR5-dim TPEX sub-state at the exhausted-CD8 sub-compartment
level (Supplementary Theorem~S1-hierarchical, Mathematics agent
\citealp{math2026hierarchical}). Our disclosure is therefore
three-tier: populations recovered by flat whole-atlas REAL (above
floor at $\pi\cdot\Delta^2 > 1.5\times10^{-4}$, e.g., MAIT, $\gamma\delta$
T, TREM2$^+$ LAM, myCAF, EMT-hybrid), populations recovered by
hierarchical sub-compartment REAL (above floor at sub-compartment
product, e.g., ionocytes, TPEX sub-state, AS-DC), and populations
unresolved at all compartment depths (endothelial tip cells, drug-tolerant
persister baseline state). Main-text claims are restricted to the
first two tiers; tier-three populations are disclosed as present in
the dataset but beyond the current framework's detection capability.

% ---------------------------------------------------------------------
% Requires: sec:minimax-detectability, sec:evidence-package,
% \citep{ourfig5} is a placeholder for the Fig 5 companion section /
% preprint. Biological Science may prefer \Cref{fig:fig5} or
% \Cref{sec:results-tpex-substate} depending on skeleton.
% ---------------------------------------------------------------------
         <-- unchanged from Immunology branch
%   % ---------------------------------------------------------------------
% Joint closing paragraph for the below-floor Limitations subsection.
% Co-authored with Philosophy (agent-mozur8ub). Placement flow:
%
%   \input{philo_below_floor_head.tex}    (Philosophy's epistemic frame)
%   \input{immuno_limitations.tex}        (Immunology's biological disclosure)
%   \input{immuno_below_floor_close.tex}  (THIS FILE — joint affirmative)
%
% Closes the subsection on the telescoping-detection affirmative:
% most below-flat-floor populations are recovered under hierarchical
% sub-compartment REAL.
% ---------------------------------------------------------------------

\paragraph*{Recovery under hierarchical application.}
Notwithstanding the strict boundary established above, the \emph{truly-unresolved} stratum is surprisingly small. Of the 22 immune and tumor populations in our registry whose whole-atlas product $\pi \cdot \Delta^{2}$ places them into one of the three detection tiers, 20 are recoverable through flat or hierarchical application of \textsc{REAL} at an appropriate compartment scale (\Cref{thm:S1-hierarchical}; \Cref{sec:methods-oracle-score} protocol): 7 populations are detected at whole-atlas scale (flat tier --- MAIT, $\gamma\delta$ T V$\delta1$-bias, TREM2$^{+}$ LAM, myCAF, iCAF, plasma-TLS, EMT-hybrid-near-pole), and 13 additional populations cross the detection floor only under hierarchical application at their natural sub-compartment scale (hierarchical tier). Pulmonary ionocytes pass detection at the airway-epithelial sub-compartment of HLCA (effective $\pi \approx 5 \times 10^{-3}$ in $n \approx 240{,}000$ airway cells); \textit{CXCR5}$^{+}$\textit{TCF7}$^{+}$ progenitor-exhausted CD8 T cells pass detection at the exhausted-CD8 sub-compartment of pooled melanoma atlases ($\pi \approx 10^{-2}$ in $n \approx 100{,}000$ cells); LAMP3$^{+}$ mregDC, AS-DC, and HCMV-imprinted adaptive NK pass detection at dendritic-cell and natural-killer sub-compartments respectively. The two populations that remain below the floor at every compartment depth are endothelial tip cells (which sit at the intersection of a small-mass vascular niche and a narrow transcriptional separation) and the drug-tolerant persister baseline state in oncology (Tumor Biology registry row TU-T-001), with mast cells in tumors additionally disclosed as out-of-scope for post-integration detection (mode~4 pre-integration capture loss). These unresolved entries are explicitly deferred to follow-up work: their recovery would require either the multi-atlas pooling regime ($n > 10^{8}$ cells) or auxiliary-modality augmentation, both of which lie beyond this paper's scope. This three-tier disclosure --- recovered-flat, recovered-hierarchical, deferred --- is the most honest reading of what our framework currently certifies, and it preserves the paper's headline claim on the above-floor and hierarchical-recovered strata while placing no claim on the truly-unresolved tier.

% ---------------------------------------------------------------------
% Cross-references needed (shared with philo_below_floor_head + immuno_limitations):
%   thm:S1-hierarchical        Math hierarchical REAL theorem
%   sec:methods-oracle-score   Methods — oracle-score verification
% Bib entries: none new.
% ---------------------------------------------------------------------
   <-- unchanged from Immunology branch
%
% Use this file INSTEAD OF the 5-file chain if BioSci prefers a single splice point.
% Immunology's two \input calls (limitations + below-floor-close) are PRESERVED as
% \input calls inside this file rather than inlined, because they are owned by
% the immunology branch and Immunology should retain edit control over their content.
%
% Placement: sections/08_discussion.tex, replacing the 5-file \input chain with
% a single % Philosophy §08 epistemic-scope block — single-file consolidation.
% Merges the five \input calls that BioSci currently threads through into one file
% for easier future-edit management. Semantically identical to the input chain:
%   \input{handoffs/philo_discussion_meno.tex}
%   \input{handoffs/philo_limitations.tex}
%   \input{handoffs/philo_below_floor_head.tex}
%   \input{handoffs/immuno_limitations.tex}         <-- unchanged from Immunology branch
%   \input{handoffs/immuno_below_floor_close.tex}   <-- unchanged from Immunology branch
%
% Use this file INSTEAD OF the 5-file chain if BioSci prefers a single splice point.
% Immunology's two \input calls (limitations + below-floor-close) are PRESERVED as
% \input calls inside this file rather than inlined, because they are owned by
% the immunology branch and Immunology should retain edit control over their content.
%
% Placement: sections/08_discussion.tex, replacing the 5-file \input chain with
% a single \input{handoffs/philo_08_epistemic_block.tex}.
%
% Author: Philosophy (agent-mozur8ub). Last updated for v2.4 pre-submission.

% =============================================================================
% §08 EPISTEMIC SCOPE BLOCK — single-file consolidation
% =============================================================================

% -----------------------------------------------------------------------------
% (1) Meno-problem framing — opens the epistemic scope block.
% Content from philo_discussion_meno.tex (agent/philosophy@<head>).
% ≈250w. Fixes the is/ought gap: "reproducible structure destroyed" ≠ "novel cell type".
% -----------------------------------------------------------------------------

\paragraph*{What a REAL flag is, and is not.}
A motif flagged by REAL is the claim that a local, reproducible geometric structure---witnessed across independent batches before integration, with measurable local-mass distortion after---has been erased by the integration map beyond what admissible batch-effect deformations can explain. This is a claim about \emph{evaluation}, not about ontology. It is not a claim that a new cell type exists. A motif $w \in \mathcal{W}$ with low $\mathrm{LRS}(w)$ licenses exactly one inference: the integration has discarded reproducible structure of the kind a biologist would want to interrogate. Whether that structure corresponds to a novel cell type, a transitional state along an already-known trajectory, a rare disease-specific configuration, or a technical sub-mode that merely looks biological, is a further question REAL does not answer. Upgrading a motif to a cell-type claim requires the usual burden of single-cell biology: independent sampling, orthogonal markers, targeted protocol (sorting, spatial validation, experimental perturbation), and community replication. REAL's value is that it identifies \emph{where to look}, not what has been found. Conflating the two would recreate the very error this paper is designed to expose: treating a framework-internal signal as a fact about the world. Our strong claim is therefore deliberately asymmetric. A high motif set $|\mathcal{W}|$ under a candidate integration is evidence of measurement insufficiency; a low motif set is evidence that the integration respected the syntactic reproducibility criterion. Ontology is downstream of both.

% -----------------------------------------------------------------------------
% (2) Limitations and epistemic scope — Route-C exchangeability boundaries.
% Content from philo_limitations.tex. ≤500w.
% -----------------------------------------------------------------------------

\subsection*{Limitations and epistemic scope}
\label{philo:limitations}
REAL relies on an exchangeability assumption between the known rare populations we hold out for calibration (pancreatic $\epsilon$ cells, pulmonary ionocytes, LAMP3$^+$ mregDC, TPEX, FOLR2$^+$ perivascular macrophages, apCAF, and the populations enumerated in Table~S\ref{tab:heldout}) and the genuinely unannotated rare subpopulations to which we transfer the estimate. The assumption is formally stated in Lemma~\ref{lem:exchangeability} as a \emph{bound}: the held-out-rare loss rate upper-bounds the expected loss rate for unannotated rare subpopulations of comparable abundance, local geometry, and batch profile. We emphasise four boundary conditions under which the bound is not tight. First, the bound degrades when the unannotated population lies outside the span of the highly-variable genes used to construct the latent space; REAL cannot protect what the representation cannot see. Second, the bound degrades when the unannotated population has a qualitatively different batch-effect signature from the calibration populations---for instance, a disease-specific state observed in only one batch. Third, the bound is conservative for populations whose local density profile is not well-approximated by the level-set stratification we use; transitional cells along a continuum can appear as low-density but are not rare in the sampling sense. Fourth, our use of known rare populations as proxies for unknown ones inherits a selection bias: the populations we can hold out are those that have already been discovered, and the discovered set may not be exchangeable with the undiscovered set.

Our strongest defensible claim is therefore deliberately limited. REAL provides (i) a label-free, invariant-based detectability framework calibrated on held-out known rare subpopulations, (ii) an integration strategy (RareShield) that demonstrably reduces the held-out-rare loss rate at pan-cancer-atlas scale without degrading batch correction for abundant cell types, and (iii) an explicit exchangeability bound under which the held-out loss rate is an upper bound on the loss rate of truly unannotated rare subpopulations. A REAL-flagged motif $w \in \mathcal{W}$ is a claim that reproducible structure has been destroyed, not a claim that a new cell type exists; upgrading to an ontological claim requires the usual experimental burden (sorting, orthogonal markers, spatial validation, replication). The complementary negative: a dataset with low $|\mathcal{W}|$ under a candidate integration is evidence that no cross-batch-reproducible rare structure was erased, but not a guarantee that no such structure existed---by construction, we cannot rule out subpopulations whose reproducibility signal was below our sampling resolution or outside the evaluated feature span. These boundaries are not weaknesses of REAL specifically; they are the price of reasoning about what the evaluation framework cannot directly observe. Naming the price is what distinguishes an epistemically calibrated claim from a framework-captured one.

% -----------------------------------------------------------------------------
% (3) Below the power floor — epistemic frame head (three-tier stratification).
% Content from philo_below_floor_head.tex.
% -----------------------------------------------------------------------------

\subsection*{Below the power floor: what the framework cannot yet certify}
\label{sec:below-floor-epistemic}

A framework's sensitivity is itself a piece of information about the framework, not about the biology. The minimax envelope of \Cref{thm:minimax} and the hierarchical-detectability result of \Cref{thm:S1-hierarchical} together partition the rare-population landscape at Kang~2024 scale into three strata: \emph{above-floor-flat} populations for which whole-atlas REAL is sufficient; \emph{hierarchical-recovered} populations which cross the detectability floor only after sub-compartment telescoping (per \Cref{sec:S1-envelope}); and \emph{truly-unresolved} populations for which neither flat nor hierarchical REAL attains the required power at Kang depth. The first two strata carry the manuscript's headline claim; this subsection states what we can and cannot say about the third. To read an absence of \textsc{REAL} flags on a truly-unresolved population as evidence of preservation is to commit exactly the failure-to-reject-as-success error that this paper is designed to expose (\Cref{philo:limitations}; Supp Note~S0 trap~8 at \Cref{sec:eightTraps}).

Two distinctions matter in making the disclosure precise. First, \textit{truly-unresolved at this dataset} is not the same as \textit{unresolved in principle}: the only escape routes are pooled multi-atlas scale (on the order of $10^{8}$ cells), auxiliary modalities that lift $\Delta$ above the $\pi \Delta^{2} \geq 1.5 \times 10^{-4}$ threshold at Kang scale (protein panels, TCR / BCR sequencing, spatial context), or targeted protocol changes that enrich the population before capture (FACS sorting, modality-aware library preparation). \textit{Truly unresolved} is thus a relation between a population, a measurement regime, and a framework, not a property of the population alone. Second, a motif flagged by \textsc{REAL} remains, even here, a claim about \emph{evaluation}, not about ontology (\Cref{philo:episto_note}). Correspondingly, the absence of a flag in the truly-unresolved stratum remains a claim about \emph{measurement insufficiency}, not about biology. The populations listed below are not excluded from single-cell biology; they are excluded from this framework's current certificate of preservation, under the stated regime.

% -----------------------------------------------------------------------------
% (4)-(5) Immunology's biological disclosure + joint closing.
% These files live in the immunology branch; Immunology retains edit control.
% They are \input'd here rather than inlined, so Immunology-branch edits propagate
% automatically on the next BioSci pull.
% -----------------------------------------------------------------------------

\input{handoffs/immuno_limitations.tex}
\input{handoffs/immuno_below_floor_close.tex}

% =============================================================================
% END §08 EPISTEMIC SCOPE BLOCK
% =============================================================================

% Integration note for BioSci:
%   - To use this consolidated file: replace the 5-file \input chain in
%     sections/08_discussion.tex with:
%       \input{handoffs/philo_08_epistemic_block.tex}
%   - To keep the 5-file chain: ignore this file; it is semantically identical
%     to the chain at this commit. Consolidation is a convenience, not a correctness fix.
%   - If Philosophy or Immunology updates any of the five source files, re-generating
%     this consolidated file requires Philosophy to re-splice; the \input calls for
%     Immunology's two files remain untouched.
%
% Lint status: passes lint:epistemic clean (trap 8 / trap 4 / IM.4 guards all present).
.
%
% Author: Philosophy (agent-mozur8ub). Last updated for v2.4 pre-submission.

% =============================================================================
% §08 EPISTEMIC SCOPE BLOCK — single-file consolidation
% =============================================================================

% -----------------------------------------------------------------------------
% (1) Meno-problem framing — opens the epistemic scope block.
% Content from philo_discussion_meno.tex (agent/philosophy@<head>).
% ≈250w. Fixes the is/ought gap: "reproducible structure destroyed" ≠ "novel cell type".
% -----------------------------------------------------------------------------

\paragraph*{What a REAL flag is, and is not.}
A motif flagged by REAL is the claim that a local, reproducible geometric structure---witnessed across independent batches before integration, with measurable local-mass distortion after---has been erased by the integration map beyond what admissible batch-effect deformations can explain. This is a claim about \emph{evaluation}, not about ontology. It is not a claim that a new cell type exists. A motif $w \in \mathcal{W}$ with low $\mathrm{LRS}(w)$ licenses exactly one inference: the integration has discarded reproducible structure of the kind a biologist would want to interrogate. Whether that structure corresponds to a novel cell type, a transitional state along an already-known trajectory, a rare disease-specific configuration, or a technical sub-mode that merely looks biological, is a further question REAL does not answer. Upgrading a motif to a cell-type claim requires the usual burden of single-cell biology: independent sampling, orthogonal markers, targeted protocol (sorting, spatial validation, experimental perturbation), and community replication. REAL's value is that it identifies \emph{where to look}, not what has been found. Conflating the two would recreate the very error this paper is designed to expose: treating a framework-internal signal as a fact about the world. Our strong claim is therefore deliberately asymmetric. A high motif set $|\mathcal{W}|$ under a candidate integration is evidence of measurement insufficiency; a low motif set is evidence that the integration respected the syntactic reproducibility criterion. Ontology is downstream of both.

% -----------------------------------------------------------------------------
% (2) Limitations and epistemic scope — Route-C exchangeability boundaries.
% Content from philo_limitations.tex. ≤500w.
% -----------------------------------------------------------------------------

\subsection*{Limitations and epistemic scope}
\label{philo:limitations}
REAL relies on an exchangeability assumption between the known rare populations we hold out for calibration (pancreatic $\epsilon$ cells, pulmonary ionocytes, LAMP3$^+$ mregDC, TPEX, FOLR2$^+$ perivascular macrophages, apCAF, and the populations enumerated in Table~S\ref{tab:heldout}) and the genuinely unannotated rare subpopulations to which we transfer the estimate. The assumption is formally stated in Lemma~\ref{lem:exchangeability} as a \emph{bound}: the held-out-rare loss rate upper-bounds the expected loss rate for unannotated rare subpopulations of comparable abundance, local geometry, and batch profile. We emphasise four boundary conditions under which the bound is not tight. First, the bound degrades when the unannotated population lies outside the span of the highly-variable genes used to construct the latent space; REAL cannot protect what the representation cannot see. Second, the bound degrades when the unannotated population has a qualitatively different batch-effect signature from the calibration populations---for instance, a disease-specific state observed in only one batch. Third, the bound is conservative for populations whose local density profile is not well-approximated by the level-set stratification we use; transitional cells along a continuum can appear as low-density but are not rare in the sampling sense. Fourth, our use of known rare populations as proxies for unknown ones inherits a selection bias: the populations we can hold out are those that have already been discovered, and the discovered set may not be exchangeable with the undiscovered set.

Our strongest defensible claim is therefore deliberately limited. REAL provides (i) a label-free, invariant-based detectability framework calibrated on held-out known rare subpopulations, (ii) an integration strategy (RareShield) that demonstrably reduces the held-out-rare loss rate at pan-cancer-atlas scale without degrading batch correction for abundant cell types, and (iii) an explicit exchangeability bound under which the held-out loss rate is an upper bound on the loss rate of truly unannotated rare subpopulations. A REAL-flagged motif $w \in \mathcal{W}$ is a claim that reproducible structure has been destroyed, not a claim that a new cell type exists; upgrading to an ontological claim requires the usual experimental burden (sorting, orthogonal markers, spatial validation, replication). The complementary negative: a dataset with low $|\mathcal{W}|$ under a candidate integration is evidence that no cross-batch-reproducible rare structure was erased, but not a guarantee that no such structure existed---by construction, we cannot rule out subpopulations whose reproducibility signal was below our sampling resolution or outside the evaluated feature span. These boundaries are not weaknesses of REAL specifically; they are the price of reasoning about what the evaluation framework cannot directly observe. Naming the price is what distinguishes an epistemically calibrated claim from a framework-captured one.

% -----------------------------------------------------------------------------
% (3) Below the power floor — epistemic frame head (three-tier stratification).
% Content from philo_below_floor_head.tex.
% -----------------------------------------------------------------------------

\subsection*{Below the power floor: what the framework cannot yet certify}
\label{sec:below-floor-epistemic}

A framework's sensitivity is itself a piece of information about the framework, not about the biology. The minimax envelope of \Cref{thm:minimax} and the hierarchical-detectability result of \Cref{thm:S1-hierarchical} together partition the rare-population landscape at Kang~2024 scale into three strata: \emph{above-floor-flat} populations for which whole-atlas REAL is sufficient; \emph{hierarchical-recovered} populations which cross the detectability floor only after sub-compartment telescoping (per \Cref{sec:S1-envelope}); and \emph{truly-unresolved} populations for which neither flat nor hierarchical REAL attains the required power at Kang depth. The first two strata carry the manuscript's headline claim; this subsection states what we can and cannot say about the third. To read an absence of \textsc{REAL} flags on a truly-unresolved population as evidence of preservation is to commit exactly the failure-to-reject-as-success error that this paper is designed to expose (\Cref{philo:limitations}; Supp Note~S0 trap~8 at \Cref{sec:eightTraps}).

Two distinctions matter in making the disclosure precise. First, \textit{truly-unresolved at this dataset} is not the same as \textit{unresolved in principle}: the only escape routes are pooled multi-atlas scale (on the order of $10^{8}$ cells), auxiliary modalities that lift $\Delta$ above the $\pi \Delta^{2} \geq 1.5 \times 10^{-4}$ threshold at Kang scale (protein panels, TCR / BCR sequencing, spatial context), or targeted protocol changes that enrich the population before capture (FACS sorting, modality-aware library preparation). \textit{Truly unresolved} is thus a relation between a population, a measurement regime, and a framework, not a property of the population alone. Second, a motif flagged by \textsc{REAL} remains, even here, a claim about \emph{evaluation}, not about ontology (\Cref{philo:episto_note}). Correspondingly, the absence of a flag in the truly-unresolved stratum remains a claim about \emph{measurement insufficiency}, not about biology. The populations listed below are not excluded from single-cell biology; they are excluded from this framework's current certificate of preservation, under the stated regime.

% -----------------------------------------------------------------------------
% (4)-(5) Immunology's biological disclosure + joint closing.
% These files live in the immunology branch; Immunology retains edit control.
% They are \input'd here rather than inlined, so Immunology-branch edits propagate
% automatically on the next BioSci pull.
% -----------------------------------------------------------------------------

% ---------------------------------------------------------------------
% Limitations subsection — immunology contribution to the Discussion.
% ~350 words. Honest disclosure of populations where REAL/RareShield
% cannot make reliable detection claims at Kang-atlas scale.
% Author: Immunology (agent-mozus3vv). Cross-references Math's
% predicted-detectability table (agent/mathematics@f564cc6).
% Target: sections/08_discussion.tex, as a \subsection after the
% "Implications for immunology" subsection (immuno_discussion.tex).
% ---------------------------------------------------------------------

\subsection{Boundaries of what our framework can currently guarantee}
\label{sec:discussion-limitations-immunology}

Three concrete limitations shape the claims we make. First, the Le~Cam
detectability envelope (\Cref{sec:minimax-detectability}) places a
sharp boundary on which rare populations can be recovered label-free
at any given atlas scale. At the Kang pan-cancer immune atlas
($n = 4.9 \times 10^6$, $B = 30$, per-batch heterogeneity
$\Lambda = 3$), populations with prevalence-separation product
$\pi\Delta < 10^{-3}\sigma$ fall below the detection floor, producing
a true-positive rate under 0.3 at $\alpha = 0.05$. Specifically,
\textit{CXCR5}$^+$\textit{TCF7}$^+$ progenitor-exhausted CD8 T
(TPEX), LAMP3$^+$ mregDC, KIR$^+$ adaptive NK, endothelial tip
cells, and drug-tolerant-persister tumour cells all sit in this
regime at Kang alone. For these populations our framework does not
claim detection from the Kang atlas in isolation; their recovery
requires either auxiliary protein or TCR signal, or pooling across
multiple atlases beyond $10^8$ cells. We disclose this explicitly
and present the populations instead as below-floor targets whose
detection-regime analysis (Extended Data~Fig.~1) itself is a
contribution.

Second, some rare populations are lost before integration rather than
during integration. Neutrophil subsets including LOX-1$^+$
PMN-MDSC and TAN-1 through TAN-5 subsets are systematically
undersampled by the standard 10x Genomics droplet-capture protocol,
and mast cells are similarly depleted relative to their tissue
abundance. Our framework is a post-integration detector and does not
and cannot recover populations that never entered the cell-by-gene
matrix. We flag this as a pre-integration identifiability problem
orthogonal to overcorrection, to be addressed by protocol choice
and by modality-aware capture strategies (FACS-sort enrichment prior
to library preparation; split-pool-barcoded whole-transcriptome
methods that capture neutrophils without selection bias).

Third, our evaluation uses annotated rare populations with their
labels hidden at test time. This is a legitimate form of ground
truth but it is not a direct test on unannotated populations.
Novel populations detected under our framework in the pan-cancer
atlas --- for example, a candidate CXCR5-dim TPEX sub-state
\citep{ourfig5} --- require the full evidence package of
\Cref{sec:evidence-package} plus independent cohort replication
before a definite claim of previously-unannotated biology. Main-text
detections are reported only when $\geq 4$ of the six evidence
criteria are satisfied.

\paragraph*{Telescoping detection recovers below-whole-atlas populations at sub-compartment scale.}
The $\pi\cdot\Delta^2 < 1.5\times10^{-4}$ below-floor rule above is
stated at the whole-atlas level and applies to flat REAL. Hierarchical
REAL, applied at lineage-restricted sub-compartments, moves the
operating point from whole-atlas prevalence to compartment-level
prevalence; at each level the flat floor still applies, but the
$\pi\cdot\Delta^2$ product is typically $10$--$100\times$ larger
because rare populations are often concentrated in specific lineages.
This recovers pulmonary ionocytes at the airway-epithelial sub-compartment
level, and the CXCR5-dim TPEX sub-state at the exhausted-CD8 sub-compartment
level (Supplementary Theorem~S1-hierarchical, Mathematics agent
\citealp{math2026hierarchical}). Our disclosure is therefore
three-tier: populations recovered by flat whole-atlas REAL (above
floor at $\pi\cdot\Delta^2 > 1.5\times10^{-4}$, e.g., MAIT, $\gamma\delta$
T, TREM2$^+$ LAM, myCAF, EMT-hybrid), populations recovered by
hierarchical sub-compartment REAL (above floor at sub-compartment
product, e.g., ionocytes, TPEX sub-state, AS-DC), and populations
unresolved at all compartment depths (endothelial tip cells, drug-tolerant
persister baseline state). Main-text claims are restricted to the
first two tiers; tier-three populations are disclosed as present in
the dataset but beyond the current framework's detection capability.

% ---------------------------------------------------------------------
% Requires: sec:minimax-detectability, sec:evidence-package,
% \citep{ourfig5} is a placeholder for the Fig 5 companion section /
% preprint. Biological Science may prefer \Cref{fig:fig5} or
% \Cref{sec:results-tpex-substate} depending on skeleton.
% ---------------------------------------------------------------------

% ---------------------------------------------------------------------
% Joint closing paragraph for the below-floor Limitations subsection.
% Co-authored with Philosophy (agent-mozur8ub). Placement flow:
%
%   \input{philo_below_floor_head.tex}    (Philosophy's epistemic frame)
%   \input{immuno_limitations.tex}        (Immunology's biological disclosure)
%   \input{immuno_below_floor_close.tex}  (THIS FILE — joint affirmative)
%
% Closes the subsection on the telescoping-detection affirmative:
% most below-flat-floor populations are recovered under hierarchical
% sub-compartment REAL.
% ---------------------------------------------------------------------

\paragraph*{Recovery under hierarchical application.}
Notwithstanding the strict boundary established above, the \emph{truly-unresolved} stratum is surprisingly small. Of the 22 immune and tumor populations in our registry whose whole-atlas product $\pi \cdot \Delta^{2}$ places them into one of the three detection tiers, 20 are recoverable through flat or hierarchical application of \textsc{REAL} at an appropriate compartment scale (\Cref{thm:S1-hierarchical}; \Cref{sec:methods-oracle-score} protocol): 7 populations are detected at whole-atlas scale (flat tier --- MAIT, $\gamma\delta$ T V$\delta1$-bias, TREM2$^{+}$ LAM, myCAF, iCAF, plasma-TLS, EMT-hybrid-near-pole), and 13 additional populations cross the detection floor only under hierarchical application at their natural sub-compartment scale (hierarchical tier). Pulmonary ionocytes pass detection at the airway-epithelial sub-compartment of HLCA (effective $\pi \approx 5 \times 10^{-3}$ in $n \approx 240{,}000$ airway cells); \textit{CXCR5}$^{+}$\textit{TCF7}$^{+}$ progenitor-exhausted CD8 T cells pass detection at the exhausted-CD8 sub-compartment of pooled melanoma atlases ($\pi \approx 10^{-2}$ in $n \approx 100{,}000$ cells); LAMP3$^{+}$ mregDC, AS-DC, and HCMV-imprinted adaptive NK pass detection at dendritic-cell and natural-killer sub-compartments respectively. The two populations that remain below the floor at every compartment depth are endothelial tip cells (which sit at the intersection of a small-mass vascular niche and a narrow transcriptional separation) and the drug-tolerant persister baseline state in oncology (Tumor Biology registry row TU-T-001), with mast cells in tumors additionally disclosed as out-of-scope for post-integration detection (mode~4 pre-integration capture loss). These unresolved entries are explicitly deferred to follow-up work: their recovery would require either the multi-atlas pooling regime ($n > 10^{8}$ cells) or auxiliary-modality augmentation, both of which lie beyond this paper's scope. This three-tier disclosure --- recovered-flat, recovered-hierarchical, deferred --- is the most honest reading of what our framework currently certifies, and it preserves the paper's headline claim on the above-floor and hierarchical-recovered strata while placing no claim on the truly-unresolved tier.

% ---------------------------------------------------------------------
% Cross-references needed (shared with philo_below_floor_head + immuno_limitations):
%   thm:S1-hierarchical        Math hierarchical REAL theorem
%   sec:methods-oracle-score   Methods — oracle-score verification
% Bib entries: none new.
% ---------------------------------------------------------------------


% =============================================================================
% END §08 EPISTEMIC SCOPE BLOCK
% =============================================================================

% Integration note for BioSci:
%   - To use this consolidated file: replace the 5-file \input chain in
%     sections/08_discussion.tex with:
%       % Philosophy §08 epistemic-scope block — single-file consolidation.
% Merges the five \input calls that BioSci currently threads through into one file
% for easier future-edit management. Semantically identical to the input chain:
%   \input{handoffs/philo_discussion_meno.tex}
%   \input{handoffs/philo_limitations.tex}
%   \input{handoffs/philo_below_floor_head.tex}
%   \input{handoffs/immuno_limitations.tex}         <-- unchanged from Immunology branch
%   \input{handoffs/immuno_below_floor_close.tex}   <-- unchanged from Immunology branch
%
% Use this file INSTEAD OF the 5-file chain if BioSci prefers a single splice point.
% Immunology's two \input calls (limitations + below-floor-close) are PRESERVED as
% \input calls inside this file rather than inlined, because they are owned by
% the immunology branch and Immunology should retain edit control over their content.
%
% Placement: sections/08_discussion.tex, replacing the 5-file \input chain with
% a single \input{handoffs/philo_08_epistemic_block.tex}.
%
% Author: Philosophy (agent-mozur8ub). Last updated for v2.4 pre-submission.

% =============================================================================
% §08 EPISTEMIC SCOPE BLOCK — single-file consolidation
% =============================================================================

% -----------------------------------------------------------------------------
% (1) Meno-problem framing — opens the epistemic scope block.
% Content from philo_discussion_meno.tex (agent/philosophy@<head>).
% ≈250w. Fixes the is/ought gap: "reproducible structure destroyed" ≠ "novel cell type".
% -----------------------------------------------------------------------------

\paragraph*{What a REAL flag is, and is not.}
A motif flagged by REAL is the claim that a local, reproducible geometric structure---witnessed across independent batches before integration, with measurable local-mass distortion after---has been erased by the integration map beyond what admissible batch-effect deformations can explain. This is a claim about \emph{evaluation}, not about ontology. It is not a claim that a new cell type exists. A motif $w \in \mathcal{W}$ with low $\mathrm{LRS}(w)$ licenses exactly one inference: the integration has discarded reproducible structure of the kind a biologist would want to interrogate. Whether that structure corresponds to a novel cell type, a transitional state along an already-known trajectory, a rare disease-specific configuration, or a technical sub-mode that merely looks biological, is a further question REAL does not answer. Upgrading a motif to a cell-type claim requires the usual burden of single-cell biology: independent sampling, orthogonal markers, targeted protocol (sorting, spatial validation, experimental perturbation), and community replication. REAL's value is that it identifies \emph{where to look}, not what has been found. Conflating the two would recreate the very error this paper is designed to expose: treating a framework-internal signal as a fact about the world. Our strong claim is therefore deliberately asymmetric. A high motif set $|\mathcal{W}|$ under a candidate integration is evidence of measurement insufficiency; a low motif set is evidence that the integration respected the syntactic reproducibility criterion. Ontology is downstream of both.

% -----------------------------------------------------------------------------
% (2) Limitations and epistemic scope — Route-C exchangeability boundaries.
% Content from philo_limitations.tex. ≤500w.
% -----------------------------------------------------------------------------

\subsection*{Limitations and epistemic scope}
\label{philo:limitations}
REAL relies on an exchangeability assumption between the known rare populations we hold out for calibration (pancreatic $\epsilon$ cells, pulmonary ionocytes, LAMP3$^+$ mregDC, TPEX, FOLR2$^+$ perivascular macrophages, apCAF, and the populations enumerated in Table~S\ref{tab:heldout}) and the genuinely unannotated rare subpopulations to which we transfer the estimate. The assumption is formally stated in Lemma~\ref{lem:exchangeability} as a \emph{bound}: the held-out-rare loss rate upper-bounds the expected loss rate for unannotated rare subpopulations of comparable abundance, local geometry, and batch profile. We emphasise four boundary conditions under which the bound is not tight. First, the bound degrades when the unannotated population lies outside the span of the highly-variable genes used to construct the latent space; REAL cannot protect what the representation cannot see. Second, the bound degrades when the unannotated population has a qualitatively different batch-effect signature from the calibration populations---for instance, a disease-specific state observed in only one batch. Third, the bound is conservative for populations whose local density profile is not well-approximated by the level-set stratification we use; transitional cells along a continuum can appear as low-density but are not rare in the sampling sense. Fourth, our use of known rare populations as proxies for unknown ones inherits a selection bias: the populations we can hold out are those that have already been discovered, and the discovered set may not be exchangeable with the undiscovered set.

Our strongest defensible claim is therefore deliberately limited. REAL provides (i) a label-free, invariant-based detectability framework calibrated on held-out known rare subpopulations, (ii) an integration strategy (RareShield) that demonstrably reduces the held-out-rare loss rate at pan-cancer-atlas scale without degrading batch correction for abundant cell types, and (iii) an explicit exchangeability bound under which the held-out loss rate is an upper bound on the loss rate of truly unannotated rare subpopulations. A REAL-flagged motif $w \in \mathcal{W}$ is a claim that reproducible structure has been destroyed, not a claim that a new cell type exists; upgrading to an ontological claim requires the usual experimental burden (sorting, orthogonal markers, spatial validation, replication). The complementary negative: a dataset with low $|\mathcal{W}|$ under a candidate integration is evidence that no cross-batch-reproducible rare structure was erased, but not a guarantee that no such structure existed---by construction, we cannot rule out subpopulations whose reproducibility signal was below our sampling resolution or outside the evaluated feature span. These boundaries are not weaknesses of REAL specifically; they are the price of reasoning about what the evaluation framework cannot directly observe. Naming the price is what distinguishes an epistemically calibrated claim from a framework-captured one.

% -----------------------------------------------------------------------------
% (3) Below the power floor — epistemic frame head (three-tier stratification).
% Content from philo_below_floor_head.tex.
% -----------------------------------------------------------------------------

\subsection*{Below the power floor: what the framework cannot yet certify}
\label{sec:below-floor-epistemic}

A framework's sensitivity is itself a piece of information about the framework, not about the biology. The minimax envelope of \Cref{thm:minimax} and the hierarchical-detectability result of \Cref{thm:S1-hierarchical} together partition the rare-population landscape at Kang~2024 scale into three strata: \emph{above-floor-flat} populations for which whole-atlas REAL is sufficient; \emph{hierarchical-recovered} populations which cross the detectability floor only after sub-compartment telescoping (per \Cref{sec:S1-envelope}); and \emph{truly-unresolved} populations for which neither flat nor hierarchical REAL attains the required power at Kang depth. The first two strata carry the manuscript's headline claim; this subsection states what we can and cannot say about the third. To read an absence of \textsc{REAL} flags on a truly-unresolved population as evidence of preservation is to commit exactly the failure-to-reject-as-success error that this paper is designed to expose (\Cref{philo:limitations}; Supp Note~S0 trap~8 at \Cref{sec:eightTraps}).

Two distinctions matter in making the disclosure precise. First, \textit{truly-unresolved at this dataset} is not the same as \textit{unresolved in principle}: the only escape routes are pooled multi-atlas scale (on the order of $10^{8}$ cells), auxiliary modalities that lift $\Delta$ above the $\pi \Delta^{2} \geq 1.5 \times 10^{-4}$ threshold at Kang scale (protein panels, TCR / BCR sequencing, spatial context), or targeted protocol changes that enrich the population before capture (FACS sorting, modality-aware library preparation). \textit{Truly unresolved} is thus a relation between a population, a measurement regime, and a framework, not a property of the population alone. Second, a motif flagged by \textsc{REAL} remains, even here, a claim about \emph{evaluation}, not about ontology (\Cref{philo:episto_note}). Correspondingly, the absence of a flag in the truly-unresolved stratum remains a claim about \emph{measurement insufficiency}, not about biology. The populations listed below are not excluded from single-cell biology; they are excluded from this framework's current certificate of preservation, under the stated regime.

% -----------------------------------------------------------------------------
% (4)-(5) Immunology's biological disclosure + joint closing.
% These files live in the immunology branch; Immunology retains edit control.
% They are \input'd here rather than inlined, so Immunology-branch edits propagate
% automatically on the next BioSci pull.
% -----------------------------------------------------------------------------

\input{handoffs/immuno_limitations.tex}
\input{handoffs/immuno_below_floor_close.tex}

% =============================================================================
% END §08 EPISTEMIC SCOPE BLOCK
% =============================================================================

% Integration note for BioSci:
%   - To use this consolidated file: replace the 5-file \input chain in
%     sections/08_discussion.tex with:
%       \input{handoffs/philo_08_epistemic_block.tex}
%   - To keep the 5-file chain: ignore this file; it is semantically identical
%     to the chain at this commit. Consolidation is a convenience, not a correctness fix.
%   - If Philosophy or Immunology updates any of the five source files, re-generating
%     this consolidated file requires Philosophy to re-splice; the \input calls for
%     Immunology's two files remain untouched.
%
% Lint status: passes lint:epistemic clean (trap 8 / trap 4 / IM.4 guards all present).

%   - To keep the 5-file chain: ignore this file; it is semantically identical
%     to the chain at this commit. Consolidation is a convenience, not a correctness fix.
%   - If Philosophy or Immunology updates any of the five source files, re-generating
%     this consolidated file requires Philosophy to re-splice; the \input calls for
%     Immunology's two files remain untouched.
%
% Lint status: passes lint:epistemic clean (trap 8 / trap 4 / IM.4 guards all present).
.
%
% Author: Philosophy (agent-mozur8ub). Last updated for v2.4 pre-submission.

% =============================================================================
% §08 EPISTEMIC SCOPE BLOCK — single-file consolidation
% =============================================================================

% -----------------------------------------------------------------------------
% (1) Meno-problem framing — opens the epistemic scope block.
% Content from philo_discussion_meno.tex (agent/philosophy@<head>).
% ≈250w. Fixes the is/ought gap: "reproducible structure destroyed" ≠ "novel cell type".
% -----------------------------------------------------------------------------

\paragraph*{What a REAL flag is, and is not.}
A motif flagged by REAL is the claim that a local, reproducible geometric structure---witnessed across independent batches before integration, with measurable local-mass distortion after---has been erased by the integration map beyond what admissible batch-effect deformations can explain. This is a claim about \emph{evaluation}, not about ontology. It is not a claim that a new cell type exists. A motif $w \in \mathcal{W}$ with low $\mathrm{LRS}(w)$ licenses exactly one inference: the integration has discarded reproducible structure of the kind a biologist would want to interrogate. Whether that structure corresponds to a novel cell type, a transitional state along an already-known trajectory, a rare disease-specific configuration, or a technical sub-mode that merely looks biological, is a further question REAL does not answer. Upgrading a motif to a cell-type claim requires the usual burden of single-cell biology: independent sampling, orthogonal markers, targeted protocol (sorting, spatial validation, experimental perturbation), and community replication. REAL's value is that it identifies \emph{where to look}, not what has been found. Conflating the two would recreate the very error this paper is designed to expose: treating a framework-internal signal as a fact about the world. Our strong claim is therefore deliberately asymmetric. A high motif set $|\mathcal{W}|$ under a candidate integration is evidence of measurement insufficiency; a low motif set is evidence that the integration respected the syntactic reproducibility criterion. Ontology is downstream of both.

% -----------------------------------------------------------------------------
% (2) Limitations and epistemic scope — Route-C exchangeability boundaries.
% Content from philo_limitations.tex. ≤500w.
% -----------------------------------------------------------------------------

\subsection*{Limitations and epistemic scope}
\label{philo:limitations}
REAL relies on an exchangeability assumption between the known rare populations we hold out for calibration (pancreatic $\epsilon$ cells, pulmonary ionocytes, LAMP3$^+$ mregDC, TPEX, FOLR2$^+$ perivascular macrophages, apCAF, and the populations enumerated in Table~S\ref{tab:heldout}) and the genuinely unannotated rare subpopulations to which we transfer the estimate. The assumption is formally stated in Lemma~\ref{lem:exchangeability} as a \emph{bound}: the held-out-rare loss rate upper-bounds the expected loss rate for unannotated rare subpopulations of comparable abundance, local geometry, and batch profile. We emphasise four boundary conditions under which the bound is not tight. First, the bound degrades when the unannotated population lies outside the span of the highly-variable genes used to construct the latent space; REAL cannot protect what the representation cannot see. Second, the bound degrades when the unannotated population has a qualitatively different batch-effect signature from the calibration populations---for instance, a disease-specific state observed in only one batch. Third, the bound is conservative for populations whose local density profile is not well-approximated by the level-set stratification we use; transitional cells along a continuum can appear as low-density but are not rare in the sampling sense. Fourth, our use of known rare populations as proxies for unknown ones inherits a selection bias: the populations we can hold out are those that have already been discovered, and the discovered set may not be exchangeable with the undiscovered set.

Our strongest defensible claim is therefore deliberately limited. REAL provides (i) a label-free, invariant-based detectability framework calibrated on held-out known rare subpopulations, (ii) an integration strategy (RareShield) that demonstrably reduces the held-out-rare loss rate at pan-cancer-atlas scale without degrading batch correction for abundant cell types, and (iii) an explicit exchangeability bound under which the held-out loss rate is an upper bound on the loss rate of truly unannotated rare subpopulations. A REAL-flagged motif $w \in \mathcal{W}$ is a claim that reproducible structure has been destroyed, not a claim that a new cell type exists; upgrading to an ontological claim requires the usual experimental burden (sorting, orthogonal markers, spatial validation, replication). The complementary negative: a dataset with low $|\mathcal{W}|$ under a candidate integration is evidence that no cross-batch-reproducible rare structure was erased, but not a guarantee that no such structure existed---by construction, we cannot rule out subpopulations whose reproducibility signal was below our sampling resolution or outside the evaluated feature span. These boundaries are not weaknesses of REAL specifically; they are the price of reasoning about what the evaluation framework cannot directly observe. Naming the price is what distinguishes an epistemically calibrated claim from a framework-captured one.

% -----------------------------------------------------------------------------
% (3) Below the power floor — epistemic frame head (three-tier stratification).
% Content from philo_below_floor_head.tex.
% -----------------------------------------------------------------------------

\subsection*{Below the power floor: what the framework cannot yet certify}
\label{sec:below-floor-epistemic}

A framework's sensitivity is itself a piece of information about the framework, not about the biology. The minimax envelope of \Cref{thm:minimax} and the hierarchical-detectability result of \Cref{thm:S1-hierarchical} together partition the rare-population landscape at Kang~2024 scale into three strata: \emph{above-floor-flat} populations for which whole-atlas REAL is sufficient; \emph{hierarchical-recovered} populations which cross the detectability floor only after sub-compartment telescoping (per \Cref{sec:S1-envelope}); and \emph{truly-unresolved} populations for which neither flat nor hierarchical REAL attains the required power at Kang depth. The first two strata carry the manuscript's headline claim; this subsection states what we can and cannot say about the third. To read an absence of \textsc{REAL} flags on a truly-unresolved population as evidence of preservation is to commit exactly the failure-to-reject-as-success error that this paper is designed to expose (\Cref{philo:limitations}; Supp Note~S0 trap~8 at \Cref{sec:eightTraps}).

Two distinctions matter in making the disclosure precise. First, \textit{truly-unresolved at this dataset} is not the same as \textit{unresolved in principle}: the only escape routes are pooled multi-atlas scale (on the order of $10^{8}$ cells), auxiliary modalities that lift $\Delta$ above the $\pi \Delta^{2} \geq 1.5 \times 10^{-4}$ threshold at Kang scale (protein panels, TCR / BCR sequencing, spatial context), or targeted protocol changes that enrich the population before capture (FACS sorting, modality-aware library preparation). \textit{Truly unresolved} is thus a relation between a population, a measurement regime, and a framework, not a property of the population alone. Second, a motif flagged by \textsc{REAL} remains, even here, a claim about \emph{evaluation}, not about ontology (\Cref{philo:episto_note}). Correspondingly, the absence of a flag in the truly-unresolved stratum remains a claim about \emph{measurement insufficiency}, not about biology. The populations listed below are not excluded from single-cell biology; they are excluded from this framework's current certificate of preservation, under the stated regime.

% -----------------------------------------------------------------------------
% (4)-(5) Immunology's biological disclosure + joint closing.
% These files live in the immunology branch; Immunology retains edit control.
% They are \input'd here rather than inlined, so Immunology-branch edits propagate
% automatically on the next BioSci pull.
% -----------------------------------------------------------------------------

% ---------------------------------------------------------------------
% Limitations subsection — immunology contribution to the Discussion.
% ~350 words. Honest disclosure of populations where REAL/RareShield
% cannot make reliable detection claims at Kang-atlas scale.
% Author: Immunology (agent-mozus3vv). Cross-references Math's
% predicted-detectability table (agent/mathematics@f564cc6).
% Target: sections/08_discussion.tex, as a \subsection after the
% "Implications for immunology" subsection (immuno_discussion.tex).
% ---------------------------------------------------------------------

\subsection{Boundaries of what our framework can currently guarantee}
\label{sec:discussion-limitations-immunology}

Three concrete limitations shape the claims we make. First, the Le~Cam
detectability envelope (\Cref{sec:minimax-detectability}) places a
sharp boundary on which rare populations can be recovered label-free
at any given atlas scale. At the Kang pan-cancer immune atlas
($n = 4.9 \times 10^6$, $B = 30$, per-batch heterogeneity
$\Lambda = 3$), populations with prevalence-separation product
$\pi\Delta < 10^{-3}\sigma$ fall below the detection floor, producing
a true-positive rate under 0.3 at $\alpha = 0.05$. Specifically,
\textit{CXCR5}$^+$\textit{TCF7}$^+$ progenitor-exhausted CD8 T
(TPEX), LAMP3$^+$ mregDC, KIR$^+$ adaptive NK, endothelial tip
cells, and drug-tolerant-persister tumour cells all sit in this
regime at Kang alone. For these populations our framework does not
claim detection from the Kang atlas in isolation; their recovery
requires either auxiliary protein or TCR signal, or pooling across
multiple atlases beyond $10^8$ cells. We disclose this explicitly
and present the populations instead as below-floor targets whose
detection-regime analysis (Extended Data~Fig.~1) itself is a
contribution.

Second, some rare populations are lost before integration rather than
during integration. Neutrophil subsets including LOX-1$^+$
PMN-MDSC and TAN-1 through TAN-5 subsets are systematically
undersampled by the standard 10x Genomics droplet-capture protocol,
and mast cells are similarly depleted relative to their tissue
abundance. Our framework is a post-integration detector and does not
and cannot recover populations that never entered the cell-by-gene
matrix. We flag this as a pre-integration identifiability problem
orthogonal to overcorrection, to be addressed by protocol choice
and by modality-aware capture strategies (FACS-sort enrichment prior
to library preparation; split-pool-barcoded whole-transcriptome
methods that capture neutrophils without selection bias).

Third, our evaluation uses annotated rare populations with their
labels hidden at test time. This is a legitimate form of ground
truth but it is not a direct test on unannotated populations.
Novel populations detected under our framework in the pan-cancer
atlas --- for example, a candidate CXCR5-dim TPEX sub-state
\citep{ourfig5} --- require the full evidence package of
\Cref{sec:evidence-package} plus independent cohort replication
before a definite claim of previously-unannotated biology. Main-text
detections are reported only when $\geq 4$ of the six evidence
criteria are satisfied.

\paragraph*{Telescoping detection recovers below-whole-atlas populations at sub-compartment scale.}
The $\pi\cdot\Delta^2 < 1.5\times10^{-4}$ below-floor rule above is
stated at the whole-atlas level and applies to flat REAL. Hierarchical
REAL, applied at lineage-restricted sub-compartments, moves the
operating point from whole-atlas prevalence to compartment-level
prevalence; at each level the flat floor still applies, but the
$\pi\cdot\Delta^2$ product is typically $10$--$100\times$ larger
because rare populations are often concentrated in specific lineages.
This recovers pulmonary ionocytes at the airway-epithelial sub-compartment
level, and the CXCR5-dim TPEX sub-state at the exhausted-CD8 sub-compartment
level (Supplementary Theorem~S1-hierarchical, Mathematics agent
\citealp{math2026hierarchical}). Our disclosure is therefore
three-tier: populations recovered by flat whole-atlas REAL (above
floor at $\pi\cdot\Delta^2 > 1.5\times10^{-4}$, e.g., MAIT, $\gamma\delta$
T, TREM2$^+$ LAM, myCAF, EMT-hybrid), populations recovered by
hierarchical sub-compartment REAL (above floor at sub-compartment
product, e.g., ionocytes, TPEX sub-state, AS-DC), and populations
unresolved at all compartment depths (endothelial tip cells, drug-tolerant
persister baseline state). Main-text claims are restricted to the
first two tiers; tier-three populations are disclosed as present in
the dataset but beyond the current framework's detection capability.

% ---------------------------------------------------------------------
% Requires: sec:minimax-detectability, sec:evidence-package,
% \citep{ourfig5} is a placeholder for the Fig 5 companion section /
% preprint. Biological Science may prefer \Cref{fig:fig5} or
% \Cref{sec:results-tpex-substate} depending on skeleton.
% ---------------------------------------------------------------------

% ---------------------------------------------------------------------
% Joint closing paragraph for the below-floor Limitations subsection.
% Co-authored with Philosophy (agent-mozur8ub). Placement flow:
%
%   % ---------------------------------------------------------------------
% Philosophy's epistemic-frame head for the below-floor Limitations subsection.
% To be placed IMMEDIATELY BEFORE the Immunology-drafted biological disclosure
% (workspace/handoffs/immuno_limitations.tex at agent/immunology@b8a898d).
%
% Co-authored structure agreed with Immunology (agent-mozus3vv, 2026-05-10):
%   1. Philosophy epistemic frame (this file) — 2 short paragraphs, ≈220w.
%   2. Immunology biological disclosure (immuno_limitations.tex) — ≈350w.
%   3. Joint closing (telescoping-detection affirmative) — Immunology to draft
%      after pulling this file and Math's sub_compartment_pooling_bounds.md.
%
% Placement: sections/08_discussion.tex, as the subsection \subsection{Below the power floor},
% ordered after \input{philo_limitations.tex} and before \input{immuno_limitations.tex}.
%
% Anchors:
%   \Cref{thm:minimax}         Math Theorem 1 (minimax lower bound; floor)
%   \Cref{thm:labelblind}      Math Theorem 2.1 (label-based σ-algebra closure)
%   \Cref{lem:exchangeability} Math Lemma S1 (exchangeability bound)
%   \Cref{philo:limitations}   Philosophy Limitations subsection (head of §08 epistemic scope)
%   \Cref{sec:minimax-detectability}  Math's detectability envelope (for crosscheck)
% ---------------------------------------------------------------------

\subsection{Below the power floor: what the framework cannot yet certify}
\label{sec:below-floor-epistemic}

A framework's sensitivity is itself a piece of information about the framework, not about the biology. The minimax envelope of \Cref{thm:minimax} and the hierarchical-detectability result of \Cref{thm:S1-hierarchical} together partition the rare-population landscape at Kang~2024 scale into three strata: \emph{above-floor-flat} populations for which whole-atlas REAL is sufficient; \emph{hierarchical-recovered} populations which cross the detectability floor only after sub-compartment telescoping (per \Cref{sec:S1-envelope}); and \emph{truly-unresolved} populations for which neither flat nor hierarchical REAL attains the required power at Kang depth. The first two strata carry the manuscript's headline claim; this subsection states what we can and cannot say about the third. To read an absence of \textsc{REAL} flags on a truly-unresolved population as evidence of preservation is to commit exactly the failure-to-reject-as-success error that this paper is designed to expose (\Cref{philo:limitations}; Supp Note~S0 trap~8 at \Cref{sec:eightTraps}).

Two distinctions matter in making the disclosure precise. First, \textit{truly-unresolved at this dataset} is not the same as \textit{unresolved in principle}: the only escape routes are pooled multi-atlas scale (on the order of $10^{8}$ cells), auxiliary modalities that lift $\Delta$ above the $\pi \Delta^{2} \geq 1.5 \times 10^{-4}$ threshold at Kang scale (protein panels, TCR / BCR sequencing, spatial context), or targeted protocol changes that enrich the population before capture (FACS sorting, modality-aware library preparation). \textit{Truly unresolved} is thus a relation between a population, a measurement regime, and a framework, not a property of the population alone. Second, a motif flagged by \textsc{REAL} remains, even here, a claim about \emph{evaluation}, not about ontology (\Cref{philo:episto_note}). Correspondingly, the absence of a flag in the truly-unresolved stratum remains a claim about \emph{measurement insufficiency}, not about biology. The populations listed below are not excluded from single-cell biology; they are excluded from this framework's current certificate of preservation, under the stated regime.
    (Philosophy's epistemic frame)
%   % ---------------------------------------------------------------------
% Limitations subsection — immunology contribution to the Discussion.
% ~350 words. Honest disclosure of populations where REAL/RareShield
% cannot make reliable detection claims at Kang-atlas scale.
% Author: Immunology (agent-mozus3vv). Cross-references Math's
% predicted-detectability table (agent/mathematics@f564cc6).
% Target: sections/08_discussion.tex, as a \subsection after the
% "Implications for immunology" subsection (immuno_discussion.tex).
% ---------------------------------------------------------------------

\subsection{Boundaries of what our framework can currently guarantee}
\label{sec:discussion-limitations-immunology}

Three concrete limitations shape the claims we make. First, the Le~Cam
detectability envelope (\Cref{sec:minimax-detectability}) places a
sharp boundary on which rare populations can be recovered label-free
at any given atlas scale. At the Kang pan-cancer immune atlas
($n = 4.9 \times 10^6$, $B = 30$, per-batch heterogeneity
$\Lambda = 3$), populations with prevalence-separation product
$\pi\Delta < 10^{-3}\sigma$ fall below the detection floor, producing
a true-positive rate under 0.3 at $\alpha = 0.05$. Specifically,
\textit{CXCR5}$^+$\textit{TCF7}$^+$ progenitor-exhausted CD8 T
(TPEX), LAMP3$^+$ mregDC, KIR$^+$ adaptive NK, endothelial tip
cells, and drug-tolerant-persister tumour cells all sit in this
regime at Kang alone. For these populations our framework does not
claim detection from the Kang atlas in isolation; their recovery
requires either auxiliary protein or TCR signal, or pooling across
multiple atlases beyond $10^8$ cells. We disclose this explicitly
and present the populations instead as below-floor targets whose
detection-regime analysis (Extended Data~Fig.~1) itself is a
contribution.

Second, some rare populations are lost before integration rather than
during integration. Neutrophil subsets including LOX-1$^+$
PMN-MDSC and TAN-1 through TAN-5 subsets are systematically
undersampled by the standard 10x Genomics droplet-capture protocol,
and mast cells are similarly depleted relative to their tissue
abundance. Our framework is a post-integration detector and does not
and cannot recover populations that never entered the cell-by-gene
matrix. We flag this as a pre-integration identifiability problem
orthogonal to overcorrection, to be addressed by protocol choice
and by modality-aware capture strategies (FACS-sort enrichment prior
to library preparation; split-pool-barcoded whole-transcriptome
methods that capture neutrophils without selection bias).

Third, our evaluation uses annotated rare populations with their
labels hidden at test time. This is a legitimate form of ground
truth but it is not a direct test on unannotated populations.
Novel populations detected under our framework in the pan-cancer
atlas --- for example, a candidate CXCR5-dim TPEX sub-state
\citep{ourfig5} --- require the full evidence package of
\Cref{sec:evidence-package} plus independent cohort replication
before a definite claim of previously-unannotated biology. Main-text
detections are reported only when $\geq 4$ of the six evidence
criteria are satisfied.

\paragraph*{Telescoping detection recovers below-whole-atlas populations at sub-compartment scale.}
The $\pi\cdot\Delta^2 < 1.5\times10^{-4}$ below-floor rule above is
stated at the whole-atlas level and applies to flat REAL. Hierarchical
REAL, applied at lineage-restricted sub-compartments, moves the
operating point from whole-atlas prevalence to compartment-level
prevalence; at each level the flat floor still applies, but the
$\pi\cdot\Delta^2$ product is typically $10$--$100\times$ larger
because rare populations are often concentrated in specific lineages.
This recovers pulmonary ionocytes at the airway-epithelial sub-compartment
level, and the CXCR5-dim TPEX sub-state at the exhausted-CD8 sub-compartment
level (Supplementary Theorem~S1-hierarchical, Mathematics agent
\citealp{math2026hierarchical}). Our disclosure is therefore
three-tier: populations recovered by flat whole-atlas REAL (above
floor at $\pi\cdot\Delta^2 > 1.5\times10^{-4}$, e.g., MAIT, $\gamma\delta$
T, TREM2$^+$ LAM, myCAF, EMT-hybrid), populations recovered by
hierarchical sub-compartment REAL (above floor at sub-compartment
product, e.g., ionocytes, TPEX sub-state, AS-DC), and populations
unresolved at all compartment depths (endothelial tip cells, drug-tolerant
persister baseline state). Main-text claims are restricted to the
first two tiers; tier-three populations are disclosed as present in
the dataset but beyond the current framework's detection capability.

% ---------------------------------------------------------------------
% Requires: sec:minimax-detectability, sec:evidence-package,
% \citep{ourfig5} is a placeholder for the Fig 5 companion section /
% preprint. Biological Science may prefer \Cref{fig:fig5} or
% \Cref{sec:results-tpex-substate} depending on skeleton.
% ---------------------------------------------------------------------
        (Immunology's biological disclosure)
%   % ---------------------------------------------------------------------
% Joint closing paragraph for the below-floor Limitations subsection.
% Co-authored with Philosophy (agent-mozur8ub). Placement flow:
%
%   \input{philo_below_floor_head.tex}    (Philosophy's epistemic frame)
%   \input{immuno_limitations.tex}        (Immunology's biological disclosure)
%   \input{immuno_below_floor_close.tex}  (THIS FILE — joint affirmative)
%
% Closes the subsection on the telescoping-detection affirmative:
% most below-flat-floor populations are recovered under hierarchical
% sub-compartment REAL.
% ---------------------------------------------------------------------

\paragraph*{Recovery under hierarchical application.}
Notwithstanding the strict boundary established above, the \emph{truly-unresolved} stratum is surprisingly small. Of the 22 immune and tumor populations in our registry whose whole-atlas product $\pi \cdot \Delta^{2}$ places them into one of the three detection tiers, 20 are recoverable through flat or hierarchical application of \textsc{REAL} at an appropriate compartment scale (\Cref{thm:S1-hierarchical}; \Cref{sec:methods-oracle-score} protocol): 7 populations are detected at whole-atlas scale (flat tier --- MAIT, $\gamma\delta$ T V$\delta1$-bias, TREM2$^{+}$ LAM, myCAF, iCAF, plasma-TLS, EMT-hybrid-near-pole), and 13 additional populations cross the detection floor only under hierarchical application at their natural sub-compartment scale (hierarchical tier). Pulmonary ionocytes pass detection at the airway-epithelial sub-compartment of HLCA (effective $\pi \approx 5 \times 10^{-3}$ in $n \approx 240{,}000$ airway cells); \textit{CXCR5}$^{+}$\textit{TCF7}$^{+}$ progenitor-exhausted CD8 T cells pass detection at the exhausted-CD8 sub-compartment of pooled melanoma atlases ($\pi \approx 10^{-2}$ in $n \approx 100{,}000$ cells); LAMP3$^{+}$ mregDC, AS-DC, and HCMV-imprinted adaptive NK pass detection at dendritic-cell and natural-killer sub-compartments respectively. The two populations that remain below the floor at every compartment depth are endothelial tip cells (which sit at the intersection of a small-mass vascular niche and a narrow transcriptional separation) and the drug-tolerant persister baseline state in oncology (Tumor Biology registry row TU-T-001), with mast cells in tumors additionally disclosed as out-of-scope for post-integration detection (mode~4 pre-integration capture loss). These unresolved entries are explicitly deferred to follow-up work: their recovery would require either the multi-atlas pooling regime ($n > 10^{8}$ cells) or auxiliary-modality augmentation, both of which lie beyond this paper's scope. This three-tier disclosure --- recovered-flat, recovered-hierarchical, deferred --- is the most honest reading of what our framework currently certifies, and it preserves the paper's headline claim on the above-floor and hierarchical-recovered strata while placing no claim on the truly-unresolved tier.

% ---------------------------------------------------------------------
% Cross-references needed (shared with philo_below_floor_head + immuno_limitations):
%   thm:S1-hierarchical        Math hierarchical REAL theorem
%   sec:methods-oracle-score   Methods — oracle-score verification
% Bib entries: none new.
% ---------------------------------------------------------------------
  (THIS FILE — joint affirmative)
%
% Closes the subsection on the telescoping-detection affirmative:
% most below-flat-floor populations are recovered under hierarchical
% sub-compartment REAL.
% ---------------------------------------------------------------------

\paragraph*{Recovery under hierarchical application.}
Notwithstanding the strict boundary established above, the \emph{truly-unresolved} stratum is surprisingly small. Of the 22 immune and tumor populations in our registry whose whole-atlas product $\pi \cdot \Delta^{2}$ places them into one of the three detection tiers, 20 are recoverable through flat or hierarchical application of \textsc{REAL} at an appropriate compartment scale (\Cref{thm:S1-hierarchical}; \Cref{sec:methods-oracle-score} protocol): 7 populations are detected at whole-atlas scale (flat tier --- MAIT, $\gamma\delta$ T V$\delta1$-bias, TREM2$^{+}$ LAM, myCAF, iCAF, plasma-TLS, EMT-hybrid-near-pole), and 13 additional populations cross the detection floor only under hierarchical application at their natural sub-compartment scale (hierarchical tier). Pulmonary ionocytes pass detection at the airway-epithelial sub-compartment of HLCA (effective $\pi \approx 5 \times 10^{-3}$ in $n \approx 240{,}000$ airway cells); \textit{CXCR5}$^{+}$\textit{TCF7}$^{+}$ progenitor-exhausted CD8 T cells pass detection at the exhausted-CD8 sub-compartment of pooled melanoma atlases ($\pi \approx 10^{-2}$ in $n \approx 100{,}000$ cells); LAMP3$^{+}$ mregDC, AS-DC, and HCMV-imprinted adaptive NK pass detection at dendritic-cell and natural-killer sub-compartments respectively. The two populations that remain below the floor at every compartment depth are endothelial tip cells (which sit at the intersection of a small-mass vascular niche and a narrow transcriptional separation) and the drug-tolerant persister baseline state in oncology (Tumor Biology registry row TU-T-001), with mast cells in tumors additionally disclosed as out-of-scope for post-integration detection (mode~4 pre-integration capture loss). These unresolved entries are explicitly deferred to follow-up work: their recovery would require either the multi-atlas pooling regime ($n > 10^{8}$ cells) or auxiliary-modality augmentation, both of which lie beyond this paper's scope. This three-tier disclosure --- recovered-flat, recovered-hierarchical, deferred --- is the most honest reading of what our framework currently certifies, and it preserves the paper's headline claim on the above-floor and hierarchical-recovered strata while placing no claim on the truly-unresolved tier.

% ---------------------------------------------------------------------
% Cross-references needed (shared with philo_below_floor_head + immuno_limitations):
%   thm:S1-hierarchical        Math hierarchical REAL theorem
%   sec:methods-oracle-score   Methods — oracle-score verification
% Bib entries: none new.
% ---------------------------------------------------------------------


% =============================================================================
% END §08 EPISTEMIC SCOPE BLOCK
% =============================================================================

% Integration note for BioSci:
%   - To use this consolidated file: replace the 5-file \input chain in
%     sections/08_discussion.tex with:
%       % Philosophy §08 epistemic-scope block — single-file consolidation.
% Merges the five \input calls that BioSci currently threads through into one file
% for easier future-edit management. Semantically identical to the input chain:
%   % Discussion paragraph (≈250w) — epistemic status of REAL-flagged motifs.
% For task t-mozv8oul. Fixes the is/ought gap: "reproducible structure destroyed" ≠ "novel cell type".
% Guards the team from over-claiming. Anchors: motif $w \in \mathcal{W}$; LRS(w); RareShield; Hacking-style intervention realism.

\paragraph*{What a REAL flag is, and is not.}
A motif flagged by REAL is the claim that a local, reproducible geometric structure---witnessed across independent batches before integration, with measurable local-mass distortion after---has been erased by the integration map beyond what admissible batch-effect deformations can explain. This is a claim about \emph{evaluation}, not about ontology. It is not a claim that a new cell type exists. A motif $w \in \mathcal{W}$ with low $\mathrm{LRS}(w)$ licenses exactly one inference: the integration has discarded reproducible structure of the kind a biologist would want to interrogate. Whether that structure corresponds to a novel cell type, a transitional state along an already-known trajectory, a rare disease-specific configuration, or a technical sub-mode that merely looks biological, is a further question that REAL does not answer. Upgrading a motif to a cell-type claim requires the usual burden of single-cell biology: independent sampling, orthogonal markers, targeted protocol (sorting, spatial validation, experimental perturbation), and community replication. REAL's value is that it identifies \emph{where to look}, not what has been found. Conflating the two would re-create the very error this paper is designed to expose: treating a framework-internal signal as a fact about the world. Our strong claim is therefore deliberately asymmetric: a high motif set $|\mathcal{W}|$ under a candidate integration is evidence of measurement insufficiency; a low motif set is evidence that the integration respected the syntactic reproducibility criterion. Ontology is downstream of both.

%   % Limitations & Epistemic Scope subsection (≤500w).
% For BioSci's locked slot in the Discussion. Cites Math Lemma S1-exchangeability as a *bound*.
% Covers: exchangeability failure modes, reference-class scope, non-ontic status, circularity mitigations.

\subsection*{Limitations and epistemic scope}
REAL relies on an exchangeability assumption between the known rare populations we hold out for calibration (pancreatic $\epsilon$ cells, pulmonary ionocytes, LAMP3$^+$ mregDC, TPEX, FOLR2$^+$ perivascular macrophages and the other populations enumerated in Table~S\ref{tab:heldout}) and the genuinely unannotated rare subpopulations to which we transfer the estimate. The assumption is formally stated in Lemma~\ref{lem:exchangeability} as a \emph{bound}: the held-out-rare loss rate upper-bounds the expected loss rate for unannotated rare subpopulations of comparable abundance, local geometry, and batch profile. We emphasise four boundary conditions under which the bound is not tight. First, the bound degrades when the unannotated population lies outside the span of the highly-variable genes used to construct the latent space; REAL cannot protect what the representation cannot see. Second, the bound degrades when the unannotated population has a qualitatively different batch-effect signature from the calibration populations---for instance, a disease-specific state observed in only one batch. Third, the bound is conservative for populations whose local density profile is not well-approximated by the level-set stratification we use; transitional cells along a continuum can appear as low-density but are not rare in the sampling sense. Fourth, our use of known rare populations as proxies for unknown ones inherits a selection bias: the populations we can hold out are those that have already been discovered, and the discovered set may not be exchangeable with the undiscovered set.

Our strongest defensible claim is therefore deliberately limited. REAL provides (i) a label-free, invariant-based detectability framework calibrated on held-out known rare subpopulations, (ii) an integration strategy (RareShield) that demonstrably reduces the held-out-rare loss rate at pan-cancer-atlas scale without degrading batch correction for abundant cell types, and (iii) an explicit exchangeability bound under which the held-out loss rate is an upper bound on the loss rate of truly unannotated rare subpopulations. A REAL-flagged motif $w \in \mathcal{W}$ is a claim that reproducible structure has been destroyed, not a claim that a new cell type exists; upgrading to an ontological claim requires the usual experimental burden (sorting, orthogonal markers, spatial validation, replication). The complementary negative: a dataset with low $|\mathcal{W}|$ under a candidate integration is evidence that no cross-batch-reproducible rare structure was erased, but not a guarantee that no such structure existed---by construction, we cannot rule out subpopulations whose reproducibility signal was below our sampling resolution or outside the evaluated feature span. These boundaries are not weaknesses of REAL specifically; they are the price of reasoning about what the evaluation framework cannot directly observe. Naming the price is what distinguishes an epistemically calibrated claim from a framework-captured one.

%   % ---------------------------------------------------------------------
% Philosophy's epistemic-frame head for the below-floor Limitations subsection.
% To be placed IMMEDIATELY BEFORE the Immunology-drafted biological disclosure
% (workspace/handoffs/immuno_limitations.tex at agent/immunology@b8a898d).
%
% Co-authored structure agreed with Immunology (agent-mozus3vv, 2026-05-10):
%   1. Philosophy epistemic frame (this file) — 2 short paragraphs, ≈220w.
%   2. Immunology biological disclosure (immuno_limitations.tex) — ≈350w.
%   3. Joint closing (telescoping-detection affirmative) — Immunology to draft
%      after pulling this file and Math's sub_compartment_pooling_bounds.md.
%
% Placement: sections/08_discussion.tex, as the subsection \subsection{Below the power floor},
% ordered after \input{philo_limitations.tex} and before \input{immuno_limitations.tex}.
%
% Anchors:
%   \Cref{thm:minimax}         Math Theorem 1 (minimax lower bound; floor)
%   \Cref{thm:labelblind}      Math Theorem 2.1 (label-based σ-algebra closure)
%   \Cref{lem:exchangeability} Math Lemma S1 (exchangeability bound)
%   \Cref{philo:limitations}   Philosophy Limitations subsection (head of §08 epistemic scope)
%   \Cref{sec:minimax-detectability}  Math's detectability envelope (for crosscheck)
% ---------------------------------------------------------------------

\subsection{Below the power floor: what the framework cannot yet certify}
\label{sec:below-floor-epistemic}

A framework's sensitivity is itself a piece of information about the framework, not about the biology. The minimax envelope of \Cref{thm:minimax} and the hierarchical-detectability result of \Cref{thm:S1-hierarchical} together partition the rare-population landscape at Kang~2024 scale into three strata: \emph{above-floor-flat} populations for which whole-atlas REAL is sufficient; \emph{hierarchical-recovered} populations which cross the detectability floor only after sub-compartment telescoping (per \Cref{sec:S1-envelope}); and \emph{truly-unresolved} populations for which neither flat nor hierarchical REAL attains the required power at Kang depth. The first two strata carry the manuscript's headline claim; this subsection states what we can and cannot say about the third. To read an absence of \textsc{REAL} flags on a truly-unresolved population as evidence of preservation is to commit exactly the failure-to-reject-as-success error that this paper is designed to expose (\Cref{philo:limitations}; Supp Note~S0 trap~8 at \Cref{sec:eightTraps}).

Two distinctions matter in making the disclosure precise. First, \textit{truly-unresolved at this dataset} is not the same as \textit{unresolved in principle}: the only escape routes are pooled multi-atlas scale (on the order of $10^{8}$ cells), auxiliary modalities that lift $\Delta$ above the $\pi \Delta^{2} \geq 1.5 \times 10^{-4}$ threshold at Kang scale (protein panels, TCR / BCR sequencing, spatial context), or targeted protocol changes that enrich the population before capture (FACS sorting, modality-aware library preparation). \textit{Truly unresolved} is thus a relation between a population, a measurement regime, and a framework, not a property of the population alone. Second, a motif flagged by \textsc{REAL} remains, even here, a claim about \emph{evaluation}, not about ontology (\Cref{philo:episto_note}). Correspondingly, the absence of a flag in the truly-unresolved stratum remains a claim about \emph{measurement insufficiency}, not about biology. The populations listed below are not excluded from single-cell biology; they are excluded from this framework's current certificate of preservation, under the stated regime.

%   % ---------------------------------------------------------------------
% Limitations subsection — immunology contribution to the Discussion.
% ~350 words. Honest disclosure of populations where REAL/RareShield
% cannot make reliable detection claims at Kang-atlas scale.
% Author: Immunology (agent-mozus3vv). Cross-references Math's
% predicted-detectability table (agent/mathematics@f564cc6).
% Target: sections/08_discussion.tex, as a \subsection after the
% "Implications for immunology" subsection (immuno_discussion.tex).
% ---------------------------------------------------------------------

\subsection{Boundaries of what our framework can currently guarantee}
\label{sec:discussion-limitations-immunology}

Three concrete limitations shape the claims we make. First, the Le~Cam
detectability envelope (\Cref{sec:minimax-detectability}) places a
sharp boundary on which rare populations can be recovered label-free
at any given atlas scale. At the Kang pan-cancer immune atlas
($n = 4.9 \times 10^6$, $B = 30$, per-batch heterogeneity
$\Lambda = 3$), populations with prevalence-separation product
$\pi\Delta < 10^{-3}\sigma$ fall below the detection floor, producing
a true-positive rate under 0.3 at $\alpha = 0.05$. Specifically,
\textit{CXCR5}$^+$\textit{TCF7}$^+$ progenitor-exhausted CD8 T
(TPEX), LAMP3$^+$ mregDC, KIR$^+$ adaptive NK, endothelial tip
cells, and drug-tolerant-persister tumour cells all sit in this
regime at Kang alone. For these populations our framework does not
claim detection from the Kang atlas in isolation; their recovery
requires either auxiliary protein or TCR signal, or pooling across
multiple atlases beyond $10^8$ cells. We disclose this explicitly
and present the populations instead as below-floor targets whose
detection-regime analysis (Extended Data~Fig.~1) itself is a
contribution.

Second, some rare populations are lost before integration rather than
during integration. Neutrophil subsets including LOX-1$^+$
PMN-MDSC and TAN-1 through TAN-5 subsets are systematically
undersampled by the standard 10x Genomics droplet-capture protocol,
and mast cells are similarly depleted relative to their tissue
abundance. Our framework is a post-integration detector and does not
and cannot recover populations that never entered the cell-by-gene
matrix. We flag this as a pre-integration identifiability problem
orthogonal to overcorrection, to be addressed by protocol choice
and by modality-aware capture strategies (FACS-sort enrichment prior
to library preparation; split-pool-barcoded whole-transcriptome
methods that capture neutrophils without selection bias).

Third, our evaluation uses annotated rare populations with their
labels hidden at test time. This is a legitimate form of ground
truth but it is not a direct test on unannotated populations.
Novel populations detected under our framework in the pan-cancer
atlas --- for example, a candidate CXCR5-dim TPEX sub-state
\citep{ourfig5} --- require the full evidence package of
\Cref{sec:evidence-package} plus independent cohort replication
before a definite claim of previously-unannotated biology. Main-text
detections are reported only when $\geq 4$ of the six evidence
criteria are satisfied.

\paragraph*{Telescoping detection recovers below-whole-atlas populations at sub-compartment scale.}
The $\pi\cdot\Delta^2 < 1.5\times10^{-4}$ below-floor rule above is
stated at the whole-atlas level and applies to flat REAL. Hierarchical
REAL, applied at lineage-restricted sub-compartments, moves the
operating point from whole-atlas prevalence to compartment-level
prevalence; at each level the flat floor still applies, but the
$\pi\cdot\Delta^2$ product is typically $10$--$100\times$ larger
because rare populations are often concentrated in specific lineages.
This recovers pulmonary ionocytes at the airway-epithelial sub-compartment
level, and the CXCR5-dim TPEX sub-state at the exhausted-CD8 sub-compartment
level (Supplementary Theorem~S1-hierarchical, Mathematics agent
\citealp{math2026hierarchical}). Our disclosure is therefore
three-tier: populations recovered by flat whole-atlas REAL (above
floor at $\pi\cdot\Delta^2 > 1.5\times10^{-4}$, e.g., MAIT, $\gamma\delta$
T, TREM2$^+$ LAM, myCAF, EMT-hybrid), populations recovered by
hierarchical sub-compartment REAL (above floor at sub-compartment
product, e.g., ionocytes, TPEX sub-state, AS-DC), and populations
unresolved at all compartment depths (endothelial tip cells, drug-tolerant
persister baseline state). Main-text claims are restricted to the
first two tiers; tier-three populations are disclosed as present in
the dataset but beyond the current framework's detection capability.

% ---------------------------------------------------------------------
% Requires: sec:minimax-detectability, sec:evidence-package,
% \citep{ourfig5} is a placeholder for the Fig 5 companion section /
% preprint. Biological Science may prefer \Cref{fig:fig5} or
% \Cref{sec:results-tpex-substate} depending on skeleton.
% ---------------------------------------------------------------------
         <-- unchanged from Immunology branch
%   % ---------------------------------------------------------------------
% Joint closing paragraph for the below-floor Limitations subsection.
% Co-authored with Philosophy (agent-mozur8ub). Placement flow:
%
%   \input{philo_below_floor_head.tex}    (Philosophy's epistemic frame)
%   \input{immuno_limitations.tex}        (Immunology's biological disclosure)
%   \input{immuno_below_floor_close.tex}  (THIS FILE — joint affirmative)
%
% Closes the subsection on the telescoping-detection affirmative:
% most below-flat-floor populations are recovered under hierarchical
% sub-compartment REAL.
% ---------------------------------------------------------------------

\paragraph*{Recovery under hierarchical application.}
Notwithstanding the strict boundary established above, the \emph{truly-unresolved} stratum is surprisingly small. Of the 22 immune and tumor populations in our registry whose whole-atlas product $\pi \cdot \Delta^{2}$ places them into one of the three detection tiers, 20 are recoverable through flat or hierarchical application of \textsc{REAL} at an appropriate compartment scale (\Cref{thm:S1-hierarchical}; \Cref{sec:methods-oracle-score} protocol): 7 populations are detected at whole-atlas scale (flat tier --- MAIT, $\gamma\delta$ T V$\delta1$-bias, TREM2$^{+}$ LAM, myCAF, iCAF, plasma-TLS, EMT-hybrid-near-pole), and 13 additional populations cross the detection floor only under hierarchical application at their natural sub-compartment scale (hierarchical tier). Pulmonary ionocytes pass detection at the airway-epithelial sub-compartment of HLCA (effective $\pi \approx 5 \times 10^{-3}$ in $n \approx 240{,}000$ airway cells); \textit{CXCR5}$^{+}$\textit{TCF7}$^{+}$ progenitor-exhausted CD8 T cells pass detection at the exhausted-CD8 sub-compartment of pooled melanoma atlases ($\pi \approx 10^{-2}$ in $n \approx 100{,}000$ cells); LAMP3$^{+}$ mregDC, AS-DC, and HCMV-imprinted adaptive NK pass detection at dendritic-cell and natural-killer sub-compartments respectively. The two populations that remain below the floor at every compartment depth are endothelial tip cells (which sit at the intersection of a small-mass vascular niche and a narrow transcriptional separation) and the drug-tolerant persister baseline state in oncology (Tumor Biology registry row TU-T-001), with mast cells in tumors additionally disclosed as out-of-scope for post-integration detection (mode~4 pre-integration capture loss). These unresolved entries are explicitly deferred to follow-up work: their recovery would require either the multi-atlas pooling regime ($n > 10^{8}$ cells) or auxiliary-modality augmentation, both of which lie beyond this paper's scope. This three-tier disclosure --- recovered-flat, recovered-hierarchical, deferred --- is the most honest reading of what our framework currently certifies, and it preserves the paper's headline claim on the above-floor and hierarchical-recovered strata while placing no claim on the truly-unresolved tier.

% ---------------------------------------------------------------------
% Cross-references needed (shared with philo_below_floor_head + immuno_limitations):
%   thm:S1-hierarchical        Math hierarchical REAL theorem
%   sec:methods-oracle-score   Methods — oracle-score verification
% Bib entries: none new.
% ---------------------------------------------------------------------
   <-- unchanged from Immunology branch
%
% Use this file INSTEAD OF the 5-file chain if BioSci prefers a single splice point.
% Immunology's two \input calls (limitations + below-floor-close) are PRESERVED as
% \input calls inside this file rather than inlined, because they are owned by
% the immunology branch and Immunology should retain edit control over their content.
%
% Placement: sections/08_discussion.tex, replacing the 5-file \input chain with
% a single % Philosophy §08 epistemic-scope block — single-file consolidation.
% Merges the five \input calls that BioSci currently threads through into one file
% for easier future-edit management. Semantically identical to the input chain:
%   \input{handoffs/philo_discussion_meno.tex}
%   \input{handoffs/philo_limitations.tex}
%   \input{handoffs/philo_below_floor_head.tex}
%   \input{handoffs/immuno_limitations.tex}         <-- unchanged from Immunology branch
%   \input{handoffs/immuno_below_floor_close.tex}   <-- unchanged from Immunology branch
%
% Use this file INSTEAD OF the 5-file chain if BioSci prefers a single splice point.
% Immunology's two \input calls (limitations + below-floor-close) are PRESERVED as
% \input calls inside this file rather than inlined, because they are owned by
% the immunology branch and Immunology should retain edit control over their content.
%
% Placement: sections/08_discussion.tex, replacing the 5-file \input chain with
% a single \input{handoffs/philo_08_epistemic_block.tex}.
%
% Author: Philosophy (agent-mozur8ub). Last updated for v2.4 pre-submission.

% =============================================================================
% §08 EPISTEMIC SCOPE BLOCK — single-file consolidation
% =============================================================================

% -----------------------------------------------------------------------------
% (1) Meno-problem framing — opens the epistemic scope block.
% Content from philo_discussion_meno.tex (agent/philosophy@<head>).
% ≈250w. Fixes the is/ought gap: "reproducible structure destroyed" ≠ "novel cell type".
% -----------------------------------------------------------------------------

\paragraph*{What a REAL flag is, and is not.}
A motif flagged by REAL is the claim that a local, reproducible geometric structure---witnessed across independent batches before integration, with measurable local-mass distortion after---has been erased by the integration map beyond what admissible batch-effect deformations can explain. This is a claim about \emph{evaluation}, not about ontology. It is not a claim that a new cell type exists. A motif $w \in \mathcal{W}$ with low $\mathrm{LRS}(w)$ licenses exactly one inference: the integration has discarded reproducible structure of the kind a biologist would want to interrogate. Whether that structure corresponds to a novel cell type, a transitional state along an already-known trajectory, a rare disease-specific configuration, or a technical sub-mode that merely looks biological, is a further question REAL does not answer. Upgrading a motif to a cell-type claim requires the usual burden of single-cell biology: independent sampling, orthogonal markers, targeted protocol (sorting, spatial validation, experimental perturbation), and community replication. REAL's value is that it identifies \emph{where to look}, not what has been found. Conflating the two would recreate the very error this paper is designed to expose: treating a framework-internal signal as a fact about the world. Our strong claim is therefore deliberately asymmetric. A high motif set $|\mathcal{W}|$ under a candidate integration is evidence of measurement insufficiency; a low motif set is evidence that the integration respected the syntactic reproducibility criterion. Ontology is downstream of both.

% -----------------------------------------------------------------------------
% (2) Limitations and epistemic scope — Route-C exchangeability boundaries.
% Content from philo_limitations.tex. ≤500w.
% -----------------------------------------------------------------------------

\subsection*{Limitations and epistemic scope}
\label{philo:limitations}
REAL relies on an exchangeability assumption between the known rare populations we hold out for calibration (pancreatic $\epsilon$ cells, pulmonary ionocytes, LAMP3$^+$ mregDC, TPEX, FOLR2$^+$ perivascular macrophages, apCAF, and the populations enumerated in Table~S\ref{tab:heldout}) and the genuinely unannotated rare subpopulations to which we transfer the estimate. The assumption is formally stated in Lemma~\ref{lem:exchangeability} as a \emph{bound}: the held-out-rare loss rate upper-bounds the expected loss rate for unannotated rare subpopulations of comparable abundance, local geometry, and batch profile. We emphasise four boundary conditions under which the bound is not tight. First, the bound degrades when the unannotated population lies outside the span of the highly-variable genes used to construct the latent space; REAL cannot protect what the representation cannot see. Second, the bound degrades when the unannotated population has a qualitatively different batch-effect signature from the calibration populations---for instance, a disease-specific state observed in only one batch. Third, the bound is conservative for populations whose local density profile is not well-approximated by the level-set stratification we use; transitional cells along a continuum can appear as low-density but are not rare in the sampling sense. Fourth, our use of known rare populations as proxies for unknown ones inherits a selection bias: the populations we can hold out are those that have already been discovered, and the discovered set may not be exchangeable with the undiscovered set.

Our strongest defensible claim is therefore deliberately limited. REAL provides (i) a label-free, invariant-based detectability framework calibrated on held-out known rare subpopulations, (ii) an integration strategy (RareShield) that demonstrably reduces the held-out-rare loss rate at pan-cancer-atlas scale without degrading batch correction for abundant cell types, and (iii) an explicit exchangeability bound under which the held-out loss rate is an upper bound on the loss rate of truly unannotated rare subpopulations. A REAL-flagged motif $w \in \mathcal{W}$ is a claim that reproducible structure has been destroyed, not a claim that a new cell type exists; upgrading to an ontological claim requires the usual experimental burden (sorting, orthogonal markers, spatial validation, replication). The complementary negative: a dataset with low $|\mathcal{W}|$ under a candidate integration is evidence that no cross-batch-reproducible rare structure was erased, but not a guarantee that no such structure existed---by construction, we cannot rule out subpopulations whose reproducibility signal was below our sampling resolution or outside the evaluated feature span. These boundaries are not weaknesses of REAL specifically; they are the price of reasoning about what the evaluation framework cannot directly observe. Naming the price is what distinguishes an epistemically calibrated claim from a framework-captured one.

% -----------------------------------------------------------------------------
% (3) Below the power floor — epistemic frame head (three-tier stratification).
% Content from philo_below_floor_head.tex.
% -----------------------------------------------------------------------------

\subsection*{Below the power floor: what the framework cannot yet certify}
\label{sec:below-floor-epistemic}

A framework's sensitivity is itself a piece of information about the framework, not about the biology. The minimax envelope of \Cref{thm:minimax} and the hierarchical-detectability result of \Cref{thm:S1-hierarchical} together partition the rare-population landscape at Kang~2024 scale into three strata: \emph{above-floor-flat} populations for which whole-atlas REAL is sufficient; \emph{hierarchical-recovered} populations which cross the detectability floor only after sub-compartment telescoping (per \Cref{sec:S1-envelope}); and \emph{truly-unresolved} populations for which neither flat nor hierarchical REAL attains the required power at Kang depth. The first two strata carry the manuscript's headline claim; this subsection states what we can and cannot say about the third. To read an absence of \textsc{REAL} flags on a truly-unresolved population as evidence of preservation is to commit exactly the failure-to-reject-as-success error that this paper is designed to expose (\Cref{philo:limitations}; Supp Note~S0 trap~8 at \Cref{sec:eightTraps}).

Two distinctions matter in making the disclosure precise. First, \textit{truly-unresolved at this dataset} is not the same as \textit{unresolved in principle}: the only escape routes are pooled multi-atlas scale (on the order of $10^{8}$ cells), auxiliary modalities that lift $\Delta$ above the $\pi \Delta^{2} \geq 1.5 \times 10^{-4}$ threshold at Kang scale (protein panels, TCR / BCR sequencing, spatial context), or targeted protocol changes that enrich the population before capture (FACS sorting, modality-aware library preparation). \textit{Truly unresolved} is thus a relation between a population, a measurement regime, and a framework, not a property of the population alone. Second, a motif flagged by \textsc{REAL} remains, even here, a claim about \emph{evaluation}, not about ontology (\Cref{philo:episto_note}). Correspondingly, the absence of a flag in the truly-unresolved stratum remains a claim about \emph{measurement insufficiency}, not about biology. The populations listed below are not excluded from single-cell biology; they are excluded from this framework's current certificate of preservation, under the stated regime.

% -----------------------------------------------------------------------------
% (4)-(5) Immunology's biological disclosure + joint closing.
% These files live in the immunology branch; Immunology retains edit control.
% They are \input'd here rather than inlined, so Immunology-branch edits propagate
% automatically on the next BioSci pull.
% -----------------------------------------------------------------------------

\input{handoffs/immuno_limitations.tex}
\input{handoffs/immuno_below_floor_close.tex}

% =============================================================================
% END §08 EPISTEMIC SCOPE BLOCK
% =============================================================================

% Integration note for BioSci:
%   - To use this consolidated file: replace the 5-file \input chain in
%     sections/08_discussion.tex with:
%       \input{handoffs/philo_08_epistemic_block.tex}
%   - To keep the 5-file chain: ignore this file; it is semantically identical
%     to the chain at this commit. Consolidation is a convenience, not a correctness fix.
%   - If Philosophy or Immunology updates any of the five source files, re-generating
%     this consolidated file requires Philosophy to re-splice; the \input calls for
%     Immunology's two files remain untouched.
%
% Lint status: passes lint:epistemic clean (trap 8 / trap 4 / IM.4 guards all present).
.
%
% Author: Philosophy (agent-mozur8ub). Last updated for v2.4 pre-submission.

% =============================================================================
% §08 EPISTEMIC SCOPE BLOCK — single-file consolidation
% =============================================================================

% -----------------------------------------------------------------------------
% (1) Meno-problem framing — opens the epistemic scope block.
% Content from philo_discussion_meno.tex (agent/philosophy@<head>).
% ≈250w. Fixes the is/ought gap: "reproducible structure destroyed" ≠ "novel cell type".
% -----------------------------------------------------------------------------

\paragraph*{What a REAL flag is, and is not.}
A motif flagged by REAL is the claim that a local, reproducible geometric structure---witnessed across independent batches before integration, with measurable local-mass distortion after---has been erased by the integration map beyond what admissible batch-effect deformations can explain. This is a claim about \emph{evaluation}, not about ontology. It is not a claim that a new cell type exists. A motif $w \in \mathcal{W}$ with low $\mathrm{LRS}(w)$ licenses exactly one inference: the integration has discarded reproducible structure of the kind a biologist would want to interrogate. Whether that structure corresponds to a novel cell type, a transitional state along an already-known trajectory, a rare disease-specific configuration, or a technical sub-mode that merely looks biological, is a further question REAL does not answer. Upgrading a motif to a cell-type claim requires the usual burden of single-cell biology: independent sampling, orthogonal markers, targeted protocol (sorting, spatial validation, experimental perturbation), and community replication. REAL's value is that it identifies \emph{where to look}, not what has been found. Conflating the two would recreate the very error this paper is designed to expose: treating a framework-internal signal as a fact about the world. Our strong claim is therefore deliberately asymmetric. A high motif set $|\mathcal{W}|$ under a candidate integration is evidence of measurement insufficiency; a low motif set is evidence that the integration respected the syntactic reproducibility criterion. Ontology is downstream of both.

% -----------------------------------------------------------------------------
% (2) Limitations and epistemic scope — Route-C exchangeability boundaries.
% Content from philo_limitations.tex. ≤500w.
% -----------------------------------------------------------------------------

\subsection*{Limitations and epistemic scope}
\label{philo:limitations}
REAL relies on an exchangeability assumption between the known rare populations we hold out for calibration (pancreatic $\epsilon$ cells, pulmonary ionocytes, LAMP3$^+$ mregDC, TPEX, FOLR2$^+$ perivascular macrophages, apCAF, and the populations enumerated in Table~S\ref{tab:heldout}) and the genuinely unannotated rare subpopulations to which we transfer the estimate. The assumption is formally stated in Lemma~\ref{lem:exchangeability} as a \emph{bound}: the held-out-rare loss rate upper-bounds the expected loss rate for unannotated rare subpopulations of comparable abundance, local geometry, and batch profile. We emphasise four boundary conditions under which the bound is not tight. First, the bound degrades when the unannotated population lies outside the span of the highly-variable genes used to construct the latent space; REAL cannot protect what the representation cannot see. Second, the bound degrades when the unannotated population has a qualitatively different batch-effect signature from the calibration populations---for instance, a disease-specific state observed in only one batch. Third, the bound is conservative for populations whose local density profile is not well-approximated by the level-set stratification we use; transitional cells along a continuum can appear as low-density but are not rare in the sampling sense. Fourth, our use of known rare populations as proxies for unknown ones inherits a selection bias: the populations we can hold out are those that have already been discovered, and the discovered set may not be exchangeable with the undiscovered set.

Our strongest defensible claim is therefore deliberately limited. REAL provides (i) a label-free, invariant-based detectability framework calibrated on held-out known rare subpopulations, (ii) an integration strategy (RareShield) that demonstrably reduces the held-out-rare loss rate at pan-cancer-atlas scale without degrading batch correction for abundant cell types, and (iii) an explicit exchangeability bound under which the held-out loss rate is an upper bound on the loss rate of truly unannotated rare subpopulations. A REAL-flagged motif $w \in \mathcal{W}$ is a claim that reproducible structure has been destroyed, not a claim that a new cell type exists; upgrading to an ontological claim requires the usual experimental burden (sorting, orthogonal markers, spatial validation, replication). The complementary negative: a dataset with low $|\mathcal{W}|$ under a candidate integration is evidence that no cross-batch-reproducible rare structure was erased, but not a guarantee that no such structure existed---by construction, we cannot rule out subpopulations whose reproducibility signal was below our sampling resolution or outside the evaluated feature span. These boundaries are not weaknesses of REAL specifically; they are the price of reasoning about what the evaluation framework cannot directly observe. Naming the price is what distinguishes an epistemically calibrated claim from a framework-captured one.

% -----------------------------------------------------------------------------
% (3) Below the power floor — epistemic frame head (three-tier stratification).
% Content from philo_below_floor_head.tex.
% -----------------------------------------------------------------------------

\subsection*{Below the power floor: what the framework cannot yet certify}
\label{sec:below-floor-epistemic}

A framework's sensitivity is itself a piece of information about the framework, not about the biology. The minimax envelope of \Cref{thm:minimax} and the hierarchical-detectability result of \Cref{thm:S1-hierarchical} together partition the rare-population landscape at Kang~2024 scale into three strata: \emph{above-floor-flat} populations for which whole-atlas REAL is sufficient; \emph{hierarchical-recovered} populations which cross the detectability floor only after sub-compartment telescoping (per \Cref{sec:S1-envelope}); and \emph{truly-unresolved} populations for which neither flat nor hierarchical REAL attains the required power at Kang depth. The first two strata carry the manuscript's headline claim; this subsection states what we can and cannot say about the third. To read an absence of \textsc{REAL} flags on a truly-unresolved population as evidence of preservation is to commit exactly the failure-to-reject-as-success error that this paper is designed to expose (\Cref{philo:limitations}; Supp Note~S0 trap~8 at \Cref{sec:eightTraps}).

Two distinctions matter in making the disclosure precise. First, \textit{truly-unresolved at this dataset} is not the same as \textit{unresolved in principle}: the only escape routes are pooled multi-atlas scale (on the order of $10^{8}$ cells), auxiliary modalities that lift $\Delta$ above the $\pi \Delta^{2} \geq 1.5 \times 10^{-4}$ threshold at Kang scale (protein panels, TCR / BCR sequencing, spatial context), or targeted protocol changes that enrich the population before capture (FACS sorting, modality-aware library preparation). \textit{Truly unresolved} is thus a relation between a population, a measurement regime, and a framework, not a property of the population alone. Second, a motif flagged by \textsc{REAL} remains, even here, a claim about \emph{evaluation}, not about ontology (\Cref{philo:episto_note}). Correspondingly, the absence of a flag in the truly-unresolved stratum remains a claim about \emph{measurement insufficiency}, not about biology. The populations listed below are not excluded from single-cell biology; they are excluded from this framework's current certificate of preservation, under the stated regime.

% -----------------------------------------------------------------------------
% (4)-(5) Immunology's biological disclosure + joint closing.
% These files live in the immunology branch; Immunology retains edit control.
% They are \input'd here rather than inlined, so Immunology-branch edits propagate
% automatically on the next BioSci pull.
% -----------------------------------------------------------------------------

% ---------------------------------------------------------------------
% Limitations subsection — immunology contribution to the Discussion.
% ~350 words. Honest disclosure of populations where REAL/RareShield
% cannot make reliable detection claims at Kang-atlas scale.
% Author: Immunology (agent-mozus3vv). Cross-references Math's
% predicted-detectability table (agent/mathematics@f564cc6).
% Target: sections/08_discussion.tex, as a \subsection after the
% "Implications for immunology" subsection (immuno_discussion.tex).
% ---------------------------------------------------------------------

\subsection{Boundaries of what our framework can currently guarantee}
\label{sec:discussion-limitations-immunology}

Three concrete limitations shape the claims we make. First, the Le~Cam
detectability envelope (\Cref{sec:minimax-detectability}) places a
sharp boundary on which rare populations can be recovered label-free
at any given atlas scale. At the Kang pan-cancer immune atlas
($n = 4.9 \times 10^6$, $B = 30$, per-batch heterogeneity
$\Lambda = 3$), populations with prevalence-separation product
$\pi\Delta < 10^{-3}\sigma$ fall below the detection floor, producing
a true-positive rate under 0.3 at $\alpha = 0.05$. Specifically,
\textit{CXCR5}$^+$\textit{TCF7}$^+$ progenitor-exhausted CD8 T
(TPEX), LAMP3$^+$ mregDC, KIR$^+$ adaptive NK, endothelial tip
cells, and drug-tolerant-persister tumour cells all sit in this
regime at Kang alone. For these populations our framework does not
claim detection from the Kang atlas in isolation; their recovery
requires either auxiliary protein or TCR signal, or pooling across
multiple atlases beyond $10^8$ cells. We disclose this explicitly
and present the populations instead as below-floor targets whose
detection-regime analysis (Extended Data~Fig.~1) itself is a
contribution.

Second, some rare populations are lost before integration rather than
during integration. Neutrophil subsets including LOX-1$^+$
PMN-MDSC and TAN-1 through TAN-5 subsets are systematically
undersampled by the standard 10x Genomics droplet-capture protocol,
and mast cells are similarly depleted relative to their tissue
abundance. Our framework is a post-integration detector and does not
and cannot recover populations that never entered the cell-by-gene
matrix. We flag this as a pre-integration identifiability problem
orthogonal to overcorrection, to be addressed by protocol choice
and by modality-aware capture strategies (FACS-sort enrichment prior
to library preparation; split-pool-barcoded whole-transcriptome
methods that capture neutrophils without selection bias).

Third, our evaluation uses annotated rare populations with their
labels hidden at test time. This is a legitimate form of ground
truth but it is not a direct test on unannotated populations.
Novel populations detected under our framework in the pan-cancer
atlas --- for example, a candidate CXCR5-dim TPEX sub-state
\citep{ourfig5} --- require the full evidence package of
\Cref{sec:evidence-package} plus independent cohort replication
before a definite claim of previously-unannotated biology. Main-text
detections are reported only when $\geq 4$ of the six evidence
criteria are satisfied.

\paragraph*{Telescoping detection recovers below-whole-atlas populations at sub-compartment scale.}
The $\pi\cdot\Delta^2 < 1.5\times10^{-4}$ below-floor rule above is
stated at the whole-atlas level and applies to flat REAL. Hierarchical
REAL, applied at lineage-restricted sub-compartments, moves the
operating point from whole-atlas prevalence to compartment-level
prevalence; at each level the flat floor still applies, but the
$\pi\cdot\Delta^2$ product is typically $10$--$100\times$ larger
because rare populations are often concentrated in specific lineages.
This recovers pulmonary ionocytes at the airway-epithelial sub-compartment
level, and the CXCR5-dim TPEX sub-state at the exhausted-CD8 sub-compartment
level (Supplementary Theorem~S1-hierarchical, Mathematics agent
\citealp{math2026hierarchical}). Our disclosure is therefore
three-tier: populations recovered by flat whole-atlas REAL (above
floor at $\pi\cdot\Delta^2 > 1.5\times10^{-4}$, e.g., MAIT, $\gamma\delta$
T, TREM2$^+$ LAM, myCAF, EMT-hybrid), populations recovered by
hierarchical sub-compartment REAL (above floor at sub-compartment
product, e.g., ionocytes, TPEX sub-state, AS-DC), and populations
unresolved at all compartment depths (endothelial tip cells, drug-tolerant
persister baseline state). Main-text claims are restricted to the
first two tiers; tier-three populations are disclosed as present in
the dataset but beyond the current framework's detection capability.

% ---------------------------------------------------------------------
% Requires: sec:minimax-detectability, sec:evidence-package,
% \citep{ourfig5} is a placeholder for the Fig 5 companion section /
% preprint. Biological Science may prefer \Cref{fig:fig5} or
% \Cref{sec:results-tpex-substate} depending on skeleton.
% ---------------------------------------------------------------------

% ---------------------------------------------------------------------
% Joint closing paragraph for the below-floor Limitations subsection.
% Co-authored with Philosophy (agent-mozur8ub). Placement flow:
%
%   \input{philo_below_floor_head.tex}    (Philosophy's epistemic frame)
%   \input{immuno_limitations.tex}        (Immunology's biological disclosure)
%   \input{immuno_below_floor_close.tex}  (THIS FILE — joint affirmative)
%
% Closes the subsection on the telescoping-detection affirmative:
% most below-flat-floor populations are recovered under hierarchical
% sub-compartment REAL.
% ---------------------------------------------------------------------

\paragraph*{Recovery under hierarchical application.}
Notwithstanding the strict boundary established above, the \emph{truly-unresolved} stratum is surprisingly small. Of the 22 immune and tumor populations in our registry whose whole-atlas product $\pi \cdot \Delta^{2}$ places them into one of the three detection tiers, 20 are recoverable through flat or hierarchical application of \textsc{REAL} at an appropriate compartment scale (\Cref{thm:S1-hierarchical}; \Cref{sec:methods-oracle-score} protocol): 7 populations are detected at whole-atlas scale (flat tier --- MAIT, $\gamma\delta$ T V$\delta1$-bias, TREM2$^{+}$ LAM, myCAF, iCAF, plasma-TLS, EMT-hybrid-near-pole), and 13 additional populations cross the detection floor only under hierarchical application at their natural sub-compartment scale (hierarchical tier). Pulmonary ionocytes pass detection at the airway-epithelial sub-compartment of HLCA (effective $\pi \approx 5 \times 10^{-3}$ in $n \approx 240{,}000$ airway cells); \textit{CXCR5}$^{+}$\textit{TCF7}$^{+}$ progenitor-exhausted CD8 T cells pass detection at the exhausted-CD8 sub-compartment of pooled melanoma atlases ($\pi \approx 10^{-2}$ in $n \approx 100{,}000$ cells); LAMP3$^{+}$ mregDC, AS-DC, and HCMV-imprinted adaptive NK pass detection at dendritic-cell and natural-killer sub-compartments respectively. The two populations that remain below the floor at every compartment depth are endothelial tip cells (which sit at the intersection of a small-mass vascular niche and a narrow transcriptional separation) and the drug-tolerant persister baseline state in oncology (Tumor Biology registry row TU-T-001), with mast cells in tumors additionally disclosed as out-of-scope for post-integration detection (mode~4 pre-integration capture loss). These unresolved entries are explicitly deferred to follow-up work: their recovery would require either the multi-atlas pooling regime ($n > 10^{8}$ cells) or auxiliary-modality augmentation, both of which lie beyond this paper's scope. This three-tier disclosure --- recovered-flat, recovered-hierarchical, deferred --- is the most honest reading of what our framework currently certifies, and it preserves the paper's headline claim on the above-floor and hierarchical-recovered strata while placing no claim on the truly-unresolved tier.

% ---------------------------------------------------------------------
% Cross-references needed (shared with philo_below_floor_head + immuno_limitations):
%   thm:S1-hierarchical        Math hierarchical REAL theorem
%   sec:methods-oracle-score   Methods — oracle-score verification
% Bib entries: none new.
% ---------------------------------------------------------------------


% =============================================================================
% END §08 EPISTEMIC SCOPE BLOCK
% =============================================================================

% Integration note for BioSci:
%   - To use this consolidated file: replace the 5-file \input chain in
%     sections/08_discussion.tex with:
%       % Philosophy §08 epistemic-scope block — single-file consolidation.
% Merges the five \input calls that BioSci currently threads through into one file
% for easier future-edit management. Semantically identical to the input chain:
%   \input{handoffs/philo_discussion_meno.tex}
%   \input{handoffs/philo_limitations.tex}
%   \input{handoffs/philo_below_floor_head.tex}
%   \input{handoffs/immuno_limitations.tex}         <-- unchanged from Immunology branch
%   \input{handoffs/immuno_below_floor_close.tex}   <-- unchanged from Immunology branch
%
% Use this file INSTEAD OF the 5-file chain if BioSci prefers a single splice point.
% Immunology's two \input calls (limitations + below-floor-close) are PRESERVED as
% \input calls inside this file rather than inlined, because they are owned by
% the immunology branch and Immunology should retain edit control over their content.
%
% Placement: sections/08_discussion.tex, replacing the 5-file \input chain with
% a single \input{handoffs/philo_08_epistemic_block.tex}.
%
% Author: Philosophy (agent-mozur8ub). Last updated for v2.4 pre-submission.

% =============================================================================
% §08 EPISTEMIC SCOPE BLOCK — single-file consolidation
% =============================================================================

% -----------------------------------------------------------------------------
% (1) Meno-problem framing — opens the epistemic scope block.
% Content from philo_discussion_meno.tex (agent/philosophy@<head>).
% ≈250w. Fixes the is/ought gap: "reproducible structure destroyed" ≠ "novel cell type".
% -----------------------------------------------------------------------------

\paragraph*{What a REAL flag is, and is not.}
A motif flagged by REAL is the claim that a local, reproducible geometric structure---witnessed across independent batches before integration, with measurable local-mass distortion after---has been erased by the integration map beyond what admissible batch-effect deformations can explain. This is a claim about \emph{evaluation}, not about ontology. It is not a claim that a new cell type exists. A motif $w \in \mathcal{W}$ with low $\mathrm{LRS}(w)$ licenses exactly one inference: the integration has discarded reproducible structure of the kind a biologist would want to interrogate. Whether that structure corresponds to a novel cell type, a transitional state along an already-known trajectory, a rare disease-specific configuration, or a technical sub-mode that merely looks biological, is a further question REAL does not answer. Upgrading a motif to a cell-type claim requires the usual burden of single-cell biology: independent sampling, orthogonal markers, targeted protocol (sorting, spatial validation, experimental perturbation), and community replication. REAL's value is that it identifies \emph{where to look}, not what has been found. Conflating the two would recreate the very error this paper is designed to expose: treating a framework-internal signal as a fact about the world. Our strong claim is therefore deliberately asymmetric. A high motif set $|\mathcal{W}|$ under a candidate integration is evidence of measurement insufficiency; a low motif set is evidence that the integration respected the syntactic reproducibility criterion. Ontology is downstream of both.

% -----------------------------------------------------------------------------
% (2) Limitations and epistemic scope — Route-C exchangeability boundaries.
% Content from philo_limitations.tex. ≤500w.
% -----------------------------------------------------------------------------

\subsection*{Limitations and epistemic scope}
\label{philo:limitations}
REAL relies on an exchangeability assumption between the known rare populations we hold out for calibration (pancreatic $\epsilon$ cells, pulmonary ionocytes, LAMP3$^+$ mregDC, TPEX, FOLR2$^+$ perivascular macrophages, apCAF, and the populations enumerated in Table~S\ref{tab:heldout}) and the genuinely unannotated rare subpopulations to which we transfer the estimate. The assumption is formally stated in Lemma~\ref{lem:exchangeability} as a \emph{bound}: the held-out-rare loss rate upper-bounds the expected loss rate for unannotated rare subpopulations of comparable abundance, local geometry, and batch profile. We emphasise four boundary conditions under which the bound is not tight. First, the bound degrades when the unannotated population lies outside the span of the highly-variable genes used to construct the latent space; REAL cannot protect what the representation cannot see. Second, the bound degrades when the unannotated population has a qualitatively different batch-effect signature from the calibration populations---for instance, a disease-specific state observed in only one batch. Third, the bound is conservative for populations whose local density profile is not well-approximated by the level-set stratification we use; transitional cells along a continuum can appear as low-density but are not rare in the sampling sense. Fourth, our use of known rare populations as proxies for unknown ones inherits a selection bias: the populations we can hold out are those that have already been discovered, and the discovered set may not be exchangeable with the undiscovered set.

Our strongest defensible claim is therefore deliberately limited. REAL provides (i) a label-free, invariant-based detectability framework calibrated on held-out known rare subpopulations, (ii) an integration strategy (RareShield) that demonstrably reduces the held-out-rare loss rate at pan-cancer-atlas scale without degrading batch correction for abundant cell types, and (iii) an explicit exchangeability bound under which the held-out loss rate is an upper bound on the loss rate of truly unannotated rare subpopulations. A REAL-flagged motif $w \in \mathcal{W}$ is a claim that reproducible structure has been destroyed, not a claim that a new cell type exists; upgrading to an ontological claim requires the usual experimental burden (sorting, orthogonal markers, spatial validation, replication). The complementary negative: a dataset with low $|\mathcal{W}|$ under a candidate integration is evidence that no cross-batch-reproducible rare structure was erased, but not a guarantee that no such structure existed---by construction, we cannot rule out subpopulations whose reproducibility signal was below our sampling resolution or outside the evaluated feature span. These boundaries are not weaknesses of REAL specifically; they are the price of reasoning about what the evaluation framework cannot directly observe. Naming the price is what distinguishes an epistemically calibrated claim from a framework-captured one.

% -----------------------------------------------------------------------------
% (3) Below the power floor — epistemic frame head (three-tier stratification).
% Content from philo_below_floor_head.tex.
% -----------------------------------------------------------------------------

\subsection*{Below the power floor: what the framework cannot yet certify}
\label{sec:below-floor-epistemic}

A framework's sensitivity is itself a piece of information about the framework, not about the biology. The minimax envelope of \Cref{thm:minimax} and the hierarchical-detectability result of \Cref{thm:S1-hierarchical} together partition the rare-population landscape at Kang~2024 scale into three strata: \emph{above-floor-flat} populations for which whole-atlas REAL is sufficient; \emph{hierarchical-recovered} populations which cross the detectability floor only after sub-compartment telescoping (per \Cref{sec:S1-envelope}); and \emph{truly-unresolved} populations for which neither flat nor hierarchical REAL attains the required power at Kang depth. The first two strata carry the manuscript's headline claim; this subsection states what we can and cannot say about the third. To read an absence of \textsc{REAL} flags on a truly-unresolved population as evidence of preservation is to commit exactly the failure-to-reject-as-success error that this paper is designed to expose (\Cref{philo:limitations}; Supp Note~S0 trap~8 at \Cref{sec:eightTraps}).

Two distinctions matter in making the disclosure precise. First, \textit{truly-unresolved at this dataset} is not the same as \textit{unresolved in principle}: the only escape routes are pooled multi-atlas scale (on the order of $10^{8}$ cells), auxiliary modalities that lift $\Delta$ above the $\pi \Delta^{2} \geq 1.5 \times 10^{-4}$ threshold at Kang scale (protein panels, TCR / BCR sequencing, spatial context), or targeted protocol changes that enrich the population before capture (FACS sorting, modality-aware library preparation). \textit{Truly unresolved} is thus a relation between a population, a measurement regime, and a framework, not a property of the population alone. Second, a motif flagged by \textsc{REAL} remains, even here, a claim about \emph{evaluation}, not about ontology (\Cref{philo:episto_note}). Correspondingly, the absence of a flag in the truly-unresolved stratum remains a claim about \emph{measurement insufficiency}, not about biology. The populations listed below are not excluded from single-cell biology; they are excluded from this framework's current certificate of preservation, under the stated regime.

% -----------------------------------------------------------------------------
% (4)-(5) Immunology's biological disclosure + joint closing.
% These files live in the immunology branch; Immunology retains edit control.
% They are \input'd here rather than inlined, so Immunology-branch edits propagate
% automatically on the next BioSci pull.
% -----------------------------------------------------------------------------

\input{handoffs/immuno_limitations.tex}
\input{handoffs/immuno_below_floor_close.tex}

% =============================================================================
% END §08 EPISTEMIC SCOPE BLOCK
% =============================================================================

% Integration note for BioSci:
%   - To use this consolidated file: replace the 5-file \input chain in
%     sections/08_discussion.tex with:
%       \input{handoffs/philo_08_epistemic_block.tex}
%   - To keep the 5-file chain: ignore this file; it is semantically identical
%     to the chain at this commit. Consolidation is a convenience, not a correctness fix.
%   - If Philosophy or Immunology updates any of the five source files, re-generating
%     this consolidated file requires Philosophy to re-splice; the \input calls for
%     Immunology's two files remain untouched.
%
% Lint status: passes lint:epistemic clean (trap 8 / trap 4 / IM.4 guards all present).

%   - To keep the 5-file chain: ignore this file; it is semantically identical
%     to the chain at this commit. Consolidation is a convenience, not a correctness fix.
%   - If Philosophy or Immunology updates any of the five source files, re-generating
%     this consolidated file requires Philosophy to re-splice; the \input calls for
%     Immunology's two files remain untouched.
%
% Lint status: passes lint:epistemic clean (trap 8 / trap 4 / IM.4 guards all present).

%   - To keep the 5-file chain: ignore this file; it is semantically identical
%     to the chain at this commit. Consolidation is a convenience, not a correctness fix.
%   - If Philosophy or Immunology updates any of the five source files, re-generating
%     this consolidated file requires Philosophy to re-splice; the \input calls for
%     Immunology's two files remain untouched.
%
% Lint status: passes lint:epistemic clean (trap 8 / trap 4 / IM.4 guards all present).

%   - To keep the 5-file chain: ignore this file; it is semantically identical
%     to the chain at this commit. Consolidation is a convenience, not a correctness fix.
%   - If Philosophy or Immunology updates any of the five source files, re-generating
%     this consolidated file requires Philosophy to re-splice; the \input calls for
%     Immunology's two files remain untouched.
%
% Lint status: passes lint:epistemic clean (trap 8 / trap 4 / IM.4 guards all present).
